\PassOptionsToPackage{unicode}{hyperref}
\PassOptionsToPackage{hyphens}{url}

\documentclass[11pt,a4paper]{article}

\usepackage{graphicx}

\usepackage{tikz}
\usetikzlibrary{shapes.geometric, arrows.meta, positioning, fit, backgrounds, calc}

\usepackage[top=2.5cm, bottom=2.5cm, left=2.5cm, right=2.5cm]{geometry}

\usepackage{xcolor}
\usepackage{colortbl}
\usepackage{amsmath,amssymb}
\setcounter{secnumdepth}{-\maxdimen}
\usepackage{iftex}

\ifPDFTeX
  \usepackage[utf8]{inputenc}
  \usepackage[T1]{fontenc}
  \usepackage{mathptmx}
\else
  \usepackage{fontspec}
  \setmainfont{Helvetica Neue}
  \setsansfont{Helvetica Neue}
  \setmonofont{Menlo}
  \usepackage{unicode-math}
\fi

\usepackage{textcomp}
\usepackage{lmodern}

\IfFileExists{upquote.sty}{\usepackage{upquote}}{}
\IfFileExists{microtype.sty}{
  \usepackage[]{microtype}
  \UseMicrotypeSet[protrusion]{basicmath}
}{}

\makeatletter
\@ifundefined{KOMAClassName}{
  \IfFileExists{parskip.sty}{
    \usepackage{parskip}
  }{
    \setlength{\parindent}{0pt}
    \setlength{\parskip}{6pt plus 2pt minus 1pt}}
}{
  \KOMAoptions{parskip=half}}
\makeatother

\usepackage{color}
\usepackage{fancyvrb}
\newcommand{\VerbBar}{|}
\newcommand{\VERB}{\Verb[commandchars=\\\{\}]}
\DefineVerbatimEnvironment{Highlighting}{Verbatim}{commandchars=\\\{\}}
\IfFileExists{fvextra.sty}{
  \usepackage{fvextra}
  \RecustomVerbatimEnvironment{Highlighting}{Verbatim}{
    commandchars=\\\{\},
    fontsize=\small,
    breaklines=true,
    breakanywhere=true}
}{}

\newenvironment{Shaded}{}{}
\newcommand{\AlertTok}[1]{#1}
\newcommand{\AnnotationTok}[1]{#1}
\newcommand{\AttributeTok}[1]{#1}
\newcommand{\BaseNTok}[1]{#1}
\newcommand{\BuiltInTok}[1]{#1}
\newcommand{\CharTok}[1]{#1}
\newcommand{\CommentTok}[1]{#1}
\newcommand{\CommentVarTok}[1]{#1}
\newcommand{\ConstantTok}[1]{#1}
\newcommand{\ControlFlowTok}[1]{#1}
\newcommand{\DataTypeTok}[1]{#1}
\newcommand{\DecValTok}[1]{#1}
\newcommand{\DocumentationTok}[1]{#1}
\newcommand{\ErrorTok}[1]{#1}
\newcommand{\ExtensionTok}[1]{#1}
\newcommand{\FloatTok}[1]{#1}
\newcommand{\FunctionTok}[1]{#1}
\newcommand{\ImportTok}[1]{#1}
\newcommand{\InformationTok}[1]{#1}
\newcommand{\KeywordTok}[1]{#1}
\newcommand{\NormalTok}[1]{#1}
\newcommand{\OperatorTok}[1]{#1}
\newcommand{\OtherTok}[1]{#1}
\newcommand{\PreprocessorTok}[1]{#1}
\newcommand{\RegionMarkerTok}[1]{#1}
\newcommand{\SpecialCharTok}[1]{#1}
\newcommand{\SpecialStringTok}[1]{#1}
\newcommand{\StringTok}[1]{#1}
\newcommand{\VariableTok}[1]{#1}
\newcommand{\VerbatimStringTok}[1]{#1}
\newcommand{\WarningTok}[1]{#1}

\usepackage{longtable,booktabs,array}
\setlength{\arrayrulewidth}{0.4pt}
\setlength{\tabcolsep}{4pt}
\AtBeginEnvironment{longtable}{\footnotesize\sloppy}
\newcounter{none}
\usepackage{calc}
\usepackage{etoolbox}
\makeatletter
\patchcmd\longtable{\par}{\if@noskipsec\mbox{}\fi\par}{}{}
\makeatother

\usepackage{enumitem}
\setlistdepth{9}
\renewlist{itemize}{itemize}{9}
\renewlist{enumerate}{enumerate}{9}
\setlist[itemize]{label=\textbullet}
\setlist[enumerate]{label=\arabic*.}

\IfFileExists{footnotehyper.sty}{\usepackage{footnotehyper}}{\usepackage{footnote}}
\makesavenoteenv{longtable}
\setlength{\emergencystretch}{3em}
\providecommand{\tightlist}{
  \setlength{\itemsep}{0pt}\setlength{\parskip}{0pt}}
\usepackage{bookmark}
\IfFileExists{xurl.sty}{\usepackage{xurl}}{}
\urlstyle{same}
\hypersetup{
  hidelinks,
  pdfcreator={LaTeX via pandoc}}

\author{}
\date{}

\usepackage{tocloft}
\renewcommand{\contentsname}{\Large\bfseries Mục lục}
\setcounter{tocdepth}{3}
\renewcommand{\cftsecleader}{\cftdotfill{\cftdotsep}}
\renewcommand{\cftsubsecleader}{\cftdotfill{\cftdotsep}}
\renewcommand{\cftsubsubsecleader}{\cftdotfill{\cftdotsep}}
\renewcommand{\cftsecpagefont}{\bfseries}
\renewcommand{\cftsubsecpagefont}{\bfseries}
\renewcommand{\cftsubsubsecpagefont}{\bfseries}
\setlength{\cftbeforesecskip}{6pt}
\setlength{\cftbeforesubsecskip}{2pt}

\usepackage{titlesec}

\titleformat{\section}
  {\LARGE\bfseries}
  {\thesection}
  {1em}
  {}
  []

\titleformat{\subsection}
  {\Large\bfseries}
  {\thesection}
  {1em}
  {}
  []

\titleformat{\subsubsection}
  {\large\bfseries}
  {\thesubsection}
  {1em}
  {}
  []

\begin{document}

\begin{titlepage}
\thispagestyle{empty}
\begin{tikzpicture}[remember picture,overlay]
  \draw[line width=1pt]
    ([xshift=1.25cm,yshift=-1.25cm]current page.north west)
    rectangle
    ([xshift=-1.25cm,yshift=1.25cm]current page.south east);
\end{tikzpicture}

\centering
\vspace*{2.2cm}
\IfFileExists{logo_ptit.png}{
  \includegraphics[width=0.42\textwidth]{logo_ptit.png}
}{
  \fbox{\parbox[c][4cm][c]{0.35\textwidth}{\centering\large Logo PTIT}}
}

\vspace{1.7cm}
 {\Huge\bfseries Software Requirements Specification and Use Case Document\par}

\vspace{1.4cm}
{\Large\bfseries Nhóm 3\par}

\vspace{0.8cm}
{\large
\begin{tabular}{c}
Lê Anh Minh -- B24DCCE180 \\
Ứng Trọng Trình -- B24DCCE271 \\
Nguyễn Đình Tiến -- B24DCCE264 \\
Nguyễn Đinh Anh Quân -- B24DCCE222
\end{tabular}
\par}

\vfill
{\small
\begin{tabular}{c}
\textbf{Personal Finance Management System (PFMS)} \\
\vspace{0.15em}
\textit{ISO/IEC/IEEE 29148:2018 \textbullet{} IREB CPRE \textbullet{} Craig Larman Use Case Format} \\
\textit{Martin Fowler Accounting Patterns}
\end{tabular}
\par}
\end{titlepage}

\subsection*{MỤC LỤC TỔNG QUAN}\label{mux1ee5c-lux1ee5c-tux1ed5ng-quan}

\begingroup
\let\originalhyperref\hyperref
\renewcommand{\hyperref}[2][]{\originalhyperref[#1]{#2}\dotfill~\pageref{#1}}
\begin{itemize}
\tightlist
\item
  \hyperref[1-giux1edbi-thiux1ec7u-tux1ed5ng-quan-introduction]{1. Giới
  thiệu tổng quan (Introduction)}

  \begin{itemize}
  \tightlist
  \item
    \hyperref[11-mux1ee5c-ux111uxedch-tuxe0i-liux1ec7u]{1.1 Mục đích tài
    liệu}
  \item
    \hyperref[12-phux1ea1m-vi-sux1ea3n-phux1ea9m]{1.2 Phạm vi sản phẩm}
  \item
    \hyperref[13-ux111ux1ed1i-tux1b0ux1ee3ng-sux1eed-dux1ee5ng-vuxe0-hux1b0ux1edbng-dux1eabn-ux111ux1ecdc-tuxe0i-liux1ec7u]{1.3
    Đối tượng sử dụng và hướng dẫn đọc tài liệu}
  \item
    \hyperref[14-tuxe0i-liux1ec7u-tham-chiux1ebfu-chuux1ea9n]{1.4 Tài
    liệu tham chiếu chuẩn}
  \item
    \hyperref[15-bux1ea3ng-thuux1eadt-ngux1eef-vuxe0-ux111ux1ecbnh-nghux129a-glossary]{1.5
    Bảng thuật ngữ và định nghĩa (Glossary)}
  \end{itemize}
\item
  \hyperref[2-muxf4-tux1ea3-tux1ed5ng-quuxe1t-hux1ec7-thux1ed1ng-overall-description]{2.
  Mô tả tổng quát hệ thống (Overall Description)}

  \begin{itemize}
  \tightlist
  \item
    \hyperref[21-bux1ed1i-cux1ea3nh-hux1ec7-thux1ed1ng]{2.1 Bối cảnh hệ
    thống}
  \item
    \hyperref[22-mux1ee5c-tiuxeau-hux1ec7-thux1ed1ng-system-goals]{2.2
    Mục tiêu hệ thống (System Goals)}
  \item
    \hyperref[23-ux111ux1eb7c-ux111iux1ec3m-tuxe1c-nhuxe2n-ngux1b0ux1eddi-duxf9ng]{2.3
    Đặc điểm tác nhân người dùng}
  \item
    \hyperref[24-ruxe0ng-buux1ed9c-kux1ef9-thuux1eadt-giux1ea3-ux111ux1ecbnh-vuxe0-sux1ef1-phux1ee5-thuux1ed9c]{2.4
    Ràng buộc kỹ thuật, giả định và sự phụ thuộc}
  \item
    \hyperref[25-muxf4i-trux1b0ux1eddng-vux1eadn-huxe0nh]{2.5 Môi trường
    vận hành}
  \end{itemize}
\item
  \hyperref[3-quy-tux1eafc-nghiux1ec7p-vux1ee5-vuxe0-ruxe0ng-buux1ed9c-tuxe0i-chuxednh-domain-invariants--business-rules]{3.
  Quy tắc nghiệp vụ và ràng buộc tài chính (Domain Invariants \&
  Business Rules)}

  \begin{itemize}
  \tightlist
  \item
    \hyperref[31-ruxe0ng-buux1ed9c-chuxednh-xuxe1c-tuxe0i-chuxednh-br-fin-01]{3.1
    Ràng buộc chính xác tài chính (BR-FIN-01)}
  \item
    \hyperref[32-quy-tux1eafc-bux1ea3o-touxe0n-sux1ed1-dux1b0-vuxed-br-fin-02]{3.2
    Quy tắc bảo toàn số dư ví (BR-FIN-02)}
  \item
    \hyperref[33-quy-tux1eafc-ghi-sux1ed5-kuxe9p-khi-chuyux1ec3n-khoux1ea3n-nux1ed9i-bux1ed9-br-fin-03]{3.3
    Quy tắc ghi sổ kép khi chuyển khoản nội bộ (BR-FIN-03)}
  \item
    \hyperref[34-quy-tux1eafc-phuxe2n-cux1ea5p-danh-mux1ee5c-thu-chi-3-cux1ea5p-br-fin-04]{3.4
    Quy tắc phân cấp danh mục thu chi 3 cấp (BR-FIN-04)}
  \item
    \hyperref[35-quy-tux1eafc-tuxednh-luxe3i-tiux1ebft-kiux1ec7m-cuxf3-kux1ef3-hux1ea1n-br-fin-05]{3.5
    Quy tắc tính lãi tiết kiệm có kỳ hạn (BR-FIN-05)}
  \item
    \hyperref[36-quy-tux1eafc-ux111ux1ecbnh-giuxe1-chux1ee9ng-chux1ec9-quux1ef9-ux111ux1ea7u-tux1b0-br-fin-06]{3.6
    Quy tắc định giá chứng chỉ quỹ đầu tư (BR-FIN-06)}
  \item
    \hyperref[37-quy-tux1eafc-theo-duxf5i-cuxf4ng-nux1ee3-br-fin-07]{3.7
    Quy tắc theo dõi công nợ (BR-FIN-07)}
  \item
    \hyperref[38-quy-tux1eafc-phuxe2n-bux1ed5-quux1ef9-mux1ee5c-tiuxeau-br-fin-08]{3.8
    Quy tắc phân bổ quỹ mục tiêu (BR-FIN-08)}
  \end{itemize}
\item
  \hyperref[4-yuxeau-cux1ea7u-chux1ee9c-nux103ng-chi-tiux1ebft-functional-requirements]{4.
  Yêu cầu chức năng chi tiết (Functional Requirements)}

  \begin{itemize}
  \tightlist
  \item
    \hyperref[41-phuxe2n-hux1ec7-quux1ea3n-luxfd-vuxed-tiux1ec1n-fr-wal]{4.1
    Phân hệ Quản lý ví tiền (FR-WAL)}
  \item
    \hyperref[42-phuxe2n-hux1ec7-ghi-nhux1eadn-thu-chi-fr-txn]{4.2 Phân
    hệ Ghi nhận thu chi (FR-TXN)}
  \item
    \hyperref[43-phuxe2n-hux1ec7-chuyux1ec3n-tiux1ec1n-nux1ed9i-bux1ed9-fr-trf]{4.3
    Phân hệ Chuyển tiền nội bộ (FR-TRF)}
  \item
    \hyperref[44-phuxe2n-hux1ec7-danh-mux1ee5c-phuxe2n-loux1ea1i-fr-cat]{4.4
    Phân hệ Danh mục phân loại (FR-CAT)}
  \item
    \hyperref[45-phuxe2n-hux1ec7-quux1ea3n-luxfd-cuxf4ng-nux1ee3-fr-debt]{4.5
    Phân hệ Quản lý công nợ (FR-DEBT)}
  \item
    \hyperref[46-phuxe2n-hux1ec7-quux1ea3n-luxfd-sux1ed5-tiux1ebft-kiux1ec7m-fr-sav]{4.6
    Phân hệ Quản lý sổ tiết kiệm (FR-SAV)}
  \item
    \hyperref[47-phuxe2n-hux1ec7-chux1ee9ng-chux1ec9-quux1ef9-ux111ux1ea7u-tux1b0-fr-inv]{4.7
    Phân hệ Chứng chỉ quỹ đầu tư (FR-INV)}
  \item
    \hyperref[48-phuxe2n-hux1ec7-quux1ea3n-luxfd-quux1ef9-mux1ee5c-tiuxeau-fr-fnd]{4.8
    Phân hệ Quản lý quỹ mục tiêu (FR-FND)}
  \item
    \hyperref[49-phuxe2n-hux1ec7-thux1ed1ng-kuxea-vuxe0-buxe1o-cuxe1o-fr-rep]{4.9
    Phân hệ Thống kê và Báo cáo (FR-REP)}
  \item
    \hyperref[410-phuxe2n-hux1ec7-xuux1ea5t-nhux1eadp-vuxe0-sao-lux1b0u-dux1eef-liux1ec7u-fr-dat]{4.10
    Phân hệ Xuất nhập và Sao lưu dữ liệu (FR-DAT)}
  \end{itemize}
\item
  \hyperref[5-yuxeau-cux1ea7u-phi-chux1ee9c-nux103ng-non-functional-requirements]{5.
  Yêu cầu phi chức năng (Non-Functional Requirements)}

  \begin{itemize}
  \tightlist
  \item
    \hyperref[51-nfr-01-ux111ux1ed9-chuxednh-xuxe1c-vuxe0-touxe0n-vux1eb9n-dux1eef-liux1ec7u-tuxe0i-chuxednh-financial-precision--integrity]{5.1
    NFR-01: Độ chính xác và toàn vẹn dữ liệu tài chính (Financial
    Precision \& Integrity)}
  \item
    \hyperref[52-nfr-02-hiux1ec7u-nux103ng-vuxe0-ux111ux1ed9-trux1ec5-phux1ea3n-hux1ed3i-performance--latency]{5.2
    NFR-02: Hiệu năng và độ trễ phản hồi (Performance \& Latency)}
  \item
    \hyperref[53-nfr-03-ux111ux1ed9-tin-cux1eady-vuxe0-lux1b0u-trux1eef-an-touxe0n-reliability--persistence]{5.3
    NFR-03: Độ tin cậy và lưu trữ an toàn (Reliability \& Persistence)}
  \item
    \hyperref[54-nfr-04-khux1ea3-nux103ng-bux1ea3o-truxec-vuxe0-kiux1ebfn-truxfac-hux1b0ux1edbng-ux111ux1ed1i-tux1b0ux1ee3ng-maintainability--oop-design]{5.4
    NFR-04: Khả năng bảo trì và kiến trúc hướng đối tượng
    (Maintainability \& OOP Design)}
  \item
    \hyperref[55-nfr-05-tuxednh-khux1ea3-dux1ee5ng-vuxe0-cuxf4ng-thuxe1i-hux1ecdc-usability--user-ergonomics]{5.5
    NFR-05: Tính khả dụng và công thái học (Usability \& User
    Ergonomics)}
  \item
    \hyperref[56-nfr-06-tuxednh-tux1b0ux1a1ng-thuxedch-vuxe0-muxf4i-trux1b0ux1eddng-thux1ef1c-thi-portability--compatibility]{5.6
    NFR-06: Tính tương thích và môi trường thực thi (Portability \&
    Compatibility)}
  \end{itemize}
\item
  \hyperref[6-ux111ux1eb7c-tux1ea3-10-ca-sux1eed-dux1ee5ng-chi-tiux1ebft-fully-dressed-use-cases]{6.
  Đặc tả 10 Ca sử dụng chi tiết (Fully Dressed Use Cases)}

  \begin{itemize}
  \tightlist
  \item
    \hyperref[61-uc01-quux1ea3n-luxfd-vuxed-tiux1ec1n-wallet-management]{6.1
    UC01: Quản lý ví tiền (Wallet Management)}
  \item
    \hyperref[62-uc02-ghi-nhux1eadn-vuxe0-quux1ea3n-luxfd-giao-dux1ecbch-thu-chi-income--expense-management]{6.2
    UC02: Ghi nhận và quản lý giao dịch thu chi (Income \& Expense
    Management)}
  \item
    \hyperref[63-uc03-chuyux1ec3n-tiux1ec1n-nux1ed9i-bux1ed9-giux1eefa-cuxe1c-vuxed-internal-fund-transfer]{6.3
    UC03: Chuyển tiền nội bộ giữa các ví (Internal Fund Transfer)}
  \item
    \hyperref[64-uc04-thiux1ebft-lux1eadp-vuxe0-quux1ea3n-luxfd-danh-mux1ee5c-phuxe2n-loux1ea1i-thu-chi-3-cux1ea5p-category-hierarchy-management]{6.4
    UC04: Thiết lập và quản lý danh mục phân loại thu chi 3 cấp
    (Category Hierarchy Management)}
  \item
    \hyperref[65-uc05-quux1ea3n-luxfd-cuxf4ng-nux1ee3-debt-management]{6.5
    UC05: Quản lý công nợ (Debt Management)}
  \item
    \hyperref[66-uc06-quux1ea3n-luxfd-sux1ed5-tiux1ebft-kiux1ec7m-cuxf3-kux1ef3-hux1ea1n-vuxe0-tuxednh-touxe1n-luxe3i-suux1ea5t-term-savings-management]{6.6
    UC06: Quản lý sổ tiết kiệm có kỳ hạn và tính toán lãi suất (Term
    Savings Management)}
  \item
    \hyperref[67-uc07-quux1ea3n-luxfd-danh-mux1ee5c-chux1ee9ng-chux1ec9-quux1ef9-ux111ux1ea7u-tux1b0-vuxe0-biux1ebfn-ux111ux1ed9ng-nav-fund-certificate-investment-management]{6.7
    UC07: Quản lý danh mục chứng chỉ quỹ đầu tư và biến động NAV (Fund
    Certificate Investment Management)}
  \item
    \hyperref[68-uc08-quux1ea3n-luxfd-quux1ef9-mux1ee5c-tiuxeau-tuxe0i-chuxednh-vuxe0-phuxe2n-bux1ed5-duxf2ng-tiux1ec1n-financial-goal--fund-management]{6.8
    UC08: Quản lý quỹ mục tiêu tài chính và phân bổ dòng tiền (Financial
    Goal \& Fund Management)}
  \item
    \hyperref[69-uc09-tra-cux1ee9u-bux1ea3ng-thux1ed1ng-kuxea-duxf2ng-tiux1ec1n-trux1b0ux1ee3t-12-thuxe1ng-vuxe0-phuxe2n-tuxedch-tux1ef7-lux1ec7-tiux1ebft-kiux1ec7m-monthly-statistics--cashflow-analysis]{6.9
    UC09: Tra cứu bảng thống kê dòng tiền trượt 12 tháng và phân tích tỷ
    lệ tiết kiệm (Monthly Statistics \& Cashflow Analysis)}
  \item
    \hyperref[610-uc10-sao-lux1b0u-xuux1ea5t-vuxe0-nhux1eadp-dux1eef-liux1ec7u-tuxe0i-chuxednh-financial-data-backup-export--import]{6.10
    UC10: Sao lưu, xuất và nhập dữ liệu tài chính (Financial Data
    Backup, Export \& Import)}
  \end{itemize}
\item
  \hyperref[7-sux1a1-ux111ux1ed3-ca-sux1eed-dux1ee5ng-hux1ec7-thux1ed1ng-use-case-diagrams]{7.
  Sơ đồ Ca sử dụng hệ thống (Use Case Diagrams)}

  \begin{itemize}
  \tightlist
  \item
    \hyperref[71-sux1a1-ux111ux1ed3-ca-sux1eed-dux1ee5ng-tux1ed5ng-thux1ec3-touxe0n-hux1ec7-thux1ed1ng-global-system-use-case-diagram]{7.1
    Sơ đồ Ca sử dụng tổng thể toàn hệ thống (Global System Use Case
    Diagram)}
  \item
    \hyperref[72-sux1a1-ux111ux1ed3-ca-sux1eed-dux1ee5ng-phuxe2n-ruxe3-theo-3-phuxe2n-vuxf9ng-nghiux1ec7p-vux1ee5-subsystem-use-case-diagrams]{7.2
    Sơ đồ Ca sử dụng phân rã theo 3 phân vùng nghiệp vụ (Subsystem Use
    Case Diagrams)}

    \begin{itemize}
    \tightlist
    \item
      \hyperref[721-phuxe2n-vuxf9ng-quux1ea3n-luxfd-duxf2ng-tiux1ec1n-vuxe0-vuxed-tiux1ec1n-cashflow--liquidity-subsystem]{7.2.1
      Phân vùng Quản lý Dòng tiền và Ví tiền (Cashflow \& Liquidity
      Subsystem)}
    \item
      \hyperref[722-phuxe2n-vuxf9ng-tuxedch-lux169y-ux111ux1ea7u-tux1b0-vuxe0-nghux129a-vux1ee5-tuxe0i-chuxednh-wealth--liabilities-subsystem]{7.2.2
      Phân vùng Tích lũy, Đầu tư và Nghĩa vụ tài chính (Wealth \&
      Liabilities Subsystem)}
    \item
      \hyperref[723-phuxe2n-vuxf9ng-buxe1o-cuxe1o-phuxe2n-tuxedch-vuxe0-quux1ea3n-trux1ecb-dux1eef-liux1ec7u-analytics--administration-subsystem]{7.2.3
      Phân vùng Báo cáo, Phân tích và Quản trị dữ liệu (Analytics \&
      Administration Subsystem)}
    \end{itemize}
  \end{itemize}
\item
  \hyperref[8-ma-trux1eadn-kiux1ec3m-chux1ee9ng-vuxe0-truy-vux1ebft-yuxeau-cux1ea7u-requirements-traceability-matrix---rtm]{8.
  Ma trận kiểm chứng và truy vết yêu cầu (Requirements Traceability
  Matrix - RTM)}

  \begin{itemize}
  \tightlist
  \item
    \hyperref[81-ma-trux1eadn-truy-vux1ebft-yuxeau-cux1ea7u-4-chiux1ec1u-requirements-traceability-matrix]{8.1
    Ma trận truy vết yêu cầu 4 chiều (Requirements Traceability Matrix)}
  \item
    \hyperref[82-ma-trux1eadn-phux1b0ux1a1ng-phuxe1p-kiux1ec3m-chux1ee9ng-yuxeau-cux1ea7u-verification-methods-matrix]{8.2
    Ma trận phương pháp kiểm chứng yêu cầu (Verification Methods
    Matrix)}

    \begin{itemize}
    \tightlist
    \item
      \hyperref[821-ma-trux1eadn-phux1b0ux1a1ng-phuxe1p-kiux1ec3m-chux1ee9ng-yuxeau-cux1ea7u-chux1ee9c-nux103ng-39-frs]{8.2.1
      Ma trận phương pháp kiểm chứng yêu cầu chức năng (39 FRs)}
    \item
      \hyperref[822-ma-trux1eadn-phux1b0ux1a1ng-phuxe1p-kiux1ec3m-chux1ee9ng-yuxeau-cux1ea7u-phi-chux1ee9c-nux103ng-6-nfrs]{8.2.2
      Ma trận phương pháp kiểm chứng yêu cầu phi chức năng (6 NFRs)}
    \end{itemize}
  \end{itemize}
\end{itemize}

\endgroup

\clearpage
\subsection{1. GIỚI THIỆU TỔNG QUAN
(INTRODUCTION)}\label{1-giux1edbi-thiux1ec7u-tux1ed5ng-quan-introduction}

\subsubsection{1.1 Mục đích tài
liệu}\label{11-mux1ee5c-ux111uxedch-tuxe0i-liux1ec7u}

Tài liệu này là bản đặc tả yêu cầu phần mềm (SRS) xác lập phạm vi chức
năng, quy tắc tài chính và tiêu chí nghiệm thu cho hệ thống Quản lý Tài
chính Cá nhân (PFMS). Tài liệu đóng vai trò là căn cứ kỹ thuật cho các
giai đoạn tiếp theo:

\begin{itemize}
\tightlist
\item
  Giai đoạn Phân tích Hướng đối tượng (OOA): Trích xuất thực thể miền,
  xác định ranh giới và phân bổ trách nhiệm theo mô hình GRASP.
\item
  Giai đoạn Thiết kế Kiến trúc và Chi tiết (OOD): Xây dựng sơ đồ lớp
  (Class Diagram), sơ đồ tuần tự (Sequence Diagram) và áp dụng các mẫu
  thiết kế GoF.
\item
  Giai đoạn Cài đặt và Kiểm thử: Hiện thực hóa mã nguồn trên nền tảng
  Java 17 và tự động hóa kiểm thử đơn vị với JUnit 5.
\end{itemize}

Tài liệu cung cấp một hợp đồng kỹ thuật nhất quán giữa bài toán quản lý
tài chính thực tế và mô hình phần mềm hướng đối tượng được cài đặt.

\subsubsection{1.2 Phạm vi sản
phẩm}\label{12-phux1ea1m-vi-sux1ea3n-phux1ea9m}

Sản phẩm PFMS là một ứng dụng máy tính cá nhân (Desktop Application)
chạy trên môi trường Java Runtime Environment. Ứng dụng được thiết kế
phục vụ một cá nhân duy nhất (mô hình Single-user / Owner) nhằm theo
dõi, phân tích và kiểm soát toàn bộ sức khỏe tài chính cá nhân.

Hệ thống triển khai đầy đủ 10 phân hệ nghiệp vụ tài chính, không phân
chia giai đoạn cốt lõi hay trì hoãn:

\begin{enumerate}
\def\labelenumi{\arabic{enumi}.}
\tightlist
\item
  Quản lý ví tiền (Wallets): Khởi tạo, theo dõi số dư khả dụng, cơ cấu
  phân bổ tỷ trọng tài sản thanh khoản giữa các ví tiền mặt và tài khoản
  ngân hàng.
\item
  Quản lý thu chi (Transactions): Ghi nhận nhật ký thu nhập và chi tiêu
  phát sinh theo thời gian thực, liên kết trực tiếp với từng ví cụ thể.
\item
  Chuyển tiền nội bộ (Transfers): Điều chuyển tiền tệ giữa các ví nội bộ
  theo nguyên tắc ghi sổ kép đối ứng, bảo toàn tổng lượng tiền mặt trong
  hệ thống.
\item
  Danh mục phân loại 3 cấp (Categories): Chuẩn hóa cây phân loại gồm
  Loại (Type), Nhóm nhu cầu (Group) và Danh mục chi tiết (Category).
\item
  Quản lý công nợ (Debts): Giám sát các khoản phải thu (cho vay) và các
  khoản phải trả (đi vay), xử lý hoàn ứng và xóa nợ khó đòi.
\item
  Quản lý sổ tiết kiệm (Savings): Theo dõi các khoản tiền gửi có kỳ hạn
  và không kỳ hạn, tính toán lãi suất dự kiến, xử lý tất toán đúng hạn
  và trước hạn.
\item
  Quản lý chứng chỉ quỹ đầu tư (Fund Certificates): Định giá danh mục
  đầu tư theo giá thị trường (Mark-to-Market) dựa trên NAV, theo dõi lợi
  nhuận chưa thực hiện và tỷ suất sinh lời hàng năm.
\item
  Quản lý quỹ mục tiêu tài chính (Financial Goals / Funds): Tổng hợp tài
  sản từ tiền mặt, sổ tiết kiệm và chứng chỉ quỹ để đo lường tiến độ
  hoàn thành các mục tiêu tài chính độc lập.
\item
  Thống kê và báo cáo (Monthly Statistics \& Reporting): Tự động trượt
  chuỗi thời gian 12 tháng gần nhất, tổng hợp dòng tiền ròng, tỷ lệ tiết
  kiệm và phân tích cơ cấu chi tiêu.
\item
  Quản lý dữ liệu và sao lưu (Data Management \& Backup): Lưu trữ cục bộ
  an toàn, hỗ trợ xuất và nhập dữ liệu để ngăn ngừa mất mát thông tin.
\end{enumerate}

Ranh giới ngoài phạm vi của hệ thống (Out of Scope):

\begin{itemize}
\tightlist
\item
  Không kết nối trực tiếp API với các cổng thanh toán hoặc ngân hàng
  thương mại.
\item
  Không hỗ trợ cơ chế đa người dùng (Multi-tenant) hoặc phân quyền truy
  cập nhiều tầng (Role-based access control).
\item
  Không thực hiện các giao dịch chuyển khoản ngân hàng thực tế mà đóng
  vai trò là sổ cái ghi nhận và phân tích thông tin.
\end{itemize}

\subsubsection{1.3 Đối tượng sử dụng và hướng dẫn đọc tài
liệu}\label{13-ux111ux1ed1i-tux1b0ux1ee3ng-sux1eed-dux1ee5ng-vuxe0-hux1b0ux1edbng-dux1eabn-ux111ux1ecdc-tuxe0i-liux1ec7u}

Tài liệu được biên soạn nhằm phục vụ các nhóm đối tượng sau:

\begin{itemize}
\tightlist
\item
  \textbf{Người phụ trách phân tích và thiết kế (Software Analysts /
  Designers):} Dùng tài liệu để xác định ranh giới hệ thống, chuyển đổi
  các ca sử dụng thành mô hình phân tích OOA và thiết kế OOD.
\item
  \textbf{Kỹ sư lập trình (Developers):} Căn cứ vào các luồng sự kiện,
  quy tắc nghiệp vụ và ràng buộc dữ liệu để cài đặt mã nguồn Java.
\item
  \textbf{Kỹ sư kiểm thử (QA / Testers):} Căn cứ vào các luồng ngoại lệ
  và tiêu chí nghiệm thu để xây dựng kịch bản kiểm thử (Test Cases).
\item
  \textbf{Người đánh giá kỹ thuật (Technical Reviewers / Stakeholders):}
  Sử dụng tài liệu để đối soát phạm vi tính năng và nghiệm thu kỹ thuật.
\end{itemize}

Hướng dẫn đọc tài liệu:

\begin{itemize}
\tightlist
\item
  Để nắm tổng quan bài toán và phạm vi: Đọc Mục 1 và Mục 2.
\item
  Để hiểu sâu về các nguyên lý và ràng buộc tài chính: Đọc kỹ Mục 3.
\item
  Để triển khai mã nguồn và thiết kế ca kiểm thử: Đọc tuần tự Mục 4, Mục
  6 và tham chiếu chéo với Mục 8.
\end{itemize}

\subsubsection{1.4 Tài liệu tham chiếu
chuẩn}\label{14-tuxe0i-liux1ec7u-tham-chiux1ebfu-chuux1ea9n}

\begin{enumerate}
\def\labelenumi{\arabic{enumi}.}
\tightlist
\item
  ISO/IEC/IEEE 29148:2018: Systems and software engineering - Life cycle
  processes - Requirements engineering.
\item
  International Requirements Engineering Board (IREB): CPRE Foundation
  Level Handbook \& Requirements Modeling Handbook.
\item
  Craig Larman (2004): Applying UML and Patterns: An Introduction to
  Object-Oriented Analysis and Design and Iterative Development, 3rd
  Edition, Prentice Hall.
\item
  Martin Fowler (1997): Accounting Patterns (Mô hình tài khoản, bút
  toán, giao dịch ghi sổ kép).
\item
  Đề bài bài toán gốc và yêu cầu kỹ thuật:
  \texttt{docs/academic/assignment-requirement.pdf}.
\end{enumerate}

\subsubsection{1.5 Bảng thuật ngữ và định nghĩa
(Glossary)}\label{15-bux1ea3ng-thuux1eadt-ngux1eef-vuxe0-ux111ux1ecbnh-nghux129a-glossary}

Bảng thuật ngữ được xây dựng theo chuẩn IREB CPRE, định nghĩa thống nhất
các khái niệm xuất hiện trong toàn bộ tài liệu:

{\def\LTcaptype{none}
\setlength{\LTleft}{\fill}
\setlength{\LTright}{\fill}
\begin{longtable}[]{|>{\raggedright\arraybackslash}p{0.16\linewidth}|>{\raggedright\arraybackslash}p{0.18\linewidth}|>{\raggedright\arraybackslash}p{0.56\linewidth}|}
\hline
\everycr{\noalign{\ifnum0=`}\fi\hrule height 0.4pt\futurelet\reserved@a\LT@@tabularcr}
\textbf{Thuật ngữ} & \textbf{Tên tiếng Anh} & \textbf{Định nghĩa kỹ thuật trong hệ thống} \\
\hline
\endfirsthead
\hline
\textbf{Thuật ngữ} & \textbf{Tên tiếng Anh} & \textbf{Định nghĩa kỹ thuật trong hệ thống} \\
\hline
\endhead
\hline
\endfoot
\hline
\endlastfoot
Tác nhân chính & Primary Actor & Người dùng duy nhất tương tác trực tiếp
với phần mềm để đạt được mục tiêu cụ thể (vai trò Owner). \\
Kịch bản & Scenario & Một chuỗi các bước tương tác tuần tự giữa tác nhân
và hệ thống trong một ngữ cảnh cụ thể. \\
Điều kiện kích hoạt & Trigger & Sự kiện khởi đầu kích thích việc thực
thi một ca sử dụng. \\
Tiền điều kiện & Precondition & Trạng thái bắt buộc của hệ thống phải
được thỏa mãn trước khi một ca sử dụng có thể bắt đầu. \\
Hậu điều kiện & Postcondition & Trạng thái bảo đảm của hệ thống sau khi
ca sử dụng kết thúc (thành công hoặc thất bại). \\
Bất biến miền & Invariant & Quy tắc nghiệp vụ bắt buộc phải luôn luôn
đúng tại mọi trạng thái hợp lệ của hệ thống. \\
Ví tiền & Wallet & Một tài khoản thanh khoản thực tế (tiền mặt, thẻ ngân
hàng, ví điện tử) lưu giữ số dư tiền khả dụng. \\
Giao dịch thu chi & Transaction & Bản ghi một lần phát sinh tài chính
làm tăng hoặc giảm số dư của một ví xác định. \\
Chuyển tiền nội bộ & Internal Transfer & Hành vi dịch chuyển giá trị
tiền tệ giữa hai ví nội bộ mà không làm thay đổi tổng tài sản ròng. \\
Cấu trúc danh mục & Category Hierarchy & Mô hình phân loại dạng cây 3
cấp gồm Loại (Type), Nhóm nhu cầu (Group) và Danh mục chi tiết
(Category). \\
Công nợ & Debt & Các nghĩa vụ tài chính phát sinh giữa Owner và các đối
tượng bên ngoài, chia thành phải thu và phải trả. \\
Khoản phải thu & Receivable & Tài sản của Owner đang bị người khác chiếm
dụng tạm thời qua hành vi cho vay (\texttt{lend}). \\
Khoản phải trả & Payable & Nghĩa vụ nợ mà Owner có trách nhiệm phải hoàn
trả cho người khác qua hành vi đi vay (\texttt{borrow}). \\
Xóa nợ & Write-off & Hành động kế toán ghi nhận tổn thất tài sản khi
không thể thu hồi khoản cho vay hoặc ghi nhận thu nhập bất thường khi
được chủ nợ xóa nợ. \\
Sổ tiết kiệm & Saving Account & Khoản tiền gửi tích lũy có kỳ hạn hoặc
không kỳ hạn tại tổ chức tín dụng để sinh lời theo lãi suất cam kết. \\
Chứng chỉ quỹ & Fund Certificate & Sản phẩm đầu tư tài chính được quản
lý bởi công ty quản lý quỹ, đại diện cho quyền sở hữu một phần danh mục
tài sản ròng. \\
Giá trị tài sản ròng & Net Asset Value (NAV) & Giá trị thị trường của
một đơn vị chứng chỉ quỹ tại thời điểm định giá. \\
Quỹ mục tiêu & Financial Goal / Fund & Một nhóm tài sản độc lập được quy
hoạch từ tiền mặt, sổ tiết kiệm và chứng chỉ quỹ để phục vụ một mục đích
tài chính cụ thể. \\
Ghi sổ kép & Double-Entry & Nguyên lý kế toán yêu cầu mỗi giao dịch
chuyển dịch giá trị phải gồm hai vế đối ứng cân bằng: ghi giảm một bên
và ghi tăng một bên tương ứng. \\
Chủ sở hữu duy nhất & Single-user Owner & Mô hình vận hành của hệ thống,
trong đó toàn bộ dữ liệu thuộc quyền kiểm soát của một người dùng duy
nhất trên máy trạm cục bộ. \\
\end{longtable}
}

\clearpage
\subsection{2. MÔ TẢ TỔNG QUÁT HỆ THỐNG (OVERALL
DESCRIPTION)}\label{2-muxf4-tux1ea3-tux1ed5ng-quuxe1t-hux1ec7-thux1ed1ng-overall-description}

\subsubsection{2.1 Bối cảnh hệ
thống}\label{21-bux1ed1i-cux1ea3nh-hux1ec7-thux1ed1ng}

Hệ thống Quản lý Tài chính Cá nhân (PFMS) được nghiên cứu và phát triển
độc lập từ đầu nhằm cung cấp một giải pháp phần mềm quản trị tài chính
cá nhân toàn diện, có tính kỷ luật cao và bảo toàn dữ liệu nghiêm ngặt
cho người dùng cá nhân (vai trò Chủ sở hữu: Owner). Trong bối cảnh tài
chính hiện đại, một cá nhân thường sở hữu nhiều kênh phân bổ tài sản
(tiền mặt, tài khoản ngân hàng, sổ tiết kiệm có kỳ hạn, danh mục chứng
chỉ quỹ đầu tư) đi kèm các nghĩa vụ công nợ và các mục tiêu tích lũy tài
chính dài hạn. Việc quản trị phân tán hoặc sử dụng các công cụ ghi chép
rời rạc dễ dẫn tới nguy cơ mất cân đối dòng tiền, sai lệch số dư và
thiếu hụt cái nhìn tổng thể về bức tranh tài chính.

PFMS được xây dựng nhằm giải quyết bài toán trên thông qua các định
hướng kỹ thuật nền tảng:

\begin{itemize}
\tightlist
\item
  Chuẩn hóa mô hình tài chính kế toán hướng đối tượng: Ứng dụng các mẫu
  kế toán chuẩn mực (Accounting Patterns của Martin Fowler), thiết lập
  tính bất biến số dư ví và áp dụng nguyên lý ghi sổ kép cho giao dịch
  chuyển tiền nhằm bảo đảm tính chính xác tuyệt đối.
\item
  Phân bổ tài sản và tổng hợp tài sản ròng theo thời gian thực: Tự động
  tổng hợp tài sản ròng (Net Worth) qua 5 cấu phần tài sản và nghĩa vụ
  nợ, phân tích xu hướng dòng tiền trượt 12 tháng.
\item
  Kiến trúc phần mềm độc lập, an toàn và toàn vẹn: Xây dựng bằng ngôn
  ngữ Java 17, thiết kế tách biệt 4 phân tầng kiến trúc, lưu trữ dữ liệu
  an toàn trên máy cục bộ với cơ chế ghi tệp nguyên tử và sao lưu tự
  động, hoàn toàn độc lập và không phụ thuộc vào bất kỳ phần mềm văn
  phòng trung gian nào.
\end{itemize}

\begin{Shaded}
\begin{Highlighting}[]
\NormalTok{flowchart TD}
\NormalTok{    subgraph CoreTx ["1. Giao dịch \& Quản lý dòng tiền cơ sở"]}
\NormalTok{        W["Quản lý Ví tiền\textless{}br/\textgreater{}(6 ví mặc định)"]}
\NormalTok{        IE["Ghi nhận Thu chi\textless{}br/\textgreater{}(Income / Expense)"]}
\NormalTok{        TR["Chuyển tiền nội bộ\textless{}br/\textgreater{}(Ghi sổ kép nguyên tử)"]}
\NormalTok{    end}

\NormalTok{    INV\{"BỘ LÕI BẢO TOÀN DÒNG TIỀN VÀ SỐ DƯ\textless{}br/\textgreater{}(Balance Invariant)"\}}

\NormalTok{    subgraph Assets ["2. Cấu phần Tài sản \& Đầu tư"]}
\NormalTok{        DEBT["Quản lý Công nợ\textless{}br/\textgreater{}(Vay, Thu, Trả)"]}
\NormalTok{        SAV["Sổ Tiết kiệm\textless{}br/\textgreater{}(Tính lãi kỳ hạn)"]}
\NormalTok{        INV\_F["Chứng chỉ quỹ Đầu tư\textless{}br/\textgreater{}(NAV, Lợi nhuận)"]}
\NormalTok{    end}

\NormalTok{    GOAL["QUỸ MỤC TIÊU TÀI CHÍNH\textless{}br/\textgreater{}(Emergency, Health, Target, Retire)"]}
\NormalTok{    REP["BÁO CÁO THỐNG KÊ 12 THÁNG \& SAO LƯU DỮ LIỆU CỤC BỘ"]}

\NormalTok{    CoreTx {-}{-}\textgreater{} INV}
\NormalTok{    INV {-}{-}\textgreater{} Assets}
\NormalTok{    Assets {-}{-}\textgreater{} GOAL}
\NormalTok{    GOAL {-}{-}\textgreater{} REP}

\NormalTok{    style INV fill:\#f9f0ea,stroke:\#d97706,stroke{-}width:2px}
\NormalTok{    style GOAL fill:\#ecfdf5,stroke:\#059669,stroke{-}width:1px}
\NormalTok{    style REP fill:\#eff6ff,stroke:\#2563eb,stroke{-}width:1px}
\end{Highlighting}
\end{Shaded}

\subsubsection{2.2 Mục tiêu hệ thống (System
Goals)}\label{22-mux1ee5c-tiuxeau-hux1ec7-thux1ed1ng-system-goals}

Hệ thống được xây dựng nhằm đáp ứng 5 mục tiêu cốt lõi:

\begin{itemize}
\item
  \textbf{G1 - Quản lý dòng tiền và phân loại thu chi minh bạch:}\\
  Cung cấp cơ chế ghi nhận tức thời mọi biến động tiền tệ phát sinh
  trong đời sống hàng ngày; phân loại chính xác từng khoản thu chi vào
  cây danh mục 3 cấp để người dùng hiểu rõ nguồn thu và cấu trúc chi phí
  sinh hoạt.
\item
  \textbf{G2 - Nắm bắt chính xác giá trị tài sản thực tế và cơ cấu phân
  bổ:}\\
  Cung cấp cái nhìn toàn diện về tài sản ròng tại bất kỳ thời điểm nào
  thông qua việc tổng hợp số dư ví thanh khoản, giá trị các sổ tiết
  kiệm, giá trị thị trường của danh mục đầu tư và chênh lệch công nợ.
\item
  \textbf{G3 - Quản lý công nợ và các khoản tích lũy, đầu tư có kỳ
  hạn:}\\
  Kiểm soát chặt chẽ các khoản cho vay và đi vay theo từng đối tượng đối
  tác; tự động tính toán lãi suất tiết kiệm có kỳ hạn và theo dõi biến
  động lãi lỗ của chứng chỉ quỹ đầu tư theo giá trị tài sản ròng cập
  nhật.
\item
  \textbf{G4 - Hỗ trợ lập kế hoạch tài chính và giám sát tiến độ mục
  tiêu:}\\
  Cho phép người dùng phân bổ các nguồn lực tài chính (tiền mặt, tiết
  kiệm, quỹ đầu tư) vào các quỹ mục tiêu độc lập (quỹ khẩn cấp, quỹ sức
  khỏe, mục tiêu mua sắm, quỹ hưu trí) và theo dõi tỷ lệ hoàn thành mục
  tiêu theo thời gian thực.
\item
  \textbf{G5 - Đảm bảo tính toán tài chính chính xác tuyệt đối:}\\
  Loại bỏ hoàn toàn các sai số làm tròn số học bằng cách chuẩn hóa toàn
  bộ các phép tính tiền tệ sang kiểu số học có độ chính xác tùy ý, bảo
  đảm nguyên tắc bảo toàn giá trị và cân bằng kép kế toán trong mọi tình
  huống.
\end{itemize}

\subsubsection{2.3 Đặc điểm tác nhân người
dùng}\label{23-ux111ux1eb7c-ux111iux1ec3m-tuxe1c-nhuxe2n-ngux1b0ux1eddi-duxf9ng}

Hệ thống phục vụ duy nhất một nhóm tác nhân:

\begin{itemize}
\tightlist
\item
  Tác nhân: Chủ sở hữu hệ thống (Owner).
\item
  Quyền hạn: Toàn quyền truy cập và thao tác trên mọi phân hệ dữ liệu.
  Không tồn tại cơ chế phân quyền nhiều cấp, không yêu cầu xác thực
  phiên đăng nhập phức tạp.
\item
  Năng lực kỹ thuật: Người dùng có kỹ năng sử dụng máy tính phổ thông,
  có nhu cầu ghi chép tài chính cá nhân nghiêm túc, có hiểu biết cơ bản
  về thu chi, tiết kiệm và quản lý nợ.
\item
  Tần suất tương tác: Tương tác hàng ngày để ghi chép các giao dịch phát
  sinh, tương tác định kỳ hàng tuần hoặc hàng tháng để cập nhật giá trị
  thị trường chứng chỉ quỹ, kiểm tra tiến độ quỹ mục tiêu và phân tích
  báo cáo tài chính.
\end{itemize}

\subsubsection{2.4 Ràng buộc kỹ thuật, giả định và sự phụ
thuộc}\label{24-ruxe0ng-buux1ed9c-kux1ef9-thuux1eadt-giux1ea3-ux111ux1ecbnh-vuxe0-sux1ef1-phux1ee5-thuux1ed9c}

Các ràng buộc kỹ thuật:

\begin{itemize}
\tightlist
\item
  Nền tảng thực thi: Ứng dụng chạy trên Java Development Kit (JDK) 17
  LTS trở lên.
\item
  Xử lý số học tài chính: Bắt buộc sử dụng lớp
  \texttt{java.math.BigDecimal} cho tất cả các trường dữ liệu biểu diễn
  số tiền, số dư, tỷ lệ phần trăm và các biến số tính toán trung gian.
  Nghiêm cấm tuyệt đối việc sử dụng kiểu số thực dấu chấm động
  (\texttt{float}, \texttt{double}) để ngăn ngừa triệt để sai số làm
  tròn theo chuẩn IEEE 754.
\item
  Chuẩn kiểm thử: Sử dụng khung kiểm thử JUnit 5 để xây dựng bộ kiểm thử
  đơn vị tự động, bảo đảm độ bao phủ mã nguồn của các quy tắc tính toán
  tài chính đạt mức tối đa.
\item
  Cơ chế lưu trữ: Dữ liệu được lưu trữ an toàn dưới dạng các tệp tin cục
  bộ có định dạng rõ ràng (JSON hoặc cấu trúc tệp hướng đối tượng), bảo
  đảm tính khả chuyển khi sao lưu và phục hồi.
\end{itemize}

Giả định của hệ thống:

\begin{itemize}
\tightlist
\item
  Hệ thống giả định một ngày có 24 giờ và một năm chuẩn trong tính toán
  tài chính ngân hàng có 365 ngày (hoặc quy đổi tháng chuẩn theo quy
  định tổ chức tín dụng).
\item
  Hệ thống giả định người dùng chỉ thao tác trên một phiên làm việc tại
  một thời điểm trên một máy tính cá nhân.
\item
  Các giá trị tỷ giá tiền tệ và giá trị tài sản ròng (NAV) được người
  dùng cập nhật thủ công dựa trên công bố thực tế của các tổ chức tài
  chính.
\end{itemize}

Sự phụ thuộc:

\begin{itemize}
\tightlist
\item
  Hệ thống phụ thuộc vào đồng hồ hệ thống của hệ điều hành (System
  Clock) để xác định mốc thời gian giao dịch, số ngày tính lãi tiết kiệm
  và trượt mốc thống kê 12 tháng.
\item
  Hệ thống phụ thuộc vào hệ thống tệp tin cục bộ của hệ điều hành để
  thực hiện việc đọc, ghi và bảo toàn dữ liệu.
\end{itemize}

\subsubsection{2.5 Môi trường vận
hành}\label{25-muxf4i-trux1b0ux1eddng-vux1eadn-huxe0nh}

\begin{itemize}
\tightlist
\item
  Hệ điều hành mục tiêu: Tương thích hoàn toàn với Microsoft Windows
  10/11 (64-bit), macOS 12 trở lên, hoặc các bản phân phối Linux phổ
  biến (Ubuntu 20.04+, Fedora 36+).
\item
  Môi trường phần mềm yêu cầu: Máy trạm đã cài đặt sẵn Java Runtime
  Environment (JRE) hoặc Java Development Kit (JDK) phiên bản 17 trở
  lên.
\item
  Yêu cầu phần cứng tối thiểu:

  \begin{itemize}
  \tightlist
  \item
    Bộ vi xử lý: Dual-core 2.0 GHz hoặc cao hơn.
  \item
    Bộ nhớ RAM: Tối thiểu 2 GB RAM khả dụng.
  \item
    Dung lượng ổ đĩa: Tối thiểu 200 MB dung lượng trống để chứa ứng dụng
    và tệp dữ liệu giao dịch.
  \item
    Màn hình: Độ phân giải tối thiểu 1280 x 720 pixels để bảo đảm hiển
    thị đầy đủ các bảng dữ liệu và biểu đồ thống kê.
  \end{itemize}
\end{itemize}

\clearpage
\subsection{3. QUY TẮC NGHIỆP VỤ VÀ RÀNG BUỘC TÀI CHÍNH (DOMAIN
INVARIANTS \& BUSINESS
RULES)}\label{3-quy-tux1eafc-nghiux1ec7p-vux1ee5-vuxe0-ruxe0ng-buux1ed9c-tuxe0i-chuxednh-domain-invariants--business-rules}

\subsubsection{3.1 Ràng buộc chính xác tài chính
(BR-FIN-01)}\label{31-ruxe0ng-buux1ed9c-chuxednh-xuxe1c-tuxe0i-chuxednh-br-fin-01}

Ràng buộc chính xác tài chính là quy tắc bất biến nền tảng chi phối toàn
bộ các phép tính số học trong hệ thống PFMS.

\begin{enumerate}
\def\labelenumi{\arabic{enumi}.}
\item
  \textbf{Cấm sử dụng số thực dấu chấm động nhị phân:}\\
  Nghiêm cấm khai báo và sử dụng các kiểu dữ liệu nguyên thủy
  \texttt{float} và \texttt{double}, cũng như các lớp bao bọc
  \texttt{java.lang.Float} và \texttt{java.lang.Double} cho bất kỳ
  trường dữ liệu hoặc biến trung gian nào liên quan đến tiền tệ, số
  lượng chứng chỉ quỹ, tỷ lệ phần trăm hoặc lãi suất. Bản chất biểu diễn
  số thực dấu chấm động theo chuẩn IEEE 754 dẫn đến sai số nhị phân
  không thể kiểm soát (ví dụ phép toán \texttt{0.1\ +\ 0.2} trong kiểu
  \texttt{double} cho kết quả \texttt{0.30000000000000004}), gây sai
  lệch cân bằng sổ cái tài chính.
\item
  \textbf{Chuẩn hóa đối tượng số học \texttt{BigDecimal}:}\\
  Mọi phép tính tiền tệ bắt buộc phải sử dụng
  \texttt{java.math.BigDecimal}. Khi khởi tạo \texttt{BigDecimal} từ các
  giá trị số, bắt buộc sử dụng hàm khởi tạo nhận chuỗi ký tự
  (\texttt{new\ BigDecimal(String)}) hoặc phương thức tĩnh
  \texttt{BigDecimal.valueOf(double/long)}. Tuyệt đối không sử dụng hàm
  khởi tạo \texttt{new\ BigDecimal(double)} vì nó giữ nguyên sai số biểu
  diễn của số thực nguyên thủy.
\item
  \textbf{Quy định về độ chính xác (Scale) và chế độ làm tròn
  (RoundingMode):}

  \begin{itemize}
  \tightlist
  \item
    Đơn vị tiền tệ chính (VNĐ): Dữ liệu lưu trữ và hiển thị cuối cùng có
    scale bằng 0 (số nguyên đồng) hoặc scale bằng 2 khi lưu trữ nội bộ
    chuẩn quốc tế.
  \item
    Các phép tính tỷ lệ, lãi suất và NAV: Các biến số trung gian (lãi
    suất năm, tỷ lệ hoàn thành mục tiêu, đơn giá NAV) phải duy trì scale
    tối thiểu từ 4 đến 6 chữ số thập phân để bảo đảm độ mịn của phép
    tính.
  \item
    Chế độ làm tròn: Khi bắt buộc phải làm tròn số trong các phép chia
    không hết (\texttt{BigDecimal.divide}), hệ thống phải chỉ định tường
    minh chế độ làm tròn \texttt{RoundingMode.HALF\_EVEN}
    (Banker\textquotesingle s Rounding, ưu tiên làm tròn về số chẵn gần
    nhất nhằm triệt tiêu độ lệch thống kê tích lũy) hoặc
    \texttt{RoundingMode.HALF\_UP} (làm tròn số học truyền thống). Tuyệt
    đối cấm thực hiện phép chia \texttt{BigDecimal} mà không chỉ định
    scale và \texttt{RoundingMode} vì sẽ gây ngoại lệ
    \texttt{ArithmeticException} khi gặp số thập phân vô hạn tuần hoàn.
  \end{itemize}
\item
  \textbf{Tiêu chuẩn sai số tuyệt đối bằng 0:}\\
  Tại mọi thời điểm, tổng các thành phần phân rã phải khớp chính xác đến
  từng đơn vị nhỏ nhất so với giá trị tổng thể. Không chấp nhận bất kỳ
  độ lệch sai số nào trong các phép toán cộng trừ số dư ví và chuyển
  khoản.
\end{enumerate}

\subsubsection{3.2 Quy tắc bảo toàn số dư ví
(BR-FIN-02)}\label{32-quy-tux1eafc-bux1ea3o-touxe0n-sux1ed1-dux1b0-vuxed-br-fin-02}

Quy tắc bảo toàn số dư ví (Balance Invariant) xác lập rằng số dư hiển
thị của bất kỳ ví tiền nào cũng phải phản ánh chính xác kết quả tổng hợp
của toàn bộ lịch sử các dòng tiền tác động trực tiếp đến ví đó.

\begin{enumerate}
\def\labelenumi{\arabic{enumi}.}
\item
  \textbf{Công thức bất biến số dư ví:}\\
  Xét một ví tiền \(w\) bất kỳ tại thời điểm quan sát \(T\), số dư khả
  dụng \(Balance_w(T)\) được xác định duy nhất thông qua công thức toán
  học sau:
  \[
  \begin{aligned}
  Balance_w(T) &= Balance_w(0)
  + \sum_{t \le T} Income_w(t)
  - \sum_{t \le T} Expense_w(t)
  + \sum_{t \le T} TransferIn_w(t) \\
  &\quad - \sum_{t \le T} TransferOut_w(t)
  + \sum_{t \le T} Borrow_w(t)
  - \sum_{t \le T} Repay_w(t) \\
  &\quad - \sum_{t \le T} Lend_w(t)
  + \sum_{t \le T} Collect_w(t)
  \end{aligned}
  \]
  Trong đó:

  \begin{itemize}
  \tightlist
  \item
    \(Balance_w(0)\): Số dư khởi tạo ban đầu của ví \(w\).
  \item
    \(Income_w(t)\): Các khoản thu nhập phát sinh trực tiếp vào ví
    \(w\).
  \item
    \(Expense_w(t)\): Các khoản chi tiêu chi trả trực tiếp từ ví \(w\).
  \item
    \(TransferIn_w(t)\): Các khoản tiền chuyển từ ví khác đến ví \(w\).
  \item
    \(TransferOut_w(t)\): Các khoản tiền chuyển từ ví \(w\) sang ví
    khác.
  \item
    \(Borrow_w(t)\): Các khoản tiền đi vay nhận về nhập vào ví \(w\)
    (dòng tiền dương).
  \item
    \(Repay_w(t)\): Các khoản tiền trả nợ đã vay trích xuất từ ví \(w\)
    (dòng tiền âm).
  \item
    \(Lend_w(t)\): Các khoản tiền cho người khác vay trích xuất từ ví
    \(w\) (dòng tiền âm).
  \item
    \(Collect_w(t)\): Các khoản tiền thu hồi nợ cho vay nhận về nhập vào
    ví \(w\) (dòng tiền dương).
  \end{itemize}
\item
  \textbf{Ranh giới dòng tiền của các nghiệp vụ phi tiền mặt:}\\
  Hệ thống phân biệt rạch ròi giữa biến động công nợ kế toán và dòng
  tiền vật lý của ví:

  \begin{itemize}
  \tightlist
  \item
    Nghiệp vụ xóa nợ cho vay (\texttt{write-off\ lend}): Là thao tác ghi
    giảm khoản phải thu do xác định đối tác mất khả năng trả nợ. Nghiệp
    vụ này không làm phát sinh dòng tiền vào ví \(w\), do đó không tham
    gia vào công thức tính số dư ví.
  \item
    Nghiệp vụ được xóa nợ đi vay (\texttt{write-off\ borrow}): Là thao
    tác ghi giảm khoản phải trả do được chủ nợ tha thứ nghĩa vụ. Nghiệp
    vụ này không làm phát sinh dòng tiền rời khỏi ví \(w\), do đó không
    làm thay đổi số dư ví.
  \end{itemize}
\item
  \textbf{Ràng buộc số dư không âm:}\\
  Mặc định đối với các ví tiền mặt và ví tiền gửi thanh khoản thông
  thường, số dư khả dụng bắt buộc phải thỏa mãn: \[Balance_w(T) \ge 0\]
  Mọi giao dịch chi tiêu, chuyển tiền đi, cho vay hoặc trả nợ khiến số
  dư của ví rơi vào trạng thái âm (\(Balance_w < 0\)) đều bị hệ thống
  chặn lại và báo lỗi thiếu hụt số dư, ngoại trừ trường hợp người dùng
  cấu hình ví thuộc loại thẻ tín dụng có hạn mức thấu chi được phê duyệt
  trước.
\end{enumerate}

\subsubsection{3.3 Quy tắc ghi sổ kép khi chuyển khoản nội bộ
(BR-FIN-03)}\label{33-quy-tux1eafc-ghi-sux1ed5-kuxe9p-khi-chuyux1ec3n-khoux1ea3n-nux1ed9i-bux1ed9-br-fin-03}

Quy tắc ghi sổ kép trong chuyển khoản nội bộ tuân thủ mô hình kế toán
của Martin Fowler, bảo đảm sự nhất quán và toàn vẹn của tổng tài sản
thanh khoản.

\begin{enumerate}
\def\labelenumi{\arabic{enumi}.}
\item
  \textbf{Cặp bút toán đối ứng cân bằng:}\\
  Mỗi hành vi chuyển tiền từ ví nguồn \(w_{src}\) sang ví đích
  \(w_{dst}\) với số tiền chuyển \(A > 0\) phải được thực hiện thông qua
  hai bút toán đối ứng:

  \begin{itemize}
  \tightlist
  \item
    Bút toán ghi giảm (Debit): Trừ chính xác số tiền \(A\) khỏi số dư
    khả dụng của ví nguồn \(w_{src}\).
  \item
    Bút toán ghi tăng (Credit): Cộng chính xác số tiền \(A\) vào số dư
    khả dụng của ví đích \(w_{dst}\).
  \end{itemize}
\item
  \textbf{Tính nguyên tử của giao dịch (Atomic Transaction):}\\
  Cả hai bút toán trên phải được liên kết chặt chẽ trong một khối giao
  dịch nguyên tử. Hệ thống bảo đảm rằng một bên không thể được cập nhật
  nếu bên kia gặp lỗi. Nếu việc trừ tiền ở ví nguồn thành công nhưng
  việc cộng tiền ở ví đích bị gián đoạn, toàn bộ giao dịch phải được
  phục hồi trạng thái ban đầu (rollback), ngăn chặn hoàn toàn nguy cơ
  hao hụt hoặc sinh thêm tiền tệ ảo.
\item
  \textbf{Bảo toàn tổng tài sản ví thanh khoản:}\\
  Giao dịch chuyển tiền nội bộ thuần túy chỉ là hành vi tái cơ cấu nơi
  lưu giữ tiền tệ, không tạo ra thu nhập cũng không tạo ra chi phí của
  chủ sở hữu. Do đó, tổng giá trị tiền mặt trên toàn bộ các ví trước và
  sau giao dịch chuyển khoản phải được bảo toàn tuyệt đối:
  \[
  \Delta\sum_{w\in Wallets} Balance_w = 0
  \]
\item
  \textbf{Quy tắc xử lý phí chuyển tiền phát sinh:}\\
  Nếu giao dịch chuyển tiền phát sinh phí dịch vụ do đơn vị trung gian
  thu:

  \begin{itemize}
  \tightlist
  \item
    Khoản phí chuyển tiền \(Fee \ge 0\) không được khấu trừ gộp vào số
    tiền chuyển \(A\) của bút toán ghi sổ kép.
  \item
    Khoản phí phải được tách thành một giao dịch chi tiêu (Expense) độc
    lập phát sinh từ ví nguồn \(w_{src}\), gắn với danh mục chi phí
    thích hợp (ví dụ: "Phí dịch vụ ngân hàng"). Khi đó, số tiền thực trừ
    khỏi ví nguồn là \(A + Fee\), số tiền thực nhận tại ví đích là
    \(A\), và tổng tài sản toàn hệ thống giảm đúng bằng khoản phí
    \(Fee\).
  \end{itemize}
\end{enumerate}

\subsubsection{3.4 Quy tắc phân cấp danh mục thu chi 3 cấp
(BR-FIN-04)}\label{34-quy-tux1eafc-phuxe2n-cux1ea5p-danh-mux1ee5c-thu-chi-3-cux1ea5p-br-fin-04}

Hệ thống thiết lập mô hình dữ liệu danh mục phân loại thu chi theo cấu
trúc phân cấp cây 3 cấp nghiêm ngặt nhằm phục vụ việc phân tích tài
chính đa chiều.

\begin{enumerate}
\def\labelenumi{\arabic{enumi}.}
\item
  \textbf{Cấu trúc phân cấp 3 cấp:}

  \begin{itemize}
  \tightlist
  \item
    Cấp 1: Loại giao dịch (Type). Chỉ nhận một trong hai giá trị định
    danh bất biến: \texttt{Income} (Thu nhập) hoặc \texttt{Expense} (Chi
    tiêu).
  \item
    Cấp 2: Nhóm nhu cầu tài chính (Group). Là tập hợp các nhóm mục đích
    lớn, trực thuộc duy nhất một \texttt{Type}.

    \begin{itemize}
    \tightlist
    \item
      Với \texttt{Type\ =\ Income}: Các nhóm chuẩn gồm
      \texttt{employment} (Thu nhập từ tiền lương, việc làm),
      \texttt{assets} (Thu nhập từ tài sản, đầu tư, lãi suất),
      \texttt{other} (Các khoản thu nhập khác).
    \item
      Với \texttt{Type\ =\ Expense}: Các nhóm chuẩn gồm
      \texttt{essentials} (Chi tiêu thiết yếu cho đời sống),
      \texttt{personal\ \&\ leisure} (Chi tiêu cá nhân và hưởng thụ),
      \texttt{obligations\ \&\ transfers} (Nghĩa vụ phụng dưỡng, thuế
      phí), \texttt{savings\ \&\ investing} (Tích lũy và đầu tư).
    \end{itemize}
  \item
    Cấp 3: Danh mục chi tiết (Category). Là đơn vị phân loại nhỏ nhất,
    phản ánh cụ thể nội dung phát sinh giao dịch (ví dụ: tiền ăn uống,
    xăng xe, tiền điện nước, học phí), trực thuộc duy nhất một
    \texttt{Group}.
  \end{itemize}
\item
  \textbf{Tính toàn vẹn quan hệ cha con:}

  \begin{itemize}
  \tightlist
  \item
    Mỗi Category bắt buộc phải có đúng một Group cha làm định danh tham
    chiếu.
  \item
    Mỗi Group bắt buộc phải có đúng một Type làm phân loại gốc.
  \item
    Hệ thống nghiêm cấm sự tồn tại của các danh mục mồ côi (Category
    không có Group, hoặc Group không có Type).
  \end{itemize}
\item
  \textbf{Ràng buộc khi sửa đổi hoặc xóa danh mục:}

  \begin{itemize}
  \tightlist
  \item
    Không được phép xóa một Group nếu vẫn còn Category trực thuộc Group
    đó.
  \item
    Không được phép xóa một Category hoặc một Group nếu đã có ít nhất
    một giao dịch trong lịch sử đang tham chiếu tới danh mục đó. Nếu
    muốn ngừng sử dụng, danh mục phải được chuyển sang trạng thái ngưng
    hoạt động (Archived/Inactive).
  \item
    Không được phép thay đổi thuộc tính \texttt{Type} của một Group hoặc
    Category khi đã phát sinh dữ liệu giao dịch liên kết, nhằm tránh làm
    đảo lộn toàn bộ báo cáo thu chi quá khứ.
  \end{itemize}
\end{enumerate}

\subsubsection{3.5 Quy tắc tính lãi tiết kiệm có kỳ hạn
(BR-FIN-05)}\label{35-quy-tux1eafc-tuxednh-luxe3i-tiux1ebft-kiux1ec7m-cuxf3-kux1ef3-hux1ea1n-br-fin-05}

Hệ thống chuẩn hóa quy tắc tính toán lãi suất cho các sổ tiết kiệm theo
đúng nghiệp vụ thực tế của các ngân hàng thương mại.

\begin{enumerate}
\def\labelenumi{\arabic{enumi}.}
\item
  \textbf{Quy đổi kỳ hạn và mốc thời gian:}\\
  Kỳ hạn gửi tiền (\texttt{term}) được quy đổi thành số ngày thực tế dựa
  trên bảng tham chiếu chuẩn:

  \begin{itemize}
  \tightlist
  \item
    Không kỳ hạn (\texttt{unlimited}): Kỳ hạn quy ước là 0 ngày.
  \item
    Có kỳ hạn: 1 tháng (30 ngày), 3 tháng (90 ngày), 6 tháng (180 ngày),
    12 tháng (365 ngày), 24 tháng (730 ngày), 36 tháng (1095 ngày), 60
    tháng (1825 ngày).
  \item
    Số ngày gửi thực tế (\(Days\)) được tính từ ngày bắt đầu gửi
    (\(Start\)) đến ngày quan sát hoặc ngày tất toán (\(Current\)):
    \[Days = \max(0, Current - Start)\]
  \end{itemize}
\item
  \textbf{Công thức tính lãi cho khoản gửi không kỳ hạn:}\\
  Với sổ tiết kiệm không kỳ hạn (\texttt{term\ =\ unlimited}), lãi suất
  được tính theo lãi đơn tích lũy hàng ngày:
  \[Interest = \frac{Principal \times Rate \times Days}{365}\] Trong đó
  \(Principal\) là số vốn gốc, \(Rate\) là lãi suất năm (\%/năm).
\item
  \textbf{Công thức tính lãi và giá trị hiện tại của sổ có kỳ hạn:}\\
  Với sổ tiết kiệm có kỳ hạn xác định (\(TermDays > 0\)):

  \begin{itemize}
  \tightlist
  \item
    Tiến độ thời gian của kỳ hạn (\(Process\)) được xác định:
    \[Process = \min\left(1.0, \frac{Days}{TermDays}\right)\]
  \item
    Lãi dự kiến toàn kỳ (\(ExpectedInterest\)):
    \[ExpectedInterest = \frac{Principal \times Rate \times TermDays}{365}\]
  \item
    Giá trị sổ tiết kiệm ghi nhận tại thời điểm quan sát (\(Balance\)):
    \[Balance = Principal + (ExpectedInterest \times Process)\]
  \end{itemize}
\item
  \textbf{Quy tắc tất toán sổ tiết kiệm:}

  \begin{itemize}
  \tightlist
  \item
    Tất toán đúng hạn (\(Days \ge TermDays\)): Chủ tài khoản được nhận
    toàn bộ Vốn gốc cùng 100\% Lãi dự kiến theo lãi suất cam kết ban
    đầu: \[Payout = Principal + ExpectedInterest\]
  \item
    Tất toán trước hạn (\(Days < TermDays\)): Theo quy chế ngân hàng,
    người gửi không được hưởng mức lãi suất có kỳ hạn ban đầu. Toàn bộ
    tiền lãi sẽ được tính toán lại hồi tố theo mức lãi suất tiền gửi
    không kỳ hạn (\(DemandRate\), thường dao động từ 0.1\% đến
    0.5\%/năm) dựa trên số ngày thực gửi:
    \[Payout = Principal + \frac{Principal \times DemandRate \times Days}{365}\]
    Toàn bộ phần chênh lệch lãi suất kỳ hạn dự kiến trước đó bị hủy bỏ
    hoàn toàn.
  \end{itemize}
\end{enumerate}

\subsubsection{3.6 Quy tắc định giá chứng chỉ quỹ đầu tư
(BR-FIN-06)}\label{36-quy-tux1eafc-ux111ux1ecbnh-giuxe1-chux1ee9ng-chux1ec9-quux1ef9-ux111ux1ea7u-tux1b0-br-fin-06}

Quy tắc định giá chứng chỉ quỹ đầu tư thiết lập cơ chế phản ánh giá trị
tài sản đầu tư theo biến động thị trường thực tế.

\begin{enumerate}
\def\labelenumi{\arabic{enumi}.}
\item
  \textbf{Phân loại danh mục chứng chỉ quỹ:}\\
  Hệ thống hỗ trợ 3 nhóm tài sản đầu tư:

  \begin{itemize}
  \tightlist
  \item
    \texttt{stock}: Quỹ đầu tư cổ phiếu (mức độ biến động và kỳ vọng
    sinh lời cao).
  \item
    \texttt{bond}: Quỹ đầu tư trái phiếu (thu nhập cố định, độ an toàn
    cao).
  \item
    \texttt{balance}: Quỹ đầu tư cân bằng (kết hợp phân bổ giữa cổ phiếu
    và trái phiếu).
  \end{itemize}
\item
  \textbf{Nguyên tắc định giá theo giá thị trường (Mark-to-Market):}\\
  Giá trị hiện tại của một khoản đầu tư chứng chỉ quỹ (\(Balance\)) tại
  ngày định giá được xác định dựa trên số lượng đơn vị quỹ nắm giữ
  (\(Units\)) và Giá trị tài sản ròng trên một đơn vị quỹ (\(NAV\)) do
  công ty quản lý quỹ công bố: \[Balance = Units \times NAV\] Trường hợp
  ghi nhận đơn giản theo giá trị cập nhật từng lô mua: \(Balance\) là
  giá trị thị trường được ghi nhận tại ngày đánh giá gần nhất.
\item
  \textbf{Xác định lợi nhuận chưa thực hiện (Unrealized Profit/Loss):}\\
  Chênh lệch giữa giá trị thị trường hiện tại và số vốn gốc đầu tư ban
  đầu (\(Principal\)) phản ánh khoản lãi hoặc lỗ chưa thực hiện của danh
  mục: \[UnrealizedPL = Balance - Principal\]

  \begin{itemize}
  \tightlist
  \item
    Khi \(UnrealizedPL > 0\): Khoản đầu tư đang có lãi trên danh nghĩa.
  \item
    Khi \(UnrealizedPL < 0\): Khoản đầu tư đang chịu lỗ trên danh nghĩa.
  \end{itemize}
\item
  \textbf{Suất sinh lời quy đổi năm (Annualized Rate of Return):}\\
  Để đánh giá hiệu quả đầu tư một cách khách quan không phụ thuộc vào
  thời gian nắm giữ, hệ thống áp dụng công thức lãi kép hàng năm (tương
  đương hàm tài chính \texttt{RATE} hoặc CAGR):
  \[AnnualizedRate = \left(\frac{Balance}{Principal}\right)^{\frac{365}{HoldDays}} - 1\]
  Quy tắc áp dụng:

  \begin{itemize}
  \tightlist
  \item
    Công thức chỉ có hiệu lực khi \(HoldDays \ge 30\) ngày và
    \(Principal > 0\).
  \item
    Nếu \(HoldDays < 30\) ngày, hệ thống chỉ hiển thị tỷ suất sinh lời
    tuyệt đối chưa quy đổi năm:
    \(\frac{Balance - Principal}{Principal} \times 100\%\), tránh hiện
    tượng méo mó số liệu thống kê do số mũ quá lớn.
  \end{itemize}
\end{enumerate}

\subsubsection{3.7 Quy tắc theo dõi công nợ
(BR-FIN-07)}\label{37-quy-tux1eafc-theo-duxf5i-cuxf4ng-nux1ee3-br-fin-07}

Hệ thống quản lý công nợ cá nhân chặt chẽ nhằm phân định rõ quyền sở hữu
tài sản và nghĩa vụ nợ đối ứng với từng đối tượng cụ thể.

\begin{enumerate}
\def\labelenumi{\arabic{enumi}.}
\item
  \textbf{Phân loại 2 nhánh công nợ:}

  \begin{itemize}
  \tightlist
  \item
    Khoản phải thu (\texttt{Receivable}): Đại diện cho quyền đòi nợ của
    Owner đối với người khác khi phát sinh hành vi cho vay tiền. Khoản
    phải thu được ghi nhận là một bộ phận của tổng tài sản cá nhân.
  \item
    Khoản phải trả (\texttt{Payable}): Đại diện cho nghĩa vụ tài chính
    Owner phải thực hiện đối với chủ nợ khi phát sinh hành vi đi vay
    tiền. Khoản phải trả được ghi nhận là một khoản nợ phải trừ khi tính
    tài sản ròng.
  \end{itemize}
\item
  \textbf{Định nghĩa 6 loại giao dịch nghiệp vụ công nợ:}

  \begin{itemize}
  \tightlist
  \item
    \texttt{lend} (Cho vay): Tiền xuất khỏi ví của Owner để chuyển cho
    đối tác mượn. Hậu quả: Giảm số dư ví xuất tiền, tăng số dư phải thu
    từ đối tác.
  \item
    \texttt{borrow} (Đi vay): Tiền nhập vào ví của Owner từ chủ nợ. Hậu
    quả: Tăng số dư ví nhận tiền, tăng số dư phải trả cho chủ nợ.
  \item
    \texttt{collect} (Thu hồi nợ): Đối tác hoàn trả tiền đã mượn trước
    đó. Tiền nhập vào ví của Owner. Hậu quả: Tăng số dư ví nhận tiền,
    giảm số dư phải thu tương ứng.
  \item
    \texttt{repay} (Trả nợ đã vay): Owner hoàn trả tiền đã vay cho chủ
    nợ. Tiền xuất khỏi ví của Owner. Hậu quả: Giảm số dư ví xuất tiền,
    giảm số dư phải trả tương ứng.
  \item
    \texttt{write-off\ lend} (Xóa nợ cho vay): Owner chính thức xác nhận
    khoản nợ cho vay không thể thu hồi được (nợ xấu). Hậu quả: Giảm số
    dư phải thu, ghi nhận tổn thất tài chính, hoàn toàn không làm biến
    động số dư ví tiền mặt.
  \item
    \texttt{write-off\ borrow} (Được xóa nợ đi vay): Chủ nợ quyết định
    miễn trừ nghĩa vụ nợ cho Owner. Hậu quả: Giảm số dư phải trả, ghi
    nhận một khoản thu nhập bất thường, hoàn toàn không làm biến động số
    dư ví tiền mặt.
  \end{itemize}
\item
  \textbf{Bất biến số dư nợ hợp lệ:}\\
  Tại mọi thời điểm, số dư nợ tổng thể và số dư nợ chi tiết theo từng cá
  nhân (\(person\)) phải bảo đảm không âm:
  \[Receivable = \sum Lend - \sum Collect - \sum WriteOffLend \ge 0\]
  \[Payable = \sum Borrow - \sum Repay - \sum WriteOffBorrow \ge 0\] Quy
  tắc kiểm tra biên:

  \begin{itemize}
  \tightlist
  \item
    Tổng số tiền thu nợ từ một cá nhân không được vượt quá số dư nợ phải
    thu hiện tại của cá nhân đó
    (\(\sum Collect_{person} \le \sum Lend_{person} - \sum WriteOffLend_{person}\)).
  \item
    Tổng số tiền trả nợ cho một chủ nợ không được vượt quá số dư nợ phải
    trả hiện tại đối với chủ nợ đó
    (\(\sum Repay_{person} \le \sum Borrow_{person} - \sum WriteOffBorrow_{person}\)).
  \item
    Mọi thao tác thu nợ hoặc trả nợ vượt quá giá trị nợ gốc còn lại đều
    bị hệ thống từ chối thực thi.
  \end{itemize}
\end{enumerate}

\subsubsection{3.8 Quy tắc phân bổ quỹ mục tiêu
(BR-FIN-08)}\label{38-quy-tux1eafc-phuxe2n-bux1ed5-quux1ef9-mux1ee5c-tiuxeau-br-fin-08}

Quỹ mục tiêu tài chính là cơ chế quy hoạch và gắn kết các nguồn lực tài
sản rời rạc vào các kế hoạch dài hạn cụ thể của cá nhân.

\begin{enumerate}
\def\labelenumi{\arabic{enumi}.}
\item
  \textbf{Các quỹ mục tiêu tiêu chuẩn:}\\
  Hệ thống khởi tạo 4 quỹ mục tiêu tiêu chuẩn phục vụ các trụ cột an
  toàn tài chính:

  \begin{itemize}
  \tightlist
  \item
    \texttt{emergency} (Quỹ khẩn cấp): Dự phòng tài chính cho các biến
    cố bất ngờ (thất nghiệp, hư hỏng tài sản lớn), yêu cầu tính thanh
    khoản cao.
  \item
    \texttt{health} (Quỹ y tế và sức khỏe): Dự phòng chi phí khám chữa
    bệnh, bảo hiểm và chăm sóc sức khỏe.
  \item
    \texttt{target} (Quỹ mục tiêu cụ thể): Tích lũy mua sắm tài sản lớn
    ngắn hạn và trung hạn (mua xe, mua sắm thiết bị, du lịch).
  \item
    \texttt{retirement} (Quỹ hưu trí): Tích lũy và đầu tư dài hạn phục
    vụ giai đoạn ngừng lao động.
  \end{itemize}
\item
  \textbf{Cấu trúc tài sản hợp thành quỹ mục tiêu:}\\
  Quỹ mục tiêu là thực thể kế toán tổng hợp giá trị từ 3 cấu phần tài
  sản (không phải tài khoản ngân hàng vật lý riêng biệt):

  \begin{itemize}
  \tightlist
  \item
    Tiền mặt phân bổ (\texttt{Cash}): Số tiền được chỉ định phân bổ từ
    bảng tiền mặt dành riêng cho quỹ (\texttt{funds\_cash}).
  \item
    Sổ tiết kiệm liên kết (\texttt{Savings}): Tổng số dư hiện tại của
    tất cả các sổ tiết kiệm có gắn mã tham chiếu tới quỹ đó.
  \item
    Chứng chỉ quỹ liên kết (\texttt{FundCerts}): Tổng giá trị thị trường
    hiện tại của tất cả các khoản chứng chỉ quỹ có gắn mã tham chiếu tới
    quỹ đó. Tổng giá trị tích lũy của một quỹ (\(Total_{fund}\)) được
    xác định:
    \[
    Total_{fund}=Cash_{fund}+\sum_{s\in Savings_{fund}}Balance_s+\sum_{c\in FundCerts_{fund}}Balance_c
    \]
  \end{itemize}
\item
  \textbf{Tính toán tiến độ hoàn thành mục tiêu:}\\
  Mỗi quỹ mục tiêu có một số tiền đích cần đạt (\(Target \ge 0\)). Tiến
  độ hoàn thành (\(Process\)) được xác định theo tỷ lệ phần trăm:
  \[Process = \begin{cases} \frac{Total_{fund}}{Target} \times 100\%, & \text{khi } Target > 0 \\ 100\%, & \text{khi } Target = 0 \text{ và } Total_{fund} \ge 0 \end{cases}\]
  Khi \(Process \ge 100\%\), quỹ mục tiêu được hệ thống đánh dấu là đã
  hoàn thành mục tiêu tài chính.
\item
  \textbf{Tính toàn vẹn của việc phân bổ tài sản:}\\
  Nhằm ngăn chặn hiện tượng tính trùng giá trị tài sản
  (double-counting), mỗi sổ tiết kiệm cụ thể và mỗi lô chứng chỉ quỹ cụ
  thể chỉ được phép gán vào duy nhất một quỹ mục tiêu tại một thời điểm.
  Khi một tài sản được chuyển quyền sở hữu sang quỹ mục tiêu khác, giá
  trị của quỹ cũ sẽ giảm tương ứng và giá trị của quỹ mới sẽ tăng tương
  ứng, bảo đảm tổng giá trị tài sản toàn hệ thống không thay đổi.
\end{enumerate}

\clearpage
\subsection{4. YÊU CẦU CHỨC NĂNG CHI TIẾT (FUNCTIONAL
REQUIREMENTS)}\label{4-yuxeau-cux1ea7u-chux1ee9c-nux103ng-chi-tiux1ebft-functional-requirements}

Phần này đặc tả toàn bộ 39 yêu cầu chức năng (Functional Requirements)
cho 10 phân hệ nghiệp vụ của Hệ thống Quản lý Tài chính Cá nhân (PFMS),
tuân thủ tiêu chuẩn ISO/IEC/IEEE 29148:2018. Mỗi yêu cầu được chuẩn hóa
theo cấu trúc 8 trường thông tin: Mã yêu cầu, Tên chức năng, Mô tả, Dữ
liệu đầu vào, Quy tắc xử lý và ràng buộc, Dữ liệu đầu ra, Ca sử dụng
tham chiếu, và Tiêu chí nghiệm thu.

\subsubsection{4.1 Phân hệ Quản lý ví tiền
(FR-WAL)}\label{41-phuxe2n-hux1ec7-quux1ea3n-luxfd-vuxed-tiux1ec1n-fr-wal}

Phân hệ Quản lý ví tiền chịu trách nhiệm theo dõi toàn bộ các nguồn tiền
thanh khoản của người dùng, bao gồm tiền mặt, tài khoản thanh toán ngân
hàng và ví điện tử. Phân hệ bảo đảm tính bất biến số dư ví theo quy tắc
BR-FIN-02 và phản ánh cơ cấu thanh khoản tức thời.

\paragraph{4.1.1 Yêu cầu FR-WAL-01: Khởi tạo ví tiền
mới}\label{411-yuxeau-cux1ea7u-fr-wal-01-khux1edfi-tux1ea1o-vuxed-tiux1ec1n-mux1edbi}

\begin{itemize}
\tightlist
\item
  \textbf{Mã yêu cầu:} \texttt{FR-WAL-01}
\item
  \textbf{Tên chức năng:} Khởi tạo ví tiền mới
\item
  \textbf{Mô tả:} Cho phép người dùng tạo một ví tài chính mới trong hệ
  thống để quản lý tiền mặt, tài khoản thanh toán ngân hàng hoặc ví điện
  tử với số dư khởi tạo xác định.
\item
  \textbf{Dữ liệu đầu vào:}

  \begin{itemize}
  \tightlist
  \item
    Tên ví (\texttt{name}): Chuỗi ký tự (String), độ dài từ 1 đến 50 ký
    tự, bắt buộc, không được trùng lặp với tên ví đã tồn tại (không phân
    biệt chữ hoa, chữ thường).
  \item
    Loại ví (\texttt{type}): Kiểu liệt kê (Enum) gồm \texttt{CASH} (tiền
    mặt), \texttt{BANK} (tài khoản ngân hàng), \texttt{E\_WALLET} (ví
    điện tử), \texttt{CREDIT} (thẻ tín dụng). Bắt buộc.
  \item
    Số dư khởi tạo (\texttt{initialBalance}): Số học chính xác
    \texttt{BigDecimal}, mặc định bằng 0, bắt buộc lớn hơn hoặc bằng 0
    đối với ví thông thường, cho phép nhỏ hơn 0 nếu loại ví là
    \texttt{CREDIT}.
  \item
    Ghi chú (\texttt{description}): Chuỗi ký tự (String), tối đa 255 ký
    tự, tùy chọn.
  \end{itemize}
\item
  \textbf{Quy tắc xử lý và ràng buộc:}

  \begin{itemize}
  \tightlist
  \item
    Kiểm tra tính duy nhất của tên ví so với danh sách các ví hiện có
    trong hệ thống (sau khi loại bỏ khoảng trắng thừa ở hai đầu và
    chuyển về chữ thường).
  \item
    Xác lập trạng thái mặc định của ví là \texttt{ACTIVE} (đang hoạt
    động).
  \item
    Ghi nhận số dư ban đầu \(Balance_w(0) = initialBalance\) theo quy
    tắc \texttt{BR-FIN-02}. Khoản số dư ban đầu này không tạo ra bản ghi
    thu chi giả mạo trong sổ nhật ký giao dịch mà đóng vai trò là số dư
    gốc của sổ cái ví.
  \item
    Chuẩn hóa số dư khởi tạo về kiểu \texttt{BigDecimal} với scale chuẩn
    là 2 chữ số thập phân theo \texttt{BR-FIN-01}.
  \end{itemize}
\item
  \textbf{Dữ liệu đầu ra:}

  \begin{itemize}
  \tightlist
  \item
    Đối tượng ví (\texttt{Wallet}) mới được lưu thành công vào cơ sở dữ
    liệu kèm mã định danh duy nhất (\texttt{walletId}).
  \item
    Thông báo tạo ví thành công hiển thị cho người dùng.
  \end{itemize}
\item
  \textbf{Ca sử dụng tham chiếu:} \texttt{UC01\ -\ Quản\ lý\ ví\ tiền}.
\item
  \textbf{Tiêu chí nghiệm thu:}

  \begin{itemize}
  \tightlist
  \item
    TC-WAL-01.1: Khi người dùng nhập tên ví "MB Bank", loại
    \texttt{BANK}, số dư ban đầu 5,000,000 VNĐ, hệ thống tạo thành công
    ví mới có số dư khả dụng bằng đúng 5,000,000 VNĐ và trạng thái
    \texttt{ACTIVE}.
  \item
    TC-WAL-01.2: Khi người dùng nhập tên ví trùng lặp với ví đã có (ví
    dụ nhập "mb bank" khi đã tồn tại "MB Bank"), hệ thống từ chối tạo và
    hiển thị thông báo lỗi tên ví đã tồn tại.
  \item
    TC-WAL-01.3: Khi người dùng nhập số dư khởi tạo là số âm cho ví loại
    \texttt{CASH} hoặc \texttt{BANK}, hệ thống từ chối và thông báo số
    dư ban đầu không được nhỏ hơn 0.
  \end{itemize}
\end{itemize}

\paragraph{4.1.2 Yêu cầu FR-WAL-02: Xem danh sách và chi tiết số dư
động}\label{412-yuxeau-cux1ea7u-fr-wal-02-xem-danh-suxe1ch-vuxe0-chi-tiux1ebft-sux1ed1-dux1b0-ux111ux1ed9ng}

\begin{itemize}
\tightlist
\item
  \textbf{Mã yêu cầu:} \texttt{FR-WAL-02}
\item
  \textbf{Tên chức năng:} Xem danh sách và chi tiết số dư động
\item
  \textbf{Mô tả:} Cung cấp danh sách toàn bộ các ví tiền đang quản lý
  cùng số dư khả dụng thực tế được tính toán động tại thời điểm truy vấn
  dựa trên toàn bộ lịch sử dòng tiền.
\item
  \textbf{Dữ liệu đầu vào:}

  \begin{itemize}
  \tightlist
  \item
    Bộ lọc trạng thái ví (\texttt{statusFilter}): Tùy chọn, gồm
    \texttt{ACTIVE} (mặc định), \texttt{ARCHIVED} (đã lưu trữ), hoặc
    \texttt{ALL} (toàn bộ).
  \end{itemize}
\item
  \textbf{Quy tắc xử lý và ràng buộc:}

  \begin{itemize}
  \tightlist
  \item
    Tính toán số dư động \(Balance_w(T)\) cho từng ví theo đúng công
    thức bất biến \texttt{BR-FIN-02}:
    \[
    \begin{aligned}
    Balance_w(T) &= Balance_w(0) + \sum Income - \sum Expense + \sum TransferIn \\
    &\quad - \sum TransferOut + \sum Borrow - \sum Repay \\
    &\quad - \sum Lend + \sum Collect
    \end{aligned}
    \]
  \item
    Tuyệt đối không lưu giá trị số dư tĩnh bị ghi đè trực tiếp trong cơ
    sở dữ liệu; mọi số dư hiển thị phải là kết quả tổng hợp tức thời từ
    các giao dịch hợp lệ.
  \item
    Tính tổng số dư khả dụng của toàn bộ các ví đang kích hoạt bằng tổng
    số học chính xác các số dư thành phần thông qua
    \texttt{BigDecimal.add()}.
  \item
    Hiển thị danh sách ví kèm các thông tin nhận diện (tên, loại, trạng
    thái, số dư khả dụng, ngày cập nhật giao dịch gần nhất).
  \end{itemize}
\item
  \textbf{Dữ liệu đầu ra:}

  \begin{itemize}
  \tightlist
  \item
    Danh sách các đối tượng ví kèm số dư động tương ứng.
  \item
    Giá trị tổng số dư tiền mặt thanh khoản toàn hệ thống.
  \end{itemize}
\item
  \textbf{Ca sử dụng tham chiếu:} \texttt{UC01\ -\ Quản\ lý\ ví\ tiền}.
\item
  \textbf{Tiêu chí nghiệm thu:}

  \begin{itemize}
  \tightlist
  \item
    TC-WAL-02.1: Cho ví "Timo" có số dư ban đầu 1,000,000 VNĐ, phát sinh
    1 giao dịch thu 500,000 VNĐ và 1 giao dịch chi 200,000 VNĐ; khi
    người dùng truy vấn danh sách ví, số dư hiển thị của ví Timo phải là
    chính xác 1,300,000 VNĐ.
  \item
    TC-WAL-02.2: Tổng số dư toàn hệ thống hiển thị ở chân danh sách phải
    khớp chính xác đến từng đơn vị đồng với tổng số dư của các ví thành
    phần mà không có sai số làm tròn.
  \end{itemize}
\end{itemize}

\paragraph{4.1.3 Yêu cầu FR-WAL-03: Cập nhật thông tin ví
tiền}\label{413-yuxeau-cux1ea7u-fr-wal-03-cux1eadp-nhux1eadt-thuxf4ng-tin-vuxed-tiux1ec1n}

\begin{itemize}
\tightlist
\item
  \textbf{Mã yêu cầu:} \texttt{FR-WAL-03}
\item
  \textbf{Tên chức năng:} Cập nhật thông tin ví tiền
\item
  \textbf{Mô tả:} Cho phép người dùng chỉnh sửa các thuộc tính phi tài
  chính của ví tiền (tên ví, mô tả, trạng thái hoạt động); nghiêm cấm
  chỉnh sửa trực tiếp trường số dư khả dụng.
\item
  \textbf{Dữ liệu đầu vào:}

  \begin{itemize}
  \tightlist
  \item
    Mã định danh ví (\texttt{walletId}): Chuỗi ký tự, bắt buộc.
  \item
    Tên ví mới (\texttt{name}): Chuỗi ký tự, từ 1 đến 50 ký tự, bắt
    buộc.
  \item
    Trạng thái mới (\texttt{status}): Enum \texttt{ACTIVE} hoặc
    \texttt{ARCHIVED}.
  \item
    Ghi chú mới (\texttt{description}): Chuỗi ký tự, tối đa 255 ký tự.
  \end{itemize}
\item
  \textbf{Quy tắc xử lý và ràng buộc:}

  \begin{itemize}
  \tightlist
  \item
    Kiểm tra tên ví mới không được trùng lặp với tên của các ví khác
    hiện có trong hệ thống.
  \item
    Nếu người dùng chuyển trạng thái ví từ \texttt{ACTIVE} sang
    \texttt{ARCHIVED}, hệ thống kiểm tra số dư hiện tại của ví:

    \begin{itemize}
    \tightlist
    \item
      Nếu số dư khác 0, hệ thống đưa ra cảnh báo khuyến nghị điều chuyển
      hết tiền trước khi lưu trữ, nhưng vẫn cho phép lưu trữ nếu người
      dùng xác nhận.
    \item
      Ví ở trạng thái \texttt{ARCHIVED} sẽ bị khóa không cho phép chọn
      làm ví nguồn hoặc ví đích trong các giao dịch phát sinh mới.
    \end{itemize}
  \item
    Hệ thống tuyệt đối không cung cấp giao diện hoặc phương thức cho
    phép sửa đổi trường \texttt{currentBalance} hay
    \texttt{initialBalance} sau khi ví đã phát sinh giao dịch lịch sử.
  \end{itemize}
\item
  \textbf{Dữ liệu đầu ra:}

  \begin{itemize}
  \tightlist
  \item
    Thông tin ví được cập nhật thành công trong hệ thống.
  \item
    Thông báo kết quả cập nhật hiển thị cho người dùng.
  \end{itemize}
\item
  \textbf{Ca sử dụng tham chiếu:} \texttt{UC01\ -\ Quản\ lý\ ví\ tiền}.
\item
  \textbf{Tiêu chí nghiệm thu:}

  \begin{itemize}
  \tightlist
  \item
    TC-WAL-03.1: Khi người dùng đổi tên ví "Timo Bank" thành "Timo
    Digital", hệ thống lưu tên mới và hiển thị tên mới trên toàn bộ giao
    diện mà số dư ví không đổi.
  \item
    TC-WAL-03.2: Khi người dùng chuyển một ví có số dư về trạng thái
    \texttt{ARCHIVED}, ví đó không còn xuất hiện trong danh sách lựa
    chọn ví khi tạo giao dịch mới.
  \item
    TC-WAL-03.3: Mọi yêu cầu cố ý chỉnh sửa số dư trực tiếp qua API hoặc
    giao diện đều bị chặn và trả về ngoại lệ bảo toàn dữ liệu.
  \end{itemize}
\end{itemize}

\paragraph{4.1.4 Yêu cầu FR-WAL-04: Thống kê cơ cấu tài sản theo
ví}\label{414-yuxeau-cux1ea7u-fr-wal-04-thux1ed1ng-kuxea-cux1a1-cux1ea5u-tuxe0i-sux1ea3n-theo-vuxed}

\begin{itemize}
\tightlist
\item
  \textbf{Mã yêu cầu:} \texttt{FR-WAL-04}
\item
  \textbf{Tên chức năng:} Thống kê cơ cấu tài sản theo ví
\item
  \textbf{Mô tả:} Phân tích và hiển thị tỷ trọng phần trăm giá trị tiền
  tệ lưu giữ tại từng ví so với tổng lượng tiền mặt thanh khoản trong
  toàn bộ hệ thống.
\item
  \textbf{Dữ liệu đầu vào:}

  \begin{itemize}
  \tightlist
  \item
    Thời điểm tính toán (\(T\)): Mặc định là thời điểm hiện tại.
  \end{itemize}
\item
  \textbf{Quy tắc xử lý và ràng buộc:}

  \begin{itemize}
  \tightlist
  \item
    Tính tổng số dư của toàn bộ các ví đang ở trạng thái
    \texttt{ACTIVE}:
    \(TotalWallets = \sum_{w \in ActiveWallets} Balance_w(T)\).
  \item
    Với từng ví \(w\), tính tỷ trọng theo công thức:
    \[Rate_w = \begin{cases} \frac{Balance_w(T)}{TotalWallets} \times 100\%, & \text{khi } TotalWallets > 0 \\ 0\%, & \text{khi } TotalWallets = 0 \end{cases}\]
  \item
    Phép chia tỷ trọng sử dụng \texttt{BigDecimal.divide()} với scale
    tối thiểu là 4 chữ số thập phân và chỉ định tường minh chế độ làm
    tròn \texttt{RoundingMode.HALF\_EVEN}.
  \item
    Giá trị hiển thị ra màn hình được định dạng 2 chữ số thập phân (ví
    dụ: 45.25\%).
  \item
    Sắp xếp danh sách ví theo thứ tự tỷ trọng giảm dần để làm nổi bật ví
    chiếm tỷ trọng thanh khoản cao nhất.
  \end{itemize}
\item
  \textbf{Dữ liệu đầu ra:}

  \begin{itemize}
  \tightlist
  \item
    Bảng cơ cấu phân bổ ví gồm: Tên ví, Loại ví, Số dư khả dụng, Tỷ
    trọng phần trăm (\%), Xếp hạng tỷ trọng.
  \end{itemize}
\item
  \textbf{Ca sử dụng tham chiếu:} \texttt{UC01\ -\ Quản\ lý\ ví\ tiền}.
\item
  \textbf{Tiêu chí nghiệm thu:}

  \begin{itemize}
  \tightlist
  \item
    TC-WAL-04.1: Cho ví A có số dư 6,000,000 VNĐ và ví B có số dư
    4,000,000 VNĐ (tổng 10,000,000 VNĐ); hệ thống tính ra tỷ trọng của
    ví A là 60.00\% và ví B là 40.00\%.
  \item
    TC-WAL-04.2: Khi tổng số dư toàn bộ các ví bằng 0, tỷ trọng hiển thị
    của từng ví là 0.00\% và hệ thống không phát sinh lỗi chia cho 0
    (\texttt{ArithmeticException}).
  \end{itemize}
\end{itemize}

\subsubsection{4.2 Phân hệ Ghi nhận thu chi
(FR-TXN)}\label{42-phuxe2n-hux1ec7-ghi-nhux1eadn-thu-chi-fr-txn}

Phân hệ Ghi nhận thu chi đóng vai trò là nhật ký ghi chép các luồng tiền
phát sinh hàng ngày của người dùng. Phân hệ kết nối trực tiếp biến động
tiền tệ với cấu trúc phân loại thu chi 3 cấp và số dư thực tế của từng
ví.

\paragraph{4.2.1 Yêu cầu FR-TXN-01: Tạo mới giao dịch thu
nhập}\label{421-yuxeau-cux1ea7u-fr-txn-01-tux1ea1o-mux1edbi-giao-dux1ecbch-thu-nhux1eadp}

\begin{itemize}
\tightlist
\item
  \textbf{Mã yêu cầu:} \texttt{FR-TXN-01}
\item
  \textbf{Tên chức năng:} Tạo mới giao dịch thu nhập
\item
  \textbf{Mô tả:} Ghi nhận một khoản thu nhập thực tế làm tăng số dư của
  một ví chỉ định và phân loại vào danh mục thu nhập tương ứng.
\item
  \textbf{Dữ liệu đầu vào:}

  \begin{itemize}
  \tightlist
  \item
    Ngày giao dịch (\texttt{date}): Kiểu \texttt{LocalDate}, bắt buộc,
    mặc định là ngày hiện tại.
  \item
    Số tiền thu (\texttt{amount}): Kiểu \texttt{BigDecimal}, bắt buộc,
    giá trị phải lớn hơn 0.
  \item
    Ví nhận tiền (\texttt{walletId}): Chuỗi ký tự, bắt buộc, ví phải ở
    trạng thái \texttt{ACTIVE}.
  \item
    Danh mục thu nhập (\texttt{categoryId}): Chuỗi ký tự, bắt buộc, danh
    mục phải thuộc nhóm có \texttt{Type\ =\ Income}.
  \item
    Ghi chú (\texttt{note}): Chuỗi ký tự, tối đa 255 ký tự, tùy chọn.
  \end{itemize}
\item
  \textbf{Quy tắc xử lý và ràng buộc:}

  \begin{itemize}
  \tightlist
  \item
    Kiểm tra tính hợp lệ của số tiền thu: \(amount > 0\), chuẩn hóa về
    kiểu \texttt{BigDecimal} với scale chuẩn là 2 theo
    \texttt{BR-FIN-01}.
  \item
    Kiểm tra tính toàn vẹn của danh mục: \texttt{categoryId} bắt buộc
    phải tồn tại và có thuộc tính cấp 1 \texttt{Type\ =\ Income} theo
    \texttt{BR-FIN-04}. Nếu người dùng chọn danh mục thuộc
    \texttt{Type\ =\ Expense}, hệ thống từ chối ghi nhận.
  \item
    Tạo bản ghi mới trong sổ nhật ký giao dịch (\texttt{transactions})
    với mã định danh duy nhất (\texttt{txnId}).
  \item
    Cập nhật số dư động của ví nhận tiền:
    \(Balance_{wallet}(t) = Balance_{wallet}(t-1) + amount\) theo đúng
    quy tắc bảo toàn số dư \texttt{BR-FIN-02}.
  \item
    Kích hoạt cập nhật tự động các bảng thống kê dòng tiền tháng liên
    quan.
  \end{itemize}
\item
  \textbf{Dữ liệu đầu ra:}

  \begin{itemize}
  \tightlist
  \item
    Bản ghi giao dịch thu nhập mới được lưu vào cơ sở dữ liệu kèm mã
    \texttt{txnId}.
  \item
    Số dư ví nhận tiền tăng thêm đúng bằng số tiền thu nhập.
  \item
    Thông báo ghi nhận giao dịch thành công.
  \end{itemize}
\item
  \textbf{Ca sử dụng tham chiếu:}
  \texttt{UC02\ -\ Ghi\ nhận\ giao\ dịch\ thu\ chi}.
\item
  \textbf{Tiêu chí nghiệm thu:}

  \begin{itemize}
  \tightlist
  \item
    TC-TXN-01.1: Cho ví "Vietcombank" có số dư 10,000,000 VNĐ; khi người
    dùng ghi nhận khoản thu nhập tiền lương 15,000,000 VNĐ vào ví
    Vietcombank, hệ thống ghi nhận thành công và số dư ví tăng lên đúng
    25,000,000 VNĐ.
  \item
    TC-TXN-01.2: Khi người dùng nhập số tiền thu nhập bằng 0 hoặc là số
    âm, hệ thống chặn lại và thông báo lỗi số tiền phải lớn hơn 0.
  \item
    TC-TXN-01.3: Khi người dùng chọn danh mục "Ăn uống" (thuộc Type
    Expense) để ghi nhận thu nhập, hệ thống từ chối và yêu cầu chọn danh
    mục thuộc nhóm thu nhập.
  \end{itemize}
\end{itemize}

\paragraph{4.2.2 Yêu cầu FR-TXN-02: Tạo mới giao dịch chi
tiêu}\label{422-yuxeau-cux1ea7u-fr-txn-02-tux1ea1o-mux1edbi-giao-dux1ecbch-chi-tiuxeau}

\begin{itemize}
\tightlist
\item
  \textbf{Mã yêu cầu:} \texttt{FR-TXN-02}
\item
  \textbf{Tên chức năng:} Tạo mới giao dịch chi tiêu
\item
  \textbf{Mô tả:} Ghi nhận một khoản chi tiêu tiêu dùng hoặc thanh toán
  sinh hoạt làm giảm số dư của một ví chỉ định và phân loại vào danh mục
  chi tiêu tương ứng.
\item
  \textbf{Dữ liệu đầu vào:}

  \begin{itemize}
  \tightlist
  \item
    Ngày giao dịch (\texttt{date}): Kiểu \texttt{LocalDate}, bắt buộc,
    mặc định là ngày hiện tại.
  \item
    Số tiền chi (\texttt{amount}): Kiểu \texttt{BigDecimal}, bắt buộc,
    giá trị phải lớn hơn 0.
  \item
    Ví chi tiền (\texttt{walletId}): Chuỗi ký tự, bắt buộc, ví phải ở
    trạng thái \texttt{ACTIVE}.
  \item
    Danh mục chi tiêu (\texttt{categoryId}): Chuỗi ký tự, bắt buộc, danh
    mục phải thuộc nhóm có \texttt{Type\ =\ Expense}.
  \item
    Ghi chú (\texttt{note}): Chuỗi ký tự, tối đa 255 ký tự, tùy chọn.
  \end{itemize}
\item
  \textbf{Quy tắc xử lý và ràng buộc:}

  \begin{itemize}
  \tightlist
  \item
    Kiểm tra tính hợp lệ của số tiền chi: \(amount > 0\).
  \item
    Kiểm tra tính toàn vẹn của danh mục: \texttt{categoryId} bắt buộc
    phải tồn tại và có \texttt{Type\ =\ Expense} theo
    \texttt{BR-FIN-04}.
  \item
    Kiểm tra ràng buộc số dư không âm theo quy tắc \texttt{BR-FIN-02}:
    Đối với các ví không phải loại thẻ tín dụng (\texttt{CREDIT}), số dư
    khả dụng hiện tại của ví phải thỏa mãn:
    \[Balance_{wallet} - amount \ge 0\] Nếu số dư không đủ để chi trả,
    hệ thống từ chối giao dịch và đưa ra thông báo thiếu hụt số dư.
  \item
    Nếu thỏa mãn mọi điều kiện, lưu bản ghi giao dịch mới vào sổ nhật ký
    và trừ số dư của ví chi tiền:
    \(Balance_{wallet}(t) = Balance_{wallet}(t-1) - amount\).
  \end{itemize}
\item
  \textbf{Dữ liệu đầu ra:}

  \begin{itemize}
  \tightlist
  \item
    Bản ghi giao dịch chi tiêu mới được lưu thành công kèm mã
    \texttt{txnId}.
  \item
    Số dư của ví chi tiền giảm đúng bằng số tiền chi tiêu.
  \item
    Thông báo ghi nhận giao dịch chi tiêu thành công.
  \end{itemize}
\item
  \textbf{Ca sử dụng tham chiếu:}
  \texttt{UC02\ -\ Ghi\ nhận\ giao\ dịch\ thu\ chi}.
\item
  \textbf{Tiêu chí nghiệm thu:}

  \begin{itemize}
  \tightlist
  \item
    TC-TXN-02.1: Cho ví "Tiền mặt" có số dư 500,000 VNĐ; khi người dùng
    tạo giao dịch chi tiêu 120,000 VNĐ cho danh mục "Ăn trưa", hệ thống
    ghi nhận thành công và số dư ví còn lại là 380,000 VNĐ.
  \item
    TC-TXN-02.2: Cho ví "Tiền mặt" có số dư 200,000 VNĐ; khi người dùng
    nhập khoản chi 350,000 VNĐ, hệ thống từ chối giao dịch và thông báo
    số dư ví không đủ để thực hiện chi trả.
  \item
    TC-TXN-02.3: Với ví loại \texttt{CREDIT} có hạn mức thấu chi, cho
    phép ghi nhận chi tiêu khiến số dư trở thành số âm trong phạm vi hạn
    mức cho phép.
  \end{itemize}
\end{itemize}

\paragraph{4.2.3 Yêu cầu FR-TXN-03: Tra cứu và lọc lịch sử giao
dịch}\label{423-yuxeau-cux1ea7u-fr-txn-03-tra-cux1ee9u-vuxe0-lux1ecdc-lux1ecbch-sux1eed-giao-dux1ecbch}

\begin{itemize}
\tightlist
\item
  \textbf{Mã yêu cầu:} \texttt{FR-TXN-03}
\item
  \textbf{Tên chức năng:} Tra cứu và lọc lịch sử giao dịch
\item
  \textbf{Mô tả:} Cung cấp công cụ tìm kiếm, lọc đa tiêu chí và tổng hợp
  số liệu trên sổ nhật ký giao dịch thu chi nhằm hỗ trợ người dùng phân
  tích lịch sử tài chính.
\item
  \textbf{Dữ liệu đầu vào:}

  \begin{itemize}
  \tightlist
  \item
    Khoảng thời gian (\texttt{fromDate}, \texttt{toDate}): Kiểu
    \texttt{LocalDate}, tùy chọn (mặc định là từ ngày đầu tháng đến ngày
    cuối tháng hiện tại).
  \item
    Danh sách mã ví (\texttt{walletIds}): Danh sách chuỗi ký tự, tùy
    chọn (mặc định lấy tất cả các ví).
  \item
    Phân loại luồng tiền (\texttt{type}): Tùy chọn gồm \texttt{ALL},
    \texttt{INCOME}, \texttt{EXPENSE}.
  \item
    Nhóm danh mục (\texttt{groupId}): Chuỗi ký tự, tùy chọn.
  \item
    Danh mục chi tiết (\texttt{categoryId}): Chuỗi ký tự, tùy chọn.
  \item
    Từ khóa tìm kiếm (\texttt{keyword}): Chuỗi ký tự, tìm kiếm không
    phân biệt hoa thường trong trường ghi chú (\texttt{note}).
  \item
    Thông số phân trang (\texttt{pageIndex}, \texttt{pageSize}): Số
    nguyên, mặc định trang 1, 50 bản ghi trên một trang.
  \end{itemize}
\item
  \textbf{Quy tắc xử lý và ràng buộc:}

  \begin{itemize}
  \tightlist
  \item
    Áp dụng phép toán kết hợp logic AND giữa các tiêu chí lọc được chọn.
  \item
    Tính toán tổng thu nhập, tổng chi tiêu và dòng tiền thuần (Cash Flow
    = Tổng thu nhập - Tổng chi tiêu) của tập kết quả thỏa mãn điều kiện
    lọc.
  \item
    Mọi phép tính tổng hợp trên tập kết quả lọc phải sử dụng
    \texttt{BigDecimal} để bảo đảm độ chính xác tuyệt đối theo
    \texttt{BR-FIN-01}.
  \item
    Sắp xếp kết quả mặc định theo ngày giao dịch giảm dần (giao dịch mới
    nhất hiển thị trước).
  \end{itemize}
\item
  \textbf{Dữ liệu đầu ra:}

  \begin{itemize}
  \tightlist
  \item
    Danh sách các bản ghi giao dịch thỏa mãn điều kiện lọc (kèm thông
    tin phân trang).
  \item
    Tổng số bản ghi tìm thấy.
  \item
    Thống kê tổng hợp: Tổng thu nhập, Tổng chi tiêu, Chênh lệch dòng
    tiền thuần của tập kết quả.
  \end{itemize}
\item
  \textbf{Ca sử dụng tham chiếu:}
  \texttt{UC02\ -\ Ghi\ nhận\ giao\ dịch\ thu\ chi}.
\item
  \textbf{Tiêu chí nghiệm thu:}

  \begin{itemize}
  \tightlist
  \item
    TC-TXN-03.1: Khi người dùng lọc theo khoảng ngày từ 01/08/2026 đến
    31/08/2026 và chọn loại giao dịch \texttt{EXPENSE}, hệ thống trả về
    chính xác tất cả các giao dịch chi tiêu trong tháng 8/2026 kèm tổng
    số tiền chi tiêu khớp với sổ cái.
  \item
    TC-TXN-03.2: Khi người dùng nhập từ khóa "tiền điện", hệ thống trả
    về toàn bộ các giao dịch có chứa cụm từ "tiền điện" trong phần ghi
    chú, không phân biệt hoa thường.
  \end{itemize}
\end{itemize}

\paragraph{4.2.4 Yêu cầu FR-TXN-04: Chỉnh sửa giao dịch thu
chi}\label{424-yuxeau-cux1ea7u-fr-txn-04-chux1ec9nh-sux1eeda-giao-dux1ecbch-thu-chi}

\begin{itemize}
\tightlist
\item
  \textbf{Mã yêu cầu:} \texttt{FR-TXN-04}
\item
  \textbf{Tên chức năng:} Chỉnh sửa giao dịch thu chi
\item
  \textbf{Mô tả:} Cho phép người dùng cập nhật thông tin của một giao
  dịch đã lưu trong quá khứ, tự động tính toán hoàn nguyên số dư ví cũ
  và cập nhật số dư ví mới nhằm bảo toàn tính toàn vẹn của sổ cái.
\item
  \textbf{Dữ liệu đầu vào:}

  \begin{itemize}
  \tightlist
  \item
    Mã giao dịch cần sửa (\texttt{txnId}): Chuỗi ký tự, bắt buộc.
  \item
    Ngày mới (\texttt{date}): Kiểu \texttt{LocalDate}, bắt buộc.
  \item
    Số tiền mới (\texttt{amount}): Kiểu \texttt{BigDecimal}, bắt buộc,
    lớn hơn 0.
  \item
    Ví mới (\texttt{walletId}): Chuỗi ký tự, bắt buộc, ví ở trạng thái
    \texttt{ACTIVE}.
  \item
    Danh mục mới (\texttt{categoryId}): Chuỗi ký tự, bắt buộc, phù hợp
    với \texttt{Type} của giao dịch.
  \item
    Ghi chú mới (\texttt{note}): Chuỗi ký tự, tối đa 255 ký tự.
  \end{itemize}
\item
  \textbf{Quy tắc xử lý và ràng buộc:}

  \begin{itemize}
  \tightlist
  \item
    Tìm bản ghi giao dịch gốc dựa trên \texttt{txnId}. Nếu không tìm
    thấy, báo lỗi không tồn tại giao dịch.
  \item
    Thực hiện hoàn nguyên (revert) tác động của giao dịch cũ lên ví cũ
    theo \texttt{BR-FIN-02}:

    \begin{itemize}
    \tightlist
    \item
      Nếu giao dịch cũ là \texttt{Income}: Trừ số tiền cũ khỏi ví cũ.
    \item
      Nếu giao dịch cũ là \texttt{Expense}: Cộng số tiền cũ trả lại ví
      cũ.
    \end{itemize}
  \item
    Kiểm tra điều kiện áp dụng cho ví mới:

    \begin{itemize}
    \tightlist
    \item
      Nếu giao dịch mới là \texttt{Expense}: Kiểm tra số dư ví mới sau
      khi hoàn nguyên có thỏa mãn \(Balance_{newWallet} - amount \ge 0\)
      hay không.
    \item
      Nếu không đủ số dư, hủy bỏ toàn bộ thao tác (rollback) và báo lỗi
      không đủ số dư tại ví mới.
    \end{itemize}
  \item
    Nếu điều kiện hợp lệ, áp dụng tác động của giao dịch mới:

    \begin{itemize}
    \tightlist
    \item
      Nếu là \texttt{Income}: Cộng số tiền mới vào ví mới.
    \item
      Nếu là \texttt{Expense}: Trừ số tiền mới khỏi ví mới.
    \end{itemize}
  \item
    Cập nhật thông tin bản ghi giao dịch trong cơ sở dữ liệu trong cùng
    một khối giao dịch nguyên tử.
  \end{itemize}
\item
  \textbf{Dữ liệu đầu ra:}

  \begin{itemize}
  \tightlist
  \item
    Bản ghi giao dịch đã cập nhật.
  \item
    Số dư của các ví liên quan được đồng bộ hóa chính xác.
  \item
    Thông báo cập nhật giao dịch thành công.
  \end{itemize}
\item
  \textbf{Ca sử dụng tham chiếu:}
  \texttt{UC02\ -\ Ghi\ nhận\ giao\ dịch\ thu\ chi}.
\item
  \textbf{Tiêu chí nghiệm thu:}

  \begin{itemize}
  \tightlist
  \item
    TC-TXN-04.1: Cho giao dịch chi 100,000 VNĐ từ ví A; khi người dùng
    sửa thành chi 150,000 VNĐ từ ví B, hệ thống hoàn lại 100,000 VNĐ cho
    ví A và trừ 150,000 VNĐ từ ví B (với điều kiện ví B có đủ tiền).
  \item
    TC-TXN-04.2: Khi người dùng sửa số tiền chi tiêu tăng lên vượt quá
    số dư khả dụng của ví, hệ thống từ chối thao tác, thông báo lỗi số
    dư và khôi phục trạng thái ban đầu của giao dịch.
  \end{itemize}
\end{itemize}

\paragraph{4.2.5 Yêu cầu FR-TXN-05: Xóa giao dịch thu
chi}\label{425-yuxeau-cux1ea7u-fr-txn-05-xuxf3a-giao-dux1ecbch-thu-chi}

\begin{itemize}
\tightlist
\item
  \textbf{Mã yêu cầu:} \texttt{FR-TXN-05}
\item
  \textbf{Tên chức năng:} Xóa giao dịch thu chi
\item
  \textbf{Mô tả:} Cho phép xóa bỏ một bản ghi giao dịch thu chi khỏi sổ
  nhật ký, đồng thời tự động hoàn nguyên số dư của ví tiền liên quan về
  trạng thái như trước khi giao dịch phát sinh.
\item
  \textbf{Dữ liệu đầu vào:}

  \begin{itemize}
  \tightlist
  \item
    Mã giao dịch cần xóa (\texttt{txnId}): Chuỗi ký tự, bắt buộc.
  \item
    Xác nhận xóa từ người dùng (\texttt{confirmed}): Kiểu
    \texttt{Boolean}, bắt buộc là \texttt{true}.
  \end{itemize}
\item
  \textbf{Quy tắc xử lý và ràng buộc:}

  \begin{itemize}
  \tightlist
  \item
    Kiểm tra sự tồn tại của giao dịch theo \texttt{txnId}.
  \item
    Thực hiện hoàn nguyên số dư ví theo quy tắc \texttt{BR-FIN-02}:

    \begin{itemize}
    \tightlist
    \item
      Nếu giao dịch cần xóa là \texttt{Income}: Hệ thống kiểm tra số dư
      hiện tại của ví có lớn hơn hoặc bằng số tiền của giao dịch hay
      không (\(Balance_{wallet} - amount \ge 0\)). Nếu việc trừ tiền làm
      ví bị âm số dư, hệ thống đưa ra cảnh báo chặn và từ chối xóa để
      tránh vi phạm bất biến số dư không âm.
    \item
      Nếu thỏa mãn, trừ \(amount\) khỏi ví tương ứng.
    \item
      Nếu giao dịch cần xóa là \texttt{Expense}: Cộng hoàn trả
      \(amount\) vào số dư của ví tương ứng.
    \end{itemize}
  \item
    Xóa bản ghi giao dịch khỏi sổ nhật ký \texttt{transactions}.
  \item
    Cập nhật lại các bảng tổng hợp thống kê tháng có liên quan.
  \end{itemize}
\item
  \textbf{Dữ liệu đầu ra:}

  \begin{itemize}
  \tightlist
  \item
    Bản ghi giao dịch bị xóa khỏi hệ thống.
  \item
    Số dư ví liên quan được hoàn nguyên chính xác.
  \item
    Thông báo xóa giao dịch thành công.
  \end{itemize}
\item
  \textbf{Ca sử dụng tham chiếu:}
  \texttt{UC02\ -\ Ghi\ nhận\ giao\ dịch\ thu\ chi}.
\item
  \textbf{Tiêu chí nghiệm thu:}

  \begin{itemize}
  \tightlist
  \item
    TC-TXN-05.1: Cho ví "MB Bank" có số dư 4,500,000 VNĐ sau khi ghi
    nhận một khoản chi 500,000 VNĐ; khi người dùng xóa giao dịch chi
    này, số dư ví MB Bank tự động phục hồi về mức 5,000,000 VNĐ.
  \item
    TC-TXN-05.2: Khi người dùng cố gắng xóa một giao dịch thu nhập nhưng
    số dư ví hiện tại nhỏ hơn số tiền thu nhập đó (do tiền đã được chi
    tiêu hết), hệ thống từ chối xóa và thông báo việc xóa sẽ làm âm số
    dư ví.
  \end{itemize}
\end{itemize}

\subsubsection{4.3 Phân hệ Chuyển tiền nội bộ
(FR-TRF)}\label{43-phuxe2n-hux1ec7-chuyux1ec3n-tiux1ec1n-nux1ed9i-bux1ed9-fr-trf}

Phân hệ Chuyển tiền nội bộ thực hiện việc điều chuyển tiền tệ giữa các
ví nội bộ của người dùng theo nguyên tắc ghi sổ kép đối ứng của Martin
Fowler, bảo đảm tính bất biến tổng tài sản thanh khoản theo quy tắc
BR-FIN-03.

\paragraph{4.3.1 Yêu cầu FR-TRF-01: Chuyển tiền giữa hai ví nội
bộ}\label{431-yuxeau-cux1ea7u-fr-trf-01-chuyux1ec3n-tiux1ec1n-giux1eefa-hai-vuxed-nux1ed9i-bux1ed9}

\begin{itemize}
\tightlist
\item
  \textbf{Mã yêu cầu:} \texttt{FR-TRF-01}
\item
  \textbf{Tên chức năng:} Chuyển tiền giữa hai ví nội bộ
\item
  \textbf{Mô tả:} Thực hiện chuyển dịch một lượng tiền tệ từ ví nguồn
  sang ví đích trong cùng hệ thống, ghi nhận đồng thời hai bút toán đối
  ứng cân bằng mà không làm thay đổi tổng tài sản ròng của người dùng.
\item
  \textbf{Dữ liệu đầu vào:}

  \begin{itemize}
  \tightlist
  \item
    Ngày chuyển tiền (\texttt{date}): Kiểu \texttt{LocalDate}, bắt buộc,
    mặc định hôm nay.
  \item
    Ví nguồn (\texttt{fromWalletId}): Chuỗi ký tự, bắt buộc, ví ở trạng
    thái \texttt{ACTIVE}.
  \item
    Ví đích (\texttt{toWalletId}): Chuỗi ký tự, bắt buộc, ví ở trạng
    thái \texttt{ACTIVE}.
  \item
    Số tiền chuyển (\texttt{amount}): Kiểu \texttt{BigDecimal}, bắt
    buộc, lớn hơn 0.
  \item
    Ghi chú (\texttt{note}): Chuỗi ký tự, tối đa 255 ký tự, tùy chọn.
  \end{itemize}
\item
  \textbf{Quy tắc xử lý và ràng buộc:}

  \begin{itemize}
  \tightlist
  \item
    Kiểm tra ràng buộc ví nguồn phải khác ví đích:
    \(fromWalletId \ne toWalletId\). Nếu trùng lặp, hệ thống từ chối và
    báo lỗi không thể chuyển tiền cho chính ví đó.
  \item
    Kiểm tra số dư khả dụng của ví nguồn theo quy tắc
    \texttt{BR-FIN-02}: \(Balance_{from} - amount \ge 0\). Nếu số dư
    không đủ, hệ thống từ chối thực hiện.
  \item
    Thực hiện ghi sổ kép nguyên tử (Atomic Double-Entry) theo quy tắc
    \texttt{BR-FIN-03}:

    \begin{itemize}
    \tightlist
    \item
      Trừ chính xác \(amount\) khỏi số dư ví nguồn \(fromWalletId\).
    \item
      Cộng chính xác \(amount\) vào số dư ví đích \(toWalletId\).
    \end{itemize}
  \item
    Cả hai thao tác trừ và cộng phải nằm trong cùng một giao dịch cơ sở
    dữ liệu duy nhất. Nếu xảy ra sự cố tại bất kỳ bước nào, toàn bộ giao
    dịch phải được rollback về trạng thái ban đầu.
  \item
    Giao dịch chuyển tiền nội bộ thuần túy không được tính vào tổng thu
    nhập (Income) cũng như tổng chi tiêu (Expense) của hệ thống; tổng số
    dư toàn bộ các ví trước và sau giao dịch bảo toàn tuyệt đối:
    \(\Delta \sum Balance = 0\).
  \end{itemize}
\item
  \textbf{Dữ liệu đầu ra:}

  \begin{itemize}
  \tightlist
  \item
    Bản ghi chuyển tiền nội bộ mới được lưu vào bảng \texttt{transfers}
    với mã định danh duy nhất (\texttt{transferId}).
  \item
    Số dư của hai ví liên quan được cập nhật đồng thời.
  \item
    Thông báo chuyển tiền thành công.
  \end{itemize}
\item
  \textbf{Ca sử dụng tham chiếu:}
  \texttt{UC03\ -\ Chuyển\ tiền\ nội\ bộ\ giữa\ các\ ví}.
\item
  \textbf{Tiêu chí nghiệm thu:}

  \begin{itemize}
  \tightlist
  \item
    TC-TRF-01.1: Cho ví "MB Bank" có 5,000,000 VNĐ và ví "Tiền mặt" có
    1,000,000 VNĐ; khi người dùng chuyển 2,000,000 VNĐ từ MB Bank sang
    Tiền mặt, số dư sau giao dịch của MB Bank là 3,000,000 VNĐ, Tiền mặt
    là 3,000,000 VNĐ, và tổng số dư hai ví vẫn giữ nguyên là 6,000,000
    VNĐ.
  \item
    TC-TRF-01.2: Khi người dùng chọn ví nguồn trùng với ví đích, hệ
    thống từ chối và thông báo lỗi ví nguồn và ví đích không được trùng
    nhau.
  \item
    TC-TRF-01.3: Khi số dư ví nguồn nhỏ hơn số tiền muốn chuyển, hệ
    thống từ chối giao dịch và đưa ra thông báo số dư không đủ.
  \end{itemize}
\end{itemize}

\paragraph{4.3.2 Yêu cầu FR-TRF-02: Ghi nhận phí chuyển tiền phát
sinh}\label{432-yuxeau-cux1ea7u-fr-trf-02-ghi-nhux1eadn-phuxed-chuyux1ec3n-tiux1ec1n-phuxe1t-sinh}

\begin{itemize}
\tightlist
\item
  \textbf{Mã yêu cầu:} \texttt{FR-TRF-02}
\item
  \textbf{Tên chức năng:} Ghi nhận phí chuyển tiền phát sinh
\item
  \textbf{Mô tả:} Tự động phát sinh và ghi nhận giao dịch chi tiêu độc
  lập cho khoản phí chuyển tiền do ngân hàng hoặc đơn vị trung gian thu,
  gắn liền với ví nguồn theo quy tắc BR-FIN-03.
\item
  \textbf{Dữ liệu đầu vào:}

  \begin{itemize}
  \tightlist
  \item
    Khoản phí chuyển tiền (\texttt{fee}): Kiểu \texttt{BigDecimal}, tùy
    chọn, mặc định bằng 0, giá trị phải lớn hơn hoặc bằng 0.
  \item
    Danh mục phí dịch vụ (\texttt{feeCategoryId}): Chuỗi ký tự, tùy chọn
    (nếu có phí, mặc định chọn danh mục "Phí ngân hàng / chuyển tiền"
    thuộc nhóm \texttt{obligations\ \&\ transfers}).
  \end{itemize}
\item
  \textbf{Quy tắc xử lý và ràng buộc:}

  \begin{itemize}
  \tightlist
  \item
    Nếu \(fee > 0\), hệ thống kiểm tra tổng năng lực chi trả của ví
    nguồn: \[Balance_{from} \ge amount + fee\] Nếu số dư ví nguồn nhỏ
    hơn tổng số tiền chuyển cộng phí, giao dịch bị từ chối toàn bộ.
  \item
    Khi điều kiện hợp lệ, hệ thống thực hiện hai bước đồng bộ:

    \begin{enumerate}
    \def\labelenumi{\arabic{enumi}.}
    \tightlist
    \item
      Ghi nhận giao dịch chuyển tiền số tiền \(amount\) từ ví nguồn sang
      ví đích như quy tắc \texttt{FR-TRF-01}.
    \item
      Tự động sinh một bản ghi giao dịch chi tiêu (\texttt{Expense}) độc
      lập với số tiền đúng bằng \(fee\), ví chi là ví nguồn
      \(fromWalletId\), danh mục là \(feeCategoryId\), và lưu trường
      tham chiếu \texttt{linkedTransferId\ =\ transferId}.
    \end{enumerate}
  \item
    Số tiền thực trừ khỏi ví nguồn là \(amount + fee\); số tiền thực
    nhận tại ví đích là \(amount\). Khoản phí \(fee\) được tính vào tổng
    chi tiêu của hệ thống trong kỳ.
  \end{itemize}
\item
  \textbf{Dữ liệu đầu ra:}

  \begin{itemize}
  \tightlist
  \item
    Bản ghi chuyển tiền nội bộ và bản ghi giao dịch chi tiêu phí liên
    kết.
  \item
    Số dư ví nguồn giảm \(amount + fee\), số dư ví đích tăng \(amount\).
  \item
    Thông báo hoàn thành chuyển tiền kèm thông tin chi tiết về khoản phí
    đã trừ.
  \end{itemize}
\item
  \textbf{Ca sử dụng tham chiếu:}
  \texttt{UC03\ -\ Chuyển\ tiền\ nội\ bộ\ giữa\ các\ ví}.
\item
  \textbf{Tiêu chí nghiệm thu:}

  \begin{itemize}
  \tightlist
  \item
    TC-TRF-02.1: Cho ví nguồn có 2,000,000 VNĐ, chuyển 1,000,000 VNĐ
    sang ví đích với phí chuyển tiền là 5,500 VNĐ; sau khi thực hiện, ví
    nguồn còn lại chính xác 994,500 VNĐ, ví đích tăng 1,000,000 VNĐ, và
    hệ thống xuất hiện một giao dịch chi tiêu phí ngân hàng 5,500 VNĐ.
  \item
    TC-TRF-02.2: Cho ví nguồn có đúng 1,000,000 VNĐ, chuyển 1,000,000
    VNĐ với phí 5,500 VNĐ; hệ thống từ chối do tổng số tiền cần trừ
    (1,005,500 VNĐ) vượt quá số dư hiện có.
  \end{itemize}
\end{itemize}

\paragraph{4.3.3 Yêu cầu FR-TRF-03: Lịch sử và hủy chuyển tiền nội
bộ}\label{433-yuxeau-cux1ea7u-fr-trf-03-lux1ecbch-sux1eed-vuxe0-hux1ee7y-chuyux1ec3n-tiux1ec1n-nux1ed9i-bux1ed9}

\begin{itemize}
\tightlist
\item
  \textbf{Mã yêu cầu:} \texttt{FR-TRF-03}
\item
  \textbf{Tên chức năng:} Lịch sử và hủy chuyển tiền nội bộ
\item
  \textbf{Mô tả:} Xem lại danh sách các giao dịch chuyển tiền nội bộ đã
  thực hiện và cho phép hủy bỏ giao dịch, tự động hoàn nguyên đồng thời
  số dư của cả hai ví đối ứng và khoản phí đi kèm nếu có.
\item
  \textbf{Dữ liệu đầu vào:}

  \begin{itemize}
  \tightlist
  \item
    Mã giao dịch chuyển tiền (\texttt{transferId}): Chuỗi ký tự, bắt
    buộc.
  \item
    Xác nhận hủy (\texttt{confirmed}): Kiểu \texttt{Boolean}, bắt buộc
    là \texttt{true}.
  \end{itemize}
\item
  \textbf{Quy tắc xử lý và ràng buộc:}

  \begin{itemize}
  \tightlist
  \item
    Tìm kiếm bản ghi chuyển tiền theo \texttt{transferId}. Nếu không tìm
    thấy hoặc giao dịch đã ở trạng thái hủy, báo lỗi.
  \item
    Kiểm tra điều kiện hoàn nguyên tại ví đích: Ví đích phải có số dư
    khả dụng tối thiểu bằng số tiền đã nhận từ giao dịch chuyển tiền
    (\(Balance_{to} - amount \ge 0\)). Nếu ví đích đã chi tiêu hết tiền
    dẫn đến số dư không đủ để trừ lại, hệ thống từ chối hủy và thông báo
    lỗi thiếu hụt số dư hoàn nguyên tại ví đích.
  \item
    Thực hiện hoàn nguyên đối ứng nguyên tử:

    \begin{itemize}
    \tightlist
    \item
      Cộng trả lại số tiền \(amount\) vào số dư của ví nguồn.
    \item
      Trừ số tiền \(amount\) khỏi số dư của ví đích.
    \item
      Nếu giao dịch có phát sinh phí liên kết, hoàn trả số tiền \(fee\)
      vào ví nguồn và xóa hoặc đánh dấu hủy bản ghi giao dịch chi tiêu
      phí tương ứng.
    \end{itemize}
  \item
    Chuyển trạng thái bản ghi chuyển tiền sang \texttt{CANCELLED} hoặc
    xóa khỏi sổ nhật ký.
  \end{itemize}
\item
  \textbf{Dữ liệu đầu ra:}

  \begin{itemize}
  \tightlist
  \item
    Trạng thái giao dịch chuyển tiền được cập nhật thành đã hủy.
  \item
    Số dư của hai ví đối ứng được hoàn nguyên chính xác về mức trước khi
    chuyển tiền.
  \item
    Thông báo hủy chuyển tiền thành công.
  \end{itemize}
\item
  \textbf{Ca sử dụng tham chiếu:}
  \texttt{UC03\ -\ Chuyển\ tiền\ nội\ bộ\ giữa\ các\ ví}.
\item
  \textbf{Tiêu chí nghiệm thu:}

  \begin{itemize}
  \tightlist
  \item
    TC-TRF-03.1: Cho giao dịch chuyển 2,000,000 VNĐ từ ví A sang ví B
    kèm phí 3,300 VNĐ; khi người dùng thực hiện hủy giao dịch thành
    công, ví A được cộng lại 2,003,300 VNĐ, ví B bị trừ đi 2,000,000
    VNĐ, và khoản phí chi tiêu bị xóa bỏ.
  \item
    TC-TRF-03.2: Cho ví B sau khi nhận 2,000,000 VNĐ đã chi tiêu hết
    tiền và chỉ còn lại 500,000 VNĐ; khi người dùng yêu cầu hủy giao
    dịch chuyển tiền, hệ thống từ chối thao tác và cảnh báo số dư ví B
    không đủ để thực hiện hoàn nguyên.
  \end{itemize}
\end{itemize}

\subsubsection{4.4 Phân hệ Danh mục phân loại
(FR-CAT)}\label{44-phuxe2n-hux1ec7-danh-mux1ee5c-phuxe2n-loux1ea1i-fr-cat}

Phân hệ Danh mục phân loại quản lý cấu trúc phân cấp cây 3 cấp chuẩn hóa
của hệ thống, phục vụ việc phân loại chi tiết các giao dịch thu chi và
làm nền tảng cho các báo cáo phân tích tài chính theo quy tắc BR-FIN-04.

\paragraph{4.4.1 Yêu cầu FR-CAT-01: Khởi tạo và quản lý cấu trúc phân
loại 3
cấp}\label{441-yuxeau-cux1ea7u-fr-cat-01-khux1edfi-tux1ea1o-vuxe0-quux1ea3n-luxfd-cux1ea5u-truxfac-phuxe2n-loux1ea1i-3-cux1ea5p}

\begin{itemize}
\tightlist
\item
  \textbf{Mã yêu cầu:} \texttt{FR-CAT-01}
\item
  \textbf{Tên chức năng:} Khởi tạo và quản lý cấu trúc phân loại 3 cấp
\item
  \textbf{Mô tả:} Thiết lập và duy trì cây phân loại tài chính chuẩn 3
  cấp gồm Cấp 1: Loại luồng tiền (Type), Cấp 2: Nhóm nhu cầu (Group),
  Cấp 3: Danh mục cụ thể (Category) theo quy tắc BR-FIN-04.
\item
  \textbf{Dữ liệu đầu vào:}

  \begin{itemize}
  \tightlist
  \item
    Dữ liệu khởi tạo hệ thống chuẩn từ bảng tham chiếu
    \texttt{database.classification}:

    \begin{itemize}
    \tightlist
    \item
      Cấp 1 (\texttt{Type}): Gồm 2 giá trị cố định \texttt{Income} và
      \texttt{Expense}.
    \item
      Cấp 2 (\texttt{Group}): 7 nhóm nhu cầu chuẩn gồm
      \texttt{employment}, \texttt{assets}, \texttt{other} (thuộc
      Income) và \texttt{essentials}, \texttt{personal\ \&\ leisure},
      \texttt{obligations\ \&\ transfers},
      \texttt{savings\ \&\ investing} (thuộc Expense).
    \item
      Cấp 3 (\texttt{Category}): Hơn 20 danh mục chi tiết mặc định gắn
      với từng nhóm tương ứng.
    \end{itemize}
  \end{itemize}
\item
  \textbf{Quy tắc xử lý và ràng buộc:}

  \begin{itemize}
  \tightlist
  \item
    Duy trì tính toàn vẹn quan hệ cha con: Mỗi Category bắt buộc phải
    trỏ tới duy nhất một Group cha hợp lệ; mỗi Group bắt buộc phải trỏ
    tới duy nhất một Type hợp lệ.
  \item
    Nghiêm cấm sự tồn tại của các danh mục mồ côi (Category không có
    Group hoặc Group không có Type).
  \item
    Tên của Group trong cùng một Type và tên của Category trong cùng một
    Group phải là duy nhất (không phân biệt hoa thường).
  \item
    Cung cấp phương thức duyệt cây danh mục để hiển thị dưới dạng
    dropdown phân cấp động khi người dùng thao tác nhập liệu giao dịch.
  \end{itemize}
\item
  \textbf{Dữ liệu đầu ra:}

  \begin{itemize}
  \tightlist
  \item
    Cây danh mục phân cấp 3 tầng hoàn chỉnh được nạp vào bộ nhớ hệ
    thống.
  \item
    Danh sách danh mục sẵn sàng cho việc lựa chọn và phân tích báo cáo.
  \end{itemize}
\item
  \textbf{Ca sử dụng tham chiếu:}
  \texttt{UC04\ -\ Quản\ lý\ danh\ mục\ phân\ loại\ thu\ chi}.
\item
  \textbf{Tiêu chí nghiệm thu:}

  \begin{itemize}
  \tightlist
  \item
    TC-CAT-01.1: Khi hệ thống khởi tạo lần đầu, cơ sở dữ liệu danh mục
    được điền đầy đủ 2 Type, 7 Group chuẩn và toàn bộ danh mục chi tiết
    khớp với bảng \texttt{database.classification}.
  \item
    TC-CAT-01.2: Khi người dùng chọn Type là \texttt{Expense}, hệ thống
    chỉ cung cấp danh sách 4 Group chi tiêu tương ứng
    (\texttt{essentials}, \texttt{personal\ \&\ leisure},
    \texttt{obligations\ \&\ transfers},
    \texttt{savings\ \&\ investing}).
  \end{itemize}
\end{itemize}

\paragraph{4.4.2 Yêu cầu FR-CAT-02: Tùy biến danh mục người
dùng}\label{442-yuxeau-cux1ea7u-fr-cat-02-tuxf9y-biux1ebfn-danh-mux1ee5c-ngux1b0ux1eddi-duxf9ng}

\begin{itemize}
\tightlist
\item
  \textbf{Mã yêu cầu:} \texttt{FR-CAT-02}
\item
  \textbf{Tên chức năng:} Tùy biến danh mục người dùng
\item
  \textbf{Mô tả:} Cho phép người dùng bổ sung các danh mục chi tiết mới
  theo nhu cầu cá nhân hoặc ẩn (vô hiệu hóa) các danh mục không còn sử
  dụng mà không làm gián đoạn hệ thống.
\item
  \textbf{Dữ liệu đầu vào:}

  \begin{itemize}
  \tightlist
  \item
    Tên danh mục mới (\texttt{name}): Chuỗi ký tự, từ 1 đến 50 ký tự,
    bắt buộc.
  \item
    Mã nhóm cha (\texttt{groupId}): Chuỗi ký tự, bắt buộc, nhóm cha phải
    tồn tại.
  \item
    Biểu tượng hoặc màu sắc nhận diện (\texttt{icon}, \texttt{color}):
    Chuỗi ký tự, tùy chọn.
  \item
    Trạng thái hoạt động (\texttt{isActive}): Kiểu \texttt{Boolean}, mặc
    định là \texttt{true}.
  \end{itemize}
\item
  \textbf{Quy tắc xử lý và ràng buộc:}

  \begin{itemize}
  \tightlist
  \item
    Kiểm tra tên danh mục mới không được trùng lặp với các danh mục khác
    trong cùng nhóm cha.
  \item
    Tạo mới bản ghi Category và liên kết chặt chẽ với Group cha đã chỉ
    định.
  \item
    Cho phép người dùng chuyển cờ \texttt{isActive} sang \texttt{false}
    để ẩn danh mục.
  \item
    Danh mục bị ẩn (\texttt{isActive\ =\ false}) sẽ tự động bị loại khỏi
    danh sách lựa chọn khi tạo giao dịch mới, nhưng vẫn tiếp tục xuất
    hiện trong các báo cáo lịch sử chứa giao dịch cũ liên quan đến danh
    mục đó.
  \end{itemize}
\item
  \textbf{Dữ liệu đầu ra:}

  \begin{itemize}
  \tightlist
  \item
    Danh mục mới được tạo và kích hoạt thành công kèm mã định danh
    (\texttt{categoryId}).
  \item
    Trạng thái hiển thị của danh mục được cập nhật.
  \item
    Thông báo kết quả thao tác hiển thị cho người dùng.
  \end{itemize}
\item
  \textbf{Ca sử dụng tham chiếu:}
  \texttt{UC04\ -\ Quản\ lý\ danh\ mục\ phân\ loại\ thu\ chi}.
\item
  \textbf{Tiêu chí nghiệm thu:}

  \begin{itemize}
  \tightlist
  \item
    TC-CAT-02.1: Khi người dùng thêm danh mục "Thú cưng" vào nhóm
    "personal \& leisure", danh mục mới được tạo thành công và xuất hiện
    trong danh sách lựa chọn khi nhập chi tiêu.
  \item
    TC-CAT-02.2: Khi người dùng đặt danh mục "Chơi gôn" sang trạng thái
    ẩn (\texttt{isActive\ =\ false}), danh mục này không còn hiển thị
    trong danh sách lựa chọn khi thêm giao dịch mới.
  \end{itemize}
\end{itemize}

\paragraph{4.4.3 Yêu cầu FR-CAT-03: Ràng buộc toàn vẹn khi sửa và xóa
danh
mục}\label{443-yuxeau-cux1ea7u-fr-cat-03-ruxe0ng-buux1ed9c-touxe0n-vux1eb9n-khi-sux1eeda-vuxe0-xuxf3a-danh-mux1ee5c}

\begin{itemize}
\tightlist
\item
  \textbf{Mã yêu cầu:} \texttt{FR-CAT-03}
\item
  \textbf{Tên chức năng:} Ràng buộc toàn vẹn khi sửa và xóa danh mục
\item
  \textbf{Mô tả:} Kiểm soát nghiêm ngặt hành vi chỉnh sửa hoặc xóa danh
  mục phân loại nhằm bảo vệ tính toàn vẹn lịch sử và tính bất biến của
  các báo cáo tài chính quá khứ theo quy tắc BR-FIN-04.
\item
  \textbf{Dữ liệu đầu vào:}

  \begin{itemize}
  \tightlist
  \item
    Mã danh mục hoặc mã nhóm cần thao tác (\texttt{categoryId} hoặc
    \texttt{groupId}): Chuỗi ký tự, bắt buộc.
  \item
    Hành động yêu cầu: Sửa tên, Chuyển nhóm cha, hoặc Xóa.
  \end{itemize}
\item
  \textbf{Quy tắc xử lý và ràng buộc:}

  \begin{itemize}
  \tightlist
  \item
    Ràng buộc đối với hành vi XÓA:

    \begin{itemize}
    \tightlist
    \item
      Hệ thống kiểm tra xem danh mục (hoặc nhóm) đã từng phát sinh bất
      kỳ giao dịch nào trong bảng nhật ký \texttt{transactions} hay
      chưa.
    \item
      Nếu ĐÃ PHÁT SINH GIAO DỊCH: Nghiêm cấm xóa vật lý bản ghi khỏi cơ
      sở dữ liệu. Hệ thống từ chối thao tác và hướng dẫn người dùng
      chuyển sang trạng thái ngưng hoạt động
      (\texttt{isActive\ =\ false}).
    \item
      Nếu CHƯA TỪNG PHÁT SINH GIAO DỊCH: Cho phép xóa vật lý bản ghi
      khỏi hệ thống.
    \item
      Không cho phép xóa một Group nếu vẫn còn chứa Category con trực
      thuộc Group đó.
    \end{itemize}
  \item
    Ràng buộc đối với hành vi SỬA:

    \begin{itemize}
    \tightlist
    \item
      Tuyệt đối cấm thay đổi thuộc tính \texttt{Type} (từ Income sang
      Expense hoặc ngược lại) của một Group hoặc Category nếu đã có giao
      dịch liên kết, nhằm tránh làm sai lệch toàn bộ cấu trúc báo cáo
      thu chi quá khứ.
    \item
      Cho phép đổi tên danh mục; việc đổi tên sẽ áp dụng cho các hiển
      thị mới mà không làm thay đổi mã định danh nội bộ.
    \end{itemize}
  \end{itemize}
\item
  \textbf{Dữ liệu đầu ra:}

  \begin{itemize}
  \tightlist
  \item
    Kết quả thực hiện thao tác xóa hoặc chỉnh sửa thành công.
  \item
    Thông báo từ chối thao tác nêu rõ lý do ràng buộc dữ liệu nếu vi
    phạm quy tắc toàn vẹn.
  \end{itemize}
\item
  \textbf{Ca sử dụng tham chiếu:}
  \texttt{UC04\ -\ Quản\ lý\ danh\ mục\ phân\ loại\ thu\ chi}.
\item
  \textbf{Tiêu chí nghiệm thu:}

  \begin{itemize}
  \tightlist
  \item
    TC-CAT-03.1: Cho danh mục "Ăn uống" đã gắn với 50 giao dịch chi tiêu
    trong quá khứ; khi người dùng yêu cầu xóa danh mục này, hệ thống từ
    chối xóa và hiển thị thông báo danh mục đã có dữ liệu giao dịch, chỉ
    được phép ẩn.
  \item
    TC-CAT-03.2: Cho danh mục "Học vẽ" vừa được tạo và chưa có giao dịch
    nào; khi người dùng chọn xóa, hệ thống xóa hoàn toàn bản ghi khỏi cơ
    sở dữ liệu.
  \item
    TC-CAT-03.3: Khi người dùng cố gắng chuyển một danh mục chi tiêu
    sang một nhóm thu nhập, hệ thống chặn lại và báo lỗi không tương
    thích loại luồng tiền (Type mismatch).
  \end{itemize}
\end{itemize}

\subsubsection{4.5 Phân hệ Quản lý công nợ
(FR-DEBT)}\label{45-phuxe2n-hux1ec7-quux1ea3n-luxfd-cuxf4ng-nux1ee3-fr-debt}

Phân hệ Quản lý công nợ giám sát chặt chẽ các quyền đòi nợ (khoản phải
thu) và nghĩa vụ tài chính (khoản phải trả) phát sinh giữa người dùng
với các cá nhân, tổ chức bên ngoài theo quy tắc BR-FIN-07.

\paragraph{4.5.1 Yêu cầu FR-DEBT-01: Ghi nhận khoản cho vay
(Lend)}\label{451-yuxeau-cux1ea7u-fr-debt-01-ghi-nhux1eadn-khoux1ea3n-cho-vay-lend}

\begin{itemize}
\tightlist
\item
  \textbf{Mã yêu cầu:} \texttt{FR-DEBT-01}
\item
  \textbf{Tên chức năng:} Ghi nhận khoản cho vay (Lend)
\item
  \textbf{Mô tả:} Ghi nhận một khoản tiền người dùng trích từ ví để cho
  đối tác hoặc bạn bè vay, làm giảm số dư tiền mặt của ví và làm tăng
  giá trị tài sản khoản phải thu (Receivable) theo quy tắc BR-FIN-07.
\item
  \textbf{Dữ liệu đầu vào:}

  \begin{itemize}
  \tightlist
  \item
    Ngày cho vay (\texttt{date}): Kiểu \texttt{LocalDate}, bắt buộc, mặc
    định hôm nay.
  \item
    Mã đợt vay (\texttt{batch}): Chuỗi ký tự, tùy chọn (dùng để gộp các
    đợt vay liên quan).
  \item
    Tên người vay (\texttt{person}): Chuỗi ký tự, từ 1 đến 100 ký tự,
    bắt buộc.
  \item
    Số tiền cho vay (\texttt{amount}): Kiểu \texttt{BigDecimal}, bắt
    buộc, lớn hơn 0.
  \item
    Ví xuất tiền (\texttt{walletId}): Chuỗi ký tự, bắt buộc, ví ở trạng
    thái \texttt{ACTIVE}.
  \item
    Ghi chú (\texttt{note}): Chuỗi ký tự, tối đa 255 ký tự, tùy chọn.
  \end{itemize}
\item
  \textbf{Quy tắc xử lý và ràng buộc:}

  \begin{itemize}
  \tightlist
  \item
    Kiểm tra số dư khả dụng của ví xuất tiền theo \texttt{BR-FIN-02}:
    \(Balance_{wallet} - amount \ge 0\). Nếu số dư không đủ, hệ thống từ
    chối giao dịch.
  \item
    Trừ số tiền \(amount\) khỏi số dư của ví xuất tiền.
  \item
    Tăng số dư nợ phải thu đối với cá nhân \texttt{person} thêm
    \(amount\):
    \[Receivable_{person}(t) = Receivable_{person}(t-1) + amount\]
  \item
    Cập nhật tổng tài sản khoản phải thu toàn hệ thống:
    \[
    Receivable = \sum Lend - \sum Collect - \sum WriteOffLend
    \]
  \item
    Giao dịch này đại diện cho sự biến chuyển hình thái tài sản (từ tiền
    mặt thanh khoản sang quyền đòi nợ), không tính vào chi phí tiêu dùng
    sinh hoạt (\texttt{Expense}) của kỳ.
  \end{itemize}
\item
  \textbf{Dữ liệu đầu ra:}

  \begin{itemize}
  \tightlist
  \item
    Bản ghi công nợ mới được lưu vào bảng \texttt{debt} với loại nghiệp
    vụ \texttt{lend} kèm mã định danh duy nhất (\texttt{debtId}).
  \item
    Số dư ví xuất tiền giảm đúng bằng số tiền cho vay.
  \item
    Số dư phải thu đối với người vay tăng tương ứng.
  \item
    Thông báo ghi nhận khoản cho vay thành công.
  \end{itemize}
\item
  \textbf{Ca sử dụng tham chiếu:}
  \texttt{UC05\ -\ Quản\ lý\ và\ theo\ dõi\ công\ nợ}.
\item
  \textbf{Tiêu chí nghiệm thu:}

  \begin{itemize}
  \tightlist
  \item
    TC-DEBT-01.1: Cho ví "MB Bank" có 10,000,000 VNĐ; khi người dùng ghi
    nhận cho "Nguyễn Văn A" vay 3,000,000 VNĐ từ ví MB Bank, số dư ví MB
    Bank còn lại 7,000,000 VNĐ và khoản phải thu đối với A là 3,000,000
    VNĐ.
  \item
    TC-DEBT-01.2: Cho ví có số dư 1,000,000 VNĐ; khi người dùng cho vay
    1,500,000 VNĐ, hệ thống từ chối do không đủ tiền trong ví.
  \end{itemize}
\end{itemize}

\paragraph{4.5.2 Yêu cầu FR-DEBT-02: Ghi nhận khoản đi vay
(Borrow)}\label{452-yuxeau-cux1ea7u-fr-debt-02-ghi-nhux1eadn-khoux1ea3n-ux111i-vay-borrow}

\begin{itemize}
\tightlist
\item
  \textbf{Mã yêu cầu:} \texttt{FR-DEBT-02}
\item
  \textbf{Tên chức năng:} Ghi nhận khoản đi vay (Borrow)
\item
  \textbf{Mô tả:} Ghi nhận một khoản tiền người dùng đi vay từ chủ nợ,
  nhận tiền vào một ví chỉ định và làm tăng nghĩa vụ nợ phải trả
  (Payable) theo quy tắc BR-FIN-07.
\item
  \textbf{Dữ liệu đầu vào:}

  \begin{itemize}
  \tightlist
  \item
    Ngày đi vay (\texttt{date}): Kiểu \texttt{LocalDate}, bắt buộc, mặc
    định hôm nay.
  \item
    Mã đợt vay (\texttt{batch}): Chuỗi ký tự, tùy chọn.
  \item
    Tên chủ nợ (\texttt{person}): Chuỗi ký tự, từ 1 đến 100 ký tự, bắt
    buộc.
  \item
    Số tiền đi vay (\texttt{amount}): Kiểu \texttt{BigDecimal}, bắt
    buộc, lớn hơn 0.
  \item
    Ví nhận tiền (\texttt{walletId}): Chuỗi ký tự, bắt buộc, ví ở trạng
    thái \texttt{ACTIVE}.
  \item
    Ghi chú (\texttt{note}): Chuỗi ký tự, tối đa 255 ký tự, tùy chọn.
  \end{itemize}
\item
  \textbf{Quy tắc xử lý và ràng buộc:}

  \begin{itemize}
  \tightlist
  \item
    Kiểm tra tính hợp lệ của số tiền đi vay: \(amount > 0\).
  \item
    Cộng số tiền \(amount\) vào số dư khả dụng của ví nhận tiền theo
    \texttt{BR-FIN-02}.
  \item
    Tăng nghĩa vụ nợ phải trả đối với chủ nợ \texttt{person} thêm
    \(amount\): \[Payable_{person}(t) = Payable_{person}(t-1) + amount\]
  \item
    Cập nhật tổng nợ phải trả toàn hệ thống:
    \[
    Payable = \sum Borrow - \sum Repay - \sum WriteOffBorrow
    \]
  \item
    Giao dịch này đại diện cho sự gia tăng tiền mặt đi kèm với nghĩa vụ
    nợ tương ứng, không được ghi nhận vào doanh thu thu nhập
    (\texttt{Income}) sinh hoạt trong kỳ.
  \end{itemize}
\item
  \textbf{Dữ liệu đầu ra:}

  \begin{itemize}
  \tightlist
  \item
    Bản ghi công nợ mới với loại nghiệp vụ \texttt{borrow} được lưu vào
    bảng \texttt{debt}.
  \item
    Số dư ví nhận tiền tăng thêm \(amount\).
  \item
    Khoản phải trả cho chủ nợ tăng tương ứng.
  \item
    Thông báo ghi nhận khoản đi vay thành công.
  \end{itemize}
\item
  \textbf{Ca sử dụng tham chiếu:}
  \texttt{UC05\ -\ Quản\ lý\ và\ theo\ dõi\ công\ nợ}.
\item
  \textbf{Tiêu chí nghiệm thu:}

  \begin{itemize}
  \tightlist
  \item
    TC-DEBT-02.1: Cho ví "Tiền mặt" có 1,000,000 VNĐ; khi người dùng ghi
    nhận vay của "Trần Thị B" 5,000,000 VNĐ nhận vào ví Tiền mặt, số dư
    ví Tiền mặt tăng lên 6,000,000 VNĐ và khoản phải trả đối với B là
    5,000,000 VNĐ.
  \item
    TC-DEBT-02.2: Khi người dùng nhập số tiền đi vay bằng 0 hoặc bỏ
    trống tên chủ nợ, hệ thống từ chối và báo lỗi dữ liệu đầu vào.
  \end{itemize}
\end{itemize}

\paragraph{4.5.3 Yêu cầu FR-DEBT-03: Thu hồi nợ cho vay
(Collect)}\label{453-yuxeau-cux1ea7u-fr-debt-03-thu-hux1ed3i-nux1ee3-cho-vay-collect}

\begin{itemize}
\tightlist
\item
  \textbf{Mã yêu cầu:} \texttt{FR-DEBT-03}
\item
  \textbf{Tên chức năng:} Thu hồi nợ cho vay (Collect)
\item
  \textbf{Mô tả:} Ghi nhận việc người vay hoàn trả tiền nợ cho người
  dùng, làm tăng số dư ví nhận tiền và giảm trừ khoản phải thu đối ứng
  theo quy tắc BR-FIN-07.
\item
  \textbf{Dữ liệu đầu vào:}

  \begin{itemize}
  \tightlist
  \item
    Ngày thu nợ (\texttt{date}): Kiểu \texttt{LocalDate}, bắt buộc.
  \item
    Tên người trả nợ (\texttt{person}): Chuỗi ký tự, bắt buộc.
  \item
    Số tiền thu hồi (\texttt{amount}): Kiểu \texttt{BigDecimal}, bắt
    buộc, lớn hơn 0.
  \item
    Ví nhận tiền (\texttt{walletId}): Chuỗi ký tự, bắt buộc, ví ở trạng
    thái \texttt{ACTIVE}.
  \item
    Ghi chú (\texttt{note}): Chuỗi ký tự, tối đa 255 ký tự, tùy chọn.
  \end{itemize}
\item
  \textbf{Quy tắc xử lý và ràng buộc:}

  \begin{itemize}
  \tightlist
  \item
    Tính toán số dư nợ phải thu hiện tại của cá nhân \texttt{person}:
    \[Receivable_{person} = \sum Lend_{person} - \sum Collect_{person} - \sum WriteOffLend_{person}\]
  \item
    Kiểm tra ràng buộc trần thu nợ theo quy tắc \texttt{BR-FIN-07}: Số
    tiền thu nợ không được vượt quá số dư nợ phải thu hiện tại
    (\(amount \le Receivable_{person}\)). Nếu
    \(amount > Receivable_{person}\), hệ thống chặn lại và báo lỗi thu
    vượt quá dư nợ gốc.
  \item
    Cộng số tiền \(amount\) vào số dư khả dụng của ví nhận tiền theo
    \texttt{BR-FIN-02}.
  \item
    Giảm số dư nợ phải thu của người vay:
    \[
    Receivable_{person}(t) = Receivable_{person}(t-1) - amount
    \]
  \end{itemize}
\item
  \textbf{Dữ liệu đầu ra:}

  \begin{itemize}
  \tightlist
  \item
    Bản ghi công nợ với loại nghiệp vụ \texttt{collect} được lưu vào cơ
    sở dữ liệu.
  \item
    Số dư ví nhận tiền tăng \(amount\).
  \item
    Khoản nợ phải thu của đối tượng giảm tương ứng.
  \item
    Thông báo thu hồi nợ thành công kèm số dư nợ còn lại của đối tượng.
  \end{itemize}
\item
  \textbf{Ca sử dụng tham chiếu:}
  \texttt{UC05\ -\ Quản\ lý\ và\ theo\ dõi\ công\ nợ}.
\item
  \textbf{Tiêu chí nghiệm thu:}

  \begin{itemize}
  \tightlist
  \item
    TC-DEBT-03.1: Cho đối tác A đang nợ 3,000,000 VNĐ; khi người dùng
    thu hồi 1,500,000 VNĐ vào ví MB Bank, số dư ví MB Bank tăng
    1,500,000 VNĐ và nợ còn lại của A giảm xuống đúng 1,500,000 VNĐ.
  \item
    TC-DEBT-03.2: Cho đối tác A đang nợ 1,000,000 VNĐ; khi người dùng
    nhập số tiền thu nợ 1,200,000 VNĐ, hệ thống từ chối thao tác và đưa
    ra cảnh báo số tiền thu nợ vượt quá số nợ gốc hiện tại.
  \end{itemize}
\end{itemize}

\paragraph{4.5.4 Yêu cầu FR-DEBT-04: Trả nợ đi vay
(Repay)}\label{454-yuxeau-cux1ea7u-fr-debt-04-trux1ea3-nux1ee3-ux111i-vay-repay}

\begin{itemize}
\tightlist
\item
  \textbf{Mã yêu cầu:} \texttt{FR-DEBT-04}
\item
  \textbf{Tên chức năng:} Trả nợ đi vay (Repay)
\item
  \textbf{Mô tả:} Ghi nhận việc người dùng hoàn trả một phần hoặc toàn
  bộ số nợ đã vay cho chủ nợ, trích tiền từ ví thanh toán và giảm trừ
  nghĩa vụ nợ phải trả theo quy tắc BR-FIN-07.
\item
  \textbf{Dữ liệu đầu vào:}

  \begin{itemize}
  \tightlist
  \item
    Ngày trả nợ (\texttt{date}): Kiểu \texttt{LocalDate}, bắt buộc.
  \item
    Tên chủ nợ (\texttt{person}): Chuỗi ký tự, bắt buộc.
  \item
    Số tiền hoàn trả (\texttt{amount}): Kiểu \texttt{BigDecimal}, bắt
    buộc, lớn hơn 0.
  \item
    Ví trích tiền (\texttt{walletId}): Chuỗi ký tự, bắt buộc, ví ở trạng
    thái \texttt{ACTIVE}.
  \item
    Ghi chú (\texttt{note}): Chuỗi ký tự, tối đa 255 ký tự, tùy chọn.
  \end{itemize}
\item
  \textbf{Quy tắc xử lý và ràng buộc:}

  \begin{itemize}
  \tightlist
  \item
    Tính toán nghĩa vụ nợ phải trả hiện tại đối với chủ nợ
    \texttt{person}:
    \[Payable_{person} = \sum Borrow_{person} - \sum Repay_{person} - \sum WriteOffBorrow_{person}\]
  \item
    Kiểm tra ràng buộc trần trả nợ theo \texttt{BR-FIN-07}: Số tiền trả
    không được vượt quá nghĩa vụ nợ hiện tại
    (\(amount \le Payable_{person}\)).
  \item
    Kiểm tra số dư ví trích tiền theo \texttt{BR-FIN-02}:
    \(Balance_{wallet} - amount \ge 0\). Nếu số dư không đủ, từ chối
    giao dịch.
  \item
    Trừ số tiền \(amount\) khỏi ví trích tiền và giảm nghĩa vụ nợ phải
    trả tương ứng:
    \[Payable_{person}(t) = Payable_{person}(t-1) - amount\]
  \end{itemize}
\item
  \textbf{Dữ liệu đầu ra:}

  \begin{itemize}
  \tightlist
  \item
    Bản ghi công nợ với loại nghiệp vụ \texttt{repay} được ghi sổ thành
    công.
  \item
    Số dư ví trích tiền giảm \(amount\).
  \item
    Nghĩa vụ nợ đối với chủ nợ giảm tương ứng.
  \item
    Thông báo trả nợ thành công kèm dư nợ còn lại.
  \end{itemize}
\item
  \textbf{Ca sử dụng tham chiếu:}
  \texttt{UC05\ -\ Quản\ lý\ và\ theo\ dõi\ công\ nợ}.
\item
  \textbf{Tiêu chí nghiệm thu:}

  \begin{itemize}
  \tightlist
  \item
    TC-DEBT-04.1: Cho khoản nợ chủ nợ B là 5,000,000 VNĐ và ví Timo có
    4,000,000 VNĐ; khi người dùng trả nợ cho B 2,000,000 VNĐ từ ví Timo,
    ví Timo còn lại 2,000,000 VNĐ và nợ còn lại đối với B là 3,000,000
    VNĐ.
  \item
    TC-DEBT-04.2: Khi người dùng trả nợ vượt quá số dư ví thanh toán
    hoặc vượt quá dư nợ thực tế với chủ nợ B, hệ thống từ chối và hiển
    thị thông báo lỗi cụ thể.
  \end{itemize}
\end{itemize}

\paragraph{4.5.5 Yêu cầu FR-DEBT-05: Xóa nợ khó đòi và miễn trừ nợ
(Write-off)}\label{455-yuxeau-cux1ea7u-fr-debt-05-xuxf3a-nux1ee3-khuxf3-ux111uxf2i-vuxe0-miux1ec5n-trux1eeb-nux1ee3-write-off}

\begin{itemize}
\tightlist
\item
  \textbf{Mã yêu cầu:} \texttt{FR-DEBT-05}
\item
  \textbf{Tên chức năng:} Xóa nợ khó đòi và miễn trừ nợ (Write-off)
\item
  \textbf{Mô tả:} Thực hiện xử lý kế toán đối với các khoản cho vay
  không thể thu hồi được (xóa nợ xấu) hoặc các khoản nợ đi vay được chủ
  nợ tha thứ (miễn trừ nghĩa vụ), hoàn toàn không làm biến động số dư
  các ví tiền mặt theo quy tắc BR-FIN-07.
\item
  \textbf{Dữ liệu đầu vào:}

  \begin{itemize}
  \tightlist
  \item
    Ngày xóa nợ (\texttt{date}): Kiểu \texttt{LocalDate}, bắt buộc.
  \item
    Đối tượng nợ (\texttt{person}): Chuỗi ký tự, bắt buộc.
  \item
    Loại xóa nợ (\texttt{type}): Phải là \texttt{write-off\ lend} (xóa
    nợ cho vay) hoặc \texttt{write-off\ borrow} (được xóa nợ đi vay).
  \item
    Số tiền xóa nợ (\texttt{amount}): Kiểu \texttt{BigDecimal}, bắt
    buộc, lớn hơn 0.
  \item
    Lý do xóa nợ (\texttt{note}): Chuỗi ký tự, từ 1 đến 255 ký tự, bắt
    buộc nhập lý do.
  \end{itemize}
\item
  \textbf{Quy tắc xử lý và ràng buộc:}

  \begin{itemize}
  \tightlist
  \item
    Nếu là \texttt{write-off\ lend}:

    \begin{itemize}
    \tightlist
    \item
      Kiểm tra số tiền xóa nợ không được vượt quá số dư phải thu còn lại
      của đối tượng: \(amount \le Receivable_{person}\).
    \item
      Giảm khoản phải thu \(Receivable_{person}\) và tổng \(Receivable\)
      của hệ thống.
    \item
      Tuyệt đối không thay đổi số dư của bất kỳ ví tiền mặt nào theo
      \texttt{BR-FIN-02} (nghiệp vụ phi tiền mặt).
    \item
      Ghi nhận khoản tổn thất tài sản trong báo cáo tổng quan.
    \end{itemize}
  \item
    Nếu là \texttt{write-off\ borrow}:

    \begin{itemize}
    \tightlist
    \item
      Kiểm tra số tiền xóa nợ không được vượt quá số dư phải trả còn
      lại: \(amount \le Payable_{person}\).
    \item
      Giảm khoản phải trả \(Payable_{person}\) và tổng \(Payable\) của
      hệ thống.
    \item
      Tuyệt đối không xuất tiền từ ví tiền mặt nào (nghiệp vụ phi tiền
      mặt).
    \item
      Ghi nhận khoản thu nhập bất thường do được xóa nợ.
    \end{itemize}
  \end{itemize}
\item
  \textbf{Dữ liệu đầu ra:}

  \begin{itemize}
  \tightlist
  \item
    Bản ghi công nợ với loại \texttt{write-off\ lend} hoặc
    \texttt{write-off\ borrow} được lưu trữ.
  \item
    Dư nợ công nợ của đối tượng giảm tương ứng.
  \item
    Số dư toàn bộ các ví tiền giữ nguyên tuyệt đối.
  \item
    Thông báo hoàn tất xử lý xóa nợ.
  \end{itemize}
\item
  \textbf{Ca sử dụng tham chiếu:}
  \texttt{UC05\ -\ Quản\ lý\ và\ theo\ dõi\ công\ nợ}.
\item
  \textbf{Tiêu chí nghiệm thu:}

  \begin{itemize}
  \tightlist
  \item
    TC-DEBT-05.1: Cho đối tác A còn nợ 1,000,000 VNĐ nhưng mất khả năng
    thanh toán; khi người dùng thực hiện \texttt{write-off\ lend}
    1,000,000 VNĐ, nợ phải thu của A về 0, tổng Receivable giảm
    1,000,000 VNĐ, và số dư toàn bộ các ví tiền không có bất kỳ sự thay
    đổi nào.
  \item
    TC-DEBT-05.2: Khi người dùng thực hiện xóa nợ với số tiền lớn hơn số
    dư nợ còn lại của đối tượng, hệ thống từ chối và thông báo số tiền
    xóa nợ vượt quá dư nợ thực tế.
  \end{itemize}
\end{itemize}

\subsubsection{4.6 Phân hệ Quản lý sổ tiết kiệm
(FR-SAV)}\label{46-phuxe2n-hux1ec7-quux1ea3n-luxfd-sux1ed5-tiux1ebft-kiux1ec7m-fr-sav}

Phân hệ Quản lý sổ tiết kiệm theo dõi các khoản tiền gửi tiết kiệm có kỳ
hạn và không kỳ hạn tại các ngân hàng thương mại, tự động tính toán tiền
lãi dự kiến và xử lý tất toán theo đúng quy chế ngân hàng tại quy tắc
BR-FIN-05.

\paragraph{4.6.1 Yêu cầu FR-SAV-01: Mở sổ tiết kiệm
mới}\label{461-yuxeau-cux1ea7u-fr-sav-01-mux1edf-sux1ed5-tiux1ebft-kiux1ec7m-mux1edbi}

\begin{itemize}
\tightlist
\item
  \textbf{Mã yêu cầu:} \texttt{FR-SAV-01}
\item
  \textbf{Tên chức năng:} Mở sổ tiết kiệm mới
\item
  \textbf{Mô tả:} Khởi tạo một sổ tiền gửi tiết kiệm mới, trích xuất vốn
  gốc từ một ví tiền thanh khoản chỉ định và liên kết sổ với một quỹ tài
  chính mục tiêu theo quy tắc BR-FIN-05 và BR-FIN-08.
\item
  \textbf{Dữ liệu đầu vào:}

  \begin{itemize}
  \tightlist
  \item
    Tên ngân hàng hoặc tổ chức tín dụng (\texttt{bank}): Chuỗi ký tự, từ
    1 đến 100 ký tự, bắt buộc.
  \item
    Số vốn gốc gửi (\texttt{principal}): Kiểu \texttt{BigDecimal}, bắt
    buộc, lớn hơn 0.
  \item
    Kỳ hạn gửi (\texttt{term}): Chuỗi ký tự chuẩn hóa (ví dụ:
    \texttt{unlimited}, \texttt{1M}, \texttt{3M}, \texttt{6M},
    \texttt{12M}, \texttt{24M}, \texttt{36M}, \texttt{60M}), bắt buộc.
  \item
    Lãi suất cam kết (\texttt{rate}): Kiểu \texttt{BigDecimal}, đơn vị
    \%/năm, bắt buộc, lớn hơn 0.
  \item
    Ngày bắt đầu gửi (\texttt{start}): Kiểu \texttt{LocalDate}, bắt
    buộc, mặc định hôm nay.
  \item
    Ví trích tiền (\texttt{walletId}): Chuỗi ký tự, bắt buộc, ví ở trạng
    thái \texttt{ACTIVE}.
  \item
    Quỹ mục tiêu liên kết (\texttt{fund}): Chuỗi ký tự chuẩn hóa
    (\texttt{emergency}, \texttt{health}, \texttt{target},
    \texttt{retirement} hoặc tên quỹ hợp lệ), tùy chọn.
  \end{itemize}
\item
  \textbf{Quy tắc xử lý và ràng buộc:}

  \begin{itemize}
  \tightlist
  \item
    Kiểm tra số dư khả dụng của ví trích tiền theo \texttt{BR-FIN-02}:
    \(Balance_{wallet} - principal \ge 0\). Nếu không đủ tiền, từ chối
    mở sổ.
  \item
    Trừ số tiền vốn gốc \(principal\) khỏi số dư ví trích tiền.
  \item
    Tra cứu số ngày kỳ hạn (\(TermDays\)) từ bảng chuẩn hóa
    \texttt{database.term} dựa trên mã \texttt{term}:

    \begin{itemize}
    \tightlist
    \item
      \texttt{unlimited}: 0 ngày.
    \item
      \texttt{1M}: 30 ngày; \texttt{3M}: 90 ngày; \texttt{6M}: 180 ngày;
      \texttt{12M}: 365 ngày; \texttt{24M}: 730 ngày; \texttt{36M}: 1095
      ngày; \texttt{60M}: 1825 ngày.
    \end{itemize}
  \item
    Xác định ngày đáo hạn: \(MaturityDate = start + TermDays\) (nếu có
    kỳ hạn).
  \item
    Khởi tạo sổ tiết kiệm với trạng thái \texttt{OPEN}, ghi nhận số dư
    ban đầu \(Balance = principal\).
  \item
    Nếu có gắn mã quỹ mục tiêu \texttt{fund}, ghi nhận liên kết tài sản
    này vào quỹ theo \texttt{BR-FIN-08}.
  \end{itemize}
\item
  \textbf{Dữ liệu đầu ra:}

  \begin{itemize}
  \tightlist
  \item
    Bản ghi sổ tiết kiệm mới được lưu vào bảng \texttt{savings} kèm mã
    định danh duy nhất (\texttt{savingId}).
  \item
    Số dư ví trích tiền giảm đúng bằng số vốn gốc.
  \item
    Thông báo mở sổ tiết kiệm thành công.
  \end{itemize}
\item
  \textbf{Ca sử dụng tham chiếu:}
  \texttt{UC06\ -\ Quản\ lý\ sổ\ tiết\ kiệm\ có\ kỳ\ hạn}.
\item
  \textbf{Tiêu chí nghiệm thu:}

  \begin{itemize}
  \tightlist
  \item
    TC-SAV-01.1: Cho ví "Vietcombank" có 50,000,000 VNĐ; khi người dùng
    mở sổ tiết kiệm 30,000,000 VNĐ kỳ hạn 6 tháng (180 ngày) lãi suất
    6.0\%/năm tại Vietcombank, số dư ví VCB giảm còn 20,000,000 VNĐ và
    sổ tiết kiệm mới được tạo với vốn gốc 30,000,000 VNĐ.
  \item
    TC-SAV-01.2: Khi số dư ví trích tiền nhỏ hơn số tiền vốn gốc mở sổ,
    hệ thống từ chối mở sổ và thông báo số dư không đủ.
  \end{itemize}
\end{itemize}

\paragraph{4.6.2 Yêu cầu FR-SAV-02: Tính toán tiền lãi dự kiến và theo
dõi kỳ
hạn}\label{462-yuxeau-cux1ea7u-fr-sav-02-tuxednh-touxe1n-tiux1ec1n-luxe3i-dux1ef1-kiux1ebfn-vuxe0-theo-duxf5i-kux1ef3-hux1ea1n}

\begin{itemize}
\tightlist
\item
  \textbf{Mã yêu cầu:} \texttt{FR-SAV-02}
\item
  \textbf{Tên chức năng:} Tính toán tiền lãi dự kiến và theo dõi kỳ hạn
\item
  \textbf{Mô tả:} Tự động tính toán số ngày đã gửi thực tế, tiến độ thời
  gian kỳ hạn, tiền lãi dự kiến toàn kỳ và giá trị sổ tích lũy tại thời
  điểm quan sát theo quy tắc BR-FIN-05.
\item
  \textbf{Dữ liệu đầu vào:}

  \begin{itemize}
  \tightlist
  \item
    Mã sổ tiết kiệm (\texttt{savingId}): Chuỗi ký tự, bắt buộc.
  \item
    Ngày quan sát (\texttt{currentDate}): Kiểu \texttt{LocalDate}, mặc
    định là ngày hiện tại.
  \end{itemize}
\item
  \textbf{Quy tắc xử lý và ràng buộc:}

  \begin{itemize}
  \tightlist
  \item
    Tính số ngày gửi thực tế: \(Days = \max(0, currentDate - start)\).
  \item
    Trường hợp gửi không kỳ hạn (\texttt{term\ =\ unlimited} hoặc
    \(TermDays = 0\)):

    \begin{itemize}
    \tightlist
    \item
      Tiến độ quy ước: \(Process = 1.0\) (100\%).
    \item
      Tiền lãi tích lũy:
      \(Interest = \frac{Principal \times Rate \times Days}{365}\).
    \item
      Giá trị sổ hiện tại: \(Balance = Principal + Interest\).
    \end{itemize}
  \item
    Trường hợp gửi có kỳ hạn (\(TermDays > 0\)):

    \begin{itemize}
    \tightlist
    \item
      Tiến độ kỳ hạn:
      \(Process = \min\left(1.0, \frac{Days}{TermDays}\right)\).
    \item
      Lãi dự kiến toàn kỳ:
      \(ExpectedInterest = \frac{Principal \times Rate \times TermDays}{365}\).
    \item
      Giá trị sổ ghi nhận tại thời điểm quan sát:
      \[Balance = Principal + (ExpectedInterest \times Process)\]
    \item
      Số ngày còn lại đến đáo hạn:
      \(RemainingDays = \max(0, TermDays - Days)\).
    \end{itemize}
  \item
    Sử dụng \texttt{BigDecimal} với scale trung gian tối thiểu là 4 chữ
    số thập phân và chỉ định \texttt{RoundingMode.HALF\_EVEN} cho các
    phép chia.
  \end{itemize}
\item
  \textbf{Dữ liệu đầu ra:}

  \begin{itemize}
  \tightlist
  \item
    Thông tin chi tiết sổ tiết kiệm gồm: Vốn gốc, Lãi suất cam kết, Ngày
    gửi, Ngày đáo hạn, Số ngày đã gửi, Số ngày còn lại, Tiến độ hoàn
    thành kỳ hạn (\%), Lãi dự kiến đến hạn, Giá trị ghi nhận hiện tại.
  \end{itemize}
\item
  \textbf{Ca sử dụng tham chiếu:}
  \texttt{UC06\ -\ Quản\ lý\ sổ\ tiết\ kiệm\ có\ kỳ\ hạn}.
\item
  \textbf{Tiêu chí nghiệm thu:}

  \begin{itemize}
  \tightlist
  \item
    TC-SAV-02.1: Cho sổ tiết kiệm vốn gốc 100,000,000 VNĐ, kỳ hạn 12
    tháng (365 ngày), lãi suất 6.0\%/năm; tại thời điểm đã gửi được 180
    ngày, hệ thống tính lãi dự kiến cả kỳ là đúng 6,000,000 VNĐ, tiến độ
    kỳ hạn đạt 49.32\%, và giá trị ghi nhận là 102,958,904 VNĐ.
  \item
    TC-SAV-02.2: Khi số ngày gửi thực tế lớn hơn hoặc bằng số ngày kỳ
    hạn (\(Days \ge TermDays\)), tiến độ kỳ hạn hiển thị 100.00\% và giá
    trị ghi nhận là toàn bộ Gốc + Lãi toàn kỳ (106,000,000 VNĐ).
  \end{itemize}
\end{itemize}

\paragraph{4.6.3 Yêu cầu FR-SAV-03: Tất toán sổ tiết kiệm đúng
hạn}\label{463-yuxeau-cux1ea7u-fr-sav-03-tux1ea5t-touxe1n-sux1ed5-tiux1ebft-kiux1ec7m-ux111uxfang-hux1ea1n}

\begin{itemize}
\tightlist
\item
  \textbf{Mã yêu cầu:} \texttt{FR-SAV-03}
\item
  \textbf{Tên chức năng:} Tất toán sổ tiết kiệm đúng hạn
\item
  \textbf{Mô tả:} Cho phép người dùng đóng sổ tiết kiệm khi đã đến hoặc
  quá ngày đáo hạn, nhận trọn vẹn toàn bộ vốn gốc và 100\% tiền lãi cam
  kết chuyển về ví thanh khoản chỉ định theo quy tắc BR-FIN-05.
\item
  \textbf{Dữ liệu đầu vào:}

  \begin{itemize}
  \tightlist
  \item
    Mã sổ tiết kiệm (\texttt{savingId}): Chuỗi ký tự, bắt buộc.
  \item
    Ngày tất toán (\texttt{settleDate}): Kiểu \texttt{LocalDate}, bắt
    buộc, điều kiện \(settleDate \ge MaturityDate\).
  \item
    Ví nhận tiền (\texttt{walletId}): Chuỗi ký tự, bắt buộc, ví ở trạng
    thái \texttt{ACTIVE}.
  \end{itemize}
\item
  \textbf{Quy tắc xử lý và ràng buộc:}

  \begin{itemize}
  \tightlist
  \item
    Kiểm tra sổ tiết kiệm đang ở trạng thái \texttt{OPEN}.
  \item
    Xác thực điều kiện tất toán đúng hạn: Ngày tất toán phải bằng hoặc
    sau ngày đáo hạn (\(settleDate \ge MaturityDate\)). Nếu chưa đến
    ngày đáo hạn, hệ thống chặn thao tác tất toán đúng hạn và thông báo
    sổ chưa đến hạn.
  \item
    Tính tổng số tiền thực nhận (Payout):
    \[Payout = Principal + ExpectedInterest = Principal + \frac{Principal \times Rate \times TermDays}{365}\]
  \item
    Cộng số tiền \(Payout\) vào số dư của ví nhận tiền theo
    \texttt{BR-FIN-02}.
  \item
    Cập nhật trạng thái sổ tiết kiệm sang \texttt{SETTLED\_MATURED} (đã
    tất toán đúng hạn), đóng sổ và ngừng theo dõi tính lãi.
  \item
    Nếu sổ có gắn với quỹ mục tiêu, ngắt liên kết tài sản và khấu trừ
    giá trị sổ khỏi quỹ mục tiêu đó theo \texttt{BR-FIN-08}.
  \end{itemize}
\item
  \textbf{Dữ liệu đầu ra:}

  \begin{itemize}
  \tightlist
  \item
    Số dư ví nhận tiền tăng thêm đúng bằng số tiền Payout (Gốc + Lãi cam
    kết).
  \item
    Trạng thái sổ tiết kiệm chuyển thành \texttt{SETTLED\_MATURED}.
  \item
    Thông báo tất toán đúng hạn thành công kèm chi tiết số tiền gốc và
    lãi đã nhận.
  \end{itemize}
\item
  \textbf{Ca sử dụng tham chiếu:}
  \texttt{UC06\ -\ Quản\ lý\ sổ\ tiết\ kiệm\ có\ kỳ\ hạn}.
\item
  \textbf{Tiêu chí nghiệm thu:}

  \begin{itemize}
  \tightlist
  \item
    TC-SAV-03.1: Cho sổ tiết kiệm gốc 50,000,000 VNĐ, kỳ hạn 6 tháng,
    lãi cam kết 1,500,000 VNĐ đã đến hạn; khi người dùng tất toán đúng
    hạn vào ví MB Bank, số dư ví MB Bank tăng thêm đúng 51,500,000 VNĐ
    và sổ chuyển sang trạng thái đã đóng.
  \item
    TC-SAV-03.2: Khi người dùng cố gắng chọn chức năng tất toán đúng hạn
    vào ngày trước ngày đáo hạn, hệ thống từ chối và thông báo sổ chưa
    đến ngày đáo hạn cam kết.
  \end{itemize}
\end{itemize}

\paragraph{4.6.4 Yêu cầu FR-SAV-04: Tất toán sổ tiết kiệm trước
hạn}\label{464-yuxeau-cux1ea7u-fr-sav-04-tux1ea5t-touxe1n-sux1ed5-tiux1ebft-kiux1ec7m-trux1b0ux1edbc-hux1ea1n}

\begin{itemize}
\tightlist
\item
  \textbf{Mã yêu cầu:} \texttt{FR-SAV-04}
\item
  \textbf{Tên chức năng:} Tất toán sổ tiết kiệm trước hạn
\item
  \textbf{Mô tả:} Xử lý tất toán rút tiền trước khi sổ tiết kiệm đến
  ngày đáo hạn, áp dụng cơ chế hồi tố lãi suất không kỳ hạn theo quy
  định ngân hàng và chuyển số tiền thực nhận về ví thanh khoản theo quy
  tắc BR-FIN-05.
\item
  \textbf{Dữ liệu đầu vào:}

  \begin{itemize}
  \tightlist
  \item
    Mã sổ tiết kiệm (\texttt{savingId}): Chuỗi ký tự, bắt buộc.
  \item
    Ngày tất toán trước hạn (\texttt{settleDate}): Kiểu
    \texttt{LocalDate}, bắt buộc, \(settleDate < MaturityDate\).
  \item
    Lãi suất không kỳ hạn hồi tố (\texttt{demandRate}): Kiểu
    \texttt{BigDecimal}, đơn vị \%/năm, bắt buộc, lớn hơn hoặc bằng 0
    (thường từ 0.1\% đến 0.5\%/năm).
  \item
    Ví nhận tiền (\texttt{walletId}): Chuỗi ký tự, bắt buộc, ví ở trạng
    thái \texttt{ACTIVE}.
  \item
    Xác nhận chấp nhận mất lãi kỳ hạn (\texttt{confirmed}): Kiểu
    \texttt{Boolean}, bắt buộc là \texttt{true}.
  \end{itemize}
\item
  \textbf{Quy tắc xử lý và ràng buộc:}

  \begin{itemize}
  \tightlist
  \item
    Kiểm tra sổ tiết kiệm đang ở trạng thái \texttt{OPEN}.
  \item
    Kiểm tra điều kiện trước hạn:
    \(Days = settleDate - start < TermDays\).
  \item
    Hiển thị bảng so sánh tài chính trước khi thực thi:

    \begin{itemize}
    \tightlist
    \item
      Tiền lãi cam kết dự kiến bị mất: \(ExpectedInterest\).
    \item
      Tiền lãi thực tế được nhận theo lãi suất không kỳ hạn:
      \[ActualInterest = \frac{Principal \times demandRate \times Days}{365}\]
    \item
      Chênh lệch thiệt hại lãi do rút trước hạn:
      \[
      \Delta Loss = ExpectedInterest - ActualInterest
      \]
    \end{itemize}
  \item
    Khi người dùng bấm xác nhận chấp nhận điều kiện:

    \begin{itemize}
    \tightlist
    \item
      Tính tổng số tiền thực nhận:
      \(Payout = Principal + ActualInterest\).
    \item
      Cộng \(Payout\) vào số dư của ví nhận tiền theo
      \texttt{BR-FIN-02}.
    \item
      Cập nhật trạng thái sổ tiết kiệm sang \texttt{SETTLED\_EARLY} (tất
      toán trước hạn), lưu lại thông tin số lãi đã mất.
    \item
      Cập nhật ngắt liên kết tài sản khỏi quỹ mục tiêu theo
      \texttt{BR-FIN-08}.
    \end{itemize}
  \end{itemize}
\item
  \textbf{Dữ liệu đầu ra:}

  \begin{itemize}
  \tightlist
  \item
    Số dư ví nhận tiền tăng thêm số tiền \(Payout\).
  \item
    Trạng thái sổ tiết kiệm chuyển thành \texttt{SETTLED\_EARLY}.
  \item
    Thông báo tất toán trước hạn hoàn tất kèm thống kê chi tiết số tiền
    gốc, lãi thực nhận và số lãi bị mất.
  \end{itemize}
\item
  \textbf{Ca sử dụng tham chiếu:}
  \texttt{UC06\ -\ Quản\ lý\ sổ\ tiết\ kiệm\ có\ kỳ\ hạn}.
\item
  \textbf{Tiêu chí nghiệm thu:}

  \begin{itemize}
  \tightlist
  \item
    TC-SAV-04.1: Cho sổ 100,000,000 VNĐ kỳ hạn 1 năm lãi 6.0\%, gửi được
    90 ngày thì tất toán trước hạn với lãi không kỳ hạn 0.2\%/năm; hệ
    thống tính tiền lãi thực nhận là (100,000,000 * 0.002 * 90) / 365 =
    49,315 VNĐ, ví nhận được cộng 100,049,315 VNĐ và thông báo đã mất
    hơn 5.95 triệu VNĐ tiền lãi kỳ hạn.
  \item
    TC-SAV-04.2: Khi người dùng chưa bấm xác nhận cảnh báo mất lãi kỳ
    hạn, hệ thống không thực hiện lệnh chuyển tiền và giữ nguyên trạng
    thái sổ tiết kiệm.
  \end{itemize}
\end{itemize}

\subsubsection{4.7 Phân hệ Chứng chỉ quỹ đầu tư
(FR-INV)}\label{47-phuxe2n-hux1ec7-chux1ee9ng-chux1ec9-quux1ef9-ux111ux1ea7u-tux1b0-fr-inv}

Phân hệ Chứng chỉ quỹ đầu tư quản lý các tài sản đầu tư tài chính của
người dùng, định giá danh mục theo giá thị trường (Mark-to-Market) dựa
trên NAV và theo dõi tỷ suất sinh lời thực tế theo quy tắc BR-FIN-06.

\paragraph{4.7.1 Yêu cầu FR-INV-01: Ghi nhận lệnh mua chứng chỉ
quỹ}\label{471-yuxeau-cux1ea7u-fr-inv-01-ghi-nhux1eadn-lux1ec7nh-mua-chux1ee9ng-chux1ec9-quux1ef9}

\begin{itemize}
\tightlist
\item
  \textbf{Mã yêu cầu:} \texttt{FR-INV-01}
\item
  \textbf{Tên chức năng:} Ghi nhận lệnh mua chứng chỉ quỹ
\item
  \textbf{Mô tả:} Ghi nhận việc trích tiền từ ví thanh khoản để mua
  chứng chỉ quỹ mở (cổ phiếu, trái phiếu, cân bằng), tính toán số lượng
  đơn vị quỹ sở hữu và ghi nhận vào danh mục đầu tư theo quy tắc
  BR-FIN-06.
\item
  \textbf{Dữ liệu đầu vào:}

  \begin{itemize}
  \tightlist
  \item
    Mã chứng chỉ quỹ (\texttt{name}): Chuỗi ký tự, từ 1 đến 50 ký tự,
    bắt buộc (ví dụ: \texttt{DCDS}, \texttt{SSISCA}, \texttt{VESAF}).
  \item
    Loại tài sản quỹ (\texttt{type}): Kiểu liệt kê (Enum) gồm
    \texttt{STOCK} (quỹ cổ phiếu), \texttt{BOND} (quỹ trái phiếu),
    \texttt{BALANCE} (quỹ cân bằng), bắt buộc.
  \item
    Ngày giao dịch (\texttt{start}): Kiểu \texttt{LocalDate}, bắt buộc,
    mặc định hôm nay.
  \item
    Số vốn gốc đầu tư (\texttt{principal}): Kiểu \texttt{BigDecimal},
    bắt buộc, lớn hơn 0.
  \item
    Đơn giá mua NAV (\texttt{buyNav}): Kiểu \texttt{BigDecimal}, bắt
    buộc, lớn hơn 0.
  \item
    Ví trích tiền (\texttt{walletId}): Chuỗi ký tự, bắt buộc, ví ở trạng
    thái \texttt{ACTIVE}.
  \item
    Quỹ mục tiêu liên kết (\texttt{fund}): Chuỗi ký tự, tùy chọn theo
    \texttt{BR-FIN-08}.
  \end{itemize}
\item
  \textbf{Quy tắc xử lý và ràng buộc:}

  \begin{itemize}
  \tightlist
  \item
    Kiểm tra số dư ví trích tiền theo \texttt{BR-FIN-02}:
    \(Balance_{wallet} - principal \ge 0\). Nếu không đủ số dư, từ chối
    lệnh mua.
  \item
    Trừ số vốn gốc \(principal\) khỏi số dư ví trích tiền.
  \item
    Tính toán số lượng chứng chỉ quỹ mua được (\(Units\)):
    \[Units = \frac{principal}{buyNav}\] Lưu trữ số lượng đơn vị quỹ với
    scale tối thiểu là 4 chữ số thập phân, sử dụng
    \texttt{RoundingMode.HALF\_EVEN}.
  \item
    Khởi tạo bản ghi đầu tư mới trong bảng \texttt{fund\_certs} với giá
    trị ban đầu \(Balance = principal\), trạng thái \texttt{HOLDING}.
  \item
    Gắn liên kết tài sản này vào quỹ mục tiêu chỉ định nếu có.
  \end{itemize}
\item
  \textbf{Dữ liệu đầu ra:}

  \begin{itemize}
  \tightlist
  \item
    Bản ghi chứng chỉ quỹ mới được lưu vào bảng \texttt{fund\_certs} kèm
    mã định danh duy nhất (\texttt{certId}).
  \item
    Số lượng đơn vị quỹ nắm giữ được xác lập.
  \item
    Số dư ví trích tiền giảm đúng bằng số tiền vốn đầu tư.
  \item
    Thông báo mua chứng chỉ quỹ thành công.
  \end{itemize}
\item
  \textbf{Ca sử dụng tham chiếu:}
  \texttt{UC07\ -\ Quản\ lý\ và\ định\ giá\ chứng\ chỉ\ quỹ\ đầu\ tư}.
\item
  \textbf{Tiêu chí nghiệm thu:}

  \begin{itemize}
  \tightlist
  \item
    TC-INV-01.1: Cho ví "Timo" có 20,000,000 VNĐ; khi người dùng mua
    chứng chỉ quỹ DCDS với số tiền 10,000,000 VNĐ tại NAV 50,000
    VNĐ/CCQ, số dư ví Timo giảm còn 10,000,000 VNĐ và hệ thống ghi nhận
    sở hữu đúng 200.0000 CCQ DCDS.
  \item
    TC-INV-01.2: Khi số tiền đầu tư vượt quá số dư khả dụng của ví trích
    tiền, hệ thống chặn lại và thông báo số dư không đủ.
  \end{itemize}
\end{itemize}

\paragraph{4.7.2 Yêu cầu FR-INV-02: Cập nhật giá trị tài sản ròng (NAV
Update)}\label{472-yuxeau-cux1ea7u-fr-inv-02-cux1eadp-nhux1eadt-giuxe1-trux1ecb-tuxe0i-sux1ea3n-ruxf2ng-nav-update}

\begin{itemize}
\tightlist
\item
  \textbf{Mã yêu cầu:} \texttt{FR-INV-02}
\item
  \textbf{Tên chức năng:} Cập nhật giá trị tài sản ròng (NAV Update)
\item
  \textbf{Mô tả:} Cho phép người dùng cập nhật đơn giá NAV thị trường
  mới nhất của các chứng chỉ quỹ đang nắm giữ, tự động tính lại giá trị
  hiện tại của toàn bộ danh mục theo quy tắc Mark-to-Market trong quy
  tắc BR-FIN-06.
\item
  \textbf{Dữ liệu đầu vào:}

  \begin{itemize}
  \tightlist
  \item
    Mã chứng chỉ quỹ (\texttt{name}): Chuỗi ký tự, bắt buộc.
  \item
    Ngày cập nhật (\texttt{updateDate}): Kiểu \texttt{LocalDate}, bắt
    buộc, mặc định hôm nay.
  \item
    Đơn giá NAV mới nhất (\texttt{latestNav}): Kiểu \texttt{BigDecimal},
    bắt buộc, lớn hơn 0.
  \end{itemize}
\item
  \textbf{Quy tắc xử lý và ràng buộc:}

  \begin{itemize}
  \tightlist
  \item
    Kiểm tra giá trị NAV mới: \(latestNav > 0\). Nếu nhỏ hơn hoặc bằng
    0, báo lỗi đơn giá không hợp lệ.
  \item
    Tìm tất cả các bản ghi đầu tư đang mở thuộc mã chứng chỉ quỹ
    \texttt{name}.
  \item
    Với từng bản ghi, tính toán lại giá trị thị trường hiện tại
    (\(Balance\)) theo nguyên tắc Mark-to-Market:
    \[Balance = Units \times latestNav\]
  \item
    Hoặc trong trường hợp người dùng cập nhật trực tiếp tổng giá trị thị
    trường của danh mục (\(Balance\)), hệ thống cập nhật giá trị trường
    \(balance\) tương ứng.
  \item
    Ghi nhận lịch sử biến động NAV vào bảng chuỗi thời gian hiệu suất
    quỹ (\texttt{fund\_performance}) để phục vụ việc vẽ biểu đồ tăng
    trưởng danh mục.
  \end{itemize}
\item
  \textbf{Dữ liệu đầu ra:}

  \begin{itemize}
  \tightlist
  \item
    Đơn giá NAV mới được cập nhật trong hệ thống.
  \item
    Giá trị thị trường hiện tại của toàn bộ các khoản đầu tư cùng mã quỹ
    được tính lại đồng bộ.
  \item
    Thông báo cập nhật NAV thành công.
  \end{itemize}
\item
  \textbf{Ca sử dụng tham chiếu:}
  \texttt{UC07\ -\ Quản\ lý\ và\ định\ giá\ chứng\ chỉ\ quỹ\ đầu\ tư}.
\item
  \textbf{Tiêu chí nghiệm thu:}

  \begin{itemize}
  \tightlist
  \item
    TC-INV-02.1: Cho lô chứng chỉ quỹ DCDS sở hữu 200 CCQ với NAV mua
    50,000 VNĐ (giá trị ban đầu 10,000,000 VNĐ); khi người dùng cập nhật
    NAV mới là 55,000 VNĐ, giá trị thị trường hiện tại của lô quỹ tự
    động tăng lên 11,000,000 VNĐ.
  \item
    TC-INV-02.2: Khi người dùng nhập NAV bằng 0 hoặc số âm, hệ thống từ
    chối cập nhật và báo lỗi đơn giá NAV phải lớn hơn 0.
  \end{itemize}
\end{itemize}

\paragraph{4.7.3 Yêu cầu FR-INV-03: Định giá thị trường và tính lãi lỗ
chưa thực
hiện}\label{473-yuxeau-cux1ea7u-fr-inv-03-ux111ux1ecbnh-giuxe1-thux1ecb-trux1b0ux1eddng-vuxe0-tuxednh-luxe3i-lux1ed7-chux1b0a-thux1ef1c-hiux1ec7n}

\begin{itemize}
\tightlist
\item
  \textbf{Mã yêu cầu:} \texttt{FR-INV-03}
\item
  \textbf{Tên chức năng:} Định giá thị trường và tính lãi lỗ chưa thực
  hiện
\item
  \textbf{Mô tả:} Tự động tính toán lãi/lỗ chưa thực hiện (danh nghĩa),
  tỷ suất lợi nhuận tuyệt đối và suất sinh lời quy đổi năm (Annualized
  Rate of Return) cho từng khoản đầu tư và toàn bộ danh mục theo quy tắc
  BR-FIN-06.
\item
  \textbf{Dữ liệu đầu vào:}

  \begin{itemize}
  \tightlist
  \item
    Mã khoản đầu tư (\texttt{certId}) hoặc yêu cầu tổng hợp toàn bộ danh
    mục.
  \item
    Ngày định giá (\texttt{currentDate}): Kiểu \texttt{LocalDate}, mặc
    định hôm nay.
  \end{itemize}
\item
  \textbf{Quy tắc xử lý và ràng buộc:}

  \begin{itemize}
  \tightlist
  \item
    Tính lãi/lỗ chưa thực hiện (Unrealized Profit/Loss):
    \[UnrealizedPL = Balance - Principal\]

    \begin{itemize}
    \tightlist
    \item
      Nếu \(UnrealizedPL > 0\): Khoản đầu tư có lãi danh nghĩa.
    \item
      Nếu \(UnrealizedPL < 0\): Khoản đầu tư chịu lỗ danh nghĩa.
    \end{itemize}
  \item
    Tính tỷ suất sinh lời tuyệt đối:
    \(AbsoluteReturn = \frac{UnrealizedPL}{Principal} \times 100\%\).
  \item
    Tính số ngày nắm giữ thực tế:
    \(HoldDays = \max(1, currentDate - start)\).
  \item
    Tính suất sinh lời quy đổi năm (Annualized Rate) theo quy tắc
    \texttt{BR-FIN-06}:

    \begin{itemize}
    \tightlist
    \item
      Nếu \(HoldDays \ge 30\) ngày và \(Principal > 0\):
      \[AnnualizedRate = \left[ \left( \frac{Balance}{Principal} \right)^{\frac{365}{HoldDays}} - 1 \right] \times 100\%\]
    \item
      Nếu \(HoldDays < 30\) ngày: Hệ thống chỉ hiển thị tỷ suất sinh lời
      tuyệt đối \(AbsoluteReturn\) và đánh dấu chưa đủ thời gian tối
      thiểu để quy đổi năm, nhằm tránh méo mó số liệu thống kê.
    \end{itemize}
  \item
    Sử dụng \texttt{BigDecimal} với scale tối thiểu 4 chữ số thập phân
    cho các phép tính số mũ và làm tròn
    \texttt{RoundingMode.HALF\_EVEN}.
  \end{itemize}
\item
  \textbf{Dữ liệu đầu ra:}

  \begin{itemize}
  \tightlist
  \item
    Bảng tổng hợp hiệu suất đầu tư gồm: Vốn gốc, Giá trị thị trường hiện
    tại, Lãi/lỗ chưa thực hiện, Tỷ suất lợi nhuận tuyệt đối (\%), Suất
    sinh lời quy đổi năm (\%/năm).
  \end{itemize}
\item
  \textbf{Ca sử dụng tham chiếu:}
  \texttt{UC07\ -\ Quản\ lý\ và\ định\ giá\ chứng\ chỉ\ quỹ\ đầu\ tư}.
\item
  \textbf{Tiêu chí nghiệm thu:}

  \begin{itemize}
  \tightlist
  \item
    TC-INV-03.1: Cho khoản đầu tư vốn gốc 10,000,000 VNĐ, giá trị hiện
    tại 12,000,000 VNĐ sau đúng 365 ngày nắm giữ; hệ thống tính ra lãi
    chưa thực hiện là +2,000,000 VNĐ và suất sinh lời quy đổi năm là
    đúng 20.00\%/năm.
  \item
    TC-INV-03.2: Cho khoản đầu tư mới mua được 15 ngày có lãi 200,000
    VNĐ trên vốn 10,000,000 VNĐ; hệ thống hiển thị tỷ suất lợi nhuận
    tuyệt đối là 2.00\% và không áp dụng công thức số mũ quy đổi năm.
  \end{itemize}
\end{itemize}

\paragraph{4.7.4 Yêu cầu FR-INV-04: Ghi nhận bán chứng chỉ
quỹ}\label{474-yuxeau-cux1ea7u-fr-inv-04-ghi-nhux1eadn-buxe1n-chux1ee9ng-chux1ec9-quux1ef9}

\begin{itemize}
\tightlist
\item
  \textbf{Mã yêu cầu:} \texttt{FR-INV-04}
\item
  \textbf{Tên chức năng:} Ghi nhận bán chứng chỉ quỹ
\item
  \textbf{Mô tả:} Ghi nhận giao dịch bán một phần hoặc toàn bộ số lượng
  chứng chỉ quỹ nắm giữ, kết chuyển tiền bán về ví thanh khoản chỉ định
  và ghi nhận lợi nhuận hoặc thua lỗ thực tế đã hiện thực hóa.
\item
  \textbf{Dữ liệu đầu vào:}

  \begin{itemize}
  \tightlist
  \item
    Mã khoản đầu tư (\texttt{certId}): Chuỗi ký tự, bắt buộc.
  \item
    Ngày bán (\texttt{sellDate}): Kiểu \texttt{LocalDate}, bắt buộc, mặc
    định hôm nay.
  \item
    Số lượng chứng chỉ quỹ bán (\texttt{sellUnits}): Kiểu
    \texttt{BigDecimal}, bắt buộc, lớn hơn 0.
  \item
    Đơn giá bán NAV (\texttt{sellNav}): Kiểu \texttt{BigDecimal}, bắt
    buộc, lớn hơn 0.
  \item
    Ví nhận tiền (\texttt{walletId}): Chuỗi ký tự, bắt buộc, ví ở trạng
    thái \texttt{ACTIVE}.
  \end{itemize}
\item
  \textbf{Quy tắc xử lý và ràng buộc:}

  \begin{itemize}
  \tightlist
  \item
    Kiểm tra số lượng bán hợp lệ: \(sellUnits \le Units\) hiện có. Nếu
    bán vượt quá số lượng sở hữu, từ chối lệnh bán.
  \item
    Tính tổng số tiền thu về từ lệnh bán:
    \(Proceeds = sellUnits \times sellNav\).
  \item
    Tính vốn gốc tương ứng của lượng chứng chỉ quỹ bán ra theo phương
    pháp tỷ trọng bình quân:
    \[CostBasis = \frac{sellUnits}{Units} \times Principal\]
  \item
    Tính lợi nhuận thực tế đã thực hiện (Realized Profit/Loss):
    \[RealizedPL = Proceeds - CostBasis\]
  \item
    Cộng số tiền \(Proceeds\) vào số dư của ví nhận tiền theo
    \texttt{BR-FIN-02}.
  \item
    Cập nhật khoản đầu tư:

    \begin{itemize}
    \tightlist
    \item
      Giảm số lượng CCQ nắm giữ: \(Units = Units - sellUnits\).
    \item
      Giảm vốn gốc còn lại: \(Principal = Principal - CostBasis\).
    \item
      Nếu \(Units = 0\), chuyển trạng thái khoản đầu tư sang
      \texttt{CLOSED} (đã đóng).
    \end{itemize}
  \item
    Ghi nhận \(RealizedPL\) vào báo cáo thu nhập/chi phí tài chính của
    kỳ.
  \end{itemize}
\item
  \textbf{Dữ liệu đầu ra:}

  \begin{itemize}
  \tightlist
  \item
    Số tiền bán \(Proceeds\) được ghi có vào ví nhận tiền.
  \item
    Số lượng CCQ và vốn gốc còn lại của khoản đầu tư được cập nhật.
  \item
    Giá trị lợi nhuận/thua lỗ thực tế được ghi nhận vào sổ cái.
  \item
    Thông báo hoàn tất giao dịch bán chứng chỉ quỹ.
  \end{itemize}
\item
  \textbf{Ca sử dụng tham chiếu:}
  \texttt{UC07\ -\ Quản\ lý\ và\ định\ giá\ chứng\ chỉ\ quỹ\ đầu\ tư}.
\item
  \textbf{Tiêu chí nghiệm thu:}

  \begin{itemize}
  \tightlist
  \item
    TC-INV-04.1: Cho khoản đầu tư sở hữu 200 CCQ với tổng giá vốn
    10,000,000 VNĐ; khi người dùng bán 100 CCQ tại giá 60,000 VNĐ/CCQ về
    ví MB Bank, ví MB Bank tăng 6,000,000 VNĐ, lợi nhuận thực hiện ghi
    nhận là +1,000,000 VNĐ, và còn lại 100 CCQ với giá vốn 5,000,000
    VNĐ.
  \item
    TC-INV-04.2: Khi người dùng cố gắng nhập số lượng bán lớn hơn số
    lượng CCQ hiện có trong tài khoản, hệ thống từ chối thao tác và báo
    lỗi số lượng bán vượt quá số lượng nắm giữ.
  \end{itemize}
\end{itemize}

\subsubsection{4.8 Phân hệ Quản lý quỹ mục tiêu
(FR-FND)}\label{48-phuxe2n-hux1ec7-quux1ea3n-luxfd-quux1ef9-mux1ee5c-tiuxeau-fr-fnd}

Phân hệ Quản lý quỹ mục tiêu thiết lập cơ chế quy hoạch và tổng hợp các
nguồn lực tài chính đa dạng (tiền mặt thanh khoản, sổ tiết kiệm và chứng
chỉ quỹ đầu tư) để phục vụ các mục tiêu dài hạn của người dùng theo quy
tắc BR-FIN-08.

\paragraph{4.8.1 Yêu cầu FR-FND-01: Thiết lập quỹ tài chính mục
tiêu}\label{481-yuxeau-cux1ea7u-fr-fnd-01-thiux1ebft-lux1eadp-quux1ef9-tuxe0i-chuxednh-mux1ee5c-tiuxeau}

\begin{itemize}
\tightlist
\item
  \textbf{Mã yêu cầu:} \texttt{FR-FND-01}
\item
  \textbf{Tên chức năng:} Thiết lập quỹ tài chính mục tiêu
\item
  \textbf{Mô tả:} Khởi tạo một mục tiêu tài chính mới hoặc cấu hình 4
  nhóm quỹ chuẩn hóa (emergency, health, target, retirement), xác định
  số tiền đích cần tích lũy và thời hạn hoàn thành theo quy tắc
  BR-FIN-08.
\item
  \textbf{Dữ liệu đầu vào:}

  \begin{itemize}
  \tightlist
  \item
    Mã định danh quỹ (\texttt{fundCode}): Chuỗi ký tự, độ dài từ 1 đến
    30 ký tự, bắt buộc, không trùng lặp (ví dụ: \texttt{emergency},
    \texttt{health}, \texttt{target}, \texttt{retirement} hoặc mã do
    người dùng tự tạo).
  \item
    Tên quỹ mục tiêu (\texttt{fundName}): Chuỗi ký tự, từ 1 đến 100 ký
    tự, bắt buộc.
  \item
    Số tiền mục tiêu cần tích lũy (\texttt{targetAmount}): Kiểu
    \texttt{BigDecimal}, bắt buộc, lớn hơn hoặc bằng 0.
  \item
    Hạn chót hoàn thành mục tiêu (\texttt{deadline}): Kiểu
    \texttt{LocalDate}, tùy chọn.
  \item
    Mô tả hoặc ghi chú định hướng (\texttt{description}): Chuỗi ký tự,
    tối đa 255 ký tự, tùy chọn.
  \end{itemize}
\item
  \textbf{Quy tắc xử lý và ràng buộc:}

  \begin{itemize}
  \tightlist
  \item
    Kiểm tra tính duy nhất của \texttt{fundCode} trong cơ sở dữ liệu hệ
    thống.
  \item
    Khởi tạo cấu trúc dữ liệu quỹ với 3 cấu phần tài sản ban đầu đều
    bằng 0: Tiền mặt phân bổ (\(Cash = 0\)), Sổ tiết kiệm
    (\(Savings = 0\)), Chứng chỉ quỹ (\(FundCerts = 0\)).
  \item
    Lưu bản ghi quỹ mới vào bảng \texttt{funds} với tiến độ hoàn thành
    khởi tạo \(Process = 0\%\).
  \item
    Nếu người dùng nhập \texttt{targetAmount\ =\ 0}, hệ thống coi quỹ
    này là quỹ theo dõi tích lũy mở (không giới hạn trần mục tiêu).
  \end{itemize}
\item
  \textbf{Dữ liệu đầu ra:}

  \begin{itemize}
  \tightlist
  \item
    Đối tượng quỹ mục tiêu mới được lưu thành công kèm mã
    \texttt{fundCode}.
  \item
    Thông báo thiết lập quỹ mục tiêu thành công.
  \end{itemize}
\item
  \textbf{Ca sử dụng tham chiếu:}
  \texttt{UC08\ -\ Quản\ lý\ quỹ\ tài\ chính\ mục\ tiêu}.
\item
  \textbf{Tiêu chí nghiệm thu:}

  \begin{itemize}
  \tightlist
  \item
    TC-FND-01.1: Khi người dùng tạo quỹ "emergency" với tên "Quỹ dự
    phòng khẩn cấp" và mục tiêu 60,000,000 VNĐ, hệ thống lưu thành công
    với tổng tài sản ban đầu bằng 0 VNĐ và tiến độ hiển thị 0.00\%.
  \item
    TC-FND-01.2: Khi người dùng nhập số tiền mục tiêu là số âm, hệ thống
    từ chối và đưa ra thông báo số tiền mục tiêu phải lớn hơn hoặc bằng
    0.
  \end{itemize}
\end{itemize}

\paragraph{4.8.2 Yêu cầu FR-FND-02: Phân bổ tài sản vào quỹ mục
tiêu}\label{482-yuxeau-cux1ea7u-fr-fnd-02-phuxe2n-bux1ed5-tuxe0i-sux1ea3n-vuxe0o-quux1ef9-mux1ee5c-tiuxeau}

\begin{itemize}
\tightlist
\item
  \textbf{Mã yêu cầu:} \texttt{FR-FND-02}
\item
  \textbf{Tên chức năng:} Phân bổ tài sản vào quỹ mục tiêu
\item
  \textbf{Mô tả:} Cho phép người dùng quy hoạch và phân bổ các nguồn tài
  sản cụ thể gồm tiền mặt dành riêng, sổ tiết kiệm và chứng chỉ quỹ đầu
  tư vào từng quỹ mục tiêu theo quy tắc BR-FIN-08.
\item
  \textbf{Dữ liệu đầu vào:}

  \begin{itemize}
  \tightlist
  \item
    Mã quỹ mục tiêu (\texttt{fundCode}): Chuỗi ký tự, bắt buộc.
  \item
    Số tiền mặt phân bổ (\texttt{cashAmount}): Kiểu \texttt{BigDecimal},
    tùy chọn, lớn hơn hoặc bằng 0 (lưu vào bảng \texttt{funds\_cash}).
  \item
    Danh sách mã sổ tiết kiệm liên kết (\texttt{savingIds}): Danh sách
    chuỗi ký tự, tùy chọn.
  \item
    Danh sách mã chứng chỉ quỹ liên kết (\texttt{certIds}): Danh sách
    chuỗi ký tự, tùy chọn.
  \end{itemize}
\item
  \textbf{Quy tắc xử lý và ràng buộc:}

  \begin{itemize}
  \tightlist
  \item
    Phân bổ tiền mặt: Cập nhật giá trị tiền mặt phân bổ cho quỹ trong
    bảng \texttt{funds\_cash}; kiểm tra tổng lượng tiền mặt phân bổ cho
    mọi quỹ không vượt quá tổng số dư tiền mặt khả dụng trên tất cả các
    ví của hệ thống.
  \item
    Phân bổ sổ tiết kiệm và chứng chỉ quỹ: Tuân thủ nghiêm ngặt nguyên
    tắc không tính trùng tài sản (No Double-Counting) tại quy tắc
    \texttt{BR-FIN-08}. Mỗi sổ tiết kiệm cụ thể và mỗi lô chứng chỉ quỹ
    cụ thể chỉ được phép thuộc về duy nhất một quỹ mục tiêu tại một thời
    điểm.
  \item
    Nếu tài sản trước đó đang gắn với một quỹ khác, hệ thống tự động gỡ
    liên kết khỏi quỹ cũ và gán vào quỹ mới, đồng thời cập nhật lại giá
    trị tích lũy của cả hai quỹ đối ứng.
  \end{itemize}
\item
  \textbf{Dữ liệu đầu ra:}

  \begin{itemize}
  \tightlist
  \item
    Các liên kết tài sản được xác lập thành công.
  \item
    Giá trị các cấu phần của quỹ mục tiêu được tính toán lại.
  \item
    Thông báo phân bổ tài sản thành công.
  \end{itemize}
\item
  \textbf{Ca sử dụng tham chiếu:}
  \texttt{UC08\ -\ Quản\ lý\ quỹ\ tài\ chính\ mục\ tiêu}.
\item
  \textbf{Tiêu chí nghiệm thu:}

  \begin{itemize}
  \tightlist
  \item
    TC-FND-02.1: Cho quỹ "emergency" được gán 10,000,000 VNĐ tiền mặt
    phân bổ và một sổ tiết kiệm giá trị 30,000,000 VNĐ; hệ thống hiển
    thị cấu phần tiền mặt là 10,000,000 VNĐ, cấu phần tiết kiệm là
    30,000,000 VNĐ, và tổng giá trị quỹ là 40,000,000 VNĐ.
  \item
    TC-FND-02.2: Cho sổ tiết kiệm S1 đang liên kết với quỹ "emergency";
    khi người dùng gán S1 sang quỹ "target", hệ thống tự động loại S1
    khỏi "emergency" và cộng vào "target", bảo đảm tổng tài sản toàn hệ
    thống không đổi.
  \end{itemize}
\end{itemize}

\paragraph{4.8.3 Yêu cầu FR-FND-03: Theo dõi tiến độ tích lũy
quỹ}\label{483-yuxeau-cux1ea7u-fr-fnd-03-theo-duxf5i-tiux1ebfn-ux111ux1ed9-tuxedch-lux169y-quux1ef9}

\begin{itemize}
\tightlist
\item
  \textbf{Mã yêu cầu:} \texttt{FR-FND-03}
\item
  \textbf{Tên chức năng:} Theo dõi tiến độ tích lũy quỹ
\item
  \textbf{Mô tả:} Tự động tổng hợp giá trị hiện tại của 3 cấu phần tài
  sản và tính toán tỷ lệ phần trăm hoàn thành mục tiêu theo thời gian
  thực theo quy tắc BR-FIN-08.
\item
  \textbf{Dữ liệu đầu vào:}

  \begin{itemize}
  \tightlist
  \item
    Mã quỹ mục tiêu (\texttt{fundCode}) hoặc yêu cầu tổng hợp toàn bộ
    danh sách quỹ.
  \end{itemize}
\item
  \textbf{Quy tắc xử lý và ràng buộc:}

  \begin{itemize}
  \tightlist
  \item
    Tính tổng tài sản tích lũy thực tế của quỹ mục tiêu
    (\(Total_{fund}\)):
    \[Total_{fund} = Cash_{fund} + \sum_{s \in Savings_{fund}} Balance_s + \sum_{c \in FundCerts_{fund}} Balance_c\]
    Trong đó:

    \begin{itemize}
    \tightlist
    \item
      \(Balance_s\): Số dư ghi nhận của sổ tiết kiệm tại thời điểm tính
      toán theo \texttt{FR-SAV-02}.
    \item
      \(Balance_c\): Giá trị thị trường hiện tại của chứng chỉ quỹ theo
      \texttt{FR-INV-03}.
    \end{itemize}
  \item
    Tính tỷ lệ tiến độ hoàn thành (\(Process\)):
    \[Process = \begin{cases} \frac{Total_{fund}}{Target} \times 100\%, & \text{khi } Target > 0 \\ 100\%, & \text{khi } Target = 0 \text{ và } Total_{fund} \ge 0 \end{cases}\]
  \item
    Tính số tiền còn thiếu để đạt mục tiêu:
    \(Remaining = \max(0, Target - Total_{fund})\).
  \item
    Sử dụng \texttt{BigDecimal} với scale 2 và
    \texttt{RoundingMode.HALF\_EVEN}.
  \item
    Tự động gắn nhãn trạng thái "ĐÃ HOÀN THÀNH" khi
    \(Process \ge 100\%\).
  \end{itemize}
\item
  \textbf{Dữ liệu đầu ra:}

  \begin{itemize}
  \tightlist
  \item
    Báo cáo chi tiết tiến độ quỹ gồm: Tổng tài sản thực tế, Phân rã 3
    cấu phần (Tiền mặt, Tiết kiệm, Đầu tư), Số tiền mục tiêu, Tỷ lệ tiến
    độ (\%), Số tiền còn thiếu, Trạng thái hoàn thành.
  \end{itemize}
\item
  \textbf{Ca sử dụng tham chiếu:}
  \texttt{UC08\ -\ Quản\ lý\ quỹ\ tài\ chính\ mục\ tiêu}.
\item
  \textbf{Tiêu chí nghiệm thu:}

  \begin{itemize}
  \tightlist
  \item
    TC-FND-03.1: Cho quỹ có số tiền mục tiêu Target = 100,000,000 VNĐ;
    tổng tài sản thực tế tích lũy hiện tại là 65,000,000 VNĐ; hệ thống
    tính ra tiến độ đạt chính xác 65.00\% và số tiền còn thiếu là
    35,000,000 VNĐ.
  \item
    TC-FND-03.2: Khi tổng tài sản tích lũy đạt 102,000,000 VNĐ (vượt mục
    tiêu 100,000,000 VNĐ), tiến độ hiển thị là 102.00\% và trạng thái
    quỹ chuyển sang "Đã hoàn thành mục tiêu".
  \end{itemize}
\end{itemize}

\paragraph{4.8.4 Yêu cầu FR-FND-04: Điều chỉnh và tái phân bổ quỹ mục
tiêu}\label{484-yuxeau-cux1ea7u-fr-fnd-04-ux111iux1ec1u-chux1ec9nh-vuxe0-tuxe1i-phuxe2n-bux1ed5-quux1ef9-mux1ee5c-tiuxeau}

\begin{itemize}
\tightlist
\item
  \textbf{Mã yêu cầu:} \texttt{FR-FND-04}
\item
  \textbf{Tên chức năng:} Điều chỉnh và tái phân bổ quỹ mục tiêu
\item
  \textbf{Mô tả:} Cho phép người dùng thay đổi số tiền mục tiêu, điều
  chỉnh hạn chót hoặc tái cơ cấu phân bổ tỷ trọng tài sản giữa các quỹ
  nhằm phù hợp với thay đổi trong kế hoạch tài chính cá nhân.
\item
  \textbf{Dữ liệu đầu vào:}

  \begin{itemize}
  \tightlist
  \item
    Mã quỹ mục tiêu (\texttt{fundCode}): Chuỗi ký tự, bắt buộc.
  \item
    Số tiền mục tiêu mới (\texttt{newTarget}): Kiểu \texttt{BigDecimal},
    bắt buộc, lớn hơn hoặc bằng 0.
  \item
    Hạn chót mới (\texttt{newDeadline}): Kiểu \texttt{LocalDate}, tùy
    chọn.
  \item
    Tỷ trọng tiền mặt điều chỉnh trong bảng \texttt{funds\_cash}: Kiểu
    \texttt{BigDecimal}.
  \end{itemize}
\item
  \textbf{Quy tắc xử lý và ràng buộc:}

  \begin{itemize}
  \tightlist
  \item
    Cập nhật số tiền mục tiêu mới và thời hạn mới vào bảng
    \texttt{funds}.
  \item
    Tính toán lại ngay lập tức tỷ lệ tiến độ \(Process\) tương ứng với
    hạn mức mới.
  \item
    Thao tác điều chỉnh mục tiêu thuần túy là việc thay đổi kế hoạch
    phân bổ trên sổ theo dõi, tuyệt đối không làm phát sinh giao dịch
    rút tiền hoặc làm thay đổi số dư thực tế của các ví tiền mặt và sổ
    tiết kiệm liên kết.
  \end{itemize}
\item
  \textbf{Dữ liệu đầu ra:}

  \begin{itemize}
  \tightlist
  \item
    Thông tin quỹ mục tiêu được cập nhật với hạn mức và tiến độ mới.
  \item
    Thông báo điều chỉnh quỹ mục tiêu thành công.
  \end{itemize}
\item
  \textbf{Ca sử dụng tham chiếu:}
  \texttt{UC08\ -\ Quản\ lý\ quỹ\ tài\ chính\ mục\ tiêu}.
\item
  \textbf{Tiêu chí nghiệm thu:}

  \begin{itemize}
  \tightlist
  \item
    TC-FND-04.1: Cho quỹ đang có mục tiêu 50,000,000 VNĐ và đã tích lũy
    được 40,000,000 VNĐ (đạt 80.00\%); khi người dùng nâng mục tiêu lên
    80,000,000 VNĐ, hệ thống cập nhật tiến độ mới thành đúng 50.00\% mà
    số tiền tích lũy 40,000,000 VNĐ không bị thay đổi.
  \item
    TC-FND-04.2: Khi người dùng nhập hạn chót là một ngày đã qua trong
    quá khứ, hệ thống hiển thị cảnh báo hạn chót đã quá hạn để người
    dùng cân nhắc trước khi lưu.
  \end{itemize}
\end{itemize}

\subsubsection{4.9 Phân hệ Thống kê và Báo cáo
(FR-REP)}\label{49-phuxe2n-hux1ec7-thux1ed1ng-kuxea-vuxe0-buxe1o-cuxe1o-fr-rep}

Phân hệ Thống kê và Báo cáo tự động tổng hợp toàn bộ dữ liệu giao dịch
và tài sản để cung cấp chuỗi thời gian dòng tiền trượt 12 tháng, cơ cấu
thu chi, báo cáo tài sản ròng và sao kê công nợ đối tác.

\paragraph{4.9.1 Yêu cầu FR-REP-01: Bảng thống kê dòng tiền trượt 12
tháng}\label{491-yuxeau-cux1ea7u-fr-rep-01-bux1ea3ng-thux1ed1ng-kuxea-duxf2ng-tiux1ec1n-trux1b0ux1ee3t-12-thuxe1ng}

\begin{itemize}
\tightlist
\item
  \textbf{Mã yêu cầu:} \texttt{FR-REP-01}
\item
  \textbf{Tên chức năng:} Bảng thống kê dòng tiền trượt 12 tháng
\item
  \textbf{Mô tả:} Tự động tổng hợp và hiển thị chuỗi thời gian 12 tháng
  gần nhất (tương ứng bảng \texttt{months\_stat}) tự động trượt theo
  ngày hệ thống, tính toán tổng thu nhập, tổng chi tiêu, dòng tiền ròng
  và tỷ lệ tiết kiệm từng tháng.
\item
  \textbf{Dữ liệu đầu vào:}

  \begin{itemize}
  \tightlist
  \item
    Mốc thời gian quan sát từ đồng hồ hệ thống (System Clock), xác định
    tháng hiện tại \(T = (Year, Month)\).
  \end{itemize}
\item
  \textbf{Quy tắc xử lý và ràng buộc:}

  \begin{itemize}
  \tightlist
  \item
    Tự động sinh danh sách 12 mốc tháng liên tục lùi từ tháng hiện tại
    (\(T, T-1, T-2, ..., T-11\)), neo dòng cuối cùng vào tháng hiện tại.
    Khi sang tháng mới, bảng dữ liệu tự động dịch chuyển thêm một tháng
    và bỏ tháng cũ nhất mà không cần can thiệp thủ công.
  \item
    Với mỗi tháng \(m\) trong chuỗi 12 tháng:

    \begin{itemize}
    \tightlist
    \item
      Tổng thu nhập (\(Income_m\)): Tổng số tiền của các giao dịch
      \texttt{Income} trong tháng \(m\).
    \item
      Tổng chi tiêu (\(Expense_m\)): Tổng số tiền của các giao dịch
      \texttt{Expense} trong tháng \(m\).
    \item
      Dòng tiền thuần (\(CashFlow_m\)):
      \[CashFlow_m = Income_m - Expense_m\]
    \item
      Tỷ lệ tiết kiệm (\(SavingsRate_m\)):
      \[SavingsRate_m = \begin{cases} \frac{CashFlow_m}{Income_m} \times 100\%, & \text{khi } Income_m > 0 \\ 0\%, & \text{khi } Income_m \le 0 \end{cases}\]
    \end{itemize}
  \item
    Toàn bộ phép tính sử dụng \texttt{BigDecimal} với
    \texttt{RoundingMode.HALF\_EVEN}.
  \end{itemize}
\item
  \textbf{Dữ liệu đầu ra:}

  \begin{itemize}
  \tightlist
  \item
    Bảng thống kê \texttt{months\_stat} gồm 12 dòng, các cột: Tháng
    (\texttt{month}), Thu nhập (\texttt{income}), Chi tiêu
    (\texttt{expense}), Dòng tiền ròng (\texttt{cash\ flow}), Tỷ lệ tiết
    kiệm (\texttt{savings\ rate}).
  \item
    Dữ liệu chuỗi thời gian sẵn sàng cung cấp cho biểu đồ xu hướng dòng
    tiền.
  \end{itemize}
\item
  \textbf{Ca sử dụng tham chiếu:}
  \texttt{UC09\ -\ Xem\ thống\ kê\ và\ báo\ cáo\ tài\ chính\ trượt\ 12\ tháng}.
\item
  \textbf{Tiêu chí nghiệm thu:}

  \begin{itemize}
  \tightlist
  \item
    TC-REP-01.1: Khi tháng hệ thống là tháng 09/2026, bảng
    \texttt{months\_stat} tự động hiển thị chuỗi 12 tháng từ tháng
    10/2025 đến tháng 09/2026.
  \item
    TC-REP-01.2: Cho tháng 08/2026 có tổng thu nhập 30,000,000 VNĐ và
    tổng chi tiêu 18,000,000 VNĐ; hệ thống tính ra Cash Flow là
    +12,000,000 VNĐ và Savings Rate là đúng 40.00\%.
  \item
    TC-REP-01.3: Cho tháng không có thu nhập phát sinh (\(Income = 0\))
    nhưng có chi tiêu 5,000,000 VNĐ; hệ thống hiển thị Cash Flow là
    -5,000,000 VNĐ, Savings Rate là 0.00\% và không gây ra ngoại lệ chia
    cho 0.
  \end{itemize}
\end{itemize}

\paragraph{4.9.2 Yêu cầu FR-REP-02: Báo cáo cơ cấu thu chi theo nhóm và
danh
mục}\label{492-yuxeau-cux1ea7u-fr-rep-02-buxe1o-cuxe1o-cux1a1-cux1ea5u-thu-chi-theo-nhuxf3m-vuxe0-danh-mux1ee5c}

\begin{itemize}
\tightlist
\item
  \textbf{Mã yêu cầu:} \texttt{FR-REP-02}
\item
  \textbf{Tên chức năng:} Báo cáo cơ cấu thu chi theo nhóm và danh mục
\item
  \textbf{Mô tả:} Tổng hợp, phân tích tỷ trọng và so sánh cơ cấu thu
  nhập hoặc chi tiêu theo từng Nhóm nhu cầu (Group) và Danh mục cụ thể
  (Category) của tháng quan sát so với các tháng trước.
\item
  \textbf{Dữ liệu đầu vào:}

  \begin{itemize}
  \tightlist
  \item
    Tháng phân tích (\texttt{targetMonth}): Kiểu \texttt{YearMonth}, mặc
    định tháng hiện tại.
  \item
    Loại phân tích (\texttt{flowType}): Enum \texttt{INCOME} hoặc
    \texttt{EXPENSE}.
  \end{itemize}
\item
  \textbf{Quy tắc xử lý và ràng buộc:}

  \begin{itemize}
  \tightlist
  \item
    Quét toàn bộ các giao dịch trong tháng \texttt{targetMonth} có
    \texttt{Type\ =\ flowType}.
  \item
    Tổng hợp số tiền theo 2 cấp độ: Cấp nhóm (\texttt{Group}) và Cấp
    danh mục chi tiết (\texttt{Category}).
  \item
    Tính tỷ trọng phần trăm của từng nhóm trên tổng thu/chi của tháng.
  \item
    Tính tỷ trọng phần trăm của từng danh mục trên tổng của nhóm cha và
    trên tổng toàn bộ thu/chi.
  \item
    Tính mức độ biến động so với tháng liền trước (\(targetMonth - 1\)):

    \begin{itemize}
    \tightlist
    \item
      Chênh lệch tuyệt đối:
      \(\Delta Amount = Amount_{thisMonth} - Amount_{lastMonth}\).
    \item
      Tỷ lệ phần trăm tăng/giảm:
      \(\% Change = \frac{\Delta Amount}{Amount_{lastMonth}} \times 100\%\).
    \end{itemize}
  \item
    Sắp xếp các danh mục theo số tiền giảm dần để làm nổi bật các khoản
    mục chiếm tỷ trọng lớn nhất.
  \end{itemize}
\item
  \textbf{Dữ liệu đầu ra:}

  \begin{itemize}
  \tightlist
  \item
    Bảng tổng hợp Pivot 2 tầng phân cấp (Nhóm -\textgreater{} Danh mục)
    kèm số tiền, tỷ trọng \%, so sánh chênh lệch với tháng trước.
  \item
    Dữ liệu sẵn sàng phục vụ biểu đồ tròn (cơ cấu) và biểu đồ cột (so
    sánh).
  \end{itemize}
\item
  \textbf{Ca sử dụng tham chiếu:}
  \texttt{UC09\ -\ Xem\ thống\ kê\ và\ báo\ cáo\ tài\ chính\ trượt\ 12\ tháng}.
\item
  \textbf{Tiêu chí nghiệm thu:}

  \begin{itemize}
  \tightlist
  \item
    TC-REP-02.1: Cho tháng 08/2026 có tổng chi tiêu 20,000,000 VNĐ,
    trong đó danh mục "Ăn uống" chi 6,000,000 VNĐ; báo cáo hiển thị danh
    mục "Ăn uống" chiếm tỷ trọng chính xác 30.00\% tổng chi tiêu.
  \item
    TC-REP-02.2: Cho chi phí danh mục "Ăn uống" tháng 07/2026 là
    5,000,000 VNĐ; báo cáo chỉ ra mức tăng +1,000,000 VNĐ (tương ứng
    tăng +20.00\%) so với tháng trước.
  \end{itemize}
\end{itemize}

\paragraph{4.9.3 Yêu cầu FR-REP-03: Báo cáo tổng tài sản và giá trị tài
sản
ròng}\label{493-yuxeau-cux1ea7u-fr-rep-03-buxe1o-cuxe1o-tux1ed5ng-tuxe0i-sux1ea3n-vuxe0-giuxe1-trux1ecb-tuxe0i-sux1ea3n-ruxf2ng}

\begin{itemize}
\tightlist
\item
  \textbf{Mã yêu cầu:} \texttt{FR-REP-03}
\item
  \textbf{Tên chức năng:} Báo cáo tổng tài sản và giá trị tài sản ròng
\item
  \textbf{Mô tả:} Tổng hợp 5 cấu phần tài sản của người dùng tại thời
  điểm quan sát để xác định Giá trị tài sản ròng (Net Worth) và phân
  tích cơ cấu phân bổ tài sản tổng thể.
\item
  \textbf{Dữ liệu đầu vào:}

  \begin{itemize}
  \tightlist
  \item
    Thời điểm quan sát: Mặc định là thời điểm hiện tại.
  \end{itemize}
\item
  \textbf{Quy tắc xử lý và ràng buộc:}

  \begin{itemize}
  \tightlist
  \item
    Tổng hợp giá trị 5 cấu phần tài sản theo chuẩn kế toán:

    \begin{enumerate}
    \def\labelenumi{\arabic{enumi}.}
    \tightlist
    \item
      Tổng số dư ví tiền thanh khoản:
      \(Wallets = \sum_{w \in Wallets} Balance_w\).
    \item
      Tổng giá trị sổ tiết kiệm:
      \(Savings = \sum_{s \in Savings} Balance_s\).
    \item
      Tổng giá trị thị trường chứng chỉ quỹ:
      \(Investments = \sum_{c \in FundCerts} Balance_c\).
    \item
      Tổng các khoản phải thu:
      \(Receivable = \sum Lend - \sum Collect - \sum WriteOffLend\).
    \item
      Tổng các khoản phải trả:
      \(Payable = \sum Borrow - \sum Repay - \sum WriteOffBorrow\).
    \end{enumerate}
  \item
    Tính Giá trị tài sản ròng (Net Worth):
    \[NetWorth = Wallets + Savings + Investments + Receivable - Payable\]
  \item
    Tính tổng tài sản gộp:
    \[
    TotalAssets = Wallets + Savings + Investments + Receivable
    \]
  \item
    Tính tỷ trọng phân bổ của từng cấu phần trên tổng tài sản gộp.
  \item
    Tất cả các phép tính thực hiện bằng \texttt{BigDecimal} bảo đảm sai
    số bằng 0.
  \end{itemize}
\item
  \textbf{Dữ liệu đầu ra:}

  \begin{itemize}
  \tightlist
  \item
    Bảng tổng hợp tài sản gồm 5 dòng: Wallets, Savings, Investments,
    Receivable, Payable.
  \item
    Giá trị Net Worth cuối cùng và tỷ trọng phần trăm của từng cấu phần.
  \end{itemize}
\item
  \textbf{Ca sử dụng tham chiếu:}
  \texttt{UC09\ -\ Xem\ thống\ kê\ và\ báo\ cáo\ tài\ chính\ trượt\ 12\ tháng}.
\item
  \textbf{Tiêu chí nghiệm thu:}

  \begin{itemize}
  \tightlist
  \item
    TC-REP-03.1: Cho Wallets = 50,000,000 VNĐ; Savings = 100,000,000
    VNĐ; Investments = 50,000,000 VNĐ; Receivable = 20,000,000 VNĐ;
    Payable = 10,000,000 VNĐ; hệ thống tính ra Giá trị tài sản ròng Net
    Worth đúng bằng 210,000,000 VNĐ.
  \item
    TC-REP-03.2: Khi có bất kỳ giao dịch phát sinh nào trong hệ thống
    (như ghi nhận chi tiêu hoặc trả nợ), báo cáo Net Worth phản ánh sự
    biến động tức thời khi người dùng làm mới màn hình.
  \end{itemize}
\end{itemize}

\paragraph{4.9.4 Yêu cầu FR-REP-04: Sao kê công nợ chi tiết theo đối
tượng}\label{494-yuxeau-cux1ea7u-fr-rep-04-sao-kuxea-cuxf4ng-nux1ee3-chi-tiux1ebft-theo-ux111ux1ed1i-tux1b0ux1ee3ng}

\begin{itemize}
\tightlist
\item
  \textbf{Mã yêu cầu:} \texttt{FR-REP-04}
\item
  \textbf{Tên chức năng:} Sao kê công nợ chi tiết theo đối tượng
\item
  \textbf{Mô tả:} Lập bảng sao kê chi tiết toàn bộ lịch sử phát sinh và
  số dư công nợ lũy kế đối ứng với từng cá nhân đối tác (Pivot/Statement
  theo từng cá nhân).
\item
  \textbf{Dữ liệu đầu vào:}

  \begin{itemize}
  \tightlist
  \item
    Tên cá nhân đối tác (\texttt{person}): Chuỗi ký tự, tùy chọn (nếu để
    trống, lập danh sách toàn bộ các đối tác có dư nợ).
  \item
    Bộ lọc trạng thái nợ (\texttt{debtStatus}): Tùy chọn gồm
    \texttt{ALL}, \texttt{OUTSTANDING} (còn dư nợ), \texttt{SETTLED} (đã
    tất toán xong).
  \end{itemize}
\item
  \textbf{Quy tắc xử lý và ràng buộc:}

  \begin{itemize}
  \tightlist
  \item
    Lọc tất cả các bản ghi trong bảng \texttt{debt} có trường
    \texttt{person} tương ứng.
  \item
    Sắp xếp các giao dịch theo trình tự thời gian tăng dần
    (chronological order).
  \item
    Với từng giao dịch, tính toán số dư nợ lũy kế sau giao dịch:

    \begin{itemize}
    \tightlist
    \item
      Nếu là \texttt{lend}: Tăng dư nợ phải thu.
    \item
      Nếu là \texttt{collect} hoặc \texttt{write-off\ lend}: Giảm dư nợ
      phải thu.
    \item
      Nếu là \texttt{borrow}: Tăng dư nợ phải trả.
    \item
      Nếu là \texttt{repay} hoặc \texttt{write-off\ borrow}: Giảm dư nợ
      phải trả.
    \end{itemize}
  \item
    Tính tổng các chỉ số tổng hợp của đối tác: Tổng đã cho vay, Tổng đã
    thu hồi, Số nợ còn phải thu, Tổng đã vay, Tổng đã hoàn trả, Số nợ
    còn phải trả.
  \item
    Tổng số nợ phải thu và phải trả của tất cả các cá nhân cộng lại phải
    khớp chính xác với chỉ số \(Receivable\) và \(Payable\) trên báo cáo
    Net Worth toàn hệ thống.
  \end{itemize}
\item
  \textbf{Dữ liệu đầu ra:}

  \begin{itemize}
  \tightlist
  \item
    Bảng sao kê công nợ chi tiết gồm: Ngày phát sinh, Loại nghiệp vụ, Số
    tiền, Ví liên kết, Ghi chú, Dư nợ lũy kế sau giao dịch.
  \item
    Bảng tổng kết số dư nợ hiện tại với đối tượng.
  \end{itemize}
\item
  \textbf{Ca sử dụng tham chiếu:}
  \texttt{UC09\ -\ Xem\ thống\ kê\ và\ báo\ cáo\ tài\ chính\ trượt\ 12\ tháng}.
\item
  \textbf{Tiêu chí nghiệm thu:}

  \begin{itemize}
  \tightlist
  \item
    TC-REP-04.1: Cho đối tác "Nguyễn Văn A" có 1 lần vay 5,000,000 VNĐ
    và 1 lần trả 2,000,000 VNĐ; sao kê của A hiển thị chi tiết 2 dòng
    giao dịch và kết luận nợ phải thu còn lại là đúng 3,000,000 VNĐ.
  \item
    TC-REP-04.2: Khi người dùng lọc danh sách theo điều kiện
    \texttt{OUTSTANDING}, hệ thống chỉ hiển thị những cá nhân có số dư
    nợ phải thu hoặc nợ phải trả lớn hơn 0.
  \end{itemize}
\end{itemize}

\subsubsection{4.10 Phân hệ Xuất nhập và Sao lưu dữ liệu
(FR-DAT)}\label{410-phuxe2n-hux1ec7-xuux1ea5t-nhux1eadp-vuxe0-sao-lux1b0u-dux1eef-liux1ec7u-fr-dat}

Phân hệ Xuất nhập và Sao lưu dữ liệu bảo đảm tính độc lập, khả năng bảo
toàn dữ liệu và phục hồi an toàn của toàn bộ cơ sở dữ liệu hệ thống trên
môi trường máy tính cục bộ.

\paragraph{4.10.1 Yêu cầu FR-DAT-01: Xuất dữ liệu tài chính
(Export)}\label{4101-yuxeau-cux1ea7u-fr-dat-01-xuux1ea5t-dux1eef-liux1ec7u-tuxe0i-chuxednh-export}

\begin{itemize}
\tightlist
\item
  \textbf{Mã yêu cầu:} \texttt{FR-DAT-01}
\item
  \textbf{Tên chức năng:} Xuất dữ liệu tài chính (Export)
\item
  \textbf{Mô tả:} Xuất toàn bộ hoặc từng phần dữ liệu sổ cái giao dịch,
  danh mục và cấu hình hệ thống ra các tệp tin cục bộ theo định dạng
  tiêu chuẩn (JSON hoặc CSV) phục vụ lưu trữ an toàn hoặc chia sẻ dữ
  liệu.
\item
  \textbf{Dữ liệu đầu vào:}

  \begin{itemize}
  \tightlist
  \item
    Phạm vi dữ liệu cần xuất (\texttt{scope}): Kiểu liệt kê gồm
    \texttt{FULL} (toàn bộ hệ thống), \texttt{TRANSACTIONS} (chỉ sổ thu
    chi), \texttt{DEBTS} (chỉ công nợ), \texttt{SAVINGS} (chỉ sổ tiết
    kiệm), \texttt{INVESTMENTS} (chỉ chứng chỉ quỹ).
  \item
    Định dạng tệp xuất (\texttt{format}): Enum \texttt{JSON} hoặc
    \texttt{CSV}.
  \item
    Thư mục lưu trữ trên máy tính (\texttt{destinationPath}): Chuỗi ký
    tự đường dẫn thư mục hợp lệ.
  \end{itemize}
\item
  \textbf{Quy tắc xử lý và ràng buộc:}

  \begin{itemize}
  \tightlist
  \item
    Đọc dữ liệu hiện tại từ bộ nhớ hoặc cơ sở dữ liệu cục bộ.
  \item
    Chuyển đổi các thực thể miền sang cấu trúc JSON có định dạng chuẩn
    hoặc các bảng CSV có dòng tiêu đề cột rõ ràng.
  \item
    Ràng buộc xử lý số học: Tất cả các trường dữ liệu
    \texttt{BigDecimal} khi xuất ra tệp phải được biểu diễn dưới dạng
    chuỗi số học chuẩn xác (ví dụ: \texttt{"15000000.00"}), nghiêm cấm
    biểu diễn dưới dạng số thực dấu chấm động hoặc ký hiệu khoa học (như
    \texttt{1.5E7}) để tránh mất độ chính xác khi nhập lại.
  \item
    Đặt tên tệp tự động theo quy tắc:
    \texttt{pfms\_export\_\{scope\}\_\{yyyyMMdd\_HHmmss\}.\{ext\}}.
  \item
    Ghi tệp an toàn vào đĩa cứng với mã hóa UTF-8.
  \end{itemize}
\item
  \textbf{Dữ liệu đầu ra:}

  \begin{itemize}
  \tightlist
  \item
    Tệp dữ liệu được tạo thành công tại đường dẫn chỉ định.
  \item
    Thông báo xuất dữ liệu thành công kèm thông tin về đường dẫn tệp,
    kích thước tệp và số lượng bản ghi đã xuất.
  \end{itemize}
\item
  \textbf{Ca sử dụng tham chiếu:}
  \texttt{UC10\ -\ Xuất\ nhập\ và\ sao\ lưu\ dữ\ liệu\ hệ\ thống}.
\item
  \textbf{Tiêu chí nghiệm thu:}

  \begin{itemize}
  \tightlist
  \item
    TC-DAT-01.1: Khi người dùng chọn xuất toàn bộ hệ thống sang định
    dạng JSON, tệp JSON được tạo ra chứa đầy đủ 10 phân hệ với số lượng
    bản ghi khớp hoàn toàn với cơ sở dữ liệu hiện tại.
  \item
    TC-DAT-01.2: Khi mở tệp JSON đã xuất, toàn bộ các trường số tiền đều
    là chuỗi số học đầy đủ chữ số, không bị làm tròn hoặc chuyển sang ký
    hiệu khoa học.
  \end{itemize}
\end{itemize}

\paragraph{4.10.2 Yêu cầu FR-DAT-02: Nhập dữ liệu và phục hồi (Import \&
Restore)}\label{4102-yuxeau-cux1ea7u-fr-dat-02-nhux1eadp-dux1eef-liux1ec7u-vuxe0-phux1ee5c-hux1ed3i-import--restore}

\begin{itemize}
\tightlist
\item
  \textbf{Mã yêu cầu:} \texttt{FR-DAT-02}
\item
  \textbf{Tên chức năng:} Nhập dữ liệu và phục hồi (Import \& Restore)
\item
  \textbf{Mô tả:} Đọc và phục hồi dữ liệu tài chính từ tệp sao lưu trước
  đó, thực hiện xác thực cấu trúc và kiểm tra tính toàn vẹn nghiệp vụ
  trước khi ghi vào hệ thống.
\item
  \textbf{Dữ liệu đầu vào:}

  \begin{itemize}
  \tightlist
  \item
    Đường dẫn tệp sao lưu cần nhập (\texttt{sourceFilePath}): Chuỗi
    đường dẫn tệp tin trên máy tính.
  \item
    Chế độ nhập (\texttt{importMode}): Enum gồm \texttt{OVERWRITE} (ghi
    đè toàn bộ dữ liệu hiện tại) hoặc \texttt{MERGE} (hợp nhất bổ sung
    dữ liệu mới).
  \end{itemize}
\item
  \textbf{Quy tắc xử lý và ràng buộc:}

  \begin{itemize}
  \tightlist
  \item
    Đọc tệp và phân tích cú pháp (Parse JSON hoặc CSV). Nếu tệp bị hỏng
    hoặc sai định dạng cú pháp, hủy bỏ thao tác và thông báo lỗi cấu
    trúc tệp.
  \item
    Thực hiện quy trình xác thực nghiệp vụ 4 bước:

    \begin{enumerate}
    \def\labelenumi{\arabic{enumi}.}
    \tightlist
    \item
      Kiểm tra các trường bắt buộc và định dạng ngày tháng hợp lệ.
    \item
      Kiểm tra tính toàn vẹn danh mục: Mọi \texttt{categoryId} tham
      chiếu phải thuộc về một \texttt{groupId} hợp lệ.
    \item
      Kiểm tra tính toàn vẹn ví: Mọi giao dịch phải trỏ tới một ví tồn
      tại trong danh sách ví.
    \item
      Kiểm tra bất biến số dư không âm theo \texttt{BR-FIN-02}: Đảm bảo
      sau khi nạp dữ liệu, không có ví thông thường nào bị số dư âm.
    \end{enumerate}
  \item
    Nếu phát hiện bất kỳ lỗi vi phạm toàn vẹn nào: Hủy bỏ toàn bộ quá
    trình nhập (rollback), giữ nguyên trạng thái dữ liệu hiện tại và trả
    về danh sách chi tiết các dòng dữ liệu bị lỗi.
  \item
    Nếu dữ liệu hợp lệ: Nạp dữ liệu vào cơ sở dữ liệu và kích hoạt tính
    toán lại toàn bộ số dư ví, quỹ mục tiêu và các bảng thống kê trượt
    12 tháng.
  \end{itemize}
\item
  \textbf{Dữ liệu đầu ra:}

  \begin{itemize}
  \tightlist
  \item
    Dữ liệu hệ thống được phục hồi thành công.
  \item
    Báo cáo kết quả nhập gồm: Tổng số bản ghi xử lý, số bản ghi thêm
    mới, số bản ghi cập nhật.
  \item
    Thông báo hoàn tất quá trình phục hồi dữ liệu.
  \end{itemize}
\item
  \textbf{Ca sử dụng tham chiếu:}
  \texttt{UC10\ -\ Xuất\ nhập\ và\ sao\ lưu\ dữ\ liệu\ hệ\ thống}.
\item
  \textbf{Tiêu chí nghiệm thu:}

  \begin{itemize}
  \tightlist
  \item
    TC-DAT-02.1: Cho tệp JSON sao lưu hợp lệ của toàn bộ hệ thống; khi
    người dùng thực hiện nhập với chế độ \texttt{OVERWRITE}, toàn bộ dữ
    liệu giao dịch và số dư ví được khôi phục nguyên vẹn khớp với thời
    điểm sao lưu.
  \item
    TC-DAT-02.2: Cho tệp nhập chứa một giao dịch tham chiếu tới danh mục
    không tồn tại; hệ thống từ chối nhập, dữ liệu hiện tại được bảo toàn
    100\% và hiển thị cảnh báo lỗi vi phạm khóa ngoại danh mục.
  \end{itemize}
\end{itemize}

\paragraph{4.10.3 Yêu cầu FR-DAT-03: Kiểm tra tính toàn vẹn dữ liệu
(Data Integrity
Audit)}\label{4103-yuxeau-cux1ea7u-fr-dat-03-kiux1ec3m-tra-tuxednh-touxe0n-vux1eb9n-dux1eef-liux1ec7u-data-integrity-audit}

\begin{itemize}
\tightlist
\item
  \textbf{Mã yêu cầu:} \texttt{FR-DAT-03}
\item
  \textbf{Tên chức năng:} Kiểm tra tính toàn vẹn dữ liệu (Data Integrity
  Audit)
\item
  \textbf{Mô tả:} Cung cấp công cụ tự động quét, kiểm toán và đối soát
  toàn bộ các bất biến tài chính trong hệ thống nhằm phát hiện sớm các
  sai lệch số dư, bút toán mất cân bằng hoặc lỗi tham chiếu dữ liệu.
\item
  \textbf{Dữ liệu đầu vào:}

  \begin{itemize}
  \tightlist
  \item
    Lệnh yêu cầu kiểm toán toàn vẹn dữ liệu từ người dùng hoặc tự động
    kích hoạt định kỳ khi khởi động ứng dụng.
  \end{itemize}
\item
  \textbf{Quy tắc xử lý và ràng buộc:}

  \begin{itemize}
  \tightlist
  \item
    Thực hiện kiểm tra tuần tự 4 quy tắc bất biến cốt lõi của hệ thống:

    \begin{enumerate}
    \def\labelenumi{\arabic{enumi}.}
    \tightlist
    \item
      Kiểm tra đối soát số dư ví theo \texttt{BR-FIN-02}: Tính toán lại
      số dư độc lập từ toàn bộ lịch sử giao dịch và so sánh với số dư
      ghi nhận của từng ví. Sai số cho phép tuyệt đối bằng 0
      (\(\Delta = 0\)).
    \item
      Kiểm tra tính cân bằng ghi sổ kép theo \texttt{BR-FIN-03}: Với mọi
      giao dịch chuyển tiền nội bộ, kiểm tra số tiền trừ tại ví nguồn
      phải bằng chính xác số tiền cộng tại ví đích.
    \item
      Kiểm tra tính toàn vẹn danh mục theo \texttt{BR-FIN-04}: Quét toàn
      bộ sổ nhật ký giao dịch để đảm bảo không có giao dịch nào tham
      chiếu tới danh mục không tồn tại hoặc mâu thuẫn về \texttt{Type}.
    \item
      Kiểm tra tính không âm của công nợ theo \texttt{BR-FIN-07}: Xác
      thực số dư nợ phải thu và nợ phải trả của từng cá nhân và toàn hệ
      thống đều lớn hơn hoặc bằng 0.
    \end{enumerate}
  \item
    Tổng hợp kết quả kiểm toán và phân loại mức độ nghiêm trọng (Thông
    tin, Cảnh báo, Lỗi nghiêm trọng).
  \end{itemize}
\item
  \textbf{Dữ liệu đầu ra:}

  \begin{itemize}
  \tightlist
  \item
    Báo cáo kiểm toán toàn vẹn dữ liệu (Data Integrity Audit Report)
    hiển thị trạng thái của từng quy tắc kiểm tra (ĐẠT / KHÔNG ĐẠT).
  \item
    Danh sách chi tiết các sai lệch hoặc bất thường phát hiện được kèm
    hướng dẫn khắc phục.
  \end{itemize}
\item
  \textbf{Ca sử dụng tham chiếu:}
  \texttt{UC10\ -\ Xuất\ nhập\ và\ sao\ lưu\ dữ\ liệu\ hệ\ thống}.
\item
  \textbf{Tiêu chí nghiệm thu:}

  \begin{itemize}
  \tightlist
  \item
    TC-DAT-03.1: Khi cơ sở dữ liệu hoàn toàn nhất quán, công cụ kiểm
    toán trả về kết quả 100\% PASS cho toàn bộ các quy tắc kiểm tra số
    dư, ghi sổ kép và danh mục.
  \item
    TC-DAT-03.2: Khi có một giao dịch ngoài luồng làm sai lệch số dư ví,
    công cụ kiểm toán phát hiện chính xác ví bị lệch, chỉ rõ số tiền
    chênh lệch cụ thể và cảnh báo người dùng thực hiện phục hồi dữ liệu.
  \end{itemize}
\end{itemize}

\clearpage
\subsection{5. YÊU CẦU PHI CHỨC NĂNG (NON-FUNCTIONAL
REQUIREMENTS)}\label{5-yuxeau-cux1ea7u-phi-chux1ee9c-nux103ng-non-functional-requirements}

Phần này đặc tả các yêu cầu phi chức năng (Non-Functional Requirements:
NFR) của Hệ thống Quản lý Tài chính Cá nhân (PFMS) theo khung tiêu chuẩn
chất lượng sản phẩm phần mềm ISO/IEC 25010:2011 và tiêu chuẩn kỹ thuật
đặc tả yêu cầu ISO/IEC/IEEE 29148:2018 §4.

Mỗi yêu cầu phi chức năng được cấu trúc thống nhất theo 5 thành phần kỹ
thuật:

\begin{enumerate}
\def\labelenumi{\arabic{enumi}.}
\tightlist
\item
  \textbf{Mã và Tên NFR:} Định danh duy nhất theo cú pháp
  \texttt{NFR-xx}.
\item
  \textbf{Phân loại thuộc tính chất lượng (Quality Attribute Category):}
  Chiếu theo tiêu chuẩn ISO/IEC 25010.
\item
  \textbf{Mô tả yêu cầu chi tiết (Detailed Description):} Xác định phạm
  vi kỹ thuật, nguyên tắc triển khai và bối cảnh áp dụng.
\item
  \textbf{Chỉ số đo lường định lượng (Quantitative Metrics /
  Threshold):} Xác lập các ngưỡng kỹ thuật có thể đo kiểm khách quan
  bằng công cụ định lượng cụ thể.
\item
  \textbf{Phương pháp kiểm chứng (Verification Method):} Xác định phương
  thức nghiệm thu theo ISO/IEC/IEEE 29148 §7 (Inspection, Test,
  Demonstration, Analysis).
\end{enumerate}

\subsubsection{5.1 NFR-01: Độ chính xác và toàn vẹn dữ liệu tài chính
(Financial Precision \&
Integrity)}\label{51-nfr-01-ux111ux1ed9-chuxednh-xuxe1c-vuxe0-touxe0n-vux1eb9n-dux1eef-liux1ec7u-tuxe0i-chuxednh-financial-precision--integrity}

\begin{itemize}
\tightlist
\item
  \textbf{Mã yêu cầu:} \texttt{NFR-01}
\item
  \textbf{Tên yêu cầu:} Độ chính xác và toàn vẹn dữ liệu tài chính
  (Financial Precision \& Integrity)
\item
  \textbf{Phân loại chất lượng:} Mức độ phù hợp chức năng: Tính chính
  xác (Functional Suitability: Accuracy) và Độ tin cậy: Tính toàn vẹn dữ
  liệu (Reliability: Data Integrity) theo ISO/IEC 25010 §4.2.1 và §4.5.
\item
  \textbf{Mô tả yêu cầu chi tiết:}

  \begin{itemize}
  \tightlist
  \item
    Chuẩn hóa kiểu dữ liệu số học: Tất cả các trường dữ liệu số, các
    biến trung gian và các phép tính số học liên quan đến số tiền giao
    dịch, số dư khả dụng, lãi suất tiết kiệm, số lượng chứng chỉ quỹ,
    đơn giá NAV và tổng tài sản bắt buộc phải sử dụng lớp
    \texttt{java.math.BigDecimal}. Nghiêm cấm tuyệt đối việc khai báo
    hoặc sử dụng các kiểu dữ liệu số thực dấu chấm động nhị phân nguyên
    thủy hoặc lớp bao bọc (\texttt{float}, \texttt{double},
    \texttt{Float}, \texttt{Double}) trong toàn bộ tầng thực thể miền
    (Domain Model), tầng dịch vụ tính toán và tầng lưu trữ. Quy định này
    nhằm loại trừ hoàn toàn sai số làm tròn biểu diễn nhị phân theo đặc
    tả IEEE 754.
  \item
    Quy chuẩn số chữ số thập phân (Scale) và chế độ làm tròn:

    \begin{itemize}
    \tightlist
    \item
      Đối với các giao dịch tiền tệ nội tệ (VND): Giá trị lưu trữ và
      hiển thị sử dụng scale mặc định là 2 chữ số thập phân (hoặc số
      nguyên trong các bảng tổng hợp không có phần thập phân).
    \item
      Đối với số lượng chứng chỉ quỹ đầu tư (\texttt{FundCertificate})
      và đơn giá NAV trên một đơn vị quỹ: Bắt buộc sử dụng scale tối
      thiểu là 4 chữ số thập phân để xử lý chính xác các lệnh mua bán lô
      lẻ và phân bổ tài sản.
    \item
      Chế độ làm tròn số học (Rounding Mode): Khi thực hiện các phép
      toán chia (\texttt{divide}), quy đổi lãi suất ngày
      (\texttt{annualRate\ /\ 365}), hoặc tính thuế phí, hệ thống bắt
      buộc phải chỉ định tường minh tham số chế độ làm tròn
      \texttt{RoundingMode.HALF\_UP} (làm tròn lên đối với nửa đơn vị)
      hoặc \texttt{RoundingMode.HALF\_EVEN} (làm tròn ngân hàng). Nghiêm
      cấm việc gọi hàm chia không chỉ định scale và chế độ làm tròn để
      tránh ngoại lệ \texttt{ArithmeticException} khi gặp số thập phân
      vô hạn tuần hoàn.
    \end{itemize}
  \item
    Kiểm soát sai số tích lũy: Sai số số học phát sinh từ mọi chuỗi phép
    tính số học trong từng giao dịch đơn lẻ cũng như trên toàn bộ sổ cái
    tổng hợp phải bằng chính xác 0.00 VNĐ.
  \item
    Toàn vẹn dữ liệu và bảo toàn số dư theo \texttt{BR-FIN-02}: Tại mọi
    thời điểm trạng thái \(t\), số dư của từng ví tiền phải luôn thỏa
    mãn công thức bảo toàn bất biến:
    \[
    \begin{aligned}
    Balance_{wallet}(t) &= Balance_{wallet}(0)
    + \sum_{k=1}^{n} Income_k
    - \sum_{j=1}^{m} Expense_j \\
    &\quad + \sum_{p=1}^{u} TransferIn_p
    - \sum_{q=1}^{v} TransferOut_q
    \end{aligned}
    \]
  \item
    Toàn vẹn chuyển khoản ghi sổ kép theo \texttt{BR-FIN-03}: Mọi thao
    tác điều chuyển tiền nội bộ giữa hai ví phải thỏa mãn phương trình
    cân bằng ghi sổ kép đối ứng (\(Debit_{src} = Credit_{dst}\)), bảo
    đảm tổng lượng tiền mặt trong hệ thống không bị thay đổi sau giao
    dịch chuyển khoản.
  \end{itemize}
\item
  \textbf{Chỉ số đo lường định lượng (Quantitative Metrics /
  Threshold):}

  \begin{itemize}
  \tightlist
  \item
    Sai số tích lũy số học cho phép: 0.00 VNĐ (tuyệt đối không sai
    lệch).
  \item
    Tỷ lệ vi phạm bất biến bảo toàn số dư \texttt{BR-FIN-02}: 0.00\%
    trên toàn bộ tập dữ liệu kiểm thử.
  \item
    Tỷ lệ xuất hiện kiểu dữ liệu \texttt{float} hoặc \texttt{double}
    trong các lớp thực thể miền (Domain Entities): 0.00\%.
  \item
    Độ chính xác số học: Scale = 2 cho số tiền giao dịch; Scale
    \(\ge 4\) cho số lượng chứng chỉ quỹ và giá trị NAV.
  \end{itemize}
\item
  \textbf{Phương pháp kiểm chứng:}

  \begin{itemize}
  \tightlist
  \item
    Phương pháp: Kiểm thử (Test) kết hợp Soát xét kỹ thuật (Inspection).
  \item
    Kỹ thuật:

    \begin{itemize}
    \tightlist
    \item
      Soát xét kỹ thuật: Sử dụng công cụ phân tích mã nguồn tĩnh
      (Checkstyle, ArchUnit) để quét toàn bộ mã nguồn trong gói
      \texttt{model} và \texttt{domain}, xác nhận không tồn tại bất kỳ
      khai báo kiểu \texttt{float} hoặc \texttt{double} nào.
    \item
      Kiểm thử tự động: Xây dựng bộ kiểm thử đơn vị JUnit 5 thực hiện
      kiểm thử giá trị biên số học (giá trị giao dịch \(10^{14}\) VNĐ,
      giao dịch chia tỷ lệ lãi suất ngày lẻ) và kiểm thử độ bền (stress
      testing) với chuỗi 100,000 giao dịch liên tiếp để đối chiếu số dư
      tính toán với số dư thực tế.
    \end{itemize}
  \end{itemize}
\end{itemize}

\subsubsection{5.2 NFR-02: Hiệu năng và độ trễ phản hồi (Performance \&
Latency)}\label{52-nfr-02-hiux1ec7u-nux103ng-vuxe0-ux111ux1ed9-trux1ec5-phux1ea3n-hux1ed3i-performance--latency}

\begin{itemize}
\tightlist
\item
  \textbf{Mã yêu cầu:} \texttt{NFR-02}
\item
  \textbf{Tên yêu cầu:} Hiệu năng và độ trễ phản hồi (Performance \&
  Latency)
\item
  \textbf{Phân loại chất lượng:} Hiệu quả vận hành: Hành vi thời gian
  (Performance Efficiency: Time Behaviour) và Mức sử dụng tài nguyên
  (Resource Utilization) theo ISO/IEC 25010 §4.4.1 và §4.4.2.
\item
  \textbf{Mô tả yêu cầu chi tiết:}

  \begin{itemize}
  \tightlist
  \item
    Độ trễ ghi nhận giao dịch (Transaction Write Latency): Thời gian từ
    lúc người dùng nhấn nút xác nhận ghi sổ một giao dịch thu chi,
    chuyển khoản hoặc cập nhật nợ cho đến khi hệ thống cập nhật hoàn tất
    vào bộ nhớ, đồng bộ số dư ví và hiển thị thông báo thành công trên
    giao diện phải đạt mức dưới 100 mili-giây (ms).
  \item
    Độ trễ truy vấn và lọc dữ liệu lịch sử (Query Latency): Thời gian
    phản hồi cho thao tác lọc danh sách giao dịch, tìm kiếm từ khóa và
    hiển thị kết quả phân trang trên tập dữ liệu có quy mô \(\ge 2,000\)
    bản ghi giao dịch phải đạt mức dưới 300 mili-giây (ms).
  \item
    Độ trễ tổng hợp báo cáo và trượt chuỗi thời gian 12 tháng: Quá trình
    tính toán bảng thống kê dòng tiền trượt 12 tháng
    (\texttt{MonthlyStat}), tổng hợp cơ cấu chi tiêu theo danh mục 3 cấp
    và tính toán tổng tài sản ròng trên tập dữ liệu 10,000 giao dịch
    phải hoàn tất trong thời gian dưới 500 mili-giây (ms).
  \item
    Thời gian khởi động ứng dụng (Cold Startup Time): Thời gian tính từ
    khi người dùng kích hoạt tệp thực thi của ứng dụng đến khi toàn bộ
    giao diện chính hiển thị đầy đủ và sẵn sàng nhận tương tác không
    được vượt quá 2.0 giây trên phần cứng tiêu chuẩn tham chiếu (CPU
    Dual-core 2.0 GHz, RAM 4GB, ổ cứng SSD).
  \item
    Mức chiếm dụng bộ nhớ RAM (Memory Footprint): Dung lượng bộ nhớ heap
    và non-heap của JVM trong trạng thái hoạt động ổn định khi nạp toàn
    bộ cơ sở dữ liệu sổ cái với 10,000 bản ghi giao dịch không được vượt
    quá 256 MB.
  \end{itemize}
\item
  \textbf{Chỉ số đo lường định lượng (Quantitative Metrics /
  Threshold):}

  \begin{itemize}
  \tightlist
  \item
    Thời gian phản hồi ghi sổ giao dịch mới (\(T_{write}\)): \(\le 100\)
    ms (đo tại bách phân vị 95th percentile).
  \item
    Thời gian phản hồi truy vấn và lọc lịch sử (\(T_{query}\) với số
    lượng giao dịch \(N \ge 2,000\)): \(\le 300\) ms.
  \item
    Thời gian tạo báo cáo dòng tiền trượt 12 tháng (\(T_{report}\) với
    \(N \ge 10,000\) giao dịch): \(\le 500\) ms.
  \item
    Thời gian khởi động nguội ứng dụng (\(T_{startup}\)): \(\le 2.0\)
    giây.
  \item
    Mức chiếm dụng bộ nhớ RAM tối đa (\(Memory_{RAM}\)): \(\le 256\) MB
    trong điều kiện làm việc liên tục.
  \end{itemize}
\item
  \textbf{Phương pháp kiểm chứng:}

  \begin{itemize}
  \tightlist
  \item
    Phương pháp: Kiểm thử (Test) kết hợp Phân tích định lượng
    (Analysis).
  \item
    Kỹ thuật:

    \begin{itemize}
    \tightlist
    \item
      Sử dụng công cụ đo chuẩn hiệu năng vi mô Java Microbenchmark
      Harness (JMH) và \texttt{System.nanoTime()} để đo thời gian thực
      thi của các phương thức cốt lõi.
    \item
      Sử dụng công cụ Java VisualVM và JConsole để giám sát chu kỳ thu
      gom rác (Garbage Collection), phân bổ không gian heap và mức chiếm
      dụng bộ nhớ trong kịch bản nạp tải 10,000 bản ghi giao dịch giả
      lập.
    \end{itemize}
  \end{itemize}
\end{itemize}

\subsubsection{5.3 NFR-03: Độ tin cậy và lưu trữ an toàn (Reliability \&
Persistence)}\label{53-nfr-03-ux111ux1ed9-tin-cux1eady-vuxe0-lux1b0u-trux1eef-an-touxe0n-reliability--persistence}

\begin{itemize}
\tightlist
\item
  \textbf{Mã yêu cầu:} \texttt{NFR-03}
\item
  \textbf{Tên yêu cầu:} Độ tin cậy và lưu trữ an toàn (Reliability \&
  Persistence)
\item
  \textbf{Phân loại chất lượng:} Độ tin cậy: Khả năng chịu lỗi
  (Reliability: Fault Tolerance) và Khả năng phục hồi (Recoverability)
  theo ISO/IEC 25010 §4.5.2 và §4.5.3.
\item
  \textbf{Mô tả yêu cầu chi tiết:}

  \begin{itemize}
  \tightlist
  \item
    Cơ chế ghi tệp an toàn nguyên tử (Atomic File Write): Nhằm ngăn ngừa
    triệt để tình trạng tệp dữ liệu lưu trữ bị hỏng hóc hoặc mất dữ liệu
    khi xảy ra các sự cố đột ngột của môi trường vận hành (mất nguồn
    điện, hệ điều hành bị treo, tiến trình bị cưỡng bức chấm dứt), hệ
    thống bắt buộc phải áp dụng quy trình ghi tệp an toàn 3 bước:

    \begin{enumerate}
    \def\labelenumi{\arabic{enumi}.}
    \tightlist
    \item
      Bước 1: Ghi toàn bộ nội dung dữ liệu cập nhật mới ra một tệp tạm
      thời (staging temporary file với phần mở rộng \texttt{.tmp}) nằm
      trong cùng thư mục lưu trữ.
    \item
      Bước 2: Đồng bộ cưỡng bức toàn bộ khối dữ liệu từ bộ đệm của hệ
      điều hành xuống đĩa vật lý thông qua lời gọi hàm
      \texttt{FileChannel.force(true)} hoặc sử dụng tùy chọn
      \texttt{StandardOpenOption.SYNC}.
    \item
      Bước 3: Thực hiện thao tác tráo đổi nguyên tử (Atomic Rename /
      Move) để thay thế tệp cơ sở dữ liệu chính thức bằng tệp tạm thông
      qua lệnh
      \texttt{Files.move(tempPath,\ targetPath,\ StandardCopyOption.ATOMIC\_MOVE,\ StandardCopyOption.REPLACE\_EXISTING)}.
      Nếu hệ điều hành không hỗ trợ cờ \texttt{ATOMIC\_MOVE}, hệ thống
      thực hiện sao lưu tệp cũ sang tệp dự phòng trước khi thay thế.
    \end{enumerate}
  \item
    Cơ chế tự động sao lưu dữ liệu (Auto-backup): Hệ thống tự động tạo
    bản sao lưu dữ liệu có gắn nhãn thời gian theo định dạng
    \texttt{backup\_YYYYMMDD\_HHMMSS.json} tại các thời điểm:

    \begin{itemize}
    \tightlist
    \item
      Trước khi thực hiện các thao tác có nguy cơ ghi đè lớn như phục
      hồi dữ liệu từ tệp nhập ngoài (\texttt{UC10}) hoặc thiết lập lại
      hệ thống.
    \item
      Tự động duy trì tối thiểu 5 bản sao lưu cục bộ gần nhất tại thư
      mục lưu trữ sao lưu được cô lập, cho phép người dùng quay lui về
      trạng thái an toàn trước đó khi cần thiết.
    \end{itemize}
  \item
    Xử lý ngoại lệ I/O: Mọi thao tác đọc ghi hệ thống tệp bắt buộc phải
    được bọc trong các khối xử lý ngoại lệ chặt chẽ
    (\texttt{IOException}, \texttt{SecurityException}). Trong trường hợp
    ổ đĩa bị đầy hoặc thư mục bị khóa quyền ghi, hệ thống phải giữ
    nguyên trạng tệp dữ liệu gốc và thông báo lỗi rõ ràng tới người
    dùng.
  \end{itemize}
\item
  \textbf{Chỉ số đo lường định lượng (Quantitative Metrics /
  Threshold):}

  \begin{itemize}
  \tightlist
  \item
    Tỷ lệ hỏng tệp dữ liệu chính khi ngắt tiến trình đột ngột trong lúc
    ghi: 0.00\%.
  \item
    Thời gian thực hiện thao tác tráo đổi tệp nguyên tử: \(\le 50\) ms.
  \item
    Số lượng bản sao lưu tự động tối thiểu được duy trì: 5 bản ghi gần
    nhất.
  \item
    Thời gian tạo bản sao lưu tự động: \(\le 500\) ms đối với kích thước
    tệp dữ liệu lên đến 10 MB.
  \item
    Tỷ lệ thất thoát dữ liệu ngoài giao dịch đang thực thi dở dang:
    0.00\%.
  \end{itemize}
\item
  \textbf{Phương pháp kiểm chứng:}

  \begin{itemize}
  \tightlist
  \item
    Phương pháp: Kiểm thử (Test) kết hợp Soát xét kỹ thuật (Inspection).
  \item
    Kỹ thuật:

    \begin{itemize}
    \tightlist
    \item
      Thực hiện kiểm thử tiêm lỗi (Fault Injection Testing): Viết chương
      trình mô phỏng quá trình ghi tệp dữ liệu quy mô lớn, đồng thời gửi
      tín hiệu cưỡng bức ngắt tiến trình (lệnh \texttt{taskkill\ /F}
      trên Windows hoặc \texttt{kill\ -9} trên Linux) tại thời điểm ghi
      tệp tạm; sau đó khởi động lại ứng dụng để xác minh tệp dữ liệu
      chính vẫn nguyên vẹn và tải dữ liệu thành công.
    \item
      Soát xét mã nguồn kiểm tra sự hiện diện đầy đủ của cờ
      \texttt{ATOMIC\_MOVE} và cơ chế đồng bộ đĩa trong tầng lưu trữ dữ
      liệu.
    \end{itemize}
  \end{itemize}
\end{itemize}

\subsubsection{5.4 NFR-04: Khả năng bảo trì và kiến trúc hướng đối tượng
(Maintainability \& OOP
Design)}\label{54-nfr-04-khux1ea3-nux103ng-bux1ea3o-truxec-vuxe0-kiux1ebfn-truxfac-hux1b0ux1edbng-ux111ux1ed1i-tux1b0ux1ee3ng-maintainability--oop-design}

\begin{itemize}
\tightlist
\item
  \textbf{Mã yêu cầu:} \texttt{NFR-04}
\item
  \textbf{Tên yêu cầu:} Khả năng bảo trì và kiến trúc hướng đối tượng
  (Maintainability \& OOP Design)
\item
  \textbf{Phân loại chất lượng:} Khả năng bảo trì: Tính module hóa
  (Modularity), Tính tái sử dụng (Reusability), Tính kiểm thử được
  (Testability) và Tính phân tích được (Analysability) theo ISO/IEC
  25010 §4.7.
\item
  \textbf{Mô tả yêu cầu chi tiết:}

  \begin{itemize}
  \tightlist
  \item
    Kiến trúc hướng đối tượng đáp ứng tiêu chuẩn kỹ thuật phần mềm:

    \begin{itemize}
    \tightlist
    \item
      Tính đóng gói cao (Encapsulation): 100\% thuộc tính của các thực
      thể nghiệp vụ tài chính phải được khai báo với mức truy cập
      \texttt{private}. Dữ liệu chỉ được truy xuất qua các phương thức
      getter hoặc các phương thức thực hiện hành vi nghiệp vụ cụ thể.
      Tuyệt đối không cung cấp các phương thức setter công khai cho phép
      tự do gán lại số dư ví mà không thông qua giao dịch thu chi hoặc
      chuyển khoản.
    \item
      Sao chép phòng thủ (Defensive Copying): Đối với các đối tượng có
      khả năng biến đổi (mutable objects) như danh sách tập hợp
      (\texttt{List}, \texttt{Map}) hoặc đối tượng ngày tháng, các
      phương thức khởi tạo và getter phải trả về bản sao hoặc danh sách
      không thể sửa đổi (\texttt{Collections.unmodifiableList}) nhằm
      ngăn chặn việc sửa đổi dữ liệu trái phép từ bên ngoài lớp.
    \item
      Nguyên lý đơn trách nhiệm (Single Responsibility Principle - SRP):
      Mỗi lớp thực thể chỉ đảm nhiệm một trách nhiệm duy nhất. Các lớp
      thực thể miền (\texttt{Wallet}, \texttt{Transaction},
      \texttt{Debt}, \texttt{SavingAccount}, \texttt{FundCertificate},
      \texttt{FinancialGoal}) chỉ quản lý trạng thái nghiệp vụ tự thân,
      tách biệt hoàn toàn khỏi logic biểu diễn giao diện và cơ chế lưu
      trữ tệp.
    \item
      Kế thừa và Đa hình hợp lý (Inheritance \& Polymorphism): Áp dụng
      kế thừa và đa hình tại các phân hệ phù hợp (như kế thừa giao dịch,
      giao diện tài sản tài chính \texttt{Asset} cho Ví, Sổ tiết kiệm và
      Chứng chỉ quỹ), chuẩn bị cho giai đoạn phân tích OOA và thiết kế
      kiến trúc OOD.
    \end{itemize}
  \item
    Độ bao phủ kiểm thử tự động (Unit Test Coverage): Toàn bộ các quy
    tắc nghiệp vụ tài chính, thuật toán tính số dư và công thức tính lãi
    suất bắt buộc phải có kiểm thử đơn vị tự động với khung kiểm thử
    JUnit 5. Độ bao phủ mã nguồn theo dòng lệnh (Line Coverage) của tầng
    miền nghiệp vụ phải đạt tối thiểu 80\%, và độ bao phủ nhánh điều
    kiện (Branch Coverage) phải đạt tối thiểu 75\%.
  \item
    Độ phức tạp chu trình mã nguồn (Cyclomatic Complexity): Độ phức tạp
    chu trình của mỗi phương thức nghiệp vụ không được vượt quá 10
    (\(V(G) \le 10\)).
  \item
    Khả năng mở rộng không sửa đổi (Open/Closed Principle): Cấu trúc
    phân loại danh mục và các lớp thực thể tài sản cho phép mở rộng các
    danh mục hoặc loại tài sản mới mà không cần chỉnh sửa mã nguồn của
    công cụ tính toán số dư.
  \end{itemize}
\item
  \textbf{Chỉ số đo lường định lượng (Quantitative Metrics /
  Threshold):}

  \begin{itemize}
  \tightlist
  \item
    Độ bao phủ dòng lệnh kiểm thử đơn vị (Line Coverage) tầng nghiệp vụ:
    \(\ge 80\%\).
  \item
    Độ bao phủ nhánh điều kiện (Branch Coverage) tầng nghiệp vụ:
    \(\ge 75\%\).
  \item
    Độ phức tạp chu trình của phương thức (\(V(G)\)): \(V(G) \le 10\)
    trên 100\% các phương thức nghiệp vụ.
  \item
    Mức độ phụ thuộc vòng giữa các gói mã nguồn (Cyclic Package
    Dependencies): 0 vòng phụ thuộc.
  \item
    Tỷ lệ áp dụng sao chép phòng thủ đối với các tập hợp dữ liệu biến
    đổi: 100\%.
  \end{itemize}
\item
  \textbf{Phương pháp kiểm chứng:}

  \begin{itemize}
  \tightlist
  \item
    Phương pháp: Phân tích định lượng (Analysis) kết hợp Soát xét kỹ
    thuật (Inspection).
  \item
    Kỹ thuật:

    \begin{itemize}
    \tightlist
    \item
      Sử dụng công cụ đo độ bao phủ mã nguồn JaCoCo tích hợp trong tiến
      trình build Maven để trích xuất báo cáo tỷ lệ bao phủ dòng lệnh và
      nhánh.
    \item
      Sử dụng các công cụ phân tích tĩnh (SonarQube, Checkstyle, Metrics
      Reloaded) để quét mã nguồn kiểm tra độ phức tạp chu trình \(V(G)\)
      và phát hiện các vi phạm quy tắc đóng gói hoặc phụ thuộc vòng.
    \end{itemize}
  \end{itemize}
\end{itemize}

\subsubsection{5.5 NFR-05: Tính khả dụng và công thái học (Usability \&
User
Ergonomics)}\label{55-nfr-05-tuxednh-khux1ea3-dux1ee5ng-vuxe0-cuxf4ng-thuxe1i-hux1ecdc-usability--user-ergonomics}

\begin{itemize}
\tightlist
\item
  \textbf{Mã yêu cầu:} \texttt{NFR-05}
\item
  \textbf{Tên yêu cầu:} Tính khả dụng và công thái học (Usability \&
  User Ergonomics)
\item
  \textbf{Phân loại chất lượng:} Tính khả dụng: Tính dễ học
  (Learnability), Tính dễ vận hành (Operability) và Bảo vệ người dùng
  khỏi lỗi (User Error Protection) theo ISO/IEC 25010 §4.3.
\item
  \textbf{Mô tả yêu cầu chi tiết:}

  \begin{itemize}
  \tightlist
  \item
    Công thái học và tối ưu hóa thao tác nhập liệu:

    \begin{itemize}
    \tightlist
    \item
      Tự động điền giá trị mặc định: Trường ngày giao dịch trong các
      biểu mẫu ghi nhận thu chi, chuyển tiền mặc định điền ngày hiện tại
      của hệ thống (\texttt{LocalDate.now()}). Người dùng chỉ cần can
      thiệp khi ghi nhận giao dịch phát sinh trong quá khứ.
    \item
      Định vị con trỏ tự động: Khi người dùng mở hộp thoại ghi nhận giao
      dịch, con trỏ chuột tự động focus vào ô nhập số tiền để sẵn sàng
      nhập số liệu ngay lập tức.
    \item
      Cây danh mục gợi ý theo ngữ cảnh: Hệ thống tự động lọc danh mục
      thu chi theo phân loại ngữ cảnh (khi chọn giao dịch Thu nhập chỉ
      hiển thị các danh mục thuộc loại Thu nhập; khi chọn Chi tiêu chỉ
      hiển thị các danh mục thuộc loại Chi tiêu) giúp loại trừ lỗi phân
      loại sai.
    \item
      Định dạng số tiền trực quan theo thời gian thực: Khi người dùng gõ
      các con số vào ô nhập tiền tệ, hệ thống tự động định dạng phân
      tách hàng nghìn bằng dấu phẩy theo chuẩn Việt Nam (ví dụ: gõ
      \texttt{1500000} tự động hiển thị \texttt{1,500,000\ VNĐ}) giúp
      hạn chế sai sót nhập thừa hoặc thiếu chữ số 0.
    \end{itemize}
  \item
    Thông báo lỗi người dùng rõ ràng và có tính định hướng khắc phục:

    \begin{itemize}
    \tightlist
    \item
      Toàn bộ thông báo lỗi xác thực dữ liệu đầu vào và vi phạm quy tắc
      nghiệp vụ phải được thể hiện bằng ngôn ngữ tiếng Việt tự nhiên,
      ngắn gọn, chỉ rõ nguyên nhân phát sinh lỗi và đưa ra hướng xử lý
      khắc phục cụ thể. Ví dụ: khi số dư ví không đủ chi trả khoản chi,
      thông báo hiển thị: "Số dư ví MB Bank hiện tại là 150,000 VNĐ,
      không đủ để chi trả khoản chi 200,000 VNĐ. Vui lòng giảm số tiền
      chi tiêu hoặc chuyển tiền từ ví khác sang".
    \item
      Tuyệt đối không hiển thị các thông báo lỗi kỹ thuật thuần túy, vết
      ngăn xếp ngoại lệ (stack trace) hoặc các mã lỗi hệ thống khó hiểu
      ra giao diện người dùng.
    \end{itemize}
  \end{itemize}
\item
  \textbf{Chỉ số đo lường định lượng (Quantitative Metrics /
  Threshold):}

  \begin{itemize}
  \tightlist
  \item
    Số lượng tương tác tối thiểu để ghi nhận một giao dịch chi tiêu tiêu
    chuẩn: \(\le 3\) hành động (chọn danh mục, nhập số tiền, nhấn lưu;
    ngày giao dịch và ví mặc định được tự động điền).
  \item
    Tỷ lệ thông báo lỗi nghiệp vụ hiển thị tiếng Việt rõ ràng kèm hướng
    dẫn khắc phục: 100\%.
  \item
    Tỷ lệ hiển thị lỗi ngoại lệ chưa bắt (Unhandled exceptions/stack
    traces) trên màn hình: 0.00\%.
  \item
    Thời gian làm quen hệ thống: Người dùng mới có thể hoàn thành thao
    tác tạo ví và ghi nhận giao dịch đầu tiên trong vòng dưới 2 phút mà
    không cần tài liệu hướng dẫn sử dụng.
  \end{itemize}
\item
  \textbf{Phương pháp kiểm chứng:}

  \begin{itemize}
  \tightlist
  \item
    Phương pháp: Trình diễn (Demonstration) kết hợp Soát xét kỹ thuật
    (Inspection).
  \item
    Kỹ thuật:

    \begin{itemize}
    \tightlist
    \item
      Tiến hành thử nghiệm quan sát thao tác người dùng (Usability
      Testing) với 3 kịch bản sử dụng phổ biến, bấm giờ hoàn thành tác
      vụ và đo tỷ lệ nhập liệu chính xác.
    \item
      Soát xét toàn bộ bảng thông điệp lỗi (Resource Bundle / Error
      Catalog) trong mã nguồn để kiểm tra ngôn ngữ và tính đầy đủ của
      hướng dẫn khắc phục.
    \end{itemize}
  \end{itemize}
\end{itemize}

\subsubsection{5.6 NFR-06: Tính tương thích và môi trường thực thi
(Portability \&
Compatibility)}\label{56-nfr-06-tuxednh-tux1b0ux1a1ng-thuxedch-vuxe0-muxf4i-trux1b0ux1eddng-thux1ef1c-thi-portability--compatibility}

\begin{itemize}
\tightlist
\item
  \textbf{Mã yêu cầu:} \texttt{NFR-06}
\item
  \textbf{Tên yêu cầu:} Tính tương thích và môi trường thực thi
  (Portability \& Compatibility)
\item
  \textbf{Phân loại chất lượng:} Tính khả chuyển: Tính thích ứng
  (Portability: Adaptability) và Tính dễ cài đặt (Installability) theo
  ISO/IEC 25010 §4.8.
\item
  \textbf{Mô tả yêu cầu chi tiết:}

  \begin{itemize}
  \tightlist
  \item
    Môi trường thực thi tiêu chuẩn: Ứng dụng chạy trên nền tảng máy ảo
    Java Virtual Machine (JVM) tương thích với phiên bản Java
    Development Kit (JDK) 17 Long-Term Support (LTS) trở lên, tận dụng
    các cấu trúc ngôn ngữ hiện đại như Record classes, Pattern Matching
    for switch và thư viện thời gian chuẩn \texttt{java.time}.
  \item
    Khả năng độc lập nền tảng hệ điều hành (Cross-Platform
    Compatibility): Ứng dụng có thể đóng gói độc lập và thực thi thống
    nhất trên cả 3 hệ điều hành máy tính cá nhân phổ biến: Microsoft
    Windows 10/11 (64-bit), macOS (kiến trúc Intel x86\_64 và Apple
    Silicon ARM64), và các bản phân phối Linux thông dụng (Ubuntu
    20.04+, Fedora 36+, Debian 11+) mà không cần thay đổi bất kỳ dòng mã
    nguồn nào hoặc biên dịch lại theo từng hệ điều hành.
  \item
    Chuẩn hóa thao tác đường dẫn hệ thống tệp: Tuyệt đối không viết cứng
    (hardcode) các ký tự phân cách thư mục đặc thù của hệ điều hành như
    dấu gạch chéo ngược Windows \texttt{\textbackslash{}} hoặc dấu gạch
    chéo Unix \texttt{/}. Toàn bộ thao tác định vị tệp và thư mục phải
    sử dụng lớp trừu tượng \texttt{java.nio.file.Path} và phương thức
    \texttt{Paths.get()}, hoặc sử dụng hằng số \texttt{File.separator}.
  \item
    Tiêu chuẩn mã hóa ký tự (Character Encoding): Mọi thao tác đọc, ghi
    tệp dữ liệu cục bộ, tệp xuất dữ liệu JSON và xuất báo cáo đều phải
    chỉ định tường minh bảng mã UTF-8 thông qua
    \texttt{StandardCharsets.UTF\_8}. Điều này bảo đảm dữ liệu tiếng
    Việt có dấu luôn được lưu trữ và hiển thị chính xác, không bị lỗi vỡ
    phông chữ (mojibake) bất kể bảng mã mặc định của hệ điều hành máy
    trạm là gì.
  \item
    Không phụ thuộc vào thư viện nhị phân ngoại lai: Ứng dụng không sử
    dụng bất kỳ thư viện nhị phân phụ thuộc nền tảng nào (0 thư viện
    JNI/JNA), bảo đảm việc đóng gói và phân phối phần mềm được thực hiện
    hoàn toàn dưới dạng tệp Java Archive (\texttt{.jar}) hoặc gói ứng
    dụng độc lập jpackage.
  \end{itemize}
\item
  \textbf{Chỉ số đo lường định lượng (Quantitative Metrics /
  Threshold):}

  \begin{itemize}
  \tightlist
  \item
    Tỷ lệ vượt qua bộ kiểm thử tự động trên cả 3 hệ điều hành (Windows,
    macOS, Linux): 100\% (toàn bộ các ca kiểm thử đạt kết quả tương
    đương).
  \item
    Tỷ lệ thao tác đọc ghi tệp chỉ định tường minh bảng mã UTF-8: 100\%.
  \item
    Số lượng thư viện liên kết nhị phân nền tảng (Native library
    dependencies): 0 thư viện.
  \item
    Phiên bản Java hỗ trợ tối thiểu: Java SE 17 LTS.
  \end{itemize}
\item
  \textbf{Phương pháp kiểm chứng:}

  \begin{itemize}
  \tightlist
  \item
    Phương pháp: Kiểm thử (Test) kết hợp Soát xét kỹ thuật (Inspection).
  \item
    Kỹ thuật:

    \begin{itemize}
    \tightlist
    \item
      Thực thi bộ kiểm thử đơn vị tự động trên cả 3 môi trường hệ điều
      hành thực tế (Windows, macOS, Linux) thông qua quy trình tích hợp
      liên tục (CI Workflow) để bảo đảm tính tương đương của kết quả
      chạy thử.
    \item
      Quét soát xét mã nguồn tĩnh đối với toàn bộ các lời gọi hàm I/O
      (\texttt{Files.newBufferedReader},
      \texttt{Files.newBufferedWriter}, \texttt{FileReader},
      \texttt{FileWriter}) để bảo đảm luôn truyền tham số
      \texttt{StandardCharsets.UTF\_8} và sử dụng \texttt{Path} chuẩn.
    \end{itemize}
  \end{itemize}
\end{itemize}

\clearpage
\subsection{6. ĐẶC TẢ 10 CA SỬ DỤNG CHI TIẾT (FULLY DRESSED USE
CASES)}\label{6-ux111ux1eb7c-tux1ea3-10-ca-sux1eed-dux1ee5ng-chi-tiux1ebft-fully-dressed-use-cases}

Phần này đặc tả chi tiết 10 Ca sử dụng (Use Cases) của Hệ thống Quản lý
Tài chính Cá nhân (PFMS) theo định dạng Fully Dressed do Craig Larman đề
xuất trong tài liệu chuẩn \emph{Applying UML and Patterns}. Mỗi ca sử
dụng bao gồm đầy đủ 10 trường thông tin kỹ thuật:

\begin{enumerate}
\def\labelenumi{\arabic{enumi}.}
\tightlist
\item
  Mã và Tên ca sử dụng
\item
  Tác nhân chính (Primary Actor: Owner)
\item
  Mô tả tóm tắt (Brief Description)
\item
  Điều kiện kích hoạt (Trigger)
\item
  Tiền điều kiện (Preconditions)
\item
  Hậu điều kiện thành công (Success Postconditions)
\item
  Luồng sự kiện chính (Main Success Scenario)
\item
  Các luồng rẽ nhánh (Alternative Flows)
\item
  Các luồng ngoại lệ (Exception Flows)
\item
  Quy tắc nghiệp vụ đặc thù (Special Requirements / Business Invariants)
\end{enumerate}

Toàn bộ các ca sử dụng được thiết kế xoay quanh tác nhân duy nhất là Chủ
sở hữu hệ thống (Owner), tuân thủ nghiêm ngặt các bất biến tài chính
được xác lập tại Mục 3 và các yêu cầu chức năng tại Mục 4.

\subsubsection{6.1 UC01: Quản lý ví tiền (Wallet
Management)}\label{61-uc01-quux1ea3n-luxfd-vuxed-tiux1ec1n-wallet-management}

\begin{itemize}
\tightlist
\item
  \textbf{Mã và Tên ca sử dụng:}
  \texttt{UC01:\ Quản\ lý\ ví\ tiền\ (Wallet\ Management)}
\item
  \textbf{Tác nhân chính:} Chủ sở hữu hệ thống (Owner)
\item
  \textbf{Mô tả tóm tắt:} Ca sử dụng này cho phép Owner khởi tạo ví tài
  chính mới (tiền mặt, tài khoản ngân hàng, ví điện tử, thẻ tín dụng),
  khai báo số dư ban đầu, theo dõi danh mục số dư động được tính toán
  tức thời từ toàn bộ lịch sử dòng tiền, xem phân tích cơ cấu tỷ trọng
  tài sản thanh khoản giữa các ví, và cập nhật thông tin ví (đổi tên,
  chuyển trạng thái kích hoạt hoặc lưu trữ).
\item
  \textbf{Điều kiện kích hoạt:} Owner truy cập phân hệ Quản lý ví tiền
  trên giao diện ứng dụng để thiết lập nguồn tiền mới hoặc tra cứu, tái
  cơ cấu danh mục tài sản thanh khoản.
\item
  \textbf{Tiền điều kiện:}

  \begin{enumerate}
  \def\labelenumi{\arabic{enumi}.}
  \tightlist
  \item
    Hệ thống PFMS đã khởi chạy và kết nối thành công với cơ sở dữ liệu
    cục bộ.
  \item
    Owner đang ở màn hình phân hệ Quản lý ví tiền.
  \end{enumerate}
\item
  \textbf{Hậu điều kiện thành công:}

  \begin{enumerate}
  \def\labelenumi{\arabic{enumi}.}
  \tightlist
  \item
    Ví tiền mới được lưu trữ trong cơ sở dữ liệu với trạng thái
    \texttt{ACTIVE} và số dư ban đầu hợp lệ.
  \item
    Bất biến số dư ví theo \texttt{BR-FIN-02} được bảo đảm; tổng tài sản
    thanh khoản toàn hệ thống được cập nhật chính xác.
  \item
    Bảng cơ cấu phân bổ tài sản phản ánh đúng tỷ trọng của từng ví theo
    \texttt{FR-WAL-04}.
  \end{enumerate}
\item
  \textbf{Luồng sự kiện chính (Main Success Scenario):}

  \begin{enumerate}
  \def\labelenumi{\arabic{enumi}.}
  \tightlist
  \item
    Người dùng chọn chức năng "Thêm ví mới" trên giao diện quản lý ví.
  \item
    Hệ thống hiển thị biểu mẫu tạo ví gồm các trường: Tên ví, Loại ví
    (\texttt{CASH}, \texttt{BANK}, \texttt{E\_WALLET}, \texttt{CREDIT}),
    Số dư ban đầu, Ghi chú mô tả.
  \item
    Người dùng nhập tên ví, chọn loại ví từ danh sách thả xuống, nhập số
    dư ban đầu (hoặc giữ giá trị mặc định là 0), và nhập ghi chú mô tả
    (nếu có).
  \item
    Người dùng nhấn nút "Xác nhận tạo ví".
  \item
    Hệ thống xác thực dữ liệu đầu vào: kiểm tra tên ví không rỗng và
    không trùng lặp với tên ví đã có trong hệ thống (không phân biệt chữ
    hoa, chữ thường sau khi cắt khoảng trắng thừa), kiểm tra số dư ban
    đầu không âm đối với các ví thuộc loại \texttt{CASH}, \texttt{BANK},
    \texttt{E\_WALLET}.
  \item
    Hệ thống tạo đối tượng ví mới trong cơ sở dữ liệu với mã định danh
    duy nhất (\texttt{walletId}), gán trạng thái \texttt{ACTIVE}, chuẩn
    hóa số dư ban đầu về kiểu \texttt{BigDecimal} có scale là 2 chữ số
    thập phân theo quy tắc \texttt{BR-FIN-01}.
  \item
    Hệ thống tính toán lại số dư động của ví mới và tổng số dư thanh
    khoản toàn hệ thống theo quy tắc \texttt{BR-FIN-02}.
  \item
    Hệ thống cập nhật bảng tỷ trọng cơ cấu tài sản theo từng ví theo quy
    tắc \texttt{FR-WAL-04}.
  \item
    Hệ thống thông báo tạo ví thành công, cập nhật danh sách ví trên màn
    hình và làm mới bảng cơ cấu tài sản.
  \end{enumerate}
\item
  \textbf{Các luồng rẽ nhánh (Alternative Flows):}

  \begin{itemize}
  \tightlist
  \item
    \emph{2a. Xem danh sách ví và chi tiết số dư động:}

    \begin{itemize}
    \tightlist
    \item
      2a1. Người dùng mở màn hình danh sách ví, tùy chọn áp dụng bộ lọc
      trạng thái (\texttt{ACTIVE}, \texttt{ARCHIVED}, hoặc
      \texttt{ALL}).
    \item
      2a2. Hệ thống quét toàn bộ lịch sử dòng tiền của từng ví thỏa mãn
      điều kiện lọc, áp dụng công thức \texttt{BR-FIN-02} để tính số dư
      khả dụng thực tế tại thời điểm truy vấn.
    \item
      2a3. Hệ thống tổng hợp tổng số dư tiền mặt thanh khoản bằng phương
      thức \texttt{BigDecimal.add()}.
    \item
      2a4. Hệ thống hiển thị danh sách ví kèm số dư động, loại ví, trạng
      thái và tổng số dư khả dụng toàn hệ thống ở chân danh sách.
    \end{itemize}
  \item
    \emph{2b. Cập nhật thông tin ví tiền:}

    \begin{itemize}
    \tightlist
    \item
      2b1. Người dùng chọn một ví cụ thể trong danh sách và nhấn nút
      "Chỉnh sửa".
    \item
      2b2. Hệ thống hiển thị biểu mẫu chỉnh sửa cho phép sửa Tên ví, Ghi
      chú mô tả, và Trạng thái hoạt động (\texttt{ACTIVE} hoặc
      \texttt{ARCHIVED}). Trường số dư bị khóa hoàn toàn, không cho phép
      can thiệp trực tiếp.
    \item
      2b3. Người dùng sửa đổi thông tin mong muốn và nhấn "Lưu thay
      đổi".
    \item
      2b4. Hệ thống kiểm tra tên ví mới không trùng lặp với các ví khác
      đang tồn tại.
    \item
      2b5. Hệ thống ghi nhận thông tin cập nhật vào cơ sở dữ liệu và
      hiển thị thông báo thành công.
    \end{itemize}
  \item
    \emph{2c. Lưu trữ ví không còn sử dụng (Archived):}

    \begin{itemize}
    \tightlist
    \item
      2c1. Tại bước 2b3, Người dùng chuyển trạng thái ví từ
      \texttt{ACTIVE} sang \texttt{ARCHIVED}.
    \item
      2c2. Hệ thống kiểm tra số dư khả dụng hiện tại của ví. Nếu số dư
      khác 0, hệ thống hiển thị hộp thoại cảnh báo: ví vẫn còn số dư,
      khuyến nghị điều chuyển hết tiền trước khi lưu trữ, nhưng vẫn cho
      phép Người dùng lựa chọn tiếp tục hoặc hủy bỏ.
    \item
      2c3. Người dùng xác nhận tiếp tục lưu trữ.
    \item
      2c4. Hệ thống cập nhật trạng thái ví thành \texttt{ARCHIVED}, ẩn
      ví này khỏi danh sách chọn ví trong các chức năng tạo giao dịch
      mới, nhưng tiếp tục duy trì dữ liệu trong các báo cáo lịch sử.
    \end{itemize}
  \item
    \emph{2d. Xem phân tích cơ cấu tài sản thanh khoản:}

    \begin{itemize}
    \tightlist
    \item
      2d1. Người dùng chọn tab "Cơ cấu phân bổ ví".
    \item
      2d2. Hệ thống tính tổng số dư tất cả các ví đang \texttt{ACTIVE}:
      \(TotalWallets = \sum Balance_w\).
    \item
      2d3. Nếu \(TotalWallets > 0\), hệ thống tính tỷ trọng phần trăm
      từng ví theo công thức:
      \(Rate_w = (Balance_w / TotalWallets) \times 100\%\), sử dụng
      \texttt{BigDecimal.divide()} với \texttt{RoundingMode.HALF\_EVEN}
      và scale tối thiểu 4 chữ số thập phân, định dạng 2 chữ số thập
      phân ra giao diện.
    \item
      2d4. Hệ thống sắp xếp các ví theo thứ tự tỷ trọng giảm dần và hiển
      thị dưới dạng bảng dữ liệu kèm biểu đồ trực quan.
    \end{itemize}
  \end{itemize}
\item
  \textbf{Các luồng ngoại lệ (Exception Flows):}

  \begin{itemize}
  \tightlist
  \item
    \emph{5a. Trùng tên ví đã tồn tại:}

    \begin{itemize}
    \tightlist
    \item
      5a1. Tại bước 5 (hoặc 2b4), hệ thống phát hiện tên ví đã tồn tại
      trong cơ sở dữ liệu (so sánh không phân biệt chữ hoa, chữ thường
      sau khi cắt khoảng trắng thừa).
    \item
      5a2. Hệ thống hủy bỏ thao tác ghi dữ liệu, đánh dấu trường tên ví
      bị lỗi và hiển thị thông báo: "Tên ví này đã tồn tại trong hệ
      thống. Vui lòng chọn một tên khác".
    \item
      5a3. Người dùng nhập lại tên ví hợp lệ hoặc nhấn nút hủy thao tác.
    \end{itemize}
  \item
    \emph{5b. Số dư khởi tạo ban đầu là số âm cho ví thông thường:}

    \begin{itemize}
    \tightlist
    \item
      5b1. Tại bước 5, hệ thống phát hiện ví thuộc loại \texttt{CASH},
      \texttt{BANK} hoặc \texttt{E\_WALLET} nhưng số dư ban đầu nhập vào
      nhỏ hơn 0 (\(initialBalance < 0\)).
    \item
      5b2. Hệ thống chặn thao tác tạo ví và hiển thị thông báo lỗi: "Số
      dư khởi tạo của ví tiền mặt, ngân hàng hoặc ví điện tử không được
      nhỏ hơn 0".
    \item
      5b3. Người dùng điều chỉnh lại số dư ban đầu lớn hơn hoặc bằng 0.
    \end{itemize}
  \item
    \emph{5c. Cố ý chỉnh sửa trực tiếp số dư ví:}

    \begin{itemize}
    \tightlist
    \item
      5c1. Nếu hệ thống phát hiện yêu cầu sửa đổi trực tiếp trường số dư
      khả dụng (\texttt{currentBalance} hoặc \texttt{initialBalance})
      của một ví đã có giao dịch lịch sử.
    \item
      5c2. Hệ thống từ chối thực thi, ném ngoại lệ vi phạm bất biến dữ
      liệu tài chính và ghi nhật ký cảnh báo an toàn.
    \end{itemize}
  \item
    \emph{2d3a. Tổng số dư toàn bộ các ví bằng 0 khi tính cơ cấu:}

    \begin{itemize}
    \tightlist
    \item
      2d3a1. Tại bước 2d3, hệ thống ghi nhận \(TotalWallets = 0\).
    \item
      2d3a2. Hệ thống gán tỷ trọng của toàn bộ các ví bằng đúng 0.00\%,
      tuyệt đối không thực hiện phép chia cho 0 nhằm tránh lỗi
      \texttt{ArithmeticException}.
    \end{itemize}
  \end{itemize}
\item
  \textbf{Quy tắc nghiệp vụ đặc thù (Special Requirements /
  Invariants):}

  \begin{enumerate}
  \def\labelenumi{\arabic{enumi}.}
  \tightlist
  \item
    Bất biến bảo toàn số dư theo \texttt{BR-FIN-02}: Số dư của ví luôn
    bằng số dư ban đầu cộng tổng thu nhập, chuyển đến, đi vay, thu nợ
    trừ tổng chi tiêu, chuyển đi, trả nợ, cho vay.
  \item
    Bất biến tính toán tài chính \texttt{BR-FIN-01}: Bắt buộc sử dụng
    \texttt{BigDecimal}, cấm tuyệt đối dùng kiểu số thực \texttt{double}
    hoặc \texttt{float}.
  \item
    Bất biến toàn vẹn dữ liệu: Không hỗ trợ xóa vật lý (hard delete) ví
    tiền nếu ví đó đã phát sinh bất kỳ giao dịch nào trong sổ nhật ký
    giao dịch, sổ chuyển tiền hoặc sổ công nợ. Chỉ hỗ trợ chuyển trạng
    thái \texttt{ARCHIVED}.
  \item
    Phù hợp các yêu cầu chức năng: \texttt{FR-WAL-01},
    \texttt{FR-WAL-02}, \texttt{FR-WAL-03}, \texttt{FR-WAL-04}.
  \end{enumerate}
\end{itemize}

\subsubsection{6.2 UC02: Ghi nhận và quản lý giao dịch thu chi (Income
\& Expense
Management)}\label{62-uc02-ghi-nhux1eadn-vuxe0-quux1ea3n-luxfd-giao-dux1ecbch-thu-chi-income--expense-management}

\begin{itemize}
\tightlist
\item
  \textbf{Mã và Tên ca sử dụng:}
  \texttt{UC02:\ Ghi\ nhận\ và\ quản\ lý\ giao\ dịch\ thu\ chi\ (Income\ \&\ Expense\ Management)}
\item
  \textbf{Tác nhân chính:} Chủ sở hữu hệ thống (Owner)
\item
  \textbf{Mô tả tóm tắt:} Ca sử dụng này cho phép Owner ghi nhận các
  phát sinh thu nhập và chi tiêu trong đời sống hàng ngày, phân loại
  chuẩn hóa theo cấu trúc cây danh mục 3 cấp, liên kết với một ví thanh
  toán cụ thể; tra cứu, lọc lịch sử giao dịch đa tiêu chí; chỉnh sửa
  hoặc xóa giao dịch với cơ chế tự động hoàn nguyên số dư ví bảo toàn
  tính toàn vẹn sổ cái.
\item
  \textbf{Điều kiện kích hoạt:} Owner phát sinh nhu cầu ghi chép một
  khoản thu chi mới hoặc tra cứu, điều chỉnh các giao dịch đã thực hiện
  trong sổ nhật ký tài chính.
\item
  \textbf{Tiền điều kiện:}

  \begin{enumerate}
  \def\labelenumi{\arabic{enumi}.}
  \tightlist
  \item
    Hệ thống đã có ít nhất một ví tiền ở trạng thái \texttt{ACTIVE}.
  \item
    Cây danh mục phân loại 3 cấp đã được thiết lập với các danh mục
    tương ứng cho Loại thu nhập hoặc chi tiêu.
  \item
    Owner đang ở màn hình Quản lý Giao dịch thu chi.
  \end{enumerate}
\item
  \textbf{Hậu điều kiện thành công:}

  \begin{enumerate}
  \def\labelenumi{\arabic{enumi}.}
  \tightlist
  \item
    Giao dịch mới được lưu vào bảng \texttt{transactions} với mã
    \texttt{txnId} duy nhất.
  \item
    Số dư của ví liên quan được cập nhật tức thời theo quy tắc
    \texttt{BR-FIN-02} (cộng tiền nếu thu nhập, trừ tiền nếu chi tiêu).
  \item
    Số dư của ví thông thường không bị rơi vào trạng thái âm sau giao
    dịch chi tiêu.
  \item
    Các bảng thống kê dòng tiền tháng và báo cáo cơ cấu thu chi được
    đồng bộ hóa dữ liệu tức thời.
  \end{enumerate}
\item
  \textbf{Luồng sự kiện chính (Main Success Scenario):}

  \begin{enumerate}
  \def\labelenumi{\arabic{enumi}.}
  \tightlist
  \item
    Người dùng chọn chức năng "Thêm giao dịch mới".
  \item
    Hệ thống hiển thị biểu mẫu nhập giao dịch gồm các trường: Ngày giao
    dịch (mặc định hôm nay), Loại giao dịch (\texttt{Income} /
    \texttt{Expense}), Ví thanh toán, Cây danh mục phân loại (Type
    -\textgreater{} Group -\textgreater{} Category), Số tiền giao dịch,
    Ghi chú mô tả.
  \item
    Người dùng chọn Loại giao dịch là \texttt{Expense} (Chi tiêu).
  \item
    Hệ thống tải động danh sách các Nhóm nhu cầu (\texttt{Group}) và
    Danh mục chi tiết (\texttt{Category}) tương ứng thuộc Loại chi tiêu
    (\texttt{Type\ =\ Expense}) theo quy tắc \texttt{BR-FIN-04}.
  \item
    Người dùng chọn ngày giao dịch, chọn ví thanh toán (đang ở trạng
    thái \texttt{ACTIVE}), chọn danh mục chi tiết, nhập số tiền giao
    dịch (\(amount > 0\)), và nhập ghi chú mô tả nội dung chi tiêu.
  \item
    Người dùng nhấn nút "Lưu giao dịch".
  \item
    Hệ thống xác thực dữ liệu:

    \begin{itemize}
    \tightlist
    \item
      Số tiền giao dịch \(amount\) phải lớn hơn 0 và được chuẩn hóa về
      kiểu \texttt{BigDecimal} scale 2 theo quy tắc \texttt{BR-FIN-01}.
    \item
      Danh mục được chọn phải đang ở trạng thái hoạt động
      (\texttt{isActive\ =\ true}) và thuộc đúng loại \texttt{Expense}
      theo quy tắc \texttt{BR-FIN-04}.
    \item
      Kiểm tra số dư ví theo \texttt{BR-FIN-02}: Nếu ví không phải loại
      thẻ tín dụng (\texttt{CREDIT}), hệ thống kiểm tra số dư khả dụng
      bảo đảm điều kiện \(Balance_{wallet} - amount \ge 0\).
    \end{itemize}
  \item
    Hệ thống tạo bản ghi giao dịch mới trong bảng \texttt{transactions}
    kèm mã \texttt{txnId} duy nhất.
  \item
    Hệ thống cập nhật số dư ví: trừ số tiền \(amount\) khỏi số dư khả
    dụng của ví được chọn theo quy tắc \texttt{BR-FIN-02}.
  \item
    Hệ thống kích hoạt làm mới tự động các chỉ số tổng hợp tài chính của
    tháng quan sát.
  \item
    Hệ thống hiển thị thông báo ghi nhận giao dịch thành công và làm mới
    bảng nhật ký giao dịch trên màn hình.
  \end{enumerate}
\item
  \textbf{Các luồng rẽ nhánh (Alternative Flows):}

  \begin{itemize}
  \tightlist
  \item
    \emph{3a. Ghi nhận giao dịch Thu nhập (Income):}

    \begin{itemize}
    \tightlist
    \item
      3a1. Tại bước 3, Người dùng chọn Loại giao dịch là
      \texttt{Income}.
    \item
      3a2. Hệ thống tải động danh sách các Nhóm nhu cầu và Danh mục
      thuộc \texttt{Type\ =\ Income} (ví dụ: tiền lương, tiền thưởng,
      lãi đầu tư).
    \item
      3a3. Người dùng chọn ví nhận tiền, chọn danh mục thu nhập, nhập số
      tiền và ghi chú mô tả.
    \item
      3a4. Tại bước 7, hệ thống kiểm tra danh mục thuộc
      \texttt{Type\ =\ Income} và số tiền \(amount > 0\) (không cần kiểm
      tra số dư ví phải đủ để trừ).
    \item
      3a5. Tại bước 9, hệ thống cộng số tiền \(amount\) vào số dư khả
      dụng của ví nhận tiền:
      \(Balance_{wallet}(t) = Balance_{wallet}(t-1) + amount\).
    \end{itemize}
  \item
    \emph{2a. Tra cứu và lọc lịch sử giao dịch đa tiêu chí:}

    \begin{itemize}
    \tightlist
    \item
      2a1. Người dùng mở màn hình "Nhật ký giao dịch" và thiết lập bộ
      lọc: khoảng thời gian (\texttt{fromDate}, \texttt{toDate}), danh
      sách ví, loại giao dịch (\texttt{ALL}, \texttt{INCOME},
      \texttt{EXPENSE}), nhóm danh mục, danh mục chi tiết, từ khóa tìm
      kiếm trong ghi chú.
    \item
      2a2. Hệ thống áp dụng kết hợp logic AND trên toàn bộ các tiêu chí
      lọc, thực hiện truy vấn và phân trang kết quả (mặc định 50 bản
      ghi/trang, sắp xếp ngày giao dịch giảm dần).
    \item
      2a3. Hệ thống tính toán tổng số tiền thu nhập, tổng số tiền chi
      tiêu và dòng tiền thuần (Cash Flow = Tổng thu nhập - Tổng chi
      tiêu) trên tập kết quả lọc bằng \texttt{BigDecimal}.
    \item
      2a4. Hệ thống hiển thị bảng danh sách giao dịch kèm tổng số bản
      ghi và khối tổng hợp tài chính tương ứng.
    \end{itemize}
  \item
    \emph{2b. Chỉnh sửa giao dịch đã lưu:}

    \begin{itemize}
    \tightlist
    \item
      2b1. Người dùng chọn một giao dịch trong danh sách và nhấn nút
      "Chỉnh sửa".
    \item
      2b2. Hệ thống nạp thông tin giao dịch hiện tại lên biểu mẫu chỉnh
      sửa.
    \item
      2b3. Người dùng sửa đổi các trường cần thiết (thay đổi số tiền,
      đổi ví thanh toán, hoặc đổi danh mục) và nhấn "Cập nhật".
    \item
      2b4. Hệ thống thực hiện hoàn nguyên số dư ví cũ theo quy tắc
      \texttt{BR-FIN-02}: nếu giao dịch cũ là \texttt{Income} thì trừ số
      tiền cũ khỏi ví cũ; nếu giao dịch cũ là \texttt{Expense} thì cộng
      hoàn trả số tiền cũ vào ví cũ.
    \item
      2b5. Hệ thống kiểm tra tính hợp lệ của giao dịch mới đối với ví
      mới: nếu là \texttt{Expense}, bảo đảm
      \(Balance_{newWallet} - newAmount \ge 0\).
    \item
      2b6. Nếu hợp lệ, hệ thống áp dụng số tiền mới vào ví mới (cộng nếu
      \texttt{Income}, trừ nếu \texttt{Expense}), cập nhật bản ghi giao
      dịch trong cơ sở dữ liệu trong cùng một khối giao dịch nguyên tử.
    \item
      2b7. Hệ thống thông báo cập nhật thành công và đồng bộ lại số dư
      các ví liên quan.
    \end{itemize}
  \item
    \emph{2c. Xóa giao dịch thu chi:}

    \begin{itemize}
    \tightlist
    \item
      2c1. Người dùng chọn một giao dịch và nhấn nút "Xóa".
    \item
      2c2. Hệ thống hiển thị hộp thoại cảnh báo xác nhận việc xóa giao
      dịch và nêu rõ tác động hoàn nguyên số dư ví.
    \item
      2c3. Người dùng xác nhận đồng ý xóa.
    \item
      2c4. Hệ thống kiểm tra điều kiện hoàn nguyên: nếu giao dịch là
      \texttt{Income}, việc trừ lại tiền khỏi ví không được làm cho số
      dư ví bị âm (\(Balance_{wallet} - amount \ge 0\)).
    \item
      2c5. Hệ thống hoàn nguyên số dư ví (trừ tiền nếu xóa giao dịch
      thu, cộng trả lại tiền nếu xóa giao dịch chi) theo quy tắc
      \texttt{BR-FIN-02}.
    \item
      2c6. Hệ thống xóa bản ghi khỏi bảng \texttt{transactions} và cập
      nhật lại các báo cáo thống kê tháng.
    \item
      2c7. Hệ thống hiển thị thông báo xóa giao dịch thành công.
    \end{itemize}
  \end{itemize}
\item
  \textbf{Các luồng ngoại lệ (Exception Flows):}

  \begin{itemize}
  \tightlist
  \item
    \emph{7a. Số tiền giao dịch nhỏ hơn hoặc bằng 0:}

    \begin{itemize}
    \tightlist
    \item
      7a1. Tại bước 7 (hoặc 2b5), hệ thống phát hiện số tiền nhập vào
      \(\le 0\) hoặc không đúng định dạng số học.
    \item
      7a2. Hệ thống từ chối lưu và hiển thị thông báo: "Số tiền giao
      dịch phải là một số dương lớn hơn 0".
    \item
      7a3. Người dùng nhập lại số tiền hợp lệ.
    \end{itemize}
  \item
    \emph{7b. Số dư ví không đủ chi trả (Thiếu hụt số dư ví):}

    \begin{itemize}
    \tightlist
    \item
      7b1. Tại bước 7, hệ thống phát hiện ví chi tiền không phải là thẻ
      tín dụng và số dư khả dụng nhỏ hơn số tiền cần chi trả
      (\(Balance_{wallet} < amount\)).
    \item
      7b2. Hệ thống từ chối thực hiện giao dịch, hiển thị thông báo: "Số
      dư khả dụng của ví không đủ để thực hiện khoản chi tiêu này".
    \item
      7b3. Người dùng có thể chọn ví thanh toán khác hoặc điều chỉnh
      giảm số tiền chi tiêu.
    \end{itemize}
  \item
    \emph{7c. Danh mục bị vô hiệu hóa hoặc không khớp Loại giao dịch:}

    \begin{itemize}
    \tightlist
    \item
      7c1. Tại bước 7, hệ thống phát hiện danh mục đã bị ẩn
      (\texttt{isActive\ =\ false}) hoặc có \texttt{Type} không khớp với
      Loại giao dịch đã chọn (ví dụ chọn danh mục Chi tiêu cho giao dịch
      Thu nhập).
    \item
      7c2. Hệ thống chặn thao tác và hiển thị cảnh báo: "Danh mục không
      hợp lệ hoặc không tương thích với loại luồng tiền đã chọn".
    \item
      7c3. Người dùng chọn lại danh mục hợp lệ từ cây danh mục.
    \end{itemize}
  \item
    \emph{2c4a. Không thể xóa giao dịch thu nhập do làm âm số dư ví:}

    \begin{itemize}
    \tightlist
    \item
      2c4a1. Tại bước 2c4, hệ thống phát hiện việc xóa giao dịch thu
      nhập sẽ làm số dư ví hiện tại rơi vào số âm (do tiền từ khoản thu
      này đã được chi tiêu hoặc chuyển đi trong các giao dịch kế tiếp).
    \item
      2c4a2. Hệ thống từ chối thao tác xóa và hiển thị thông báo: "Không
      thể xóa giao dịch thu nhập này vì sẽ làm số dư ví bị âm. Vui lòng
      kiểm tra lại các khoản chi tiêu liên quan của ví trước khi thực
      hiện".
    \end{itemize}
  \end{itemize}
\item
  \textbf{Quy tắc nghiệp vụ đặc thù (Special Requirements /
  Invariants):}

  \begin{enumerate}
  \def\labelenumi{\arabic{enumi}.}
  \tightlist
  \item
    Bảo toàn số dư ví theo \texttt{BR-FIN-02}: Mọi giao dịch thu chi đều
    phải phản ánh trực tiếp vào số dư ví tương ứng.
  \item
    Bất biến phân cấp danh mục \texttt{BR-FIN-04}: Mọi giao dịch phải
    được phân loại chuẩn xác theo cấu trúc 3 cấp và bảo đảm tính tương
    thích tuyệt đối về \texttt{Type}.
  \item
    Độ chính xác tài chính \texttt{BR-FIN-01}: Xử lý toàn bộ bằng
    \texttt{BigDecimal}, không phát sinh sai số làm tròn.
  \item
    Tính nguyên tử khi chỉnh sửa và xóa: Thao tác sửa hoặc xóa giao dịch
    phải bảo đảm đồng bộ hoàn nguyên số dư ví và cập nhật sổ cái trong
    cùng một khối giao dịch nguyên tử; nếu có lỗi ở bất kỳ bước nào, hệ
    thống phải tự động rollback toàn bộ.
  \item
    Phù hợp các yêu cầu chức năng: \texttt{FR-TXN-01},
    \texttt{FR-TXN-02}, \texttt{FR-TXN-03}, \texttt{FR-TXN-04},
    \texttt{FR-TXN-05}.
  \end{enumerate}
\end{itemize}

\subsubsection{6.3 UC03: Chuyển tiền nội bộ giữa các ví (Internal Fund
Transfer)}\label{63-uc03-chuyux1ec3n-tiux1ec1n-nux1ed9i-bux1ed9-giux1eefa-cuxe1c-vuxed-internal-fund-transfer}

\begin{itemize}
\tightlist
\item
  \textbf{Mã và Tên ca sử dụng:}
  \texttt{UC03:\ Chuyển\ tiền\ nội\ bộ\ giữa\ các\ ví\ (Internal\ Fund\ Transfer)}
\item
  \textbf{Tác nhân chính:} Chủ sở hữu hệ thống (Owner)
\item
  \textbf{Mô tả tóm tắt:} Ca sử dụng này cho phép Owner thực hiện điều
  chuyển giá trị tiền tệ giữa hai ví thanh khoản nội bộ trong hệ thống
  (ví dụ rút tiền từ tài khoản ngân hàng về ví tiền mặt, hoặc chuyển
  khoản giữa hai tài khoản ngân hàng). Hệ thống áp dụng nguyên lý ghi sổ
  kép đối ứng của Martin Fowler, tự động trừ ví nguồn, cộng ví đích và
  tách khoản phí chuyển tiền phát sinh (nếu có) thành một khoản chi phí
  độc lập, bảo toàn tuyệt đối tổng tài sản thanh khoản.
\item
  \textbf{Điều kiện kích hoạt:} Owner có nhu cầu tái cơ cấu số dư tiền
  mặt giữa các ví hoặc ghi nhận giao dịch chuyển khoản ngân hàng đã phát
  sinh trong thực tế.
\item
  \textbf{Tiền điều kiện:}

  \begin{enumerate}
  \def\labelenumi{\arabic{enumi}.}
  \tightlist
  \item
    Hệ thống có ít nhất hai ví tiền ở trạng thái \texttt{ACTIVE}.
  \item
    Số dư khả dụng của ví nguồn phải lớn hơn hoặc bằng tổng số tiền
    chuyển cộng phí chuyển tiền.
  \item
    Owner đang ở giao diện Chuyển tiền nội bộ.
  \end{enumerate}
\item
  \textbf{Hậu điều kiện thành công:}

  \begin{enumerate}
  \def\labelenumi{\arabic{enumi}.}
  \tightlist
  \item
    Bản ghi chuyển tiền mới được lưu vào bảng \texttt{transfers} với mã
    \texttt{transferId} duy nhất.
  \item
    Bút toán kép được ghi nhận nguyên tử: ví nguồn giảm chính xác
    \(amount + fee\), ví đích tăng chính xác \(amount\).
  \item
    Nếu có phí chuyển tiền (\(fee > 0\)), một giao dịch chi tiêu phí độc
    lập được tự động ghi nhận vào bảng \texttt{transactions} liên kết
    với \texttt{transferId}.
  \item
    Tổng tài sản ví thanh khoản toàn hệ thống giảm đúng bằng khoản phí
    \(fee\) (nếu không có phí thì tổng tài sản bảo toàn 100\% theo
    \texttt{BR-FIN-03}).
  \end{enumerate}
\item
  \textbf{Luồng sự kiện chính (Main Success Scenario):}

  \begin{enumerate}
  \def\labelenumi{\arabic{enumi}.}
  \tightlist
  \item
    Người dùng chọn chức năng "Chuyển tiền nội bộ".
  \item
    Hệ thống hiển thị biểu mẫu chuyển tiền gồm các trường: Ngày chuyển
    tiền (mặc định hôm nay), Ví nguồn (trích tiền), Ví đích (nhận tiền),
    Số tiền chuyển, Phí chuyển tiền (tùy chọn, mặc định 0), Danh mục phí
    (mặc định "Phí ngân hàng / chuyển tiền"), Ghi chú mô tả.
  \item
    Người dùng chọn Ví nguồn từ danh sách ví đang \texttt{ACTIVE}. Hệ
    thống hiển thị số dư khả dụng hiện tại của ví nguồn để hỗ trợ quan
    sát.
  \item
    Người dùng chọn Ví đích từ danh sách ví đang \texttt{ACTIVE}.
  \item
    Người dùng nhập số tiền cần chuyển (\(amount > 0\)).
  \item
    Người dùng nhập số tiền phí chuyển tiền (\(fee \ge 0\)) nếu giao
    dịch thực tế có tính phí.
  \item
    Người dùng nhập ghi chú mô tả (nếu có) và nhấn nút "Xác nhận chuyển
    tiền".
  \item
    Hệ thống tiến hành xác thực nghiệp vụ:

    \begin{itemize}
    \tightlist
    \item
      Kiểm tra ví nguồn phải khác ví đích
      (\(fromWalletId \ne toWalletId\)).
    \item
      Kiểm tra số tiền chuyển \(amount > 0\) và phí chuyển tiền
      \(fee \ge 0\), chuẩn hóa về kiểu \texttt{BigDecimal} scale 2 theo
      quy tắc \texttt{BR-FIN-01}.
    \item
      Kiểm tra năng lực thanh toán của ví nguồn theo \texttt{BR-FIN-02}:
      \(Balance_{from} \ge amount + fee\).
    \end{itemize}
  \item
    Hệ thống thực hiện khối ghi sổ kép nguyên tử (Atomic Transaction)
    theo quy tắc \texttt{BR-FIN-03}:

    \begin{itemize}
    \tightlist
    \item
      Trừ số tiền \(amount\) khỏi số dư ví nguồn \(fromWalletId\).
    \item
      Cộng số tiền \(amount\) vào số dư ví đích \(toWalletId\).
    \item
      Lưu bản ghi chuyển tiền vào bảng \texttt{transfers} với mã định
      danh \texttt{transferId}.
    \end{itemize}
  \item
    Nếu \(fee > 0\), hệ thống tự động sinh một bản ghi chi tiêu
    (\texttt{Expense}) độc lập trong bảng \texttt{transactions} với số
    tiền \(fee\), trừ tiếp \(fee\) khỏi ví nguồn \(fromWalletId\), gắn
    với danh mục chi phí ngân hàng và lưu mã tham chiếu
    \texttt{linkedTransferId\ =\ transferId}.
  \item
    Hệ thống đồng bộ hóa dữ liệu cơ sở dữ liệu và xác nhận giao dịch
    hoàn tất thành công.
  \item
    Hệ thống hiển thị thông báo chuyển tiền thành công kèm thông tin chi
    tiết: số tiền chuyển, phí đã trừ, số dư mới của ví nguồn và ví đích.
  \end{enumerate}
\item
  \textbf{Các luồng rẽ nhánh (Alternative Flows):}

  \begin{itemize}
  \tightlist
  \item
    \emph{6a. Chuyển tiền không mất phí (\(fee = 0\)):}

    \begin{itemize}
    \tightlist
    \item
      6a1. Tại bước 6, Người dùng giữ nguyên mức phí bằng 0 (hoặc để
      trống).
    \item
      6a2. Hệ thống chỉ thực hiện ghi sổ kép đối ứng trừ \(amount\) ở ví
      nguồn và cộng \(amount\) ở ví đích.
    \item
      6a3. Hệ thống không sinh bản ghi chi tiêu phí liên kết; tổng số dư
      toàn bộ các ví trước và sau giao dịch bảo toàn tuyệt đối
      (\(\Delta \sum Balance = 0\)).
    \end{itemize}
  \item
    \emph{2a. Xem lịch sử chuyển tiền nội bộ:}

    \begin{itemize}
    \tightlist
    \item
      2a1. Người dùng mở tab "Lịch sử chuyển tiền".
    \item
      2a2. Hệ thống hiển thị danh sách các lần chuyển tiền nội bộ đã
      thực hiện, sắp xếp theo thứ tự ngày giảm dần, hiển thị rõ ví
      nguồn, ví đích, số tiền chuyển, phí dịch vụ và trạng thái.
    \end{itemize}
  \item
    \emph{2b. Hủy bỏ giao dịch chuyển tiền nội bộ:}

    \begin{itemize}
    \tightlist
    \item
      2b1. Tại danh sách lịch sử chuyển tiền, Người dùng chọn một giao
      dịch và nhấn nút "Hủy chuyển tiền".
    \item
      2b2. Hệ thống hiển thị hộp thoại xác nhận hủy, liệt kê chi tiết
      các khoản hoàn nguyên sẽ diễn ra.
    \item
      2b3. Người dùng nhấn xác nhận hủy.
    \item
      2b4. Hệ thống kiểm tra điều kiện hoàn nguyên tại ví đích theo quy
      tắc \texttt{FR-TRF-03}: số dư khả dụng của ví đích phải lớn hơn
      hoặc bằng số tiền đã nhận (\(Balance_{to} \ge amount\)).
    \item
      2b5. Hệ thống thực hiện hoàn nguyên đối ứng nguyên tử:

      \begin{itemize}
      \tightlist
      \item
        Cộng trả lại \(amount\) vào số dư ví nguồn.
      \item
        Trừ \(amount\) khỏi số dư ví đích.
      \item
        Nếu giao dịch có phí chuyển tiền liên kết, hệ thống hoàn trả
        \(fee\) vào ví nguồn và xóa bản ghi chi tiêu phí trong bảng
        \texttt{transactions}.
      \end{itemize}
    \item
      2b6. Hệ thống cập nhật trạng thái của bản ghi chuyển tiền sang
      \texttt{CANCELLED} (hoặc xóa khỏi sổ chuyển tiền).
    \item
      2b7. Hệ thống hiển thị thông báo hủy chuyển tiền thành công và cập
      nhật lại số dư hai ví.
    \end{itemize}
  \end{itemize}
\item
  \textbf{Các luồng ngoại lệ (Exception Flows):}

  \begin{itemize}
  \tightlist
  \item
    \emph{8a. Ví nguồn trùng ví đích:}

    \begin{itemize}
    \tightlist
    \item
      8a1. Tại bước 8, hệ thống phát hiện \(fromWalletId = toWalletId\).
    \item
      8a2. Hệ thống chặn giao dịch và hiển thị thông báo lỗi: "Ví nguồn
      và ví đích không được trùng nhau. Vui lòng chọn hai ví khác nhau".
    \item
      8a3. Người dùng chọn lại ví đích phù hợp.
    \end{itemize}
  \item
    \emph{8b. Số tiền chuyển nhỏ hơn hoặc bằng 0:}

    \begin{itemize}
    \tightlist
    \item
      8b1. Tại bước 8, hệ thống phát hiện số tiền chuyển nhập vào
      \(\le 0\).
    \item
      8b2. Hệ thống thông báo: "Số tiền chuyển phải là một số dương lớn
      hơn 0".
    \item
      8b3. Người dùng nhập lại số tiền hợp lệ.
    \end{itemize}
  \item
    \emph{8c. Số dư ví nguồn không đủ thanh toán
    (\(Balance_{from} < amount + fee\)):}

    \begin{itemize}
    \tightlist
    \item
      8c1. Tại bước 8, hệ thống tính tổng nghĩa vụ cần trừ tại ví nguồn
      là \(amount + fee\) và phát hiện vượt quá số dư khả dụng hiện tại
      của ví nguồn.
    \item
      8c2. Hệ thống từ chối thực hiện và hiển thị thông báo: "Số dư ví
      nguồn không đủ để thực hiện chuyển tiền".
    \item
      8c3. Người dùng có thể nạp thêm tiền vào ví nguồn, giảm số tiền
      chuyển hoặc chọn ví nguồn khác.
    \end{itemize}
  \item
    \emph{9a. Lỗi một vế trong quá trình ghi sổ kép (Lỗi hệ thống hoặc
    cơ sở dữ liệu):}

    \begin{itemize}
    \tightlist
    \item
      9a1. Trong quá trình thực thi bước 9 hoặc bước 10, nếu thao tác
      trừ tiền ở ví nguồn thành công nhưng thao tác cộng tiền ở ví đích
      bị gián đoạn do lỗi kỹ thuật.
    \item
      9a2. Hệ thống tự động kích hoạt cơ chế Rollback toàn bộ
      transaction, hủy bỏ thao tác trừ tiền ở ví nguồn, khôi phục nguyên
      vẹn trạng thái trước giao dịch.
    \item
      9a3. Hệ thống ghi nhật ký lỗi chi tiết và hiển thị thông báo:
      "Giao dịch chuyển tiền bị gián đoạn và đã được phục hồi nguyên
      trạng để bảo đảm toàn vẹn tài sản. Vui lòng thử lại sau".
    \end{itemize}
  \item
    \emph{2b4a. Ví đích không đủ số dư để hoàn nguyên khi hủy giao
    dịch:}

    \begin{itemize}
    \tightlist
    \item
      2b4a1. Tại bước 2b4, hệ thống phát hiện số dư khả dụng hiện tại
      của ví đích nhỏ hơn số tiền đã nhận từ giao dịch chuyển tiền
      (\(Balance_{to} < amount\)) do ví đích đã dùng tiền này để chi
      tiêu trong các giao dịch tiếp theo.
    \item
      2b4a2. Hệ thống từ chối thao tác hủy và hiển thị thông báo: "Không
      thể hủy giao dịch chuyển tiền này vì ví đích không còn đủ số dư để
      hoàn trả".
    \end{itemize}
  \end{itemize}
\item
  \textbf{Quy tắc nghiệp vụ đặc thù (Special Requirements /
  Invariants):}

  \begin{enumerate}
  \def\labelenumi{\arabic{enumi}.}
  \tightlist
  \item
    Bất biến ghi sổ kép theo \texttt{BR-FIN-03}: Tổng tài sản trên toàn
    bộ các ví được bảo toàn tuyệt đối đối với giao dịch chuyển tiền
    thuần túy (\(\Delta \sum Balance = 0\)).
  \item
    Bất biến tách biệt phí dịch vụ: Khoản phí ngân hàng bắt buộc phải
    được hạch toán thành một giao dịch chi tiêu (\texttt{Expense}) riêng
    biệt, không được khấu trừ gộp làm méo mó số tiền chuyển thực nhận
    của ví đích.
  \item
    Bất biến số học \texttt{BR-FIN-01}: Toàn bộ các bước tính toán trừ
    cộng số dư đều thực hiện qua \texttt{BigDecimal}.
  \item
    Phù hợp các yêu cầu chức năng: \texttt{FR-TRF-01},
    \texttt{FR-TRF-02}, \texttt{FR-TRF-03}.
  \end{enumerate}
\end{itemize}

\subsubsection{6.4 UC04: Thiết lập và quản lý danh mục phân loại thu chi
3 cấp (Category Hierarchy
Management)}\label{64-uc04-thiux1ebft-lux1eadp-vuxe0-quux1ea3n-luxfd-danh-mux1ee5c-phuxe2n-loux1ea1i-thu-chi-3-cux1ea5p-category-hierarchy-management}

\begin{itemize}
\tightlist
\item
  \textbf{Mã và Tên ca sử dụng:}
  \texttt{UC04:\ Thiết\ lập\ và\ quản\ lý\ danh\ mục\ phân\ loại\ thu\ chi\ 3\ cấp\ (Category\ Hierarchy\ Management)}
\item
  \textbf{Tác nhân chính:} Chủ sở hữu hệ thống (Owner)
\item
  \textbf{Mô tả tóm tắt:} Ca sử dụng này cho phép Owner khởi tạo, duyệt
  cây phân cấp danh mục 3 tầng (Type -\textgreater{} Group
  -\textgreater{} Category), thêm mới các danh mục chi tiết phù hợp với
  thói quen sinh hoạt cá nhân, đổi tên danh mục, và quản lý vòng đời
  danh mục (kích hoạt hoặc ẩn). Hệ thống kiểm soát nghiêm ngặt ràng buộc
  toàn vẹn lịch sử, nghiêm cấm xóa vật lý danh mục đã phát sinh dữ liệu
  giao dịch theo quy tắc BR-FIN-04.
\item
  \textbf{Điều kiện kích hoạt:} Owner truy cập màn hình Cấu hình danh
  mục để xem xét cấu trúc phân loại hoặc bổ sung danh mục mới phục vụ
  ghi chép thu chi.
\item
  \textbf{Tiền điều kiện:}

  \begin{enumerate}
  \def\labelenumi{\arabic{enumi}.}
  \tightlist
  \item
    Hệ thống đã khởi tạo sẵn bộ danh mục chuẩn từ bảng
    \texttt{database.classification} (2 Type, 7 Group chuẩn).
  \item
    Owner đang ở giao diện Quản lý Danh mục.
  \end{enumerate}
\item
  \textbf{Hậu điều kiện thành công:}

  \begin{enumerate}
  \def\labelenumi{\arabic{enumi}.}
  \tightlist
  \item
    Danh mục mới được thêm vào đúng Group cha và thừa hưởng thuộc tính
    \texttt{Type} tương ứng.
  \item
    Tính toàn vẹn của cây danh mục 3 cấp được bảo toàn: không có danh
    mục mồ côi, không trùng lặp tên trong cùng một nhóm cha.
  \item
    Các danh mục đã gắn với giao dịch lịch sử được bảo vệ an toàn, không
    bị biến dạng loại hình hoặc bị xóa mất dấu vết kế toán.
  \end{enumerate}
\item
  \textbf{Luồng sự kiện chính (Main Success Scenario):}

  \begin{enumerate}
  \def\labelenumi{\arabic{enumi}.}
  \tightlist
  \item
    Người dùng chọn chức năng "Quản lý danh mục thu chi".
  \item
    Hệ thống hiển thị toàn bộ cây danh mục phân loại 3 cấp: Cấp 1 gồm
    hai nhánh gốc \texttt{Income} và \texttt{Expense}; Cấp 2 gồm các
    Nhóm nhu cầu (\texttt{Group}); Cấp 3 gồm các Danh mục chi tiết
    (\texttt{Category}) trực thuộc từng nhóm.
  \item
    Người dùng chọn một Nhóm nhu cầu cụ thể (ví dụ: nhóm
    \texttt{essentials} thuộc \texttt{Expense}) và nhấn nút "Thêm danh
    mục mới".
  \item
    Hệ thống hiển thị biểu mẫu tạo danh mục gồm các trường: Nhóm cha đã
    chọn, Tên danh mục mới, Biểu tượng và màu sắc nhận diện (tùy chọn),
    Trạng thái hoạt động (mặc định kích hoạt).
  \item
    Người dùng nhập tên danh mục mới, chọn biểu tượng và nhấn nút "Lưu
    danh mục".
  \item
    Hệ thống xác thực dữ liệu:

    \begin{itemize}
    \tightlist
    \item
      Tên danh mục không được để trống (độ dài từ 1 đến 50 ký tự).
    \item
      Kiểm tra tên danh mục không trùng lặp với các danh mục khác đã tồn
      tại trong cùng Nhóm cha (so sánh không phân biệt hoa thường sau
      khi cắt khoảng trắng thừa).
    \item
      Kiểm tra tính tồn tại của Nhóm cha (\texttt{groupId}).
    \end{itemize}
  \item
    Hệ thống tạo bản ghi mới cho Category trong cơ sở dữ liệu với mã
    định danh duy nhất (\texttt{categoryId}), thiết lập liên kết khóa
    ngoại với \texttt{groupId}, gán cờ \texttt{isActive\ =\ true}.
  \item
    Hệ thống làm mới cây danh mục hiển thị trên giao diện, mở rộng nhánh
    nhóm cha vừa thêm để làm nổi bật danh mục mới.
  \item
    Hệ thống cập nhật danh sách chọn động tại biểu mẫu ghi nhận giao
    dịch để sẵn sàng sử dụng danh mục mới.
  \end{enumerate}
\item
  \textbf{Các luồng rẽ nhánh (Alternative Flows):}

  \begin{itemize}
  \tightlist
  \item
    \emph{2a. Đổi tên danh mục chi tiết:}

    \begin{itemize}
    \tightlist
    \item
      2a1. Người dùng chọn một danh mục cụ thể trên cây phân cấp và nhấn
      "Đổi tên".
    \item
      2a2. Hệ thống hiển thị hộp thoại cho phép chỉnh sửa tên hiển thị.
    \item
      2a3. Người dùng nhập tên mới và xác nhận lưu.
    \item
      2a4. Hệ thống kiểm tra tên mới không trùng lặp với các danh mục
      khác trong cùng nhóm.
    \item
      2a5. Hệ thống cập nhật tên danh mục trong cơ sở dữ liệu. Toàn bộ
      các giao dịch trong quá khứ liên kết với \texttt{categoryId} này
      sẽ tự động hiển thị theo tên mới mà không bị phá vỡ liên kết khóa
      ngoại.
    \item
      2a6. Hệ thống hiển thị thông báo đổi tên thành công.
    \end{itemize}
  \item
    \emph{2b. Ẩn / Ngừng sử dụng danh mục (Archived / Inactive):}

    \begin{itemize}
    \tightlist
    \item
      2b1. Người dùng chọn một danh mục không còn nhu cầu sử dụng và
      chọn thao tác "Ngừng sử dụng (Ẩn)".
    \item
      2b2. Hệ thống cập nhật trường \texttt{isActive\ =\ false} của danh
      mục.
    \item
      2b3. Hệ thống loại bỏ danh mục này khỏi danh sách gợi ý khi tạo
      giao dịch mới, đồng thời hiển thị ký hiệu đánh dấu danh mục đã bị
      ẩn trên cây quản lý.
    \item
      2b4. Mọi báo cáo thống kê lịch sử của các tháng trước có chứa giao
      dịch thuộc danh mục này vẫn hiển thị dữ liệu bình thường, bảo đảm
      tính liên tục của số liệu tài chính.
    \end{itemize}
  \item
    \emph{2c. Kích hoạt lại danh mục đã bị ẩn:}

    \begin{itemize}
    \tightlist
    \item
      2c1. Người dùng chọn một danh mục đang ở trạng thái ẩn
      (\texttt{isActive\ =\ false}) và nhấn "Kích hoạt lại".
    \item
      2c2. Hệ thống cập nhật cờ \texttt{isActive\ =\ true}. Danh mục
      xuất hiện trở lại trong danh sách lựa chọn khi thêm giao dịch mới.
    \end{itemize}
  \item
    \emph{2d. Xóa danh mục chưa phát sinh giao dịch:}

    \begin{itemize}
    \tightlist
    \item
      2d1. Người dùng chọn một danh mục và nhấn nút "Xóa danh mục".
    \item
      2d2. Hệ thống quét bảng nhật ký \texttt{transactions} để kiểm tra
      sự tồn tại của các giao dịch tham chiếu tới \texttt{categoryId}
      này.
    \item
      2d3. Nếu hoàn toàn CHƯA CÓ giao dịch nào liên kết, hệ thống hiển
      thị hộp thoại xác nhận xóa vĩnh viễn.
    \item
      2d4. Người dùng xác nhận xóa.
    \item
      2d5. Hệ thống thực hiện xóa vật lý bản ghi danh mục khỏi cơ sở dữ
      liệu và thông báo xóa thành công.
    \end{itemize}
  \end{itemize}
\item
  \textbf{Các luồng ngoại lệ (Exception Flows):}

  \begin{itemize}
  \tightlist
  \item
    \emph{6a. Trùng tên danh mục trong cùng một nhóm:}

    \begin{itemize}
    \tightlist
    \item
      6a1. Tại bước 6 (hoặc 2a4), hệ thống phát hiện tên danh mục nhập
      vào trùng với tên của một danh mục đã tồn tại trong cùng Nhóm cha.
    \item
      6a2. Hệ thống từ chối lưu và hiển thị thông báo: "Tên danh mục này
      đã tồn tại trong nhóm. Vui lòng đặt tên khác để tránh nhầm lẫn".
    \item
      6a3. Người dùng điều chỉnh lại tên danh mục.
    \end{itemize}
  \item
    \emph{6b. Nhóm cha không tồn tại hoặc bị lỗi tham chiếu:}

    \begin{itemize}
    \tightlist
    \item
      6b1. Tại bước 6, hệ thống phát hiện mã nhóm cha \texttt{groupId}
      không tồn tại trong danh mục hệ thống.
    \item
      6b2. Hệ thống từ chối tạo và thông báo lỗi cấu trúc: "Không tìm
      thấy nhóm nhu cầu cha hợp lệ".
    \end{itemize}
  \item
    \emph{2d3a. Cố gắng xóa danh mục đã phát sinh giao dịch lịch sử:}

    \begin{itemize}
    \tightlist
    \item
      2d3a1. Tại bước 2d3, hệ thống phát hiện đã có ít nhất một bản ghi
      giao dịch trong bảng \texttt{transactions} đang liên kết với danh
      mục này.
    \item
      2d3a2. Hệ thống từ chối xóa vật lý bản ghi và hiển thị hộp thoại
      thông báo: "Danh mục này đã có dữ liệu giao dịch phát sinh trong
      lịch sử và không thể bị xóa để bảo toàn tính toàn vẹn của sổ sách
      kế toán. Bạn có muốn chuyển danh mục này sang trạng thái ẩn (ngừng
      sử dụng) hay không?".
    \item
      2d3a3. Nếu Người dùng chọn đồng ý ẩn, hệ thống chuyển hướng thực
      hiện luồng 2b; nếu Người dùng chọn hủy, hệ thống giữ nguyên trạng
      thái danh mục.
    \end{itemize}
  \item
    \emph{2a3a. Cố ý thay đổi Loại luồng tiền (Type) của danh mục đã có
    giao dịch:}

    \begin{itemize}
    \tightlist
    \item
      2a3a1. Nếu có yêu cầu chuyển đổi một danh mục từ nhóm Chi tiêu
      sang nhóm Thu nhập (hoặc ngược lại) khi danh mục đó đã có giao
      dịch trong quá khứ.
    \item
      2a3a2. Hệ thống chặn đứng thao tác, báo lỗi vi phạm bất biến phân
      loại thu chi theo quy tắc \texttt{BR-FIN-04} và từ chối thay đổi.
    \end{itemize}
  \end{itemize}
\item
  \textbf{Quy tắc nghiệp vụ đặc thù (Special Requirements /
  Invariants):}

  \begin{enumerate}
  \def\labelenumi{\arabic{enumi}.}
  \tightlist
  \item
    Bất biến cấu trúc 3 cấp theo \texttt{BR-FIN-04}: Luôn bảo đảm cây
    phân cấp Type -\textgreater{} Group -\textgreater{} Category nguyên
    vẹn, không tồn tại nút mồ côi.
  \item
    Bất biến bảo toàn lịch sử giao dịch: Danh mục đã phát sinh dữ liệu
    trong sổ cái không được phép xóa vật lý (chỉ cho phép soft-delete
    qua cờ \texttt{isActive}).
  \item
    Bất biến tính nhất quán của Loại: Không cho phép chuyển đổi nhánh
    gốc giữa Thu nhập và Chi tiêu đối với danh mục đã có giao dịch nhằm
    tránh làm xáo trộn báo cáo dòng tiền quá khứ.
  \item
    Phù hợp các yêu cầu chức năng: \texttt{FR-CAT-01},
    \texttt{FR-CAT-02}, \texttt{FR-CAT-03}.
  \end{enumerate}
\end{itemize}

\subsubsection{6.5 UC05: Quản lý công nợ (Debt
Management)}\label{65-uc05-quux1ea3n-luxfd-cuxf4ng-nux1ee3-debt-management}

\begin{itemize}
\tightlist
\item
  \textbf{Mã và Tên ca sử dụng:}
  \texttt{UC05:\ Quản\ lý\ công\ nợ\ (Debt\ Management)}
\item
  \textbf{Tác nhân chính:} Chủ sở hữu hệ thống (Owner)
\item
  \textbf{Mô tả tóm tắt:} Ca sử dụng này cho phép Owner theo dõi toàn
  diện các nghĩa vụ tài chính hai chiều gồm các khoản phải thu (cho vay
  - \texttt{lend}) và các khoản phải trả (đi vay - \texttt{borrow}) đối
  ứng với từng đối tác cụ thể; ghi nhận quá trình hoàn ứng (thu hồi nợ -
  \texttt{collect}, trả nợ - \texttt{repay}); xử lý xóa nợ khó đòi phi
  tiền mặt (\texttt{write-off}); và truy xuất sao kê công nợ chi tiết
  theo đối tượng theo quy tắc BR-FIN-07.
\item
  \textbf{Điều kiện kích hoạt:} Owner phát sinh giao dịch cho vay, đi
  vay, thu hồi nợ, trả nợ tiền mặt hoặc cần đối soát, xử lý xóa nợ với
  một đối tác bên ngoài.
\item
  \textbf{Tiền điều kiện:}

  \begin{enumerate}
  \def\labelenumi{\arabic{enumi}.}
  \tightlist
  \item
    Hệ thống có ít nhất một ví tiền ở trạng thái \texttt{ACTIVE} (đối
    với các nghiệp vụ phát sinh dòng tiền mặt).
  \item
    Đối với nghiệp vụ thu hồi nợ hoặc trả nợ: đối tác phải đang có số dư
    nợ tương ứng lớn hơn 0.
  \item
    Owner đang ở giao diện Quản lý Công nợ.
  \end{enumerate}
\item
  \textbf{Hậu điều kiện thành công:}

  \begin{enumerate}
  \def\labelenumi{\arabic{enumi}.}
  \tightlist
  \item
    Bản ghi công nợ mới được lưu vào bảng \texttt{debt} với loại nghiệp
    vụ chính xác (\texttt{lend}, \texttt{borrow}, \texttt{collect},
    \texttt{repay}, \texttt{write-off\ lend},
    \texttt{write-off\ borrow}).
  \item
    Số dư ví liên quan được đồng bộ hóa tức thời (nếu là nghiệp vụ tiền
    mặt) theo \texttt{BR-FIN-02}.
  \item
    Số dư nợ chi tiết của cá nhân đối tác và tổng chỉ số \(Receivable\),
    \(Payable\) của hệ thống được tính toán lại chính xác theo
    \texttt{BR-FIN-07}.
  \item
    Bất biến số dư nợ không âm luôn được thỏa mãn: không thu nợ vượt quá
    nợ phải thu, không trả nợ vượt quá nợ phải trả.
  \end{enumerate}
\item
  \textbf{Luồng sự kiện chính (Main Success Scenario) - Ghi nhận khoản
  Cho vay (Lend):}

  \begin{enumerate}
  \def\labelenumi{\arabic{enumi}.}
  \tightlist
  \item
    Người dùng chọn chức năng "Tạo khoản cho vay mới".
  \item
    Hệ thống hiển thị biểu mẫu ghi nhận cho vay gồm các trường: Ngày cho
    vay (mặc định hôm nay), Mã đợt vay (tùy chọn), Tên người vay
    (\texttt{person}), Số tiền cho vay, Ví xuất tiền, Ghi chú mô tả.
  \item
    Người dùng nhập tên người vay, nhập số tiền cho vay
    (\(amount > 0\)), chọn ví trích tiền (đang \texttt{ACTIVE}), nhập mã
    đợt và ghi chú (nếu có).
  \item
    Người dùng nhấn nút "Xác nhận cho vay".
  \item
    Hệ thống xác thực dữ liệu:

    \begin{itemize}
    \tightlist
    \item
      Tên người vay không được để trống (độ dài từ 1 đến 100 ký tự).
    \item
      Số tiền cho vay \(amount\) phải lớn hơn 0, chuẩn hóa về kiểu
      \texttt{BigDecimal} scale 2 theo quy tắc \texttt{BR-FIN-01}.
    \item
      Kiểm tra số dư ví xuất tiền theo \texttt{BR-FIN-02}:
      \(Balance_{wallet} - amount \ge 0\).
    \end{itemize}
  \item
    Hệ thống thực hiện ghi sổ nghiệp vụ:

    \begin{itemize}
    \tightlist
    \item
      Trừ số tiền \(amount\) khỏi số dư khả dụng của ví xuất tiền.
    \item
      Tăng số dư nợ phải thu của người vay:
      \[
      Receivable_{person}(t) = Receivable_{person}(t-1) + amount
      \]
    \item
      Lưu bản ghi mới vào bảng \texttt{debt} với loại \texttt{lend} kèm
      mã \texttt{debtId} duy nhất.
    \end{itemize}
  \item
    Hệ thống cập nhật tổng tài sản khoản phải thu toàn hệ thống
    (\(Receivable\)). Khoản tiền này chuyển từ tài sản thanh khoản sang
    tài sản quyền đòi nợ, không tính vào chi phí tiêu dùng
    (\texttt{Expense}).
  \item
    Hệ thống thông báo ghi nhận khoản cho vay thành công và cập nhật
    bảng theo dõi công nợ trên màn hình.
  \end{enumerate}
\item
  \textbf{Các luồng rẽ nhánh (Alternative Flows):}

  \begin{itemize}
  \tightlist
  \item
    \emph{2a. Ghi nhận khoản Đi vay (Borrow):}

    \begin{itemize}
    \tightlist
    \item
      2a1. Người dùng chọn chức năng "Ghi nhận khoản đi vay".
    \item
      2a2. Hệ thống hiển thị biểu mẫu nhập: Ngày vay, Tên chủ nợ
      (\texttt{person}), Số tiền vay, Ví nhận tiền, Ghi chú mô tả.
    \item
      2a3. Người dùng nhập thông tin hợp lệ (\(amount > 0\)) và nhấn xác
      nhận lưu.
    \item
      2a4. Hệ thống cộng số tiền \(amount\) vào số dư khả dụng của ví
      nhận tiền theo quy tắc \texttt{BR-FIN-02}.
    \item
      2a5. Hệ thống tăng nghĩa vụ nợ phải trả đối với chủ nợ:
      \[
      Payable_{person}(t) = Payable_{person}(t-1) + amount
      \]
    \item
      2a6. Hệ thống lưu bản ghi công nợ loại \texttt{borrow}, cập nhật
      tổng \(Payable\) toàn hệ thống và hiển thị thông báo thành công.
    \end{itemize}
  \item
    \emph{2b. Thu hồi nợ cho vay (Collect):}

    \begin{itemize}
    \tightlist
    \item
      2b1. Người dùng chọn một người vay trong danh sách nợ phải thu và
      nhấn "Thu nợ".
    \item
      2b2. Hệ thống hiển thị số dư nợ phải thu còn lại của người này và
      mở biểu mẫu thu nợ: Ngày thu, Số tiền thu, Ví nhận tiền, Ghi chú.
    \item
      2b3. Người dùng nhập số tiền thu hồi (\(amount > 0\)), chọn ví
      nhận tiền và nhấn "Xác nhận thu nợ".
    \item
      2b4. Hệ thống kiểm tra ràng buộc trần thu nợ theo quy tắc
      \texttt{BR-FIN-07}: \(amount \le Receivable_{person}\).
    \item
      2b5. Hệ thống cộng số tiền \(amount\) vào ví nhận tiền, đồng thời
      giảm số dư nợ phải thu của người vay:
      \(Receivable_{person}(t) = Receivable_{person}(t-1) - amount\).
    \item
      2b6. Hệ thống lưu bản ghi loại \texttt{collect} vào bảng
      \texttt{debt} và hiển thị thông báo thành công kèm dư nợ còn lại.
    \end{itemize}
  \item
    \emph{2c. Trả nợ đi vay (Repay):}

    \begin{itemize}
    \tightlist
    \item
      2c1. Người dùng chọn một chủ nợ trong danh sách nợ phải trả và
      nhấn "Trả nợ".
    \item
      2c2. Hệ thống hiển thị dư nợ phải trả hiện tại và biểu mẫu trả nợ:
      Ngày trả, Số tiền trả, Ví trích tiền, Ghi chú.
    \item
      2c3. Người dùng nhập số tiền hoàn trả (\(amount > 0\)), chọn ví
      trích tiền và xác nhận.
    \item
      2c4. Hệ thống kiểm tra ràng buộc trần trả nợ
      (\(amount \le Payable_{person}\)) và kiểm tra số dư ví trích tiền
      (\(Balance_{wallet} \ge amount\)).
    \item
      2c5. Hệ thống trừ \(amount\) khỏi ví trích tiền, giảm nghĩa vụ nợ
      phải trả đối với chủ nợ:
      \(Payable_{person}(t) = Payable_{person}(t-1) - amount\).
    \item
      2c6. Hệ thống lưu bản ghi loại \texttt{repay} vào bảng
      \texttt{debt} và hiển thị thông báo thành công kèm dư nợ còn lại.
    \end{itemize}
  \item
    \emph{2d. Xóa nợ khó đòi phi tiền mặt (Write-off Lend / Write-off
    Borrow):}

    \begin{itemize}
    \tightlist
    \item
      2d1. Người dùng chọn một khoản nợ không thể thu hồi hoặc được chủ
      nợ miễn trừ nghĩa vụ và chọn chức năng "Xóa nợ (Write-off)".
    \item
      2d2. Người dùng chọn loại hình xóa nợ (\texttt{write-off\ lend}
      nếu là khoản cho vay mất khả năng thu hồi, hoặc
      \texttt{write-off\ borrow} nếu được chủ nợ xóa nợ), nhập số tiền
      xóa nợ (\(amount > 0\)), và bắt buộc nhập lý do xóa nợ vào ô ghi
      chú.
    \item
      2d3. Hệ thống kiểm tra số tiền xóa nợ không vượt quá dư nợ hiện
      tại của đối tác (\(amount \le Receivable_{person}\) hoặc
      \(amount \le Payable_{person}\)).
    \item
      2d4. Hệ thống ghi giảm số dư nợ của đối tác tương ứng theo quy tắc
      \texttt{BR-FIN-07}.
    \item
      2d5. Hệ thống giữ nguyên 100\% số dư của toàn bộ các ví tiền mặt
      (nghiệp vụ phi tiền mặt theo quy tắc \texttt{BR-FIN-02}).
    \item
      2d6. Hệ thống lưu bản ghi công nợ loại \texttt{write-off\ lend}
      hoặc \texttt{write-off\ borrow} và thông báo hoàn tất xóa nợ.
    \end{itemize}
  \item
    \emph{2e. Xem sao kê công nợ chi tiết theo đối tượng:}

    \begin{itemize}
    \tightlist
    \item
      2e1. Người dùng chọn chức năng "Sao kê công nợ theo cá nhân" và
      chọn một cá nhân đối tác cụ thể.
    \item
      2e2. Hệ thống lọc toàn bộ các bản ghi trong bảng \texttt{debt} có
      trường \texttt{person} tương ứng, sắp xếp theo trình tự thời gian
      tăng dần.
    \item
      2e3. Hệ thống tính số dư nợ lũy kế sau từng giao dịch và tổng hợp:
      Tổng cho vay, Tổng đã thu hồi, Dư nợ phải thu còn lại, Tổng đã
      vay, Tổng đã trả, Dư nợ phải trả còn lại.
    \item
      2e4. Hệ thống hiển thị bảng sao kê chi tiết phục vụ việc đối soát
      và thanh toán.
    \end{itemize}
  \end{itemize}
\item
  \textbf{Các luồng ngoại lệ (Exception Flows):}

  \begin{itemize}
  \tightlist
  \item
    \emph{5a. Số dư ví không đủ để cho vay hoặc trả nợ:}

    \begin{itemize}
    \tightlist
    \item
      5a1. Tại bước 5 (khi cho vay) hoặc bước 2c4 (khi trả nợ), hệ thống
      phát hiện số dư khả dụng của ví trích tiền nhỏ hơn số tiền giao
      dịch (\(Balance_{wallet} < amount\)).
    \item
      5a2. Hệ thống từ chối thực thi và hiển thị thông báo: "Số dư ví
      trích tiền không đủ để thực hiện giao dịch công nợ".
    \item
      5a3. Người dùng chọn ví trích tiền khác hoặc giảm số tiền giao
      dịch.
    \end{itemize}
  \item
    \emph{2b4a. Thu hồi nợ vượt quá dư nợ phải thu
    (\(amount > Receivable_{person}\)):}

    \begin{itemize}
    \tightlist
    \item
      2b4a1. Tại bước 2b4, hệ thống phát hiện số tiền thu nợ lớn hơn số
      dư nợ phải thu còn lại của người vay.
    \item
      2b4a2. Hệ thống chặn lại và thông báo lỗi: "Số tiền thu hồi không
      được vượt quá số dư nợ gốc hiện tại của đối tác".
    \item
      2b4a3. Người dùng điều chỉnh lại số tiền thu nợ chính xác.
    \end{itemize}
  \item
    \emph{2c4a. Trả nợ vượt quá dư nợ phải trả
    (\(amount > Payable_{person}\)):}

    \begin{itemize}
    \tightlist
    \item
      2c4a1. Tại bước 2c4, hệ thống phát hiện số tiền trả nợ lớn hơn
      nghĩa vụ nợ phải trả còn lại đối với chủ nợ.
    \item
      2c4a2. Hệ thống chặn lại và thông báo lỗi: "Số tiền hoàn trả không
      được vượt quá nghĩa vụ nợ hiện tại đối với chủ nợ".
    \item
      2c4a3. Người dùng điều chỉnh lại số tiền trả nợ.
    \end{itemize}
  \item
    \emph{2d3a. Xóa nợ vượt quá dư nợ thực tế:}

    \begin{itemize}
    \tightlist
    \item
      2d3a1. Tại bước 2d3, số tiền xóa nợ nhập vào vượt quá số dư nợ của
      đối tác.
    \item
      2d3a2. Hệ thống từ chối và báo lỗi số tiền xóa nợ vượt quá giới
      hạn dư nợ cho phép.
    \end{itemize}
  \item
    \emph{5b. Nhập số tiền nhỏ hơn hoặc bằng 0 hoặc bỏ trống tên đối
    tác:}

    \begin{itemize}
    \tightlist
    \item
      5b1. Hệ thống phát hiện \(amount \le 0\) hoặc trường tên đối tác
      \texttt{person} để trống.
    \item
      5b2. Hệ thống đánh dấu lỗi tại các trường tương ứng và yêu cầu
      Người dùng điền đầy đủ dữ liệu hợp lệ.
    \end{itemize}
  \end{itemize}
\item
  \textbf{Quy tắc nghiệp vụ đặc thù (Special Requirements /
  Invariants):}

  \begin{enumerate}
  \def\labelenumi{\arabic{enumi}.}
  \tightlist
  \item
    Bất biến số dư nợ không âm theo \texttt{BR-FIN-07}:
    \(Receivable \ge 0\) và \(Payable \ge 0\) tại mọi cấp độ (toàn hệ
    thống và từng cá nhân).
  \item
    Bất biến dòng tiền ví theo \texttt{BR-FIN-02}: Nghiệp vụ
    \texttt{lend}, \texttt{borrow}, \texttt{collect}, \texttt{repay} tác
    động trực tiếp đến số dư ví; trong khi nghiệp vụ \texttt{write-off}
    tuyệt đối không làm thay đổi số dư ví.
  \item
    Bất biến độc lập với thu chi tiêu dùng: Cho vay và thu nợ không phải
    là chi tiêu/thu nhập; đi vay và trả nợ không phải là thu nhập/chi
    tiêu. Không được tính gộp vào báo cáo kết quả hoạt động tiêu dùng
    hàng tháng.
  \item
    Phù hợp các yêu cầu chức năng: \texttt{FR-DEBT-01},
    \texttt{FR-DEBT-02}, \texttt{FR-DEBT-03}, \texttt{FR-DEBT-04},
    \texttt{FR-DEBT-05}, \texttt{FR-REP-04}.
  \end{enumerate}
\end{itemize}

\subsubsection{6.6 UC06: Quản lý sổ tiết kiệm có kỳ hạn và tính toán lãi
suất (Term Savings
Management)}\label{66-uc06-quux1ea3n-luxfd-sux1ed5-tiux1ebft-kiux1ec7m-cuxf3-kux1ef3-hux1ea1n-vuxe0-tuxednh-touxe1n-luxe3i-suux1ea5t-term-savings-management}

\begin{itemize}
\tightlist
\item
  \textbf{Mã và Tên ca sử dụng:}
  \texttt{UC06:\ Quản\ lý\ sổ\ tiết\ kiệm\ có\ kỳ\ hạn\ và\ tính\ toán\ lãi\ suất\ (Term\ Savings\ Management)}
\item
  \textbf{Tác nhân chính:} Chủ sở hữu hệ thống (Owner)
\item
  \textbf{Mô tả tóm tắt:} Ca sử dụng này cho phép Owner mở sổ tiết kiệm
  mới trích vốn từ ví thanh khoản; tự động theo dõi tiến độ kỳ hạn, tính
  toán tiền lãi dự kiến tích lũy theo ngày và giá trị ghi nhận hiện tại;
  thực hiện tất toán đúng hạn (thu hồi trọn vẹn vốn gốc cùng 100\% lãi
  cam kết về ví chỉ định); và xử lý tất toán trước hạn với cơ chế hồi tố
  lãi suất không kỳ hạn theo quy định ngân hàng tại BR-FIN-05, giải
  phóng hoặc đồng bộ hóa liên kết với quỹ mục tiêu tài chính tại
  BR-FIN-08.
\item
  \textbf{Điều kiện kích hoạt:} Owner có khoản tiền nhàn rỗi cần gửi
  tiết kiệm tại ngân hàng, cần kiểm tra tình trạng sinh lời của các sổ
  tiết kiệm hiện có, hoặc thực hiện đóng sổ (đúng hạn hoặc trước hạn) để
  thu hồi dòng tiền.
\item
  \textbf{Tiền điều kiện:}

  \begin{enumerate}
  \def\labelenumi{\arabic{enumi}.}
  \tightlist
  \item
    Hệ thống có ít nhất một ví tiền ở trạng thái \texttt{ACTIVE} có số
    dư khả dụng đủ để trích nợ vốn gốc mở sổ.
  \item
    Bảng chuẩn hóa kỳ hạn (\texttt{database.term}) đã nạp đầy đủ các kỳ
    hạn chuẩn (\texttt{unlimited}, \texttt{1M}, \texttt{3M},
    \texttt{6M}, \texttt{12M}, \texttt{24M}, \texttt{36M},
    \texttt{60M}).
  \item
    Owner đang ở giao diện phân hệ Quản lý Sổ tiết kiệm.
  \end{enumerate}
\item
  \textbf{Hậu điều kiện thành công:}

  \begin{enumerate}
  \def\labelenumi{\arabic{enumi}.}
  \tightlist
  \item
    Khi mở sổ: Sổ tiết kiệm mới được lưu vào bảng \texttt{savings} với
    trạng thái \texttt{OPEN}, số dư ví trích tiền giảm đúng bằng số vốn
    gốc \(Principal\), bảo toàn bất biến dòng tiền theo
    \texttt{BR-FIN-02}.
  \item
    Khi tính toán: Hệ thống phản ánh chính xác số ngày gửi thực tế
    \(Days\), tiến độ kỳ hạn \(Process\), lãi dự kiến toàn kỳ
    \(ExpectedInterest\), và giá trị hiện tại \(Balance\).
  \item
    Khi tất toán đúng hạn: Sổ tiết kiệm chuyển trạng thái
    \texttt{SETTLED\_MATURED}, ví thanh khoản nhận trọn vẹn
    \(Payout = Principal + ExpectedInterest\).
  \item
    Khi tất toán trước hạn: Sổ tiết kiệm chuyển trạng thái
    \texttt{SETTLED\_EARLY}, ví thanh khoản nhận
    \(Payout = Principal + ActualInterest\) (theo lãi suất không kỳ
    hạn), phần lãi kỳ hạn chênh lệch được hạch toán là khoản lãi bị hủy
    bỏ.
  \item
    Liên kết quỹ mục tiêu tài chính (nếu có) được cập nhật đồng bộ theo
    \texttt{BR-FIN-08}.
  \end{enumerate}
\item
  \textbf{Luồng sự kiện chính (Main Success Scenario) - Mở sổ tiết kiệm
  mới:}

  \begin{enumerate}
  \def\labelenumi{\arabic{enumi}.}
  \tightlist
  \item
    Người dùng chọn chức năng "Mở sổ tiết kiệm mới".
  \item
    Hệ thống hiển thị biểu mẫu mở sổ gồm các trường: Tên ngân hàng / Tổ
    chức tín dụng (\texttt{bank}), Số vốn gốc gửi (\texttt{principal}),
    Kỳ hạn gửi (\texttt{term} - danh sách chọn từ bảng chuẩn hóa), Lãi
    suất cam kết (\texttt{rate} - \%/năm), Ngày bắt đầu gửi
    (\texttt{start} - mặc định hôm nay), Ví trích tiền
    (\texttt{walletId} - danh sách ví \texttt{ACTIVE}), Quỹ mục tiêu
    liên kết (\texttt{fund} - tùy chọn), Ghi chú mô tả.
  \item
    Người dùng nhập tên ngân hàng, nhập số tiền vốn gốc
    (\(principal > 0\)), chọn kỳ hạn gửi (ví dụ: \texttt{6M}), nhập lãi
    suất năm cam kết (ví dụ: \texttt{6.0\%}), chọn ngày bắt đầu, chọn ví
    nguồn trích tiền, và chọn quỹ mục tiêu liên kết (nếu có).
  \item
    Người dùng nhấn nút "Xác nhận mở sổ".
  \item
    Hệ thống xác thực nghiệp vụ:

    \begin{itemize}
    \tightlist
    \item
      Tên ngân hàng không được để trống (độ dài từ 1 đến 100 ký tự).
    \item
      Số vốn gốc \(principal > 0\) và lãi suất \(rate > 0\), chuẩn hóa
      về kiểu \texttt{BigDecimal} theo quy tắc \texttt{BR-FIN-01}.
    \item
      Kiểm tra kỳ hạn thuộc danh mục hợp lệ và tra cứu số ngày kỳ hạn
      (\(TermDays\)) từ bảng \texttt{database.term}.
    \item
      Kiểm tra số dư khả dụng của ví trích tiền theo quy tắc
      \texttt{BR-FIN-02}: \(Balance_{wallet} - principal \ge 0\).
    \end{itemize}
  \item
    Hệ thống thực hiện bút toán trích tiền: Trừ số vốn gốc \(principal\)
    khỏi số dư khả dụng của ví trích tiền được chọn.
  \item
    Hệ thống xác định ngày đáo hạn: \(MaturityDate = start + TermDays\)
    (đối với kỳ hạn xác định).
  \item
    Hệ thống khởi tạo bản ghi mới trong bảng \texttt{savings} với mã
    \texttt{savingId} duy nhất, lưu các thuộc tính: \texttt{bank},
    \texttt{principal}, \texttt{term}, \texttt{rate}, \texttt{start},
    \texttt{maturityDate}, \texttt{walletId}, \texttt{fund}, trạng thái
    \texttt{status\ =\ OPEN}, số dư ghi nhận ban đầu
    \(Balance = principal\).
  \item
    Nếu người dùng chỉ định quỹ mục tiêu (\texttt{fund}), hệ thống ghi
    nhận liên kết sổ tiết kiệm này vào quỹ mục tiêu tương ứng theo quy
    tắc \texttt{BR-FIN-08}, tự động cộng giá trị sổ vào tài sản quỹ.
  \item
    Hệ thống hiển thị thông báo mở sổ tiết kiệm thành công kèm thông tin
    chi tiết: mã sổ, ngày đáo hạn dự kiến, lãi dự kiến toàn kỳ và số dư
    mới của ví nguồn.
  \end{enumerate}
\item
  \textbf{Các luồng rẽ nhánh (Alternative Flows):}

  \begin{itemize}
  \tightlist
  \item
    \emph{2a. Xem danh sách sổ tiết kiệm và tính toán tiến độ, lãi dự
    kiến:}

    \begin{itemize}
    \tightlist
    \item
      2a1. Người dùng mở màn hình danh sách sổ tiết kiệm, chọn bộ lọc
      trạng thái (\texttt{OPEN}, \texttt{SETTLED\_MATURED},
      \texttt{SETTLED\_EARLY}, hoặc \texttt{ALL}).
    \item
      2a2. Với từng sổ tiết kiệm đang ở trạng thái \texttt{OPEN}, hệ
      thống lấy ngày hệ thống hiện tại (\(currentDate\)) và tự động tính
      toán theo quy tắc \texttt{BR-FIN-05}:

      \begin{itemize}
      \tightlist
      \item
        Số ngày thực gửi: \(Days = \max(0, currentDate - start)\).
      \item
        Nếu là gửi không kỳ hạn (\texttt{unlimited} hoặc
        \(TermDays = 0\)):
        \[Process = 1.0, \quad Interest = \frac{Principal \times Rate \times Days}{365}, \quad Balance = Principal + Interest\]
      \item
        Nếu là gửi có kỳ hạn (\(TermDays > 0\)):
        \[Process = \min\left(1.0, \frac{Days}{TermDays}\right)\]
        \[ExpectedInterest = \frac{Principal \times Rate \times TermDays}{365}\]
        \[Balance = Principal + (ExpectedInterest \times Process)\]
        \[RemainingDays = \max(0, TermDays - Days)\]
      \end{itemize}
    \item
      2a3. Hệ thống tổng hợp tổng vốn gốc đang gửi, tổng lãi dự kiến và
      tổng giá trị hiện tại của toàn bộ danh mục tiết kiệm bằng phương
      thức \texttt{BigDecimal.add()}.
    \item
      2a4. Hệ thống hiển thị bảng danh sách sổ tiết kiệm kèm thanh tiến
      độ thời gian, số ngày còn lại, lãi tích lũy và các nút chức năng
      thao tác tương ứng.
    \end{itemize}
  \item
    \emph{2b. Tất toán sổ tiết kiệm đúng hạn (Matured Settlement):}

    \begin{itemize}
    \tightlist
    \item
      2b1. Tại danh sách sổ tiết kiệm, Người dùng chọn một sổ đang mở có
      ngày hiện tại đã đến hoặc quá ngày đáo hạn
      (\(currentDate \ge MaturityDate\)) và nhấn nút "Tất toán đúng
      hạn".
    \item
      2b2. Hệ thống hiển thị biểu mẫu tất toán đúng hạn gồm: Mã sổ, Ngân
      hàng, Vốn gốc, Lãi suất, Số ngày gửi thực tế, Tiền lãi cam kết
      trọn kỳ (\(ExpectedInterest\)), Tổng số tiền thực nhận
      (\(Payout = Principal + ExpectedInterest\)), và danh sách chọn Ví
      thanh khoản nhận tiền.
    \item
      2b3. Người dùng chọn ví nhận tiền (trạng thái \texttt{ACTIVE}) và
      nhấn "Xác nhận tất toán".
    \item
      2b4. Hệ thống kiểm tra điều kiện ngày tất toán: bảo đảm
      \(currentDate \ge MaturityDate\).
    \item
      2b5. Hệ thống cộng số tiền \(Payout\) vào số dư khả dụng của ví
      nhận tiền theo quy tắc \texttt{BR-FIN-02}.
    \item
      2b6. Hệ thống cập nhật trạng thái sổ tiết kiệm sang
      \texttt{SETTLED\_MATURED}, ghi nhận ngày tất toán thực tế
      \texttt{settledDate\ =\ currentDate}, đóng sổ và ngừng tính lãi.
    \item
      2b7. Nếu sổ có liên kết với quỹ mục tiêu, hệ thống tự động gỡ liên
      kết sổ tiết kiệm khỏi quỹ mục tiêu đó theo quy tắc
      \texttt{BR-FIN-08}, trừ giá trị sổ khỏi quỹ tương ứng.
    \item
      2b8. Hệ thống hiển thị thông báo tất toán thành công kèm biên lai
      tóm tắt số tiền vốn gốc và tiền lãi đã nhập ví.
    \end{itemize}
  \item
    \emph{2c. Tất toán sổ tiết kiệm trước hạn (Early Settlement):}

    \begin{itemize}
    \tightlist
    \item
      2c1. Người dùng chọn một sổ tiết kiệm đang mở nhưng chưa đến ngày
      đáo hạn (\(currentDate < MaturityDate\)) và nhấn nút "Tất toán
      trước hạn".
    \item
      2c2. Hệ thống hiển thị hộp thoại cảnh báo tất toán trước hạn kèm
      biểu mẫu nhập: Ngày tất toán (mặc định hôm nay), Lãi suất không kỳ
      hạn áp dụng hồi tố (\texttt{demandRate}, mặc định 0.2\%/năm, cho
      phép chỉnh sửa), Ví nhận tiền, và ô đánh dấu bắt buộc "Tôi hiểu và
      chấp nhận mất quyền hưởng lãi suất có kỳ hạn".
    \item
      2c3. Hệ thống tính toán và hiển thị bảng đối chiếu tài chính minh
      bạch:

      \begin{itemize}
      \tightlist
      \item
        Số ngày thực gửi: \(Days = currentDate - start\).
      \item
        Lãi cam kết ban đầu bị mất:
        \(ExpectedInterest = \frac{Principal \times Rate \times TermDays}{365}\).
      \item
        Lãi không kỳ hạn thực nhận:
        \(ActualInterest = \frac{Principal \times demandRate \times Days}{365}\).
      \item
        Khoản thiệt hại lãi:
        \(\Delta Loss = ExpectedInterest - ActualInterest\).
      \item
        Tổng tiền thực nhận chuyển về ví:
        \(Payout = Principal + ActualInterest\).
      \end{itemize}
    \item
      2c4. Người dùng kiểm tra số liệu, chọn ví nhận tiền, tích chọn ô
      cam kết chấp nhận mất lãi và nhấn "Xác nhận tất toán trước hạn".
    \item
      2c5. Hệ thống cộng số tiền \(Payout\) vào số dư khả dụng của ví
      nhận tiền theo quy tắc \texttt{BR-FIN-02}.
    \item
      2c6. Hệ thống cập nhật trạng thái sổ tiết kiệm sang
      \texttt{SETTLED\_EARLY}, ghi nhận \texttt{settledDate},
      \texttt{actualInterest}, và khoản lãi bị mất
      \texttt{lostInterest\ =\ \textbackslash{}Delta\ Loss}.
    \item
      2c7. Hệ thống gỡ liên kết sổ tiết kiệm khỏi quỹ mục tiêu liên quan
      (nếu có) theo quy tắc \texttt{BR-FIN-08}.
    \item
      2c8. Hệ thống hiển thị thông báo tất toán trước hạn hoàn tất, ghi
      nhận vào sổ nhật ký tài chính.
    \end{itemize}
  \end{itemize}
\item
  \textbf{Các luồng ngoại lệ (Exception Flows):}

  \begin{itemize}
  \tightlist
  \item
    \emph{5a. Số dư ví nguồn không đủ để mở sổ
    (\(Balance_{wallet} < principal\)):}

    \begin{itemize}
    \tightlist
    \item
      5a1. Tại bước 5, hệ thống phát hiện số dư khả dụng của ví nguồn
      nhỏ hơn số vốn gốc gửi tiết kiệm.
    \item
      5a2. Hệ thống từ chối mở sổ và hiển thị thông báo lỗi: "Số dư khả
      dụng của ví nguồn không đủ để trích nợ mở sổ tiết kiệm".
    \item
      5a3. Người dùng chọn ví nguồn khác có đủ tiền hoặc điều chỉnh giảm
      số vốn gốc.
    \end{itemize}
  \item
    \emph{5b. Số vốn gốc hoặc lãi suất nhỏ hơn hoặc bằng 0:}

    \begin{itemize}
    \tightlist
    \item
      5b1. Tại bước 5, hệ thống phát hiện \(principal \le 0\) hoặc
      \(rate \le 0\).
    \item
      5b2. Hệ thống báo lỗi: "Vốn gốc và lãi suất gửi tiết kiệm phải là
      số dương lớn hơn 0".
    \item
      5b3. Người dùng nhập lại các giá trị hợp lệ.
    \end{itemize}
  \item
    \emph{5c. Ngày bắt đầu gửi nằm quá xa trong tương lai:}

    \begin{itemize}
    \tightlist
    \item
      5c1. Tại bước 5, hệ thống phát hiện ngày bắt đầu gửi \(start\) lớn
      hơn ngày hiện tại quá 30 ngày.
    \item
      5c2. Hệ thống từ chối và cảnh báo: "Ngày bắt đầu gửi không được
      vượt quá 30 ngày so với thời điểm hiện tại".
    \end{itemize}
  \item
    \emph{2b4a. Tất toán đúng hạn khi sổ chưa đến hạn đáo hạn
    (\(currentDate < MaturityDate\)):}

    \begin{itemize}
    \tightlist
    \item
      2b4a1. Tại bước 2b4, nếu có yêu cầu gọi lệnh tất toán đúng hạn
      nhưng ngày hiện tại vẫn trước ngày đáo hạn.
    \item
      2b4a2. Hệ thống chặn thao tác, thông báo: "Sổ tiết kiệm chưa đến
      ngày đáo hạn cam kết. Vui lòng sử dụng chức năng Tất toán trước
      hạn nếu có nhu cầu rút vốn ngay".
    \end{itemize}
  \item
    \emph{2c4a. Người dùng chưa xác nhận chấp nhận mất lãi khi tất toán
    trước hạn:}

    \begin{itemize}
    \tightlist
    \item
      2c4a1. Tại bước 2c4, Người dùng chưa tích chọn ô xác nhận chịu mất
      lãi kỳ hạn.
    \item
      2c4a2. Hệ thống vô hiệu hóa nút xác nhận và hiển thị nhắc nhở yêu
      cầu xác nhận trước khi tiếp tục.
    \end{itemize}
  \item
    \emph{2b1a. Thao tác trên sổ tiết kiệm đã đóng (Status khác OPEN):}

    \begin{itemize}
    \tightlist
    \item
      2b1a1. Nếu người dùng cố gắng thực hiện tất toán trên một sổ tiết
      kiệm đã có trạng thái \texttt{SETTLED\_MATURED} hoặc
      \texttt{SETTLED\_EARLY}.
    \item
      2b1a2. Hệ thống từ chối thực thi, thông báo sổ tiết kiệm này đã
      được tất toán trước đó.
    \end{itemize}
  \end{itemize}
\item
  \textbf{Quy tắc nghiệp vụ đặc thù (Special Requirements /
  Invariants):}

  \begin{enumerate}
  \def\labelenumi{\arabic{enumi}.}
  \tightlist
  \item
    Bất biến tính toán lãi tiết kiệm \texttt{BR-FIN-05}: Lãi suất tính
    trên cơ sở 365 ngày/năm; làm tròn bằng \texttt{BigDecimal} với
    \texttt{RoundingMode.HALF\_EVEN}.
  \item
    Bất biến bảo toàn số dư ví \texttt{BR-FIN-02}: Mở sổ làm giảm số dư
    ví nguồn; tất toán làm tăng số dư ví đích đúng bằng số tiền thực
    nhận \(Payout\).
  \item
    Bất biến duy nhất của liên kết quỹ mục tiêu \texttt{BR-FIN-08}: Mỗi
    sổ tiết kiệm chỉ được liên kết với tối đa một quỹ mục tiêu tại một
    thời điểm.
  \item
    Bất biến trạng thái vòng đời sổ tiết kiệm: Chuyển dịch trạng thái
    tuân thủ nghiêm ngặt máy trạng thái: \texttt{OPEN} chuyển sang
    \texttt{SETTLED\_MATURED} hoặc \texttt{OPEN} chuyển sang
    \texttt{SETTLED\_EARLY}. Không cho phép tái kích hoạt sổ đã tất
    toán.
  \item
    Phù hợp các yêu cầu chức năng: \texttt{FR-SAV-01},
    \texttt{FR-SAV-02}, \texttt{FR-SAV-03}, \texttt{FR-SAV-04}.
  \end{enumerate}
\end{itemize}

\subsubsection{6.7 UC07: Quản lý danh mục chứng chỉ quỹ đầu tư và biến
động NAV (Fund Certificate Investment
Management)}\label{67-uc07-quux1ea3n-luxfd-danh-mux1ee5c-chux1ee9ng-chux1ec9-quux1ef9-ux111ux1ea7u-tux1b0-vuxe0-biux1ebfn-ux111ux1ed9ng-nav-fund-certificate-investment-management}

\begin{itemize}
\tightlist
\item
  \textbf{Mã và Tên ca sử dụng:}
  \texttt{UC07:\ Quản\ lý\ danh\ mục\ chứng\ chỉ\ quỹ\ đầu\ tư\ và\ biến\ động\ NAV\ (Fund\ Certificate\ Investment\ Management)}
\item
  \textbf{Tác nhân chính:} Chủ sở hữu hệ thống (Owner)
\item
  \textbf{Mô tả tóm tắt:} Ca sử dụng này cho phép Owner ghi nhận lệnh
  mua chứng chỉ quỹ mở (cổ phiếu, trái phiếu, cân bằng), trích tiền từ
  ví thanh khoản và tính toán số lượng đơn vị quỹ sở hữu; cập nhật định
  kỳ giá trị tài sản ròng (NAV) theo thị trường; tự động định giá danh
  mục theo nguyên tắc Mark-to-Market, xác định lãi lỗ chưa thực hiện và
  suất sinh lời quy đổi năm (CAGR/RATE) theo BR-FIN-06; và ghi nhận lệnh
  bán chứng chỉ quỹ (một phần hoặc toàn bộ), kết chuyển tiền bán về ví
  thanh khoản và hiện thực hóa lợi nhuận đầu tư.
\item
  \textbf{Điều kiện kích hoạt:} Owner thực hiện giao dịch mua bán chứng
  chỉ quỹ trên thị trường, cần cập nhật NAV mới nhất do công ty quản lý
  quỹ công bố, hoặc tra cứu hiệu suất sinh lời của danh mục đầu tư.
\item
  \textbf{Tiền điều kiện:}

  \begin{enumerate}
  \def\labelenumi{\arabic{enumi}.}
  \tightlist
  \item
    Hệ thống có ít nhất một ví tiền ở trạng thái \texttt{ACTIVE} có đủ
    số dư khả dụng khi mua CCQ.
  \item
    Danh mục loại quỹ đầu tư hỗ trợ đầy đủ 3 loại: \texttt{STOCK},
    \texttt{BOND}, \texttt{BALANCE}.
  \item
    Owner đang ở giao diện phân hệ Quản lý Chứng chỉ quỹ đầu tư.
  \end{enumerate}
\item
  \textbf{Hậu điều kiện thành công:}

  \begin{enumerate}
  \def\labelenumi{\arabic{enumi}.}
  \tightlist
  \item
    Khi mua CCQ: Bản ghi đầu tư mới được lưu vào bảng
    \texttt{fund\_certs} với trạng thái \texttt{HOLDING}, số lượng đơn
    vị quỹ \(Units = Principal / NAV_{buy}\) được tính toán với độ chính
    xác tối thiểu 4 chữ số thập phân, số dư ví trích tiền giảm đúng bằng
    số vốn đầu tư theo \texttt{BR-FIN-02}.
  \item
    Khi cập nhật NAV: Giá trị thị trường hiện tại của toàn bộ các lô quỹ
    cùng mã được tính lại tức thời:
    \(Balance = Units \times NAV_{latest}\), lịch sử biến động NAV được
    lưu vào bảng \texttt{fund\_performance}.
  \item
    Khi tra cứu: Lãi lỗ chưa thực hiện (\(UnrealizedPL\)), tỷ suất lợi
    nhuận tuyệt đối và suất sinh lời quy đổi năm (\(AnnualizedRate\) nếu
    \(HoldDays \ge 30\)) được tính toán chính xác theo
    \texttt{BR-FIN-06}.
  \item
    Khi bán CCQ: Số tiền thu về \(Proceeds = sellUnits \times sellNav\)
    được cộng vào ví nhận tiền; số lượng CCQ nắm giữ và vốn gốc được
    khấu trừ theo phương pháp tỷ trọng bình quân; lợi nhuận thực tế
    (\(RealizedPL\)) được ghi nhận vào sổ kết quả tài chính; nếu bán hết
    thì chuyển trạng thái sang \texttt{CLOSED}.
  \end{enumerate}
\item
  \textbf{Luồng sự kiện chính (Main Success Scenario) - Mua chứng chỉ
  quỹ mới:}

  \begin{enumerate}
  \def\labelenumi{\arabic{enumi}.}
  \tightlist
  \item
    Người dùng chọn chức năng "Mua chứng chỉ quỹ".
  \item
    Hệ thống hiển thị biểu mẫu mua chứng chỉ quỹ gồm các trường: Mã
    chứng chỉ quỹ (\texttt{name}, ví dụ: \texttt{DCDS}, \texttt{VESAF},
    \texttt{SSISCA}), Loại tài sản quỹ (\texttt{type}: \texttt{STOCK},
    \texttt{BOND}, \texttt{BALANCE}), Ngày giao dịch (\texttt{start} -
    mặc định hôm nay), Số vốn đầu tư (\texttt{principal}), Đơn giá NAV
    mua (\texttt{buyNav}), Ví trích tiền (\texttt{walletId} - danh sách
    ví \texttt{ACTIVE}), Quỹ mục tiêu liên kết (\texttt{fund} - tùy
    chọn), Ghi chú mô tả.
  \item
    Người dùng nhập mã chứng chỉ quỹ, chọn loại tài sản, chọn ngày giao
    dịch, nhập số tiền đầu tư (\(principal > 0\)), nhập đơn giá NAV mua
    (\(buyNav > 0\)), chọn ví trích tiền và chọn quỹ mục tiêu (nếu có).
  \item
    Người dùng nhấn nút "Xác nhận mua CCQ".
  \item
    Hệ thống xác thực nghiệp vụ:

    \begin{itemize}
    \tightlist
    \item
      Mã chứng chỉ quỹ không được để trống (độ dài từ 1 đến 50 ký tự, tự
      động chuyển thành chữ hoa).
    \item
      Số tiền đầu tư \(principal > 0\) và đơn giá NAV \(buyNav > 0\),
      chuẩn hóa về kiểu \texttt{BigDecimal} theo quy tắc
      \texttt{BR-FIN-01}.
    \item
      Kiểm tra số dư ví trích tiền theo \texttt{BR-FIN-02}:
      \(Balance_{wallet} - principal \ge 0\).
    \end{itemize}
  \item
    Hệ thống trừ số tiền vốn đầu tư \(principal\) khỏi số dư khả dụng
    của ví trích tiền.
  \item
    Hệ thống tính số lượng chứng chỉ quỹ mua được theo quy tắc
    \texttt{BR-FIN-06}: \[Units = \frac{principal}{buyNav}\] Kết quả
    được lưu trữ với scale tối thiểu là 4 chữ số thập phân, sử dụng
    \texttt{RoundingMode.HALF\_EVEN}.
  \item
    Hệ thống tạo bản ghi đầu tư mới trong bảng \texttt{fund\_certs} với
    mã \texttt{certId} duy nhất, lưu trữ các trường: \texttt{name},
    \texttt{type}, \texttt{start}, \texttt{principal}, \texttt{buyNav},
    \texttt{units}, \texttt{currentNav\ =\ buyNav},
    \texttt{balance\ =\ principal}, trạng thái \texttt{HOLDING}, và
    \texttt{fund} liên kết.
  \item
    Hệ thống ghi nhận điểm dữ liệu đầu tiên vào bảng lịch sử hiệu suất
    \texttt{fund\_performance} (ngày giao dịch, mã quỹ, NAV mua).
  \item
    Nếu có chỉ định quỹ mục tiêu, hệ thống đồng bộ hóa giá trị tài sản
    đầu tư vào quỹ theo quy tắc \texttt{BR-FIN-08}.
  \item
    Hệ thống hiển thị thông báo mua thành công kèm chi tiết: số lượng
    CCQ sở hữu, đơn giá mua và số dư mới của ví thanh khoản.
  \end{enumerate}
\item
  \textbf{Các luồng rẽ nhánh (Alternative Flows):}

  \begin{itemize}
  \tightlist
  \item
    \emph{2a. Cập nhật đơn giá NAV thị trường định kỳ:}

    \begin{itemize}
    \tightlist
    \item
      2a1. Người dùng chọn chức năng "Cập nhật NAV thị trường".
    \item
      2a2. Hệ thống hiển thị biểu mẫu cập nhật: Danh sách chọn Mã chứng
      chỉ quỹ đang nắm giữ, Ngày cập nhật (\texttt{updateDate}, mặc định
      hôm nay), Đơn giá NAV mới nhất (\texttt{latestNav}).
    \item
      2a3. Người dùng chọn mã quỹ (ví dụ: \texttt{DCDS}), nhập giá trị
      NAV mới (ví dụ: 58,250 VNĐ/CCQ) và nhấn "Cập nhật".
    \item
      2a4. Hệ thống kiểm tra điều kiện \(latestNav > 0\).
    \item
      2a5. Hệ thống tìm tất cả các khoản đầu tư đang mở
      (\texttt{HOLDING}) thuộc mã quỹ này.
    \item
      2a6. Với từng khoản đầu tư, hệ thống định giá lại theo nguyên tắc
      Mark-to-Market tại quy tắc \texttt{BR-FIN-06}:
      \[Balance = Units \times latestNav\] Cập nhật trường
      \texttt{currentNav\ =\ latestNav} và \texttt{balance\ =\ Balance}.
    \item
      2a7. Hệ thống thêm bản ghi biến động giá vào bảng
      \texttt{fund\_performance} (mã quỹ, ngày cập nhật, đơn giá NAV
      mới).
    \item
      2a8. Nếu các khoản đầu tư này có liên kết với quỹ mục tiêu, hệ
      thống cập nhật lại tổng giá trị tích lũy của quỹ mục tiêu tương
      ứng theo \texttt{BR-FIN-08}.
    \item
      2a9. Hệ thống thông báo cập nhật NAV thành công và làm mới bảng
      định giá danh mục đầu tư trên màn hình.
    \end{itemize}
  \item
    \emph{2b. Tra cứu hiệu suất danh mục và lãi lỗ chưa thực hiện:}

    \begin{itemize}
    \tightlist
    \item
      2b1. Người dùng truy cập tab "Danh mục đầu tư \& Hiệu suất".
    \item
      2b2. Hệ thống nạp danh sách toàn bộ các khoản đầu tư đang nắm giữ
      (\texttt{HOLDING}), lấy ngày hiện tại (\(currentDate\)) và tính
      toán các chỉ số cho từng khoản theo quy tắc \texttt{BR-FIN-06}:

      \begin{itemize}
      \tightlist
      \item
        Lãi lỗ chưa thực hiện: \(UnrealizedPL = Balance - Principal\).
      \item
        Tỷ suất sinh lời tuyệt đối:
        \(AbsoluteReturn = \frac{UnrealizedPL}{Principal} \times 100\%\).
      \item
        Số ngày nắm giữ: \(HoldDays = \max(1, currentDate - start)\).
      \item
        Suất sinh lời quy đổi năm (Annualized Rate / CAGR):

        \begin{itemize}
        \tightlist
        \item
          Nếu \(HoldDays \ge 30\) ngày và \(Principal > 0\):
          \[AnnualizedRate = \left[ \left( \frac{Balance}{Principal} \right)^{\frac{365}{HoldDays}} - 1 \right] \times 100\%\]
        \item
          Nếu \(HoldDays < 30\) ngày: Hiển thị giá trị
          \(AbsoluteReturn\) kèm ghi chú "Chưa đủ 30 ngày để quy đổi
          năm" để ngăn chặn sai lệch thống kê số mũ.
        \end{itemize}
      \end{itemize}
    \item
      2b3. Hệ thống tổng hợp các chỉ tiêu toàn danh mục: Tổng vốn đầu tư
      ban đầu, Tổng giá trị thị trường hiện tại, Tổng lãi lỗ danh nghĩa
      toàn danh mục, Tỷ trọng phân bổ theo từng loại quỹ
      (\texttt{STOCK}, \texttt{BOND}, \texttt{BALANCE}).
    \item
      2b4. Hệ thống hiển thị bảng dữ liệu chi tiết kèm các đồ thị phân
      bổ tỷ trọng và đường cong tăng trưởng NAV theo thời gian.
    \end{itemize}
  \item
    \emph{2c. Ghi nhận bán chứng chỉ quỹ (Tất toán một phần hoặc toàn
    bộ):}

    \begin{itemize}
    \tightlist
    \item
      2c1. Tại bảng danh mục đầu tư, Người dùng chọn một khoản đầu tư cụ
      thể và nhấn nút "Bán chứng chỉ quỹ".
    \item
      2c2. Hệ thống hiển thị biểu mẫu bán gồm: Mã quỹ, Loại tài sản, Số
      lượng CCQ hiện có (\(Units\)), Vốn gốc ban đầu (\(Principal\)),
      Giá vốn bình quân (\(BuyNAV\)), Ngày bán (mặc định hôm nay), Số
      lượng CCQ muốn bán (\texttt{sellUnits}), Đơn giá bán NAV
      (\texttt{sellNav}), Ví nhận tiền (\texttt{walletId}), Ghi chú.
    \item
      2c3. Người dùng nhập số lượng CCQ cần bán
      (\(0 < sellUnits \le Units\)), đơn giá bán NAV (\(sellNav > 0\)),
      chọn ví nhận tiền (trạng thái \texttt{ACTIVE}) và nhấn "Xác nhận
      bán".
    \item
      2c4. Hệ thống kiểm tra tính hợp lệ của số lượng bán:
      \(sellUnits \le Units\).
    \item
      2c5. Hệ thống tính toán các giá trị kế toán theo quy tắc
      \texttt{FR-INV-04}:

      \begin{itemize}
      \tightlist
      \item
        Số tiền bán thu về: \(Proceeds = sellUnits \times sellNav\).
      \item
        Giá vốn tương ứng của lượng CCQ bán:
        \(CostBasis = \frac{sellUnits}{Units} \times Principal\).
      \item
        Lợi nhuận thực tế đã hiện thực hóa:
        \(RealizedPL = Proceeds - CostBasis\).
      \end{itemize}
    \item
      2c6. Hệ thống thực hiện các bút toán đồng bộ:

      \begin{itemize}
      \tightlist
      \item
        Cộng số tiền thu về \(Proceeds\) vào số dư khả dụng của ví nhận
        tiền theo \texttt{BR-FIN-02}.
      \item
        Giảm số lượng CCQ nắm giữ của khoản đầu tư:
        \(Units_{mới} = Units - sellUnits\).
      \item
        Giảm vốn gốc của khoản đầu tư:
        \(Principal_{mới} = Principal - CostBasis\).
      \item
        Tính lại giá trị thị trường còn lại:
        \(Balance_{mới} = Units_{mới} \times sellNav\).
      \item
        Nếu \(Units_{mới} = 0\), chuyển trạng thái khoản đầu tư sang
        \texttt{CLOSED} (đã đóng).
      \end{itemize}
    \item
      2c7. Hệ thống ghi nhận \(RealizedPL\) vào báo cáo kết quả hoạt
      động tài chính trong kỳ.
    \item
      2c8. Nếu khoản đầu tư có liên kết quỹ mục tiêu, hệ thống cập nhật
      lại giá trị quỹ theo \texttt{BR-FIN-08}.
    \item
      2c9. Hệ thống thông báo hoàn tất giao dịch bán, hiển thị biên lai
      tóm tắt gồm số tiền thu về, lãi/lỗ thực tế và số lượng CCQ còn
      lại.
    \end{itemize}
  \end{itemize}
\item
  \textbf{Các luồng ngoại lệ (Exception Flows):}

  \begin{itemize}
  \tightlist
  \item
    \emph{5a. Số dư ví không đủ để thực hiện lệnh mua:}

    \begin{itemize}
    \tightlist
    \item
      5a1. Tại bước 5, hệ thống phát hiện
      \(Balance_{wallet} < principal\).
    \item
      5a2. Hệ thống từ chối lệnh mua và hiển thị thông báo: "Số dư ví
      nguồn không đủ để mua chứng chỉ quỹ".
    \item
      5a3. Người dùng đổi ví khác hoặc giảm số vốn mua.
    \end{itemize}
  \item
    \emph{5b. Đơn giá NAV hoặc số tiền đầu tư nhỏ hơn hoặc bằng 0:}

    \begin{itemize}
    \tightlist
    \item
      5b1. Tại bước 5 (khi mua), bước 2a4 (khi cập nhật NAV) hoặc bước
      2c4 (khi bán), hệ thống ghi nhận \(principal \le 0\),
      \(buyNav \le 0\), \(latestNav \le 0\), hoặc \(sellNav \le 0\).
    \item
      5b2. Hệ thống báo lỗi: "Đơn giá NAV và số tiền giao dịch phải là
      các số dương lớn hơn 0".
    \item
      5b3. Người dùng nhập lại đơn giá hợp lệ.
    \end{itemize}
  \item
    \emph{2c4a. Số lượng bán vượt quá số lượng CCQ sở hữu
    (\(sellUnits > Units\)):}

    \begin{itemize}
    \tightlist
    \item
      2c4a1. Tại bước 2c4, Người dùng nhập số lượng bán lớn hơn số lượng
      CCQ đang có trong khoản đầu tư.
    \item
      2c4a2. Hệ thống từ chối và thông báo: "Số lượng chứng chỉ quỹ bán
      vượt quá số lượng hiện đang nắm giữ".
    \item
      2c4a3. Người dùng điều chỉnh số lượng bán nhỏ hơn hoặc bằng số
      lượng hiện có.
    \end{itemize}
  \item
    \emph{2c1a. Cố gắng bán khoản đầu tư đã đóng
    (\texttt{status\ =\ CLOSED}):}

    \begin{itemize}
    \tightlist
    \item
      2c1a1. Người dùng chọn thao tác bán trên khoản đầu tư đã có số
      lượng bằng 0 và trạng thái \texttt{CLOSED}.
    \item
      2c1a2. Hệ thống từ chối thực hiện và thông báo khoản đầu tư đã
      được tất toán toàn bộ.
    \end{itemize}
  \end{itemize}
\item
  \textbf{Quy tắc nghiệp vụ đặc thù (Special Requirements /
  Invariants):}

  \begin{enumerate}
  \def\labelenumi{\arabic{enumi}.}
  \tightlist
  \item
    Bất biến định giá theo thị trường \texttt{BR-FIN-06}: Giá trị ghi
    nhận \(Balance\) luôn phản ánh chính xác tích số giữa số lượng đơn
    vị quỹ và NAV tại ngày định giá gần nhất.
  \item
    Bất biến độ chính xác số học \texttt{BR-FIN-01}: Số lượng đơn vị quỹ
    \(Units\) được duy trì scale tối thiểu là 4 chữ số thập phân; số
    tiền quy đổi và lợi nhuận làm tròn theo
    \texttt{RoundingMode.HALF\_EVEN}.
  \item
    Bất biến bảo toàn dòng tiền \texttt{BR-FIN-02}: Mua CCQ làm giảm số
    dư ví nguồn đúng bằng \(Principal\); bán CCQ làm tăng số dư ví đích
    đúng bằng \(Proceeds\).
  \item
    Bất biến liên kết tài sản quỹ mục tiêu \texttt{BR-FIN-08}: Mỗi khoản
    đầu tư CCQ chỉ được thuộc về duy nhất một quỹ mục tiêu để bảo đảm
    không tính trùng tài sản.
  \item
    Phù hợp các yêu cầu chức năng: \texttt{FR-INV-01},
    \texttt{FR-INV-02}, \texttt{FR-INV-03}, \texttt{FR-INV-04}.
  \end{enumerate}
\end{itemize}

\subsubsection{6.8 UC08: Quản lý quỹ mục tiêu tài chính và phân bổ dòng
tiền (Financial Goal \& Fund
Management)}\label{68-uc08-quux1ea3n-luxfd-quux1ef9-mux1ee5c-tiuxeau-tuxe0i-chuxednh-vuxe0-phuxe2n-bux1ed5-duxf2ng-tiux1ec1n-financial-goal--fund-management}

\begin{itemize}
\tightlist
\item
  \textbf{Mã và Tên ca sử dụng:}
  \texttt{UC08:\ Quản\ lý\ quỹ\ mục\ tiêu\ tài\ chính\ và\ phân\ bổ\ dòng\ tiền\ (Financial\ Goal\ \&\ Fund\ Management)}
\item
  \textbf{Tác nhân chính:} Chủ sở hữu hệ thống (Owner)
\item
  \textbf{Mô tả tóm tắt:} Ca sử dụng này cho phép Owner thiết lập các
  quỹ mục tiêu tài chính độc lập (quỹ khẩn cấp, quỹ y tế sức khỏe, quỹ
  mua sắm mục tiêu, quỹ hưu trí dài hạn); phân bổ tài sản đa kênh gồm 3
  cấu phần (tiền mặt phân bổ từ bảng \texttt{funds\_cash}, sổ tiết kiệm
  liên kết, chứng chỉ quỹ liên kết) theo quy tắc BR-FIN-08; theo dõi
  tiến độ tích lũy theo thời gian thực
  (\(Total / Target \times 100\%\)); và thực hiện điều chỉnh hạn mức mục
  tiêu hoặc tái phân bổ tài sản giữa các quỹ nhằm thích ứng với biến
  động kế hoạch mà không làm biến dạng số dư sổ cái thực tế.
\item
  \textbf{Điều kiện kích hoạt:} Owner xây dựng kế hoạch tài chính cá
  nhân mới, muốn phân bổ nguồn lực tài sản hiện có vào các mục tiêu cụ
  thể, hoặc cần theo dõi, tái cân bằng danh mục quỹ mục tiêu.
\item
  \textbf{Tiền điều kiện:}

  \begin{enumerate}
  \def\labelenumi{\arabic{enumi}.}
  \tightlist
  \item
    Hệ thống đã khởi tạo sẵn cấu trúc 4 quỹ chuẩn (\texttt{emergency},
    \texttt{health}, \texttt{target}, \texttt{retirement}) hoặc cho phép
    tạo thêm quỹ tùy chỉnh.
  \item
    Owner đang ở giao diện phân hệ Quản lý Quỹ mục tiêu tài chính.
  \end{enumerate}
\item
  \textbf{Hậu điều kiện thành công:}

  \begin{enumerate}
  \def\labelenumi{\arabic{enumi}.}
  \tightlist
  \item
    Quỹ mục tiêu mới hoặc thông tin cập nhật được lưu vào bảng
    \texttt{funds} với mã định danh \texttt{fundCode} duy nhất.
  \item
    Tổng giá trị tích lũy của quỹ \(Total_{fund}\) được tính toán tự
    động từ 3 cấu phần
    (\(Cash_{fund} + \sum Balance_{savings} + \sum Balance_{fundCerts}\))
    theo quy tắc \texttt{BR-FIN-08}.
  \item
    Tiến độ hoàn thành \(Process\) được tính toán chính xác; tự động
    hiển thị trạng thái "ĐÃ HOÀN THÀNH" khi \(Process \ge 100\%\).
  \item
    Nguyên tắc không tính trùng tài sản (No Double-Counting) được bảo
    đảm tuyệt đối: không có sổ tiết kiệm hay khoản CCQ nào bị tính vào
    hai quỹ cùng lúc.
  \item
    Tổng lượng tiền mặt phân bổ trên toàn bộ các quỹ mục tiêu không vượt
    quá tổng số dư tiền mặt khả dụng trên tất cả các ví thanh khoản của
    hệ thống.
  \end{enumerate}
\item
  \textbf{Luồng sự kiện chính (Main Success Scenario) - Thiết lập và
  phân bổ quỹ mục tiêu:}

  \begin{enumerate}
  \def\labelenumi{\arabic{enumi}.}
  \tightlist
  \item
    Người dùng chọn chức năng "Thiết lập quỹ mục tiêu mới".
  \item
    Hệ thống hiển thị biểu mẫu thiết lập quỹ gồm các trường: Mã quỹ
    (\texttt{fundCode}), Tên quỹ mục tiêu (\texttt{fundName}), Số tiền
    mục tiêu cần tích lũy (\texttt{targetAmount}), Hạn chót hoàn thành
    (\texttt{deadline} - tùy chọn), Mô tả mục tiêu.
  \item
    Người dùng nhập mã quỹ (ví dụ: \texttt{study\_abroad}), tên quỹ (ví
    dụ: "Quỹ du học nâng cao trình độ"), số tiền mục tiêu (ví dụ:
    150,000,000 VNĐ), chọn hạn chót trong tương lai và nhập mô tả định
    hướng.
  \item
    Người dùng nhấn nút "Lưu thiết lập quỹ".
  \item
    Hệ thống xác thực dữ liệu:

    \begin{itemize}
    \tightlist
    \item
      Mã quỹ \texttt{fundCode} không trùng lặp với các quỹ đã có trong
      hệ thống (chuỗi không khoảng trắng, tự động chuyển chữ thường).
    \item
      Tên quỹ không được để trống (độ dài từ 1 đến 100 ký tự).
    \item
      Số tiền mục tiêu \(targetAmount \ge 0\), chuẩn hóa về kiểu
      \texttt{BigDecimal} theo quy tắc \texttt{BR-FIN-01}.
    \end{itemize}
  \item
    Hệ thống tạo bản ghi quỹ mới trong bảng \texttt{funds} với 3 cấu
    phần ban đầu bằng 0 (\(Cash = 0, Savings = 0, FundCerts = 0\)), tiến
    độ ban đầu \(Process = 0.00\%\).
  \item
    Hệ thống chuyển sang màn hình "Phân bổ tài sản vào quỹ" và hiển thị
    các nguồn lực tài sản khả dụng:

    \begin{itemize}
    \tightlist
    \item
      Số dư tiền mặt khả dụng chưa phân bổ trong hệ thống:
      \[
      UnallocatedCash = \sum Balance_{wallets} - \sum Cash_{funds}
      \]
    \item
      Danh sách các sổ tiết kiệm đang mở chưa liên kết với quỹ nào.
    \item
      Danh sách các khoản chứng chỉ quỹ đang nắm giữ chưa liên kết với
      quỹ nào.
    \end{itemize}
  \item
    Người dùng nhập số tiền mặt muốn phân bổ từ ví vào quỹ này
    (\(cashAmount\)), tích chọn các sổ tiết kiệm muốn gắn vào quỹ, và
    tích chọn các khoản CCQ muốn gắn vào quỹ.
  \item
    Người dùng nhấn nút "Xác nhận phân bổ tài sản".
  \item
    Hệ thống xác thực nghiệp vụ phân bổ:

    \begin{itemize}
    \tightlist
    \item
      Kiểm tra tiền mặt: \(cashAmount \le UnallocatedCash\).
    \item
      Kiểm tra tài sản liên kết: Bảo đảm các sổ tiết kiệm và CCQ được
      chọn không bị liên kết trùng lặp theo quy tắc \texttt{BR-FIN-08}.
    \end{itemize}
  \item
    Hệ thống cập nhật bảng \texttt{funds\_cash}, lưu trường
    \texttt{cashAmount} cho quỹ tương ứng.
  \item
    Hệ thống cập nhật trường \texttt{fund\ =\ fundCode} cho các bản ghi
    sổ tiết kiệm trong bảng \texttt{savings} và các bản ghi CCQ trong
    bảng \texttt{fund\_certs}.
  \item
    Hệ thống tính toán lại tổng tài sản tích lũy thực tế
    \(Total_{fund}\) và tỷ lệ tiến độ \(Process\) theo quy tắc
    \texttt{BR-FIN-08}:
    \[Total_{fund} = Cash_{fund} + \sum_{s \in Savings_{fund}} Balance_s + \sum_{c \in FundCerts_{fund}} Balance_c\]
    \[Process = \frac{Total_{fund}}{targetAmount} \times 100\%\]
  \item
    Hệ thống thông báo phân bổ tài sản thành công, cập nhật giao diện
    tổng quan các quỹ mục tiêu kèm thanh tiến độ trực quan.
  \end{enumerate}
\item
  \textbf{Các luồng rẽ nhánh (Alternative Flows):}

  \begin{itemize}
  \tightlist
  \item
    \emph{2a. Tra cứu tiến độ hoàn thành và cơ cấu 3 cấu phần tài sản:}

    \begin{itemize}
    \tightlist
    \item
      2a1. Người dùng mở màn hình tổng quan "Quỹ mục tiêu tài chính".
    \item
      2a2. Hệ thống quét dữ liệu và tính toán động các chỉ số cho toàn
      bộ danh sách quỹ (\texttt{emergency}, \texttt{health},
      \texttt{target}, \texttt{retirement} và các quỹ tùy chỉnh):

      \begin{itemize}
      \tightlist
      \item
        Cấu phần tiền mặt: Lấy giá trị \(Cash_{fund}\) từ bảng
        \texttt{funds\_cash}.
      \item
        Cấu phần sổ tiết kiệm: Quét bảng \texttt{savings} có
        \texttt{fund\ =\ fundCode} và trạng thái \texttt{OPEN}, lấy tổng
        \(Balance_s\) tính theo \texttt{FR-SAV-02}.
      \item
        Cấu phần đầu tư: Quét bảng \texttt{fund\_certs} có
        \texttt{fund\ =\ fundCode} và trạng thái \texttt{HOLDING}, lấy
        tổng \(Balance_c\) định giá theo \texttt{FR-INV-03}.
      \item
        Tổng tài sản tích lũy:
        \[
        Total_{fund} = Cash_{fund} + Savings_{fund} + FundCerts_{fund}
        \]
      \item
        Tỷ lệ tiến độ:
        \[
        Process = \left(\frac{Total_{fund}}{Target}\right) \times 100\%
        \]
        Nếu
        \(Target = 0\), \(Process = 100\%\).
      \item
        Số tiền còn thiếu:
        \(Remaining = \max(0, Target - Total_{fund})\).
      \end{itemize}
    \item
      2a3. Hệ thống phân loại trạng thái quỹ: Nếu \(Process \ge 100\%\),
      gắn nhãn "ĐÃ ĐẠT MỤC TIÊU"; nếu \(Process < 100\%\), gắn nhãn
      "ĐANG TÍCH LŨY".
    \item
      2a4. Hệ thống hiển thị bảng so sánh tổng hợp, biểu đồ tiến độ
      thanh ngang và phân rã màu sắc cho 3 cấu phần tài sản của từng
      quỹ.
    \end{itemize}
  \item
    \emph{2b. Điều chỉnh hạn mức mục tiêu và thời hạn hoàn thành (Adjust
    Goal):}

    \begin{itemize}
    \tightlist
    \item
      2b1. Người dùng chọn một quỹ mục tiêu và nhấn nút "Điều chỉnh mục
      tiêu".
    \item
      2b2. Hệ thống hiển thị biểu mẫu cho phép sửa: Số tiền mục tiêu mới
      (\texttt{newTarget}), Hạn chót mới (\texttt{newDeadline}), Ghi chú
      định hướng.
    \item
      2b3. Người dùng nhập các thông số mới và nhấn "Lưu điều chỉnh".
    \item
      2b4. Hệ thống kiểm tra \(newTarget \ge 0\).
    \item
      2b5. Hệ thống cập nhật bảng \texttt{funds} với hạn mức mới, tính
      toán lại tỷ lệ tiến độ
      \(Process_{mới} = (Total_{fund} / newTarget) \times 100\%\).
    \item
      2b6. Thao tác này hoàn toàn không làm biến động số dư thực tế của
      ví tiền, sổ tiết kiệm hay CCQ.
    \item
      2b7. Hệ thống thông báo điều chỉnh thành công và cập nhật tỷ lệ
      tiến độ mới.
    \end{itemize}
  \item
    \emph{2c. Tái phân bổ tài sản giữa các quỹ mục tiêu (Reallocate
    Assets):}

    \begin{itemize}
    \tightlist
    \item
      2c1. Người dùng chọn chức năng "Tái phân bổ tài sản giữa các quỹ".
    \item
      2c2. Hệ thống hiển thị bảng chuyển giao tài sản:

      \begin{itemize}
      \tightlist
      \item
        Chuyển tiền mặt phân bổ: Cho phép giảm tiền mặt ở Quỹ A và tăng
        tương ứng ở Quỹ B trong bảng \texttt{funds\_cash}.
      \item
        Chuyển sổ tiết kiệm hoặc CCQ: Cho phép chọn một sổ tiết kiệm
        hoặc khoản CCQ đang thuộc Quỹ A và gán lại sang Quỹ B.
      \end{itemize}
    \item
      2c3. Người dùng thực hiện thao tác dịch chuyển tài sản mong muốn
      và nhấn "Xác nhận tái phân bổ".
    \item
      2c4. Hệ thống kiểm tra tính hợp lệ: Số tiền mặt trích khỏi Quỹ A
      không vượt quá số tiền mặt hiện có của Quỹ A
      (\(amount \le Cash_A\)).
    \item
      2c5. Hệ thống cập nhật đồng bộ:

      \begin{itemize}
      \tightlist
      \item
        Nếu chuyển tiền mặt:
        \[
        Cash_A = Cash_A - amount;\qquad Cash_B = Cash_B + amount
        \]
      \item
        Nếu chuyển sổ tiết kiệm hoặc CCQ: Cập nhật trường
        \texttt{fund\ =\ fundCode\_B} cho tài sản đó.
      \end{itemize}
    \item
      2c6. Hệ thống tính toán lại \(Total\) và \(Process\) của cả hai
      Quỹ A và Quỹ B. Tổng tài sản toàn hệ thống được bảo toàn 100\%.
    \item
      2c7. Hệ thống hiển thị thông báo tái phân bổ hoàn tất kèm bảng so
      sánh trước và sau khi tái cơ cấu.
    \end{itemize}
  \end{itemize}
\item
  \textbf{Các luồng ngoại lệ (Exception Flows):}

  \begin{itemize}
  \tightlist
  \item
    \emph{10a. Phân bổ tiền mặt vượt quá tổng số dư tiền mặt khả dụng
    toàn hệ thống:}

    \begin{itemize}
    \tightlist
    \item
      10a1. Tại bước 10 (hoặc 2c4), hệ thống phát hiện số tiền mặt người
      dùng muốn phân bổ vào quỹ lớn hơn số dư tiền mặt chưa phân bổ
      (\(cashAmount > UnallocatedCash\)).
    \item
      10a2. Hệ thống từ chối phân bổ và thông báo: "Số tiền mặt phân bổ
      vượt quá tổng lượng tiền mặt khả dụng hiện có trên tất cả các ví".
    \item
      10a3. Người dùng giảm số tiền mặt phân bổ hoặc nạp thêm tiền vào
      các ví.
    \end{itemize}
  \item
    \emph{10b. Liên kết trùng lặp tài sản vào nhiều quỹ khác nhau:}

    \begin{itemize}
    \tightlist
    \item
      10b1. Tại bước 10, hệ thống phát hiện một sổ tiết kiệm hoặc khoản
      CCQ đang được chọn đã thuộc về một quỹ mục tiêu khác và người dùng
      chưa xác nhận chuyển quyền sở hữu.
    \item
      10b2. Hệ thống hiển thị hộp thoại cảnh báo: "Tài sản này đang
      thuộc về quỹ khác. Bạn có muốn chuyển tài sản sang quỹ này không?
      Thao tác này sẽ làm giảm giá trị của quỹ cũ".
    \item
      10b3. Nếu Người dùng chọn Đồng ý, hệ thống tự động gỡ khỏi quỹ cũ
      và gán sang quỹ mới theo quy tắc \texttt{BR-FIN-08}; nếu Người
      dùng chọn Hủy, hệ thống giữ nguyên liên kết cũ.
    \end{itemize}
  \item
    \emph{5a. Số tiền mục tiêu nhỏ hơn 0:}

    \begin{itemize}
    \tightlist
    \item
      5a1. Tại bước 5 (hoặc 2b4), Người dùng nhập số tiền mục tiêu là số
      âm (\(targetAmount < 0\)).
    \item
      5a2. Hệ thống chặn lại và thông báo lỗi: "Số tiền mục tiêu phải là
      số lớn hơn hoặc bằng 0".
    \end{itemize}
  \item
    \emph{5b. Trùng mã quỹ mục tiêu:}

    \begin{itemize}
    \tightlist
    \item
      5b1. Tại bước 5, hệ thống phát hiện mã quỹ \texttt{fundCode} đã
      tồn tại trong cơ sở dữ liệu.
    \item
      5b2. Hệ thống báo lỗi trùng mã và yêu cầu Người dùng nhập mã định
      danh khác.
    \end{itemize}
  \end{itemize}
\item
  \textbf{Quy tắc nghiệp vụ đặc thù (Special Requirements /
  Invariants):}

  \begin{enumerate}
  \def\labelenumi{\arabic{enumi}.}
  \tightlist
  \item
    Bất biến cấu trúc 3 cấu phần theo \texttt{BR-FIN-08}: Giá trị quỹ
    luôn bằng tổng đại số chính xác của 3 cấu phần
    (\(Cash + Savings + FundCerts\)).
  \item
    Bất biến không tính trùng tài sản (No Double-Counting): Mỗi sổ tiết
    kiệm và khoản CCQ chỉ được thuộc về duy nhất một quỹ mục tiêu tại
    một thời điểm.
  \item
    Bất biến hạn mức tiền mặt: Tổng tiền mặt phân bổ trên toàn bộ các
    quỹ không bao giờ được vượt quá tổng số dư tiền mặt trên toàn bộ các
    ví (\(\sum Cash_{funds} \le \sum Balance_{wallets}\)).
  \item
    Độc lập kế toán: Quỹ mục tiêu là thực thể quy hoạch kế hoạch, không
    phải là ví tiền thực tế; việc tạo, sửa, xóa quỹ mục tiêu tuyệt đối
    không làm thay đổi tổng tài sản ròng hay số dư khả dụng thực tế của
    ví.
  \item
    Phù hợp các yêu cầu chức năng: \texttt{FR-FND-01},
    \texttt{FR-FND-02}, \texttt{FR-FND-03}, \texttt{FR-FND-04}.
  \end{enumerate}
\end{itemize}

\subsubsection{6.9 UC09: Tra cứu bảng thống kê dòng tiền trượt 12 tháng
và phân tích tỷ lệ tiết kiệm (Monthly Statistics \& Cashflow
Analysis)}\label{69-uc09-tra-cux1ee9u-bux1ea3ng-thux1ed1ng-kuxea-duxf2ng-tiux1ec1n-trux1b0ux1ee3t-12-thuxe1ng-vuxe0-phuxe2n-tuxedch-tux1ef7-lux1ec7-tiux1ebft-kiux1ec7m-monthly-statistics--cashflow-analysis}

\begin{itemize}
\tightlist
\item
  \textbf{Mã và Tên ca sử dụng:}
  \texttt{UC09:\ Tra\ cứu\ bảng\ thống\ kê\ dòng\ tiền\ trượt\ 12\ tháng\ và\ phân\ tích\ tỷ\ lệ\ tiết\ kiệm\ (Monthly\ Statistics\ \&\ Cashflow\ Analysis)}
\item
  \textbf{Tác nhân chính:} Chủ sở hữu hệ thống (Owner)
\item
  \textbf{Mô tả tóm tắt:} Ca sử dụng này cho phép Owner tra cứu bảng
  thống kê chuỗi thời gian dòng tiền 12 tháng gần nhất (tương ứng bảng
  \texttt{months\_stat}) tự động trượt theo ngày hệ thống, phản ánh thu
  nhập, chi tiêu, dòng tiền thuần và tỷ lệ tiết kiệm từng tháng; xem báo
  cáo phân tích sâu cơ cấu chi tiêu theo Nhóm nhu cầu và Danh mục của
  tháng hiện tại so với tháng trước; tra cứu báo cáo tổng tài sản và giá
  trị tài sản ròng (Net Worth) hợp nhất 5 cấu phần; và tra cứu bảng sao
  kê công nợ chi tiết theo từng đối tác.
\item
  \textbf{Điều kiện kích hoạt:} Owner truy cập phân hệ Báo cáo \& Thống
  kê để đánh giá sức khỏe tài chính định kỳ, theo dõi xu hướng thu chi
  hoặc chuẩn bị kế hoạch ngân sách cho các tháng tiếp theo.
\item
  \textbf{Tiền điều kiện:}

  \begin{enumerate}
  \def\labelenumi{\arabic{enumi}.}
  \tightlist
  \item
    Hệ thống đã khởi chạy và kết nối với cơ sở dữ liệu giao dịch.
  \item
    Owner đang ở giao diện phân hệ Thống kê và Báo cáo tài chính.
  \end{enumerate}
\item
  \textbf{Hậu điều kiện thành công:}

  \begin{enumerate}
  \def\labelenumi{\arabic{enumi}.}
  \tightlist
  \item
    Bảng dữ liệu \texttt{months\_stat} gồm chính xác 12 dòng liên tục
    lùi từ tháng hiện tại được tổng hợp đầy đủ, không để trống (null)
    bất kỳ trường dữ liệu nào.
  \item
    Dòng tiền thuần \(CashFlow = Income - Expense\) và Tỷ lệ tiết kiệm
    \(SavingsRate\) được tính toán chính xác bằng \texttt{BigDecimal}
    theo quy tắc \texttt{FR-REP-01}.
  \item
    Báo cáo cơ cấu chi tiêu 2 tầng phân cấp (Group -\textgreater{}
    Category) phản ánh đúng tỷ trọng và biến động so với tháng liền kề
    theo \texttt{FR-REP-02}.
  \item
    Báo cáo Giá trị tài sản ròng (Net Worth) hợp nhất chính xác 5 cấu
    phần (Wallets + Savings + Investments + Receivable - Payable) theo
    \texttt{FR-REP-03}.
  \item
    Sao kê công nợ theo đối tác phản ánh đúng lịch sử lũy kế và số dư nợ
    còn lại theo \texttt{FR-REP-04}.
  \end{enumerate}
\item
  \textbf{Luồng sự kiện chính (Main Success Scenario) - Tra cứu bảng
  thống kê trượt 12 tháng:}

  \begin{enumerate}
  \def\labelenumi{\arabic{enumi}.}
  \tightlist
  \item
    Người dùng chọn chức năng "Bảng thống kê dòng tiền trượt 12 tháng".
  \item
    Hệ thống tự động xác định mốc thời gian hiện tại từ đồng hồ hệ thống
    (\texttt{LocalDate.now()}), trích xuất năm và tháng hiện tại
    (\(YearMonth_{now}\)).
  \item
    Hệ thống tạo danh sách 12 mốc tháng liên tục từ
    \(YearMonth_{now} - 11\) đến \(YearMonth_{now}\), sắp xếp theo trình
    tự thời gian tăng dần (neo tháng hiện tại ở dòng cuối cùng).
  \item
    Với từng tháng \(m\) trong chuỗi 12 tháng, hệ thống thực hiện truy
    vấn quét bảng \texttt{transactions}:

    \begin{itemize}
    \tightlist
    \item
      Lọc các giao dịch có ngày nằm trong khoảng từ ngày đầu tháng đến
      ngày cuối tháng của tháng \(m\).
    \item
      Tính tổng thu nhập: \(Income_m = \sum Amount\) của các giao dịch
      có \texttt{Type\ =\ Income}.
    \item
      Tính tổng chi tiêu: \(Expense_m = \sum Amount\) của các giao dịch
      có \texttt{Type\ =\ Expense}.
    \end{itemize}
  \item
    Hệ thống tính toán các chỉ số phái sinh cho từng tháng bằng
    \texttt{BigDecimal} theo quy tắc \texttt{FR-REP-01}:

    \begin{itemize}
    \tightlist
    \item
      Dòng tiền ròng: \(CashFlow_m = Income_m - Expense_m\).
    \item
      Tỷ lệ tiết kiệm:
      \[SavingsRate_m = \begin{cases} \frac{CashFlow_m}{Income_m} \times 100\%, & \text{khi } Income_m > 0 \\ 0\%, & \text{khi } Income_m \le 0 \end{cases}\]
      Sử dụng \texttt{RoundingMode.HALF\_EVEN} với scale hiển thị là 2
      chữ số thập phân.
    \end{itemize}
  \item
    Hệ thống tính tổng lũy kế và mức bình quân tháng của chuỗi 12 tháng:
    Tổng thu nhập 12 tháng, Tổng chi tiêu 12 tháng, Dòng tiền ròng tích
    lũy, và Tỷ lệ tiết kiệm bình quân.
  \item
    Hệ thống hiển thị bảng dữ liệu \texttt{months\_stat} gồm 12 dòng kèm
    dòng tổng cộng / bình quân ở chân bảng, đồng thời hiển thị biểu đồ
    trực quan dạng kết hợp (cột chồng thu chi và đường biểu diễn dòng
    tiền thuần).
  \end{enumerate}
\item
  \textbf{Các luồng rẽ nhánh (Alternative Flows):}

  \begin{itemize}
  \tightlist
  \item
    \emph{2a. Xem báo cáo cơ cấu thu chi theo Nhóm và Danh mục:}

    \begin{itemize}
    \tightlist
    \item
      2a1. Người dùng chọn tab "Cơ cấu thu chi", chọn tháng cần phân
      tích (mặc định tháng hiện tại) và chọn loại luồng tiền
      (\texttt{EXPENSE} hoặc \texttt{INCOME}).
    \item
      2a2. Hệ thống quét các giao dịch trong tháng được chọn có loại
      tương ứng.
    \item
      2a3. Hệ thống tổng hợp số liệu theo 2 tầng phân cấp theo quy tắc
      \texttt{FR-REP-02}:

      \begin{itemize}
      \tightlist
      \item
        Cấp nhóm nhu cầu (\texttt{Group}): Tính tổng số tiền và tỷ trọng
        \% trên tổng thu/chi của tháng.
      \item
        Cấp danh mục chi tiết (\texttt{Category}): Tính tổng số tiền, tỷ
        trọng \% trên nhóm cha và tỷ trọng \% trên toàn bộ tháng.
      \end{itemize}
    \item
      2a4. Hệ thống truy vấn số liệu của tháng liền trước
      (\(month - 1\)) để tính toán so sánh:

      \begin{itemize}
      \tightlist
      \item
        Chênh lệch tuyệt đối:
        \(\Delta Amount = Amount_{\text{tháng này}} - Amount_{\text{tháng trước}}\).
      \item
        Tỷ lệ phần trăm biến động:
        \[
        \% Change = \left(\frac{\Delta Amount}{Amount_{\text{tháng trước}}}\right) \times 100\%
        \]
        (nếu \(Amount_{\text{tháng trước}} > 0\)).
      \end{itemize}
    \item
      2a5. Hệ thống sắp xếp các danh mục theo số tiền chi tiêu giảm dần
      để làm nổi bật nhóm chi tiêu lớn nhất.
    \item
      2a6. Hệ thống hiển thị bảng Pivot phân cấp kèm biểu đồ tròn cơ cấu
      chi tiêu.
    \end{itemize}
  \item
    \emph{2b. Xem báo cáo tổng tài sản và giá trị tài sản ròng (Net
    Worth):}

    \begin{itemize}
    \tightlist
    \item
      2b1. Người dùng chọn tab "Báo cáo tài sản ròng (Net Worth)".
    \item
      2b2. Hệ thống truy vấn đồng thời 5 cấu phần tài sản tại thời điểm
      hiện tại theo quy tắc \texttt{FR-REP-03}:

      \begin{enumerate}
      \def\labelenumi{\arabic{enumi}.}
      \tightlist
      \item
        Tiền mặt thanh khoản (\(Wallets\)): Tổng số dư các ví đang
        \texttt{ACTIVE} theo \texttt{BR-FIN-02}.
      \item
        Tiền gửi tiết kiệm (\(Savings\)): Tổng giá trị hiện tại của các
        sổ tiết kiệm đang \texttt{OPEN} theo \texttt{BR-FIN-05}.
      \item
        Danh mục đầu tư (\(Investments\)): Tổng giá trị thị trường các
        khoản CCQ đang \texttt{HOLDING} theo \texttt{BR-FIN-06}.
      \item
        Khoản phải thu (\(Receivable\)): Tổng dư nợ cho vay chưa thu hồi
        theo \texttt{BR-FIN-07}.
      \item
        Khoản phải trả (\(Payable\)): Tổng dư nợ đi vay chưa thanh toán
        theo \texttt{BR-FIN-07}.
      \end{enumerate}
    \item
      2b3. Hệ thống tính toán:
      \[NetWorth = Wallets + Savings + Investments + Receivable - Payable\]
      \[TotalAssets = Wallets + Savings + Investments + Receivable\]
    \item
      2b4. Hệ thống tính tỷ trọng phần trăm của từng cấu phần trên tổng
      tài sản gộp \(TotalAssets\).
    \item
      2b5. Hệ thống hiển thị bảng cân đối tài sản cá nhân kèm biểu đồ
      thể hiện tỷ trọng phân bổ tài sản.
    \end{itemize}
  \item
    \emph{2c. Tra cứu sao kê công nợ chi tiết theo đối tác:}

    \begin{itemize}
    \tightlist
    \item
      2c1. Người dùng chọn tab "Sao kê công nợ đối tác" và chọn một cá
      nhân cụ thể từ danh sách đối tác có phát sinh giao dịch nợ.
    \item
      2c2. Hệ thống trích xuất toàn bộ các bản ghi trong bảng
      \texttt{debt} liên quan đến cá nhân này, sắp xếp theo thứ tự ngày
      tăng dần theo quy tắc \texttt{FR-REP-04}.
    \item
      2c3. Với từng giao dịch, hệ thống tính số dư nợ lũy kế: tăng phải
      thu nếu \texttt{lend}, giảm phải thu nếu
      \texttt{collect}/\texttt{write-off\ lend}; tăng phải trả nếu
      \texttt{borrow}, giảm phải trả nếu
      \texttt{repay}/\texttt{write-off\ borrow}.
    \item
      2c4. Hệ thống tổng kết: Tổng đã cho vay, Tổng đã thu hồi, Dư nợ
      phải thu còn lại, Tổng đã vay, Tổng đã trả, Dư nợ phải trả còn
      lại.
    \item
      2c5. Hệ thống hiển thị bảng sao kê chi tiết phục vụ đối soát minh
      bạch.
    \end{itemize}
  \end{itemize}
\item
  \textbf{Các luồng ngoại lệ (Exception Flows):}

  \begin{itemize}
  \tightlist
  \item
    \emph{4a. Tháng hoàn toàn không có giao dịch phát sinh:}

    \begin{itemize}
    \tightlist
    \item
      4a1. Tại bước 4, hệ thống ghi nhận trong một tháng \(m\) nào đó
      không có bất kỳ giao dịch thu chi nào phát sinh (\(Income_m = 0\)
      và \(Expense_m = 0\)).
    \item
      4a2. Hệ thống hiển thị giá trị số tiền là "0.00 VNĐ", không để
      trống ô dữ liệu hoặc gán giá trị null.
    \item
      4a3. Hệ thống gán \(CashFlow_m = 0.00\) và
      \(SavingsRate_m = 0.00\%\), bảo đảm bảng dữ liệu 12 tháng luôn
      liên tục và không gây lỗi giao diện.
    \end{itemize}
  \item
    \emph{5a. Tháng có thu nhập bằng 0 nhưng chi tiêu lớn hơn 0
    (\(Income = 0\) và \(Expense > 0\)):}

    \begin{itemize}
    \tightlist
    \item
      5a1. Tại bước 5, hệ thống phát hiện \(Income_m = 0\) nhưng
      \(Expense_m > 0\) (ví dụ người dùng không có thu nhập phát sinh
      trong tháng nhưng vẫn có chi tiêu sinh hoạt).
    \item
      5a2. Nếu thực hiện phép chia cho 0 để tính \(SavingsRate\), hệ
      thống sẽ gặp lỗi \texttt{ArithmeticException}.
    \item
      5a3. Theo quy tắc xử lý biên \texttt{FR-REP-01}, hệ thống gán
      \(SavingsRate_m = 0.00\%\) (hoặc đánh dấu chỉ số âm định danh
      không áp dụng), hiển thị dòng tiền ròng âm
      (\(CashFlow = -Expense\)) và đưa ra chỉ báo thâm hụt ngân sách.
    \end{itemize}
  \item
    \emph{2a4a. Tháng trước không có chi tiêu trong danh mục khi so sánh
    biến động:}

    \begin{itemize}
    \tightlist
    \item
      2a4a1. Tại bước 2a4, hệ thống phát hiện tháng trước không có chi
      tiêu trong danh mục này (\(Amount_{\text{tháng trước}} = 0\)) nhưng
      tháng này có chi tiêu (\(Amount_{\text{tháng này}} > 0\)).
    \item
      2a4a2. Hệ thống hiển thị chênh lệch tuyệt đối
      \(\Delta Amount = Amount_{\text{tháng này}}\) và hiển thị tỷ lệ biến
      động là "Mới" (hoặc N/A), tránh thực hiện phép chia cho 0.
    \end{itemize}
  \end{itemize}
\item
  \textbf{Quy tắc nghiệp vụ đặc thù (Special Requirements /
  Invariants):}

  \begin{enumerate}
  \def\labelenumi{\arabic{enumi}.}
  \tightlist
  \item
    Bất biến chuỗi thời gian trượt 12 tháng \texttt{FR-REP-01}: Bảng
    luôn gồm đúng 12 tháng liên tục lùi từ tháng hiện tại, tự động thích
    ứng với ngày hệ thống mà không cần cập nhật thủ công.
  \item
    Bất biến tính toán tài chính \texttt{BR-FIN-01}: Toàn bộ các chỉ số
    thu nhập, chi tiêu, dòng tiền ròng, tài sản ròng đều tính toán qua
    \texttt{BigDecimal}, không phát sinh sai số làm tròn.
  \item
    Bất biến Giá trị tài sản ròng (Net Worth) \texttt{FR-REP-03}: Luôn
    bảo đảm cân bằng kế toán giữa 5 cấu phần, phản ánh trung thực toàn
    bộ tài sản thanh khoản, tài sản đầu tư, tiền gửi tiết kiệm và nghĩa
    vụ công nợ.
  \item
    Phù hợp các yêu cầu chức năng: \texttt{FR-REP-01},
    \texttt{FR-REP-02}, \texttt{FR-REP-03}, \texttt{FR-REP-04}.
  \end{enumerate}
\end{itemize}

\subsubsection{6.10 UC10: Sao lưu, xuất và nhập dữ liệu tài chính
(Financial Data Backup, Export \&
Import)}\label{610-uc10-sao-lux1b0u-xuux1ea5t-vuxe0-nhux1eadp-dux1eef-liux1ec7u-tuxe0i-chuxednh-financial-data-backup-export--import}

\begin{itemize}
\tightlist
\item
  \textbf{Mã và Tên ca sử dụng:}
  \texttt{UC10:\ Sao\ lưu,\ xuất\ và\ nhập\ dữ\ liệu\ tài\ chính\ (Financial\ Data\ Backup,\ Export\ \&\ Import)}
\item
  \textbf{Tác nhân chính:} Chủ sở hữu hệ thống (Owner)
\item
  \textbf{Mô tả tóm tắt:} Ca sử dụng này cho phép Owner xuất toàn bộ
  hoặc từng phần sổ cái dữ liệu tài chính của hệ thống ra các tệp lưu
  trữ an toàn định dạng JSON hoặc CSV trên máy tính cá nhân, bảo toàn
  nguyên vẹn độ chính xác kiểu số học \texttt{BigDecimal}; thực hiện
  nhập và phục hồi dữ liệu từ tệp sao lưu với quy trình kiểm định 4 bước
  nghiêm ngặt (cú pháp, tính toàn vẹn khóa ngoại danh mục, tính toàn vẹn
  ví, và bất biến số dư không âm theo BR-FIN-02) kèm cơ chế Rollback tự
  động; và thực thi công cụ kiểm toán đối soát toàn vẹn dữ liệu (Data
  Integrity Audit) để phát hiện sớm các sai lệch số dư hoặc bút toán mất
  cân bằng.
\item
  \textbf{Điều kiện kích hoạt:} Owner có nhu cầu sao lưu định kỳ đề
  phòng sự cố máy tính, muốn xuất dữ liệu sang tệp phân tích định dạng
  mở phục vụ xử lý bên ngoài, khôi phục dữ liệu từ bản sao lưu trước đó,
  hoặc chạy kiểm toán tính toàn vẹn hệ thống.
\item
  \textbf{Tiền điều kiện:}

  \begin{enumerate}
  \def\labelenumi{\arabic{enumi}.}
  \tightlist
  \item
    Hệ thống PFMS đang hoạt động bình thường, có quyền đọc/ghi trên thư
    mục lưu trữ cục bộ được chỉ định.
  \item
    Owner đang ở giao diện phân hệ Sao lưu và Quản lý dữ liệu.
  \end{enumerate}
\item
  \textbf{Hậu điều kiện thành công:}

  \begin{enumerate}
  \def\labelenumi{\arabic{enumi}.}
  \tightlist
  \item
    Khi xuất dữ liệu: Tệp sao lưu được tạo thành công với mã hóa UTF-8,
    các trường số học \texttt{BigDecimal} được lưu dưới dạng chuỗi số
    học đầy đủ chữ số (không bị làm tròn, không dùng ký hiệu khoa học),
    tệp đặt tên tự động theo chuẩn
    \texttt{pfms\_export\_\{scope\}\_\{yyyyMMdd\_HHmmss\}.\{ext\}} theo
    quy tắc \texttt{FR-DAT-01}.
  \item
    Khi nhập dữ liệu: Dữ liệu từ tệp sao lưu được phục hồi trọn vẹn vào
    hệ thống theo chế độ đã chọn (\texttt{OVERWRITE} hoặc
    \texttt{MERGE}), toàn bộ các bất biến số dư \texttt{BR-FIN-02}, ghi
    sổ kép \texttt{BR-FIN-03} và danh mục \texttt{BR-FIN-04} được thỏa
    mãn 100\%.
  \item
    Nếu tệp nhập chứa dữ liệu lỗi: Toàn bộ thao tác bị hủy bỏ và
    rollback nguyên trạng, cơ sở dữ liệu hiện tại không bị thay đổi hoặc
    làm sai lệch.
  \item
    Khi kiểm toán: Báo cáo kiểm toán toàn vẹn dữ liệu (Data Integrity
    Audit Report) phản ánh trung thực trạng thái của 4 quy tắc bất biến
    cốt lõi theo \texttt{FR-DAT-03}.
  \end{enumerate}
\item
  \textbf{Luồng sự kiện chính (Main Success Scenario) - Xuất dữ liệu tài
  chính (Export):}

  \begin{enumerate}
  \def\labelenumi{\arabic{enumi}.}
  \tightlist
  \item
    Người dùng chọn chức năng "Xuất dữ liệu tài chính".
  \item
    Hệ thống hiển thị biểu mẫu xuất dữ liệu gồm các tùy chọn:

    \begin{itemize}
    \tightlist
    \item
      Phạm vi dữ liệu (\texttt{scope}): \texttt{FULL} (Toàn bộ hệ
      thống), \texttt{TRANSACTIONS} (Sổ nhật ký thu chi), \texttt{DEBTS}
      (Sổ công nợ), \texttt{SAVINGS} (Sổ tiết kiệm),
      \texttt{INVESTMENTS} (Chứng chỉ quỹ đầu tư).
    \item
      Định dạng tệp (\texttt{format}): \texttt{JSON} (chuẩn cấu trúc dữ
      liệu phân cấp) hoặc \texttt{CSV} (chuẩn dữ liệu dạng bảng phân
      tách bằng dấu phẩy).
    \item
      Thư mục lưu trữ trên máy tính (\texttt{destinationPath}): Người
      dùng chọn qua hộp thoại duyệt thư mục.
    \end{itemize}
  \item
    Người dùng chọn phạm vi là \texttt{FULL}, định dạng \texttt{JSON},
    và chọn thư mục lưu trữ đích.
  \item
    Người dùng nhấn nút "Bắt đầu xuất dữ liệu".
  \item
    Hệ thống trích xuất toàn bộ dữ liệu hiện tại từ cơ sở dữ liệu: ví
    tiền, danh mục phân loại, nhật ký thu chi, chuyển tiền nội bộ, công
    nợ, sổ tiết kiệm, chứng chỉ quỹ, quỹ mục tiêu.
  \item
    Hệ thống chuyển đổi các thực thể miền sang cấu trúc JSON tiêu chuẩn:

    \begin{itemize}
    \tightlist
    \item
      Chuẩn hóa toàn bộ các trường \texttt{BigDecimal} thành chuỗi số
      học chính xác (ví dụ: \texttt{"50000000.00"}), cấm tuyệt đối
      chuyển đổi sang kiểu \texttt{Double} hoặc ký hiệu số mũ khoa học
      (\texttt{5.0E7}) theo quy tắc \texttt{FR-DAT-01}.
    \item
      Định dạng các trường ngày tháng theo chuẩn ISO-8601
      (\texttt{yyyy-MM-dd}).
    \end{itemize}
  \item
    Hệ thống tạo tên tệp theo quy tắc:
    \texttt{pfms\_export\_full\_\{timestamp\}.json}.
  \item
    Hệ thống ghi tệp an toàn xuống đĩa cứng với bảng mã UTF-8.
  \item
    Hệ thống kiểm tra lại kích thước và tính toàn vẹn của tệp vừa tạo.
  \item
    Hệ thống hiển thị thông báo xuất dữ liệu thành công kèm đường dẫn
    tệp, kích thước tệp và tổng số lượng bản ghi đã xuất của từng phân
    hệ.
  \end{enumerate}
\item
  \textbf{Các luồng rẽ nhánh (Alternative Flows):}

  \begin{itemize}
  \tightlist
  \item
    \emph{2a. Nhập và phục hồi dữ liệu từ tệp sao lưu (Import \&
    Restore):}

    \begin{itemize}
    \tightlist
    \item
      2a1. Người dùng chọn chức năng "Nhập dữ liệu từ tệp sao lưu".
    \item
      2a2. Hệ thống hiển thị biểu mẫu nhập dữ liệu: Hộp thoại chọn tệp
      sao lưu từ máy tính (\texttt{.json} hoặc \texttt{.csv}), và Tùy
      chọn chế độ nhập (\texttt{OVERWRITE} - ghi đè toàn bộ, hoặc
      \texttt{MERGE} - hợp nhất bổ sung).
    \item
      2a3. Người dùng chọn tệp sao lưu JSON hợp lệ, chọn chế độ
      \texttt{OVERWRITE}.
    \item
      2a4. Hệ thống hiển thị hộp thoại cảnh báo: "Chế độ ghi đè sẽ thay
      thế toàn bộ dữ liệu tài chính hiện tại bằng dữ liệu từ tệp sao
      lưu. Bạn có chắc chắn muốn tiếp tục?".
    \item
      2a5. Người dùng nhấn xác nhận tiếp tục.
    \item
      2a6. Hệ thống phân tích cú pháp tệp (JSON Parsing) và thực hiện
      quy trình kiểm định toàn vẹn 4 bước theo quy tắc
      \texttt{FR-DAT-02}:

      \begin{enumerate}
      \def\labelenumi{\arabic{enumi}.}
      \tightlist
      \item
        Bước 1: Kiểm tra cấu trúc cú pháp tệp và tính hợp lệ của định
        dạng ngày tháng, số tiền.
      \item
        Bước 2: Kiểm tra tính toàn vẹn danh mục theo \texttt{BR-FIN-04}:
        Mọi \texttt{categoryId} tham chiếu phải thuộc về nhóm hợp lệ.
      \item
        Bước 3: Kiểm tra tính toàn vẹn ví: Mọi giao dịch phải tham chiếu
        đến một ví tồn tại trong danh sách ví.
      \item
        Bước 4: Kiểm tra bất biến số dư ví theo \texttt{BR-FIN-02}: Tính
        toán lại số dư của từng ví từ lịch sử giao dịch trong tệp, bảo
        đảm không có ví thông thường nào bị rơi vào số dư âm.
      \end{enumerate}
    \item
      2a7. Nếu toàn bộ 4 bước kiểm định đều ĐẠT:

      \begin{itemize}
      \tightlist
      \item
        Hệ thống bắt đầu một giao dịch cơ sở dữ liệu nguyên tử (Database
        Transaction).
      \item
        Xóa sạch dữ liệu cũ (nếu là \texttt{OVERWRITE}) và nạp toàn bộ
        các bản ghi từ tệp vào cơ sở dữ liệu.
      \item
        Hệ thống cam kết giao dịch (Commit).
      \item
        Kích hoạt tính toán lại toàn bộ số dư động, quỹ mục tiêu và các
        bảng thống kê trượt 12 tháng.
      \end{itemize}
    \item
      2a8. Hệ thống hiển thị báo cáo tổng kết nhập dữ liệu thành công:
      số lượng ví, giao dịch, sổ tiết kiệm, CCQ đã được phục hồi nguyên
      vẹn.
    \end{itemize}
  \item
    \emph{2b. Kiểm toán đối soát toàn vẹn dữ liệu (Data Integrity
    Audit):}

    \begin{itemize}
    \tightlist
    \item
      2b1. Người dùng chọn chức năng "Kiểm toán toàn vẹn dữ liệu
      (Audit)".
    \item
      2b2. Hệ thống tự động kích hoạt tiến trình quét toàn bộ cơ sở dữ
      liệu, thực hiện kiểm tra tuần tự 4 quy tắc bất biến cốt lõi theo
      quy tắc \texttt{FR-DAT-03}:

      \begin{enumerate}
      \def\labelenumi{\arabic{enumi}.}
      \tightlist
      \item
        Kiểm toán số dư ví: Quét toàn bộ lịch sử giao dịch của từng ví,
        tính số dư độc lập theo công thức \texttt{BR-FIN-02} và so sánh
        với số dư ghi nhận của từng ví. Kiểm tra sai số \(\Delta = 0\).
      \item
        Kiểm toán ghi sổ kép chuyển tiền: Quét toàn bộ bảng
        \texttt{transfers}, kiểm tra số tiền trừ tại ví nguồn bằng chính
        xác số tiền cộng tại ví đích theo \texttt{BR-FIN-03}.
      \item
        Kiểm toán liên kết danh mục: Quét bảng \texttt{transactions},
        kiểm tra không có giao dịch nào tham chiếu danh mục mồ côi hoặc
        không khớp \texttt{Type} theo \texttt{BR-FIN-04}.
      \item
        Kiểm toán bất biến công nợ: Kiểm tra toàn bộ các đối tác có số
        dư nợ phải thu và nợ phải trả đều \(\ge 0\) theo
        \texttt{BR-FIN-07}.
      \end{enumerate}
    \item
      2b3. Hệ thống tổng hợp kết quả thành Báo cáo Kiểm toán Toàn vẹn Dữ
      liệu (Audit Report).
    \item
      2b4. Nếu tất cả đều ĐẠT, hệ thống hiển thị dấu tích xanh và thông
      báo: "Cơ sở dữ liệu hoàn toàn nhất quán và toàn vẹn 100\%".
    \item
      2b5. Nếu phát hiện sai lệch, hệ thống đánh dấu đỏ, liệt kê chi
      tiết mã giao dịch, số tiền chênh lệch và hướng dẫn người dùng biện
      pháp xử lý.
    \end{itemize}
  \end{itemize}
\item
  \textbf{Các luồng ngoại lệ (Exception Flows):}

  \begin{itemize}
  \tightlist
  \item
    \emph{6a. Tệp sao lưu không đúng định dạng cú pháp hoặc bị hỏng
    (Parse Error):}

    \begin{itemize}
    \tightlist
    \item
      6a1. Tại bước 2a6, hệ thống đọc tệp và phát hiện tệp tin bị lỗi
      định dạng JSON (thiếu dấu ngoặc, tệp rỗng, hoặc mã hóa sai ký tự).
    \item
      6a2. Hệ thống hủy bỏ tiến trình nạp, thông báo lỗi: "Tệp sao lưu
      bị lỗi cấu trúc cú pháp hoặc bị hỏng dữ liệu. Vui lòng kiểm tra
      lại tệp nguồn".
    \item
      6a3. Dữ liệu hiện tại trong hệ thống được bảo toàn 100\%.
    \end{itemize}
  \item
    \emph{2a6a. Dữ liệu trong tệp vi phạm bất biến kế toán (Validation
    Failure):}

    \begin{itemize}
    \tightlist
    \item
      2a6a1. Tại bước 2a6, hệ thống phát hiện ít nhất một lỗi vi phạm:
      ví dụ giao dịch chi tiêu trỏ tới một ví không tồn tại, hoặc lịch
      sử giao dịch trong tệp làm cho số dư ví tiền mặt bị âm
      (\(Balance_{wallet} < 0\)).
    \item
      2a6a2. Hệ thống lập tức hủy bỏ toàn bộ quá trình nhập, không ghi
      bất kỳ bản ghi nào vào cơ sở dữ liệu (Rollback).
    \item
      2a6a3. Hệ thống hiển thị bảng thông báo lỗi chi tiết: liệt kê số
      dòng bị lỗi, loại vi phạm (vi phạm khóa ngoại hoặc vi phạm bất
      biến số dư \texttt{BR-FIN-02}) để người dùng sửa chữa tệp trước
      khi thử lại.
    \end{itemize}
  \item
    \emph{2a5a. Người dùng hủy bỏ thao tác khi được cảnh báo ghi đè:}

    \begin{itemize}
    \tightlist
    \item
      2a5a1. Tại bước 2a4, trước cảnh báo ghi đè toàn bộ dữ liệu, Người
      dùng chọn "Hủy bỏ".
    \item
      2a5a2. Hệ thống dừng thao tác nhập, đóng hộp thoại và giữ nguyên
      hiện trạng dữ liệu.
    \end{itemize}
  \item
    \emph{8a. Lỗi quyền ghi tệp hoặc ổ đĩa đầy khi xuất dữ liệu:}

    \begin{itemize}
    \tightlist
    \item
      8a1. Tại bước 8, hệ thống không thể ghi tệp xuống thư mục chỉ định
      do không đủ quyền truy cập hệ thống tệp hoặc ổ đĩa đã hết dung
      lượng lưu trữ.
    \item
      8a2. Hệ thống bắt ngoại lệ \texttt{IOException}, hiển thị thông
      báo: "Không thể lưu tệp dữ liệu vào thư mục đã chọn do lỗi quyền
      truy cập hoặc ổ đĩa đầy. Vui lòng chọn thư mục khác".
    \end{itemize}
  \end{itemize}
\item
  \textbf{Quy tắc nghiệp vụ đặc thù (Special Requirements /
  Invariants):}

  \begin{enumerate}
  \def\labelenumi{\arabic{enumi}.}
  \tightlist
  \item
    Bất biến bảo toàn độ chính xác số học \texttt{FR-DAT-01}: Tuyệt đối
    không làm tròn hoặc sử dụng ký hiệu khoa học đối với trường dữ liệu
    \texttt{BigDecimal} khi xuất và nhập tệp sao lưu.
  \item
    Bất biến giao dịch nguyên tử (Atomicity): Quá trình nhập dữ liệu
    phải tuân thủ nguyên tắc All-or-Nothing; nếu xuất hiện bất kỳ lỗi
    nào, toàn bộ quá trình phải được rollback hoàn toàn về trạng thái
    trước đó.
  \item
    Bất biến kiểm định 4 bước \texttt{FR-DAT-02}: Nghiêm cấm nạp dữ liệu
    vào cơ sở dữ liệu nếu chưa vượt qua đầy đủ 4 bước kiểm tra cấu trúc,
    danh mục, ví và số dư không âm.
  \item
    Phù hợp các yêu cầu chức năng: \texttt{FR-DAT-01},
    \texttt{FR-DAT-02}, \texttt{FR-DAT-03}.
  \end{enumerate}
\end{itemize}

\clearpage
\subsection{7. SƠ ĐỒ CA SỬ DỤNG HỆ THỐNG (USE CASE
DIAGRAMS)}\label{7-sux1a1-ux111ux1ed3-ca-sux1eed-dux1ee5ng-hux1ec7-thux1ed1ng-use-case-diagrams}

Phần này xây dựng hệ thống sơ đồ ca sử dụng (Use Case Diagrams) cho Hệ
thống Quản lý Tài chính Cá nhân (PFMS) theo chuẩn mực mô hình hóa của
Hội đồng Kỹ thuật Yêu cầu Quốc tế (International Requirements
Engineering Board - IREB CPRE Modeling Handbook) và đặc tả UML 2.5
(ISO/IEC 19505:2012).

Mô hình ca sử dụng xác lập ranh giới kỹ thuật kết nối giữa các yêu cầu
chức năng (FR) với các thực thể phân tích hướng đối tượng (OOA) trong
giai đoạn phân tích OOA. Các sơ đồ tuân thủ nghiêm ngặt các nguyên tắc
chuẩn mực sau:

\begin{enumerate}
\def\labelenumi{\arabic{enumi}.}
\tightlist
\item
  \textbf{Ranh giới hệ thống phân định rõ ràng (System Boundary):} Khối
  hộp chữ nhật biểu diễn ranh giới phần mềm PFMS phân biệt tường minh
  giữa môi trường bên ngoài và các xử lý nghiệp vụ tự động bên trong hệ
  thống.
\item
  \textbf{Mô hình tác nhân đơn nhất (Single-User Primary Actor):} Xác
  định tác nhân duy nhất là \texttt{Chủ\ sở\ hữu\ hệ\ thống\ (Owner)},
  phản ánh đúng ngữ cảnh ứng dụng cá nhân cục bộ trên máy tính để bàn
  (Desktop Application). Mọi quyền hạn tạo lập dữ liệu, điều chuyển dòng
  tiền, thiết lập danh mục và sao lưu hệ thống đều quy tụ về tác nhân
  này.
\item
  \textbf{Phân định ngữ nghĩa quan hệ \textless{}\textgreater{} và
  \textless{}\textgreater{} chuẩn UML:}

  \begin{itemize}
  \tightlist
  \item
    \textbf{Quan hệ
    \texttt{\textless{}\textless{}include\textgreater{}\textgreater{}}
    (Bao hàm):} Hướng mũi tên đứt nét từ ca sử dụng cơ sở đến ca sử dụng
    được bao hàm
    (\texttt{Base\ ..\textgreater{}\ Included\ :\ \textless{}\textless{}include\textgreater{}\textgreater{}}).
    Thể hiện hành vi nghiệp vụ bắt buộc phải thực thi tại một vị trí xác
    định trong tiến trình của ca sử dụng cơ sở, nhằm tái sử dụng các quy
    tắc kiểm tra bất biến cốt lõi (như kiểm tra số dư ví không âm theo
    \texttt{BR-FIN-02}, ghi sổ kép đối ứng theo \texttt{BR-FIN-03}, hoặc
    kiểm soát tỷ lệ phân bổ quỹ mục tiêu theo \texttt{BR-FIN-08}).
  \item
    \textbf{Quan hệ
    \texttt{\textless{}\textless{}extend\textgreater{}\textgreater{}}
    (Mở rộng):} Hướng mũi tên đứt nét từ ca sử dụng mở rộng trỏ ngược về
    ca sử dụng cơ sở
    (\texttt{Extension\ ..\textgreater{}\ Base\ :\ \textless{}\textless{}extend\textgreater{}\textgreater{}}).
    Thể hiện hành vi chỉ kích hoạt có điều kiện khi thỏa mãn điểm mở
    rộng (Extension Point) cụ thể, giúp tách biệt các luồng xử lý tùy
    chọn hoặc ngoại lệ mà không làm phức tạp hóa luồng sự kiện chính
    (như ghi nhận phí chuyển khoản khi có phí phát sinh, tất toán trước
    hạn khi người dùng có nhu cầu rút vốn trước ngày đáo hạn, hoặc xóa
    nợ khó đòi khi đối tác mất khả năng hoàn trả).
  \end{itemize}
\item
  \textbf{Kiến trúc mô hình hóa phân tầng (Two-Tier Modeling):}

  \begin{itemize}
  \tightlist
  \item
    Tầng 1: Sơ đồ Ca sử dụng tổng thể toàn hệ thống (Global System
    Diagram - Mục 7.1) cung cấp cái nhìn toàn diện về toàn bộ 10 ca sử
    dụng và các liên kết phụ thuộc liên phân hệ.
  \item
    Tầng 2: 3 Sơ đồ Ca sử dụng phân rã theo phân vùng nghiệp vụ
    (Subsystem Diagrams - Mục 7.2) làm rõ các ca sử dụng con, điểm mở
    rộng và ràng buộc kiểm định dữ liệu nội bộ từng phân vùng.
  \end{itemize}
\end{enumerate}

\subsubsection{7.1 Sơ đồ Ca sử dụng tổng thể toàn hệ thống (Global
System Use Case
Diagram)}\label{71-sux1a1-ux111ux1ed3-ca-sux1eed-dux1ee5ng-tux1ed5ng-thux1ec3-touxe0n-hux1ec7-thux1ed1ng-global-system-use-case-diagram}

Sơ đồ tổng thể bao quát toàn bộ 10 ca sử dụng cốt lõi của hệ thống PFMS
(\texttt{UC01} đến \texttt{UC10}), thể hiện mối tương tác giữa tác nhân
Owner với ranh giới hệ thống, cũng như các quan hệ phụ thuộc cấu trúc
giữa các ca sử dụng.

\begin{Shaded}
\begin{Highlighting}[]
\NormalTok{@startuml}
\NormalTok{left to right direction}
\NormalTok{skinparam handwritten false}
\NormalTok{skinparam monochrome false}
\NormalTok{skinparam shadowing false}
\NormalTok{skinparam defaultFontName "Segoe UI", Arial, sans{-}serif}
\NormalTok{skinparam defaultFontSize 12}
\NormalTok{skinparam roundCorner 8}

\NormalTok{skinparam actor \{}
\NormalTok{    BackgroundColor \#E8F4FD}
\NormalTok{    BorderColor \#2B6CB0}
\NormalTok{\}}

\NormalTok{skinparam usecase \{}
\NormalTok{    BackgroundColor \#FFFFFF}
\NormalTok{    BorderColor \#2B6CB0}
\NormalTok{    ArrowColor \#2D3748}
\NormalTok{\}}

\NormalTok{skinparam package \{}
\NormalTok{    BackgroundColor \#F7FAFC}
\NormalTok{    BorderColor \#CBD5E0}
\NormalTok{\}}

\NormalTok{skinparam rectangle \{}
\NormalTok{    BackgroundColor \#FFFFFF}
\NormalTok{    BorderColor \#A0AEC0}
\NormalTok{\}}

\NormalTok{actor :Chủ sở hữu hệ thống\textbackslash{}n(Owner): as Owner}

\NormalTok{rectangle "Personal Finance Management System (PFMS)" as PFMS \{}
\NormalTok{    package "Phân vùng Dòng tiền và Thanh khoản\textbackslash{}n(Cashflow \& Liquidity Subsystem)" as P1 \{}
\NormalTok{        usecase "UC01: Quản lý ví tiền\textbackslash{}n(Wallet Management)" as UC01}
\NormalTok{        usecase "UC02: Ghi nhận thu chi\textbackslash{}n(Income \& Expense Management)" as UC02}
\NormalTok{        usecase "UC03: Chuyển tiền nội bộ\textbackslash{}n(Internal Fund Transfer)" as UC03}
\NormalTok{        usecase "UC04: Quản lý danh mục 3 cấp\textbackslash{}n(Category Hierarchy Management)" as UC04}
        
\NormalTok{        usecase "Kiểm tra số dư ví nguồn\textbackslash{}n(Validate Source Balance)" as UC\_CheckBal}
\NormalTok{        usecase "Ghi nhận phí chuyển tiền\textbackslash{}n(Record Transfer Fee)" as UC\_Fee}
\NormalTok{    \}}

\NormalTok{    package "Phân vùng Tích lũy, Đầu tư và Nghĩa vụ\textbackslash{}n(Wealth \& Liabilities Subsystem)" as P2 \{}
\NormalTok{        usecase "UC05: Quản lý công nợ\textbackslash{}n(Debt Management)" as UC05}
\NormalTok{        usecase "UC06: Quản lý sổ tiết kiệm\textbackslash{}n(Term Savings Management)" as UC06}
\NormalTok{        usecase "UC07: Quản lý chứng chỉ quỹ\textbackslash{}n(Fund Certificate Management)" as UC07}
\NormalTok{        usecase "UC08: Quản lý quỹ mục tiêu\textbackslash{}n(Financial Goal Management)" as UC08}
        
\NormalTok{        usecase "Xóa nợ khó đòi phi tiền mặt\textbackslash{}n(Write{-}off Bad Debt)" as UC\_WriteOff}
\NormalTok{        usecase "Tất toán sổ trước hạn\textbackslash{}n(Early Withdrawal Settlement)" as UC\_EarlyWithdraw}
\NormalTok{        usecase "Tổng hợp tài sản 3 cấu phần\textbackslash{}n(Aggregate Multi{-}Asset Balance)" as UC\_AggAssets}
\NormalTok{    \}}

\NormalTok{    package "Phân vùng Báo cáo, Phân tích \& Quản trị\textbackslash{}n(Analytics \& Administration Subsystem)" as P3 \{}
\NormalTok{        usecase "UC09: Thống kê \& Phân tích tài chính\textbackslash{}n(Monthly Statistics \& Reporting)" as UC09}
\NormalTok{        usecase "UC10: Sao lưu \& Quản trị dữ liệu\textbackslash{}n(Financial Data Backup \& Admin)" as UC10}
        
\NormalTok{        usecase "Kiểm toán toàn vẹn dữ liệu\textbackslash{}n(Data Integrity Audit)" as UC\_Audit}
\NormalTok{    \}}
\NormalTok{\}}

\NormalTok{\textquotesingle{} Tác nhân kích hoạt 10 ca sử dụng cốt lõi}
\NormalTok{Owner {-}{-}\textgreater{} UC01}
\NormalTok{Owner {-}{-}\textgreater{} UC02}
\NormalTok{Owner {-}{-}\textgreater{} UC03}
\NormalTok{Owner {-}{-}\textgreater{} UC04}
\NormalTok{Owner {-}{-}\textgreater{} UC05}
\NormalTok{Owner {-}{-}\textgreater{} UC06}
\NormalTok{Owner {-}{-}\textgreater{} UC07}
\NormalTok{Owner {-}{-}\textgreater{} UC08}
\NormalTok{Owner {-}{-}\textgreater{} UC09}
\NormalTok{Owner {-}{-}\textgreater{} UC10}

\NormalTok{\textquotesingle{} Mối quan hệ include và extend nội bộ}
\NormalTok{UC03 ..\textgreater{} UC\_CheckBal : \textless{}\textless{}include\textgreater{}\textgreater{}}
\NormalTok{UC\_Fee ..\textgreater{} UC03 : \textless{}\textless{}extend\textgreater{}\textgreater{}\textbackslash{}n(Điểm mở rộng: Có phí phát sinh)}

\NormalTok{UC\_WriteOff ..\textgreater{} UC05 : \textless{}\textless{}extend\textgreater{}\textgreater{}\textbackslash{}n(Điểm mở rộng: Mất khả năng thanh toán)}
\NormalTok{UC\_EarlyWithdraw ..\textgreater{} UC06 : \textless{}\textless{}extend\textgreater{}\textgreater{}\textbackslash{}n(Điểm mở rộng: Rút trước đáo hạn)}
\NormalTok{UC08 ..\textgreater{} UC\_AggAssets : \textless{}\textless{}include\textgreater{}\textgreater{}}

\NormalTok{UC10 ..\textgreater{} UC\_Audit : \textless{}\textless{}include\textgreater{}\textgreater{}\textbackslash{}n(Quy trình kiểm định 4 bước)}

\NormalTok{\textquotesingle{} Mối quan hệ tương tác dữ liệu liên phân vùng}
\NormalTok{UC02 ..\textgreater{} UC01 : \textless{}\textless{}include\textgreater{}\textgreater{}\textbackslash{}n(Cập nhật số dư ví)}
\NormalTok{UC02 ..\textgreater{} UC04 : \textless{}\textless{}include\textgreater{}\textgreater{}\textbackslash{}n(Phân loại giao dịch)}
\NormalTok{UC03 ..\textgreater{} UC01 : \textless{}\textless{}include\textgreater{}\textgreater{}\textbackslash{}n(Ghi sổ kép ví nguồn \& ví đích)}
\NormalTok{UC05 ..\textgreater{} UC01 : \textless{}\textless{}include\textgreater{}\textgreater{}\textbackslash{}n(Biến động ví giải ngân/thu nợ)}
\NormalTok{UC06 ..\textgreater{} UC01 : \textless{}\textless{}include\textgreater{}\textgreater{}\textbackslash{}n(Trích tiền gửi/nhận gốc và lãi)}
\NormalTok{UC07 ..\textgreater{} UC01 : \textless{}\textless{}include\textgreater{}\textgreater{}\textbackslash{}n(Thanh toán mua/bán chứng chỉ quỹ)}
\NormalTok{UC08 ..\textgreater{} UC01 : \textless{}\textless{}include\textgreater{}\textgreater{}\textbackslash{}n(Phân bổ tiền mặt ví)}
\NormalTok{UC08 ..\textgreater{} UC06 : \textless{}\textless{}include\textgreater{}\textgreater{}\textbackslash{}n(Phân bổ sổ tiết kiệm)}
\NormalTok{UC08 ..\textgreater{} UC07 : \textless{}\textless{}include\textgreater{}\textgreater{}\textbackslash{}n(Phân bổ chứng chỉ quỹ)}
\NormalTok{UC09 ..\textgreater{} UC02 : \textless{}\textless{}include\textgreater{}\textgreater{}\textbackslash{}n(Truy vấn dòng tiền 12 tháng)}
\NormalTok{@enduml}
\end{Highlighting}
\end{Shaded}

\paragraph{Phân tích kiến trúc và quan hệ ca sử dụng tổng
thể}\label{phuxe2n-tuxedch-kiux1ebfn-truxfac-vuxe0-quan-hux1ec7-ca-sux1eed-dux1ee5ng-tux1ed5ng-thux1ec3}

\begin{enumerate}
\def\labelenumi{\arabic{enumi}.}
\tightlist
\item
  \textbf{Bảng tổng hợp 10 Ca sử dụng hệ thống:}
\end{enumerate}

{\def\LTcaptype{none}
\begin{longtable}[]{|>{\raggedright\arraybackslash}p{0.08\linewidth}|>{\raggedright\arraybackslash}p{0.20\linewidth}|>{\raggedright\arraybackslash}p{0.18\linewidth}|>{\raggedright\arraybackslash}p{0.14\linewidth}|>{\raggedright\arraybackslash}p{0.26\linewidth}|}
\hline
\everycr{\noalign{\ifnum0=`}\fi\hrule height 0.4pt\futurelet\reserved@a\LT@@tabularcr}
\textbf{Mã UC} & \textbf{Tên Ca sử dụng} & \textbf{Phân vùng nghiệp vụ} & \textbf{Mục tiêu phục vụ} & \textbf{Bất biến nghiệp vụ gắn liền} \\
\hline
\endfirsthead
\hline
\textbf{Mã UC} & \textbf{Tên Ca sử dụng} & \textbf{Phân vùng nghiệp vụ} & \textbf{Mục tiêu phục vụ} & \textbf{Bất biến nghiệp vụ gắn liền} \\
\hline
\endhead
\hline
\endfoot
\hline
\endlastfoot
UC01 & Quản lý ví tiền & Dòng tiền và Thanh khoản & G1, G2 & BR-FIN-01,
BR-FIN-02 \\
UC02 & Ghi nhận thu chi & Dòng tiền và Thanh khoản & G1, G5 & BR-FIN-01,
BR-FIN-02, BR-FIN-04 \\
UC03 & Chuyển tiền nội bộ & Dòng tiền và Thanh khoản & G1, G5 &
BR-FIN-01, BR-FIN-02, BR-FIN-03 \\
UC04 & Quản lý danh mục 3 cấp & Dòng tiền và Thanh khoản & G1 &
BR-FIN-04 \\
UC05 & Quản lý công nợ & Tích lũy, Đầu tư \& Nghĩa vụ & G2, G3 &
BR-FIN-01, BR-FIN-02, BR-FIN-07 \\
UC06 & Quản lý sổ tiết kiệm & Tích lũy, Đầu tư \& Nghĩa vụ & G2, G3, G5
& BR-FIN-01, BR-FIN-02, BR-FIN-05 \\
UC07 & Quản lý chứng chỉ quỹ & Tích lũy, Đầu tư \& Nghĩa vụ & G2, G3, G5
& BR-FIN-01, BR-FIN-02, BR-FIN-06 \\
UC08 & Quản lý quỹ mục tiêu & Tích lũy, Đầu tư \& Nghĩa vụ & G2, G4 &
BR-FIN-01, BR-FIN-08 \\
UC09 & Thống kê \& Phân tích tài chính & Báo cáo, Phân tích \& Quản trị
& G1, G2, G3, G5 & BR-FIN-01, FR-REP-01 đến FR-REP-04 \\
UC10 & Sao lưu \& Quản trị dữ liệu & Báo cáo, Phân tích \& Quản trị &
G2, G5 & BR-FIN-01, BR-FIN-02, FR-DAT-01 đến FR-DAT-03 \\
\end{longtable}
}

\begin{enumerate}
\def\labelenumi{\arabic{enumi}.}
\setcounter{enumi}{1}
\item
  \textbf{Cơ chế phụ thuộc
  \texttt{\textless{}\textless{}include\textgreater{}\textgreater{}}
  (Bao hàm):}

  \begin{itemize}
  \tightlist
  \item
    \texttt{UC03\ ..\textgreater{}\ UC\_CheckBal}: Mỗi thao tác chuyển
    khoản bắt buộc phải bao hàm bước kiểm tra số dư khả dụng của ví
    nguồn nhằm thực thi bất biến \texttt{BR-FIN-02}
    (\(Balance_{source} - Amount - Fee \ge 0\)). Nếu không thỏa mãn,
    tiến trình chuyển khoản lập tức dừng lại.
  \item
    \texttt{UC08\ ..\textgreater{}\ UC\_AggAssets}: Ca sử dụng quản lý
    quỹ mục tiêu bắt buộc phải bao hàm tiến trình tổng hợp giá trị tài
    sản từ cả 3 cấu phần (Ví tiền mặt UC01, Sổ tiết kiệm UC06, Chứng chỉ
    quỹ UC07) để xác định tỷ lệ \% hoàn thành tiến độ theo thời gian
    thực (\texttt{BR-FIN-08}).
  \item
    \texttt{UC10\ ..\textgreater{}\ UC\_Audit}: Tiến trình nạp dữ liệu
    từ tệp sao lưu bắt buộc phải kích hoạt quy trình kiểm định 4 bước để
    ngăn chặn nạp dữ liệu hỏng hoặc dữ liệu vi phạm bất biến vào cơ sở
    dữ liệu (\texttt{FR-DAT-02}, \texttt{FR-DAT-03}).
  \item
    Các phụ thuộc liên phân hệ: Toàn bộ các ca sử dụng có phát sinh dòng
    tiền (\texttt{UC02}, \texttt{UC03}, \texttt{UC05}, \texttt{UC06},
    \texttt{UC07}, \texttt{UC08}) đều bao hàm mối liên hệ cập nhật hoặc
    trích xuất số dư từ \texttt{UC01:\ Quản\ lý\ ví\ tiền}. Điều này
    phản ánh đúng kiến trúc kế toán: ví tiền mặt là trung tâm thanh
    khoản của toàn bộ hệ thống tài chính cá nhân.
  \end{itemize}
\item
  \textbf{Cơ chế phụ thuộc
  \texttt{\textless{}\textless{}extend\textgreater{}\textgreater{}} (Mở
  rộng):}

  \begin{itemize}
  \tightlist
  \item
    \texttt{UC\_Fee\ ..\textgreater{}\ UC03}: Điểm mở rộng "Có phí phát
    sinh" kích hoạt khi người dùng nhập khoản phí ngân hàng \(Fee > 0\).
    Khi đó hệ thống mở rộng luồng xử lý để tạo thêm một giao dịch chi
    phí phát sinh, trừ bổ sung vào ví nguồn mà không làm ảnh hưởng tới
    luồng chuyển tiền cơ sở.
  \item
    \texttt{UC\_WriteOff\ ..\textgreater{}\ UC05}: Điểm mở rộng "Mất khả
    năng thanh toán" kích hoạt khi một khoản nợ cho vay hoặc đi vay được
    xác định không thể thu hồi hoặc được miễn trừ. Hành vi xóa nợ phi
    tiền mặt này mở rộng ca sử dụng quản lý nợ mà không làm biến động số
    dư các ví tiền (\texttt{FR-DEBT-05}).
  \item
    \texttt{UC\_EarlyWithdraw\ ..\textgreater{}\ UC06}: Điểm mở rộng
    "Rút trước đáo hạn" kích hoạt khi ngày tất toán nhỏ hơn ngày đáo hạn
    thỏa thuận trên hợp đồng. Luồng mở rộng thay thế công thức lãi suất
    có kỳ hạn bằng lãi suất tiền gửi không kỳ hạn thực tế
    (\texttt{FR-SAV-04}).
  \end{itemize}
\end{enumerate}

\subsubsection{7.2 Sơ đồ Ca sử dụng phân rã theo 3 phân vùng nghiệp vụ
(Subsystem Use Case
Diagrams)}\label{72-sux1a1-ux111ux1ed3-ca-sux1eed-dux1ee5ng-phuxe2n-ruxe3-theo-3-phuxe2n-vuxf9ng-nghiux1ec7p-vux1ee5-subsystem-use-case-diagrams}

Nhằm làm rõ các luồng chức năng chi tiết, các điểm kiểm soát điều kiện
và các ca sử dụng con, hệ thống được phân rã thành 3 phân vùng nghiệp vụ
độc lập:

\begin{enumerate}
\def\labelenumi{\arabic{enumi}.}
\tightlist
\item
  \textbf{Phân vùng Quản lý Dòng tiền và Ví tiền (Cashflow \& Liquidity
  Subsystem):} Phụ trách quản trị toàn bộ các tài sản thanh khoản tức
  thời, các luồng tiền thu chi hàng ngày và hệ thống phân loại danh mục
  thu chi.
\item
  \textbf{Phân vùng Tích lũy, Đầu tư và Nghĩa vụ tài chính (Wealth \&
  Liabilities Subsystem):} Phụ trách các tài sản tài chính trung và dài
  hạn, các cam kết công nợ cá nhân và các mục tiêu tài chính tương lai.
\item
  \textbf{Phân vùng Báo cáo, Phân tích và Quản trị dữ liệu (Analytics \&
  Administration Subsystem):} Phụ trách tổng hợp thông tin quản trị,
  phân tích xu hướng chi tiêu, đo lường giá trị tài sản ròng và bảo vệ
  an toàn dữ liệu sổ cái.
\end{enumerate}

\paragraph{7.2.1 Phân vùng Quản lý Dòng tiền và Ví tiền (Cashflow \&
Liquidity
Subsystem)}\label{721-phuxe2n-vuxf9ng-quux1ea3n-luxfd-duxf2ng-tiux1ec1n-vuxe0-vuxed-tiux1ec1n-cashflow--liquidity-subsystem}

Phân vùng này bao gồm 4 ca sử dụng cốt lõi: \texttt{UC01},
\texttt{UC02}, \texttt{UC03} và \texttt{UC04}. Đây là phân vùng có tần
suất tương tác cao nhất trong hệ thống, bảo đảm quản trị chặt chẽ từng
đơn vị tiền tệ thanh khoản và phân loại chi tiêu chính xác.

\begin{Shaded}
\begin{Highlighting}[]
\NormalTok{@startuml}
\NormalTok{left to right direction}
\NormalTok{skinparam handwritten false}
\NormalTok{skinparam monochrome false}
\NormalTok{skinparam shadowing false}
\NormalTok{skinparam defaultFontName "Segoe UI", Arial, sans{-}serif}
\NormalTok{skinparam defaultFontSize 12}
\NormalTok{skinparam roundCorner 8}

\NormalTok{skinparam actor \{}
\NormalTok{    BackgroundColor \#E8F4FD}
\NormalTok{    BorderColor \#2B6CB0}
\NormalTok{\}}

\NormalTok{skinparam usecase \{}
\NormalTok{    BackgroundColor \#FFFFFF}
\NormalTok{    BorderColor \#2B6CB0}
\NormalTok{    ArrowColor \#2D3748}
\NormalTok{\}}

\NormalTok{skinparam package \{}
\NormalTok{    BackgroundColor \#F7FAFC}
\NormalTok{    BorderColor \#CBD5E0}
\NormalTok{\}}

\NormalTok{skinparam rectangle \{}
\NormalTok{    BackgroundColor \#FFFFFF}
\NormalTok{    BorderColor \#A0AEC0}
\NormalTok{\}}

\NormalTok{actor :Chủ sở hữu hệ thống\textbackslash{}n(Owner): as Owner}

\NormalTok{rectangle "Phân vùng Quản lý Dòng tiền và Ví tiền (Cashflow \& Liquidity Subsystem)" \{}
\NormalTok{    package "UC01: Quản lý ví tiền (Wallet Management)" \{}
\NormalTok{        usecase "UC01.1: Khởi tạo ví tiền mới\textbackslash{}n(Create Wallet)" as UC01\_1}
\NormalTok{        usecase "UC01.2: Tra cứu danh sách \& số dư động\textbackslash{}n(View Wallets \& Balances)" as UC01\_2}
\NormalTok{        usecase "UC01.3: Cập nhật thông tin ví tiền\textbackslash{}n(Update Wallet Info)" as UC01\_3}
\NormalTok{        usecase "UC01.4: Thống kê cơ cấu tài sản ví\textbackslash{}n(Analyze Wallet Allocation)" as UC01\_4}
\NormalTok{    \}}

\NormalTok{    package "UC02: Ghi nhận thu chi (Income \& Expense)" \{}
\NormalTok{        usecase "UC02.1: Ghi nhận giao dịch thu nhập\textbackslash{}n(Record Income)" as UC02\_1}
\NormalTok{        usecase "UC02.2: Ghi nhận giao dịch chi tiêu\textbackslash{}n(Record Expense)" as UC02\_2}
\NormalTok{        usecase "UC02.3: Tra cứu \& lọc lịch sử giao dịch\textbackslash{}n(Search \& Filter Transactions)" as UC02\_3}
\NormalTok{        usecase "UC02.4: Chỉnh sửa giao dịch thu chi\textbackslash{}n(Update Transaction)" as UC02\_4}
\NormalTok{        usecase "UC02.5: Xóa giao dịch thu chi\textbackslash{}n(Delete Transaction)" as UC02\_5}
        
\NormalTok{        usecase "Kiểm tra số dư ví khả dụng\textbackslash{}n(Validate Sufficient Balance)" as UC\_ValBal}
\NormalTok{    \}}

\NormalTok{    package "UC03: Chuyển tiền nội bộ (Internal Fund Transfer)" \{}
\NormalTok{        usecase "UC03.1: Chuyển tiền giữa hai ví\textbackslash{}n(Execute Fund Transfer)" as UC03\_1}
\NormalTok{        usecase "UC03.2: Ghi nhận phí chuyển tiền\textbackslash{}n(Record Transfer Fee)" as UC03\_2}
\NormalTok{        usecase "UC03.3: Tra cứu \& hủy lệnh chuyển tiền\textbackslash{}n(Rollback Transfer)" as UC03\_3}
        
\NormalTok{        usecase "Ghi sổ kép đối ứng đồng thời\textbackslash{}n(Atomic Double{-}Entry Posting)" as UC\_DoubleEntry}
\NormalTok{    \}}

\NormalTok{    package "UC04: Danh mục phân loại 3 cấp (Category Hierarchy)" \{}
\NormalTok{        usecase "UC04.1: Khởi tạo cây danh mục chuẩn\textbackslash{}n(Init Default Taxonomy)" as UC04\_1}
\NormalTok{        usecase "UC04.2: Tùy biến thêm danh mục con\textbackslash{}n(Add Custom Subcategory)" as UC04\_2}
\NormalTok{        usecase "UC04.3: Ràng buộc an toàn khi sửa và xóa\textbackslash{}n(Enforce Category Integrity)" as UC04\_3}
\NormalTok{    \}}
\NormalTok{\}}

\NormalTok{\textquotesingle{} Tác nhân Owner kích hoạt các ca sử dụng}
\NormalTok{Owner {-}{-}\textgreater{} UC01\_1}
\NormalTok{Owner {-}{-}\textgreater{} UC01\_2}
\NormalTok{Owner {-}{-}\textgreater{} UC01\_3}
\NormalTok{Owner {-}{-}\textgreater{} UC01\_4}

\NormalTok{Owner {-}{-}\textgreater{} UC02\_1}
\NormalTok{Owner {-}{-}\textgreater{} UC02\_2}
\NormalTok{Owner {-}{-}\textgreater{} UC02\_3}
\NormalTok{Owner {-}{-}\textgreater{} UC02\_4}
\NormalTok{Owner {-}{-}\textgreater{} UC02\_5}

\NormalTok{Owner {-}{-}\textgreater{} UC03\_1}
\NormalTok{Owner {-}{-}\textgreater{} UC03\_3}

\NormalTok{Owner {-}{-}\textgreater{} UC04\_1}
\NormalTok{Owner {-}{-}\textgreater{} UC04\_2}
\NormalTok{Owner {-}{-}\textgreater{} UC04\_3}

\NormalTok{\textquotesingle{} Mối quan hệ Include}
\NormalTok{UC02\_2 ..\textgreater{} UC\_ValBal : \textless{}\textless{}include\textgreater{}\textgreater{}}
\NormalTok{UC03\_1 ..\textgreater{} UC\_ValBal : \textless{}\textless{}include\textgreater{}\textgreater{}}
\NormalTok{UC03\_1 ..\textgreater{} UC\_DoubleEntry : \textless{}\textless{}include\textgreater{}\textgreater{}}
\NormalTok{UC02\_4 ..\textgreater{} UC\_ValBal : \textless{}\textless{}include\textgreater{}\textgreater{}\textbackslash{}n(Khi tăng giá trị chi)}

\NormalTok{UC02\_1 ..\textgreater{} UC04\_1 : \textless{}\textless{}include\textgreater{}\textgreater{}\textbackslash{}n(Chọn danh mục Thu)}
\NormalTok{UC02\_2 ..\textgreater{} UC04\_1 : \textless{}\textless{}include\textgreater{}\textgreater{}\textbackslash{}n(Chọn danh mục Chi)}

\NormalTok{UC02\_1 ..\textgreater{} UC01\_2 : \textless{}\textless{}include\textgreater{}\textgreater{}\textbackslash{}n(Cộng số dư ví nhận)}
\NormalTok{UC02\_2 ..\textgreater{} UC01\_2 : \textless{}\textless{}include\textgreater{}\textgreater{}\textbackslash{}n(Trừ số dư ví chi)}

\NormalTok{\textquotesingle{} Mối quan hệ Extend}
\NormalTok{UC03\_2 ..\textgreater{} UC03\_1 : \textless{}\textless{}extend\textgreater{}\textgreater{}\textbackslash{}n(Điểm mở rộng: Có phí chuyển khoản)}
\NormalTok{UC04\_3 ..\textgreater{} UC04\_2 : \textless{}\textless{}extend\textgreater{}\textgreater{}\textbackslash{}n(Điểm mở rộng: Chuyển giao dịch khi xóa mục)}
\NormalTok{@enduml}
\end{Highlighting}
\end{Shaded}

\subparagraph{Phân tích chi tiết Phân vùng Dòng tiền và Ví
tiền}\label{phuxe2n-tuxedch-chi-tiux1ebft-phuxe2n-vuxf9ng-duxf2ng-tiux1ec1n-vuxe0-vuxed-tiux1ec1n}

\begin{enumerate}
\def\labelenumi{\arabic{enumi}.}
\item
  \textbf{Phân rã chức năng con:}

  \begin{itemize}
  \tightlist
  \item
    \texttt{UC01:\ Quản\ lý\ ví\ tiền}: Bao gồm khởi tạo ví với số dư
    ban đầu không âm (\texttt{UC01.1}), tính toán và hiển thị số dư động
    theo thời gian thực (\texttt{UC01.2}), cập nhật tên hoặc trạng thái
    lưu trữ (\texttt{UC01.3}), và thống kê tỷ trọng tài sản thanh khoản
    giữa các ví (\texttt{UC01.4}).
  \item
    \texttt{UC02:\ Ghi\ nhận\ thu\ chi}: Cung cấp khả năng ghi nhận thu
    nhập làm tăng số dư (\texttt{UC02.1}), ghi nhận chi tiêu làm giảm số
    dư (\texttt{UC02.2}), tra cứu giao dịch đa tiêu chí
    (\texttt{UC02.3}), sửa giao dịch với cơ chế hoàn tác và ghi nhận mới
    (\texttt{UC02.4}), và xóa giao dịch với việc phục hồi số dư ví
    (\texttt{UC02.5}).
  \item
    \texttt{UC03:\ Chuyển\ tiền\ nội\ bộ}: Thực thi điều chuyển vốn giữa
    hai ví nội bộ (\texttt{UC03.1}), hạch toán phí ngân hàng đi kèm
    (\texttt{UC03.2}), và hủy lệnh chuyển khoản phục hồi nguyên trạng cả
    hai ví (\texttt{UC03.3}).
  \item
    \texttt{UC04:\ Danh\ mục\ phân\ loại\ 3\ cấp}: Thiết lập cây danh
    mục chuẩn ISO (\texttt{UC04.1}), mở rộng danh mục con cấp 3 theo nhu
    cầu người dùng (\texttt{UC04.2}), và kiểm soát toàn vẹn khóa ngoại
    chặn xóa danh mục đang có giao dịch phát sinh (\texttt{UC04.3}).
  \end{itemize}
\item
  \textbf{Ràng buộc kiểm soát bất biến:}

  \begin{itemize}
  \tightlist
  \item
    Ca sử dụng \texttt{Kiểm\ tra\ số\ dư\ ví\ khả\ dụng\ (UC\_ValBal)}
    được bao hàm bắt buộc
    (\texttt{\textless{}\textless{}include\textgreater{}\textgreater{}})
    trong mọi hành động làm giảm số dư ví (ghi chi tiêu \texttt{UC02.2},
    chuyển tiền \texttt{UC03.1}, chỉnh sửa tăng chi phí
    \texttt{UC02.4}). Việc này bảo đảm bất biến \texttt{BR-FIN-02} được
    thực thi ở mức kiến trúc phần mềm, ngăn chặn hoàn toàn số dư âm.
  \item
    Ca sử dụng
    \texttt{Ghi\ sổ\ kép\ đối\ ứng\ đồng\ thời\ (UC\_DoubleEntry)} được
    bao hàm trong \texttt{UC03.1} để bảo đảm tính nguyên tử (Atomicity):
    số tiền bị trừ ở ví nguồn bằng chính xác số tiền được cộng vào ví
    đích, bảo toàn tổng lượng tiền tệ của hệ thống theo
    \texttt{BR-FIN-03}.
  \end{itemize}
\end{enumerate}

\paragraph{7.2.2 Phân vùng Tích lũy, Đầu tư và Nghĩa vụ tài chính
(Wealth \& Liabilities
Subsystem)}\label{722-phuxe2n-vuxf9ng-tuxedch-lux169y-ux111ux1ea7u-tux1b0-vuxe0-nghux129a-vux1ee5-tuxe0i-chuxednh-wealth--liabilities-subsystem}

Phân vùng này bao gồm 4 ca sử dụng: \texttt{UC05}, \texttt{UC06},
\texttt{UC07} và \texttt{UC08}. Phân vùng chịu trách nhiệm quản lý cấu
trúc tài sản dài hạn, danh mục đầu tư tăng trưởng, các hợp đồng tiền gửi
tiết kiệm và nghĩa vụ công nợ cá nhân.

\begin{Shaded}
\begin{Highlighting}[]
\NormalTok{@startuml}
\NormalTok{left to right direction}
\NormalTok{skinparam handwritten false}
\NormalTok{skinparam monochrome false}
\NormalTok{skinparam shadowing false}
\NormalTok{skinparam defaultFontName "Segoe UI", Arial, sans{-}serif}
\NormalTok{skinparam defaultFontSize 12}
\NormalTok{skinparam roundCorner 8}

\NormalTok{skinparam actor \{}
\NormalTok{    BackgroundColor \#E8F4FD}
\NormalTok{    BorderColor \#2B6CB0}
\NormalTok{\}}

\NormalTok{skinparam usecase \{}
\NormalTok{    BackgroundColor \#FFFFFF}
\NormalTok{    BorderColor \#2B6CB0}
\NormalTok{    ArrowColor \#2D3748}
\NormalTok{\}}

\NormalTok{skinparam package \{}
\NormalTok{    BackgroundColor \#F7FAFC}
\NormalTok{    BorderColor \#CBD5E0}
\NormalTok{\}}

\NormalTok{skinparam rectangle \{}
\NormalTok{    BackgroundColor \#FFFFFF}
\NormalTok{    BorderColor \#A0AEC0}
\NormalTok{\}}

\NormalTok{actor :Chủ sở hữu hệ thống\textbackslash{}n(Owner): as Owner}

\NormalTok{rectangle "Phân vùng Tích lũy, Đầu tư và Nghĩa vụ tài chính (Wealth \& Liabilities Subsystem)" \{}
\NormalTok{    package "UC05: Quản lý công nợ (Debt Management)" \{}
\NormalTok{        usecase "UC05.1: Ghi nhận cho vay\textbackslash{}n(Record Lend)" as UC05\_1}
\NormalTok{        usecase "UC05.2: Ghi nhận đi vay\textbackslash{}n(Record Borrow)" as UC05\_2}
\NormalTok{        usecase "UC05.3: Thu hồi nợ cho vay\textbackslash{}n(Collect Debt)" as UC05\_3}
\NormalTok{        usecase "UC05.4: Trả nợ đi vay\textbackslash{}n(Repay Debt)" as UC05\_4}
\NormalTok{        usecase "UC05.5: Xóa nợ khó đòi phi tiền mặt\textbackslash{}n(Write{-}off Bad Debt)" as UC05\_5}
        
\NormalTok{        usecase "Kiểm soát dư nợ không âm\textbackslash{}n(Enforce Debt Invariant \textgreater{}= 0)" as UC\_ValDebt}
\NormalTok{    \}}

\NormalTok{    package "UC06: Quản lý sổ tiết kiệm (Term Savings Management)" \{}
\NormalTok{        usecase "UC06.1: Mở sổ tiết kiệm mới\textbackslash{}n(Open Savings Account)" as UC06\_1}
\NormalTok{        usecase "UC06.2: Tính lãi dự kiến \& theo dõi kỳ hạn\textbackslash{}n(Calculate Accrued Interest)" as UC06\_2}
\NormalTok{        usecase "UC06.3: Tất toán sổ tiết kiệm đúng hạn\textbackslash{}n(Maturity Settlement)" as UC06\_3}
\NormalTok{        usecase "UC06.4: Tất toán sổ tiết kiệm trước hạn\textbackslash{}n(Early Withdrawal Settlement)" as UC06\_4}
        
\NormalTok{        usecase "Tính lãi có kỳ hạn thỏa thuận\textbackslash{}n(Compute Term Interest)" as UC\_TermInt}
\NormalTok{        usecase "Tính lãi không kỳ hạn theo ngày\textbackslash{}n(Compute Demand Interest)" as UC\_DemandInt}
\NormalTok{    \}}

\NormalTok{    package "UC07: Quản lý chứng chỉ quỹ (Fund Investment Management)" \{}
\NormalTok{        usecase "UC07.1: Ghi nhận lệnh mua CCQ\textbackslash{}n(Buy Fund Certificates)" as UC07\_1}
\NormalTok{        usecase "UC07.2: Cập nhật NAV hàng ngày\textbackslash{}n(Update Daily NAV)" as UC07\_2}
\NormalTok{        usecase "UC07.3: Định giá thị trường \& lãi lỗ\textbackslash{}n(Mark{-}to{-}Market Valuation)" as UC07\_3}
\NormalTok{        usecase "UC07.4: Ghi nhận lệnh bán CCQ\textbackslash{}n(Sell Fund Certificates)" as UC07\_4}
        
\NormalTok{        usecase "Kiểm tra số lượng CCQ khả dụng\textbackslash{}n(Validate Holding Units)" as UC\_ValUnits}
\NormalTok{    \}}

\NormalTok{    package "UC08: Quản lý quỹ mục tiêu (Financial Goals Management)" \{}
\NormalTok{        usecase "UC08.1: Thiết lập quỹ tài chính mục tiêu\textbackslash{}n(Create Financial Goal)" as UC08\_1}
\NormalTok{        usecase "UC08.2: Phân bổ tài sản vào mục tiêu\textbackslash{}n(Allocate Assets {-} 3 Pillars)" as UC08\_2}
\NormalTok{        usecase "UC08.3: Theo dõi tiến độ tích lũy\textbackslash{}n(Track Goal Progress)" as UC08\_3}
\NormalTok{        usecase "UC08.4: Điều chỉnh \& tái phân bổ\textbackslash{}n(Rebalance \& Adjust Goal)" as UC08\_4}
        
\NormalTok{        usecase "Kiểm soát hạn mức phân bổ \textless{}= 100\%\textbackslash{}n(Enforce Allocation Cap)" as UC\_ValAlloc}
\NormalTok{    \}}
\NormalTok{\}}

\NormalTok{\textquotesingle{} Tác nhân Owner kích hoạt các ca sử dụng}
\NormalTok{Owner {-}{-}\textgreater{} UC05\_1}
\NormalTok{Owner {-}{-}\textgreater{} UC05\_2}
\NormalTok{Owner {-}{-}\textgreater{} UC05\_3}
\NormalTok{Owner {-}{-}\textgreater{} UC05\_4}

\NormalTok{Owner {-}{-}\textgreater{} UC06\_1}
\NormalTok{Owner {-}{-}\textgreater{} UC06\_2}
\NormalTok{Owner {-}{-}\textgreater{} UC06\_3}

\NormalTok{Owner {-}{-}\textgreater{} UC07\_1}
\NormalTok{Owner {-}{-}\textgreater{} UC07\_2}
\NormalTok{Owner {-}{-}\textgreater{} UC07\_3}
\NormalTok{Owner {-}{-}\textgreater{} UC07\_4}

\NormalTok{Owner {-}{-}\textgreater{} UC08\_1}
\NormalTok{Owner {-}{-}\textgreater{} UC08\_2}
\NormalTok{Owner {-}{-}\textgreater{} UC08\_3}
\NormalTok{Owner {-}{-}\textgreater{} UC08\_4}

\NormalTok{\textquotesingle{} Mối quan hệ Include}
\NormalTok{UC05\_3 ..\textgreater{} UC\_ValDebt : \textless{}\textless{}include\textgreater{}\textgreater{}}
\NormalTok{UC05\_4 ..\textgreater{} UC\_ValDebt : \textless{}\textless{}include\textgreater{}\textgreater{}}

\NormalTok{UC06\_3 ..\textgreater{} UC\_TermInt : \textless{}\textless{}include\textgreater{}\textgreater{}}
\NormalTok{UC06\_4 ..\textgreater{} UC\_DemandInt : \textless{}\textless{}include\textgreater{}\textgreater{}}

\NormalTok{UC07\_4 ..\textgreater{} UC\_ValUnits : \textless{}\textless{}include\textgreater{}\textgreater{}}

\NormalTok{UC08\_2 ..\textgreater{} UC\_ValAlloc : \textless{}\textless{}include\textgreater{}\textgreater{}}
\NormalTok{UC08\_3 ..\textgreater{} UC08\_2 : \textless{}\textless{}include\textgreater{}\textgreater{}\textbackslash{}n(Cập nhật giá trị tài sản 3 cấu phần)}

\NormalTok{\textquotesingle{} Mối quan hệ Extend}
\NormalTok{UC05\_5 ..\textgreater{} UC05\_3 : \textless{}\textless{}extend\textgreater{}\textgreater{}\textbackslash{}n(Điểm mở rộng: Đối tác mất khả năng trả nợ)}
\NormalTok{UC05\_5 ..\textgreater{} UC05\_4 : \textless{}\textless{}extend\textgreater{}\textgreater{}\textbackslash{}n(Điểm mở rộng: Được chủ nợ xóa bớt nghĩa vụ)}
\NormalTok{UC06\_4 ..\textgreater{} UC06\_3 : \textless{}\textless{}extend\textgreater{}\textgreater{}\textbackslash{}n(Điểm mở rộng: Rút vốn trước ngày đáo hạn)}
\NormalTok{@enduml}
\end{Highlighting}
\end{Shaded}

\subparagraph{Phân tích chi tiết Phân vùng Tích lũy, Đầu tư và Nghĩa vụ
tài
chính}\label{phuxe2n-tuxedch-chi-tiux1ebft-phuxe2n-vuxf9ng-tuxedch-lux169y-ux111ux1ea7u-tux1b0-vuxe0-nghux129a-vux1ee5-tuxe0i-chuxednh}

\begin{enumerate}
\def\labelenumi{\arabic{enumi}.}
\item
  \textbf{Phân rã chức năng con:}

  \begin{itemize}
  \tightlist
  \item
    \texttt{UC05:\ Quản\ lý\ công\ nợ}: Quản trị 2 sổ nợ độc lập gồm Nợ
    phải thu (cho vay \texttt{UC05.1}, thu nợ \texttt{UC05.3}) và Nợ
    phải trả (đi vay \texttt{UC05.2}, trả nợ \texttt{UC05.4}). Ca sử
    dụng \texttt{UC05.5} mở rộng kịch bản thu/trả nợ để xử lý xóa nợ khó
    đòi phi tiền mặt khi phát sinh biến cố đối tác.
  \item
    \texttt{UC06:\ Quản\ lý\ sổ\ tiết\ kiệm}: Quản lý việc gửi tiền tiết
    kiệm từ ví thanh toán (\texttt{UC06.1}), tính toán lãi tích lũy theo
    dõi thời gian đáo hạn (\texttt{UC06.2}), và xử lý tất toán. Điểm mở
    rộng \texttt{UC06.4} can thiệp vào quy trình tất toán để tính lãi
    không kỳ hạn khi người dùng rút vốn trước ngày đáo hạn hợp đồng theo
    \texttt{BR-FIN-05}.
  \item
    \texttt{UC07:\ Quản\ lý\ chứng\ chỉ\ quỹ}: Ghi nhận mua chứng chỉ
    quỹ theo giá NAV khớp lệnh (\texttt{UC07.1}), lưu trữ lịch sử biến
    động NAV theo ngày (\texttt{UC07.2}), định giá lại danh mục theo giá
    trị thị trường Mark-to-Market (\texttt{UC07.3}), và thực hiện bán
    chứng chỉ quỹ có khấu trừ thuế phí (\texttt{UC07.4}) theo
    \texttt{BR-FIN-06}.
  \item
    \texttt{UC08:\ Quản\ lý\ quỹ\ mục\ tiêu}: Khởi tạo mục tiêu tài
    chính độc lập (\texttt{UC08.1}), phân bổ tài sản từ 3 cấu phần tài
    chính độc lập (\texttt{UC08.2}), tự động đo lường tiến độ phần trăm
    tích lũy (\texttt{UC08.3}), và tái cân bằng cơ cấu phân bổ khi có sự
    thay đổi kế hoạch (\texttt{UC08.4}).
  \end{itemize}
\item
  \textbf{Ràng buộc kiểm soát bất biến:}

  \begin{itemize}
  \tightlist
  \item
    Quy tắc công nợ \texttt{BR-FIN-07}: \texttt{UC\_ValDebt} kiểm tra số
    tiền thu nợ hoặc trả nợ không được vượt quá số dư nợ còn lại, bảo
    đảm \(Balance_{debt} \ge 0\).
  \item
    Quy tắc phân bổ quỹ mục tiêu \texttt{BR-FIN-08}:
    \texttt{UC\_ValAlloc} kiểm soát tổng tỷ lệ phần trăm phân bổ của một
    tài sản (ví tiền, sổ tiết kiệm, hoặc CCQ) vào các quỹ mục tiêu không
    bao giờ vượt quá 100\%.
  \end{itemize}
\end{enumerate}

\paragraph{7.2.3 Phân vùng Báo cáo, Phân tích và Quản trị dữ liệu
(Analytics \& Administration
Subsystem)}\label{723-phuxe2n-vuxf9ng-buxe1o-cuxe1o-phuxe2n-tuxedch-vuxe0-quux1ea3n-trux1ecb-dux1eef-liux1ec7u-analytics--administration-subsystem}

Phân vùng này bao gồm 2 ca sử dụng chiến lược: \texttt{UC09} và
\texttt{UC10}. Phân vùng cung cấp thông tin quản trị dòng tiền, đo lường
toàn diện sức khỏe tài chính và thiết lập cơ chế phòng vệ dữ liệu an
toàn cho người dùng.

\begin{Shaded}
\begin{Highlighting}[]
\NormalTok{@startuml}
\NormalTok{left to right direction}
\NormalTok{skinparam handwritten false}
\NormalTok{skinparam monochrome false}
\NormalTok{skinparam shadowing false}
\NormalTok{skinparam defaultFontName "Segoe UI", Arial, sans{-}serif}
\NormalTok{skinparam defaultFontSize 12}
\NormalTok{skinparam roundCorner 8}

\NormalTok{skinparam actor \{}
\NormalTok{    BackgroundColor \#E8F4FD}
\NormalTok{    BorderColor \#2B6CB0}
\NormalTok{\}}

\NormalTok{skinparam usecase \{}
\NormalTok{    BackgroundColor \#FFFFFF}
\NormalTok{    BorderColor \#2B6CB0}
\NormalTok{    ArrowColor \#2D3748}
\NormalTok{\}}

\NormalTok{skinparam package \{}
\NormalTok{    BackgroundColor \#F7FAFC}
\NormalTok{    BorderColor \#CBD5E0}
\NormalTok{\}}

\NormalTok{skinparam rectangle \{}
\NormalTok{    BackgroundColor \#FFFFFF}
\NormalTok{    BorderColor \#A0AEC0}
\NormalTok{\}}

\NormalTok{actor :Chủ sở hữu hệ thống\textbackslash{}n(Owner): as Owner}

\NormalTok{rectangle "Phân vùng Báo cáo, Phân tích và Quản trị dữ liệu (Analytics \& Administration Subsystem)" \{}
\NormalTok{    package "UC09: Thống kê \& Phân tích tài chính (Analytics \& Reporting)" \{}
\NormalTok{        usecase "UC09.1: Bảng thống kê dòng tiền trượt 12 tháng\textbackslash{}n(Rolling 12{-}Month Cashflow)" as UC09\_1}
\NormalTok{        usecase "UC09.2: Báo cáo cơ cấu thu chi 2 tầng\textbackslash{}n(Income \& Expense Breakdown)" as UC09\_2}
\NormalTok{        usecase "UC09.3: Báo cáo tài sản ròng Net Worth 5 cấu phần\textbackslash{}n(Net Worth \& Balance Sheet)" as UC09\_3}
\NormalTok{        usecase "UC09.4: Sao kê công nợ chi tiết theo đối tác\textbackslash{}n(Partner Debt Statement)" as UC09\_4}
        
\NormalTok{        usecase "Xử lý cửa sổ trượt 12 tháng\textbackslash{}n(Sliding Window Generator)" as UC\_Sliding}
\NormalTok{        usecase "Xử lý biên chia cho 0 tỷ lệ tiết kiệm\textbackslash{}n(Safe Savings Rate Calculation)" as UC\_SafeDiv}
\NormalTok{        usecase "Hợp nhất 5 cấu phần tài sản \& công nợ\textbackslash{}n(Consolidate 5{-}Pillar Net Worth)" as UC\_Consolidate}
\NormalTok{    \}}

\NormalTok{    package "UC10: Quản trị dữ liệu \& Sao lưu (Backup, Export \& Import)" \{}
\NormalTok{        usecase "UC10.1: Xuất dữ liệu tài chính ra JSON/CSV\textbackslash{}n(Export Financial Data)" as UC10\_1}
\NormalTok{        usecase "UC10.2: Nhập và phục hồi dữ liệu từ tệp sao lưu\textbackslash{}n(Import \& Restore Data)" as UC10\_2}
\NormalTok{        usecase "UC10.3: Kiểm toán đối soát toàn vẹn dữ liệu\textbackslash{}n(Data Integrity Audit)" as UC10\_3}
        
\NormalTok{        usecase "Quy trình kiểm định an toàn 4 bước\textbackslash{}n(4{-}Step Integrity Validation)" as UC\_Val4Step}
\NormalTok{        usecase "Cơ chế giao dịch nguyên tử \& Rollback\textbackslash{}n(Atomic Transaction \& Rollback)" as UC\_Rollback}
\NormalTok{        usecase "Chuẩn hóa chuỗi số học BigDecimal\textbackslash{}n(Format Exact Decimal String)" as UC\_ExactDecimal}
\NormalTok{    \}}
\NormalTok{\}}

\NormalTok{\textquotesingle{} Tác nhân Owner kích hoạt các ca sử dụng}
\NormalTok{Owner {-}{-}\textgreater{} UC09\_1}
\NormalTok{Owner {-}{-}\textgreater{} UC09\_2}
\NormalTok{Owner {-}{-}\textgreater{} UC09\_3}
\NormalTok{Owner {-}{-}\textgreater{} UC09\_4}

\NormalTok{Owner {-}{-}\textgreater{} UC10\_1}
\NormalTok{Owner {-}{-}\textgreater{} UC10\_2}
\NormalTok{Owner {-}{-}\textgreater{} UC10\_3}

\NormalTok{\textquotesingle{} Quan hệ Include trong UC09}
\NormalTok{UC09\_1 ..\textgreater{} UC\_Sliding : \textless{}\textless{}include\textgreater{}\textgreater{}}
\NormalTok{UC09\_1 ..\textgreater{} UC\_SafeDiv : \textless{}\textless{}include\textgreater{}\textgreater{}}
\NormalTok{UC09\_3 ..\textgreater{} UC\_Consolidate : \textless{}\textless{}include\textgreater{}\textgreater{}}

\NormalTok{\textquotesingle{} Quan hệ Include trong UC10}
\NormalTok{UC10\_1 ..\textgreater{} UC\_ExactDecimal : \textless{}\textless{}include\textgreater{}\textgreater{}}
\NormalTok{UC10\_2 ..\textgreater{} UC\_Val4Step : \textless{}\textless{}include\textgreater{}\textgreater{}}
\NormalTok{UC10\_2 ..\textgreater{} UC\_Rollback : \textless{}\textless{}include\textgreater{}\textgreater{}\textbackslash{}n(Khi phát hiện vi phạm)}

\NormalTok{\textquotesingle{} Quan hệ Extend}
\NormalTok{UC10\_3 ..\textgreater{} UC10\_2 : \textless{}\textless{}extend\textgreater{}\textgreater{}\textbackslash{}n(Điểm mở rộng: Kiểm toán sau phục hồi dữ liệu)}
\NormalTok{@enduml}
\end{Highlighting}
\end{Shaded}

\subparagraph{Phân tích chi tiết Phân vùng Báo cáo, Phân tích và Quản
trị dữ
liệu}\label{phuxe2n-tuxedch-chi-tiux1ebft-phuxe2n-vuxf9ng-buxe1o-cuxe1o-phuxe2n-tuxedch-vuxe0-quux1ea3n-trux1ecb-dux1eef-liux1ec7u}

\begin{enumerate}
\def\labelenumi{\arabic{enumi}.}
\item
  \textbf{Phân rã chức năng con:}

  \begin{itemize}
  \tightlist
  \item
    \texttt{UC09:\ Thống\ kê\ \&\ Phân\ tích\ tài\ chính}: Cung cấp ma
    trận dòng tiền trượt 12 tháng (\texttt{UC09.1}), phân tích cơ cấu
    chi tiêu 2 tầng nhóm và danh mục con (\texttt{UC09.2}), báo cáo bảng
    cân đối kế toán cá nhân và giá trị tài sản ròng Net Worth 5 cấu phần
    (\texttt{UC09.3}), và trích xuất sao kê lịch sử công nợ chi tiết
    theo từng đối tác giao dịch (\texttt{UC09.4}).
  \item
    \texttt{UC10:\ Quản\ trị\ dữ\ liệu\ \&\ Sao\ lưu}: Kết xuất toàn bộ
    hoặc từng phần sổ cái ra định dạng JSON/CSV có cấu trúc an toàn
    (\texttt{UC10.1}), nạp và khôi phục dữ liệu từ tệp tin
    (\texttt{UC10.2}), và cung cấp công cụ tự động kiểm toán tính toàn
    vẹn dữ liệu sổ cái (\texttt{UC10.3}).
  \end{itemize}
\item
  \textbf{Ràng buộc kiểm soát bất biến:}

  \begin{itemize}
  \tightlist
  \item
    Bảo toàn chuỗi thời gian liên tục \texttt{FR-REP-01}:
    \texttt{UC\_Sliding} bảo đảm bảng dữ liệu luôn có đủ đúng 12 tháng
    liên tục lùi từ tháng hiện tại, tự động bù đắp tháng không phát sinh
    giao dịch bằng số tiền \texttt{0.00\ VNĐ}.
  \item
    Kiểm soát an toàn phép chia \texttt{FR-REP-01}: \texttt{UC\_SafeDiv}
    xử lý an toàn phép tính tỷ lệ tiết kiệm
    \(SavingsRate = (CashFlow / Income) \times 100\%\), gán giá trị
    \(0.00\%\) khi \(Income = 0\) để loại bỏ hoàn toàn nguy cơ phát sinh
    ngoại lệ \texttt{ArithmeticException}.
  \item
    Cân bằng tài sản ròng \texttt{FR-REP-03}: \texttt{UC\_Consolidate}
    thực hiện hợp nhất đồng thời 5 cấu phần tài sản:
    \[NetWorth = Wallets + Savings + Investments + Receivable - Payable\]
  \item
    Quy trình kiểm định 4 bước \texttt{FR-DAT-02}: \texttt{UC\_Val4Step}
    bắt buộc thực hiện kiểm tra tuần tự: (1) Cú pháp tệp và kiểu dữ
    liệu; (2) Toàn vẹn khóa ngoại danh mục theo \texttt{BR-FIN-04}; (3)
    Toàn vẹn tham chiếu ví tiền; và (4) Bất biến số dư ví không âm theo
    \texttt{BR-FIN-02}.
  \item
    Giao dịch nguyên tử và bảo toàn số học \texttt{FR-DAT-01},
    \texttt{NFR-03}: \texttt{UC\_Rollback} bảo đảm nguyên tắc
    All-or-Nothing khi nạp tệp sao lưu; \texttt{UC\_ExactDecimal} bảo
    đảm dữ liệu số học \texttt{BigDecimal} được duy trì nguyên bản, loại
    bỏ hoàn toàn hiện tượng mất chữ số có nghĩa hoặc sai số làm tròn.
  \end{itemize}
\end{enumerate}

\clearpage
\subsection{8. MA TRẬN KIỂM CHỨNG VÀ TRUY VẾT YÊU CẦU (REQUIREMENTS
TRACEABILITY MATRIX -
RTM)}\label{8-ma-trux1eadn-kiux1ec3m-chux1ee9ng-vuxe0-truy-vux1ebft-yuxeau-cux1ea7u-requirements-traceability-matrix---rtm}

Phần này thiết lập ma trận truy vết và phương pháp kiểm chứng yêu cầu
theo tiêu chuẩn ISO/IEC/IEEE 29148:2018 §7 (Requirements Verification
and Validation) và §8.3 (Traceability Matrix). Ma trận cung cấp cơ sở kỹ
thuật nhằm:

\begin{itemize}
\tightlist
\item
  Bảo đảm tính đầy đủ: Mọi mục tiêu hệ thống (G1 đến G5) đều được chuyển
  hóa đầy đủ thành các yêu cầu chức năng (FR-WAL đến FR-DAT) và được
  hiện thực hóa qua các ca sử dụng (UC01 đến UC10).
\item
  Định hướng phân tích hướng đối tượng (OOA): Thiết lập mối liên kết
  trực tiếp giữa các yêu cầu chức năng với các lớp đối tượng thực thể
  miền (Domain Model Entities) được phân tích chi tiết trong giai đoạn
  phân tích OOA.
\item
  Kiểm soát chất lượng nghiệm thu: Xác định phương pháp kiểm chứng cụ
  thể (Inspection, Test, Demonstration, Analysis) cùng tiêu chí nghiệm
  thu tương ứng cho toàn bộ 39 yêu cầu chức năng và 6 yêu cầu phi chức
  năng.
\end{itemize}

\subsubsection{8.1 Ma trận truy vết yêu cầu 4 chiều (Requirements
Traceability
Matrix)}\label{81-ma-trux1eadn-truy-vux1ebft-yuxeau-cux1ea7u-4-chiux1ec1u-requirements-traceability-matrix}

Bảng dưới đây đối chiếu 4 chiều kỹ thuật của hệ thống:

\begin{enumerate}
\def\labelenumi{\arabic{enumi}.}
\tightlist
\item
  \textbf{Mục tiêu hệ thống (Goal ID):} 5 mục tiêu cốt lõi từ Mục 2.2
  (\texttt{G1} đến \texttt{G5}).
\item
  \textbf{Phân hệ nghiệp vụ \& Yêu cầu chức năng (FR ID):} 39 yêu cầu
  chức năng từ Mục 4 (\texttt{FR-WAL-01} đến \texttt{FR-DAT-03}).
\item
  \textbf{Ca sử dụng tham chiếu (UC ID):} 10 ca sử dụng chi tiết từ Mục
  6 (\texttt{UC01} đến \texttt{UC10}).
\item
  \textbf{Lớp thực thể miền OOA (Domain Model Entities):} Các thực thể
  nghiệp vụ xuất phát từ ngữ cảnh tài chính thực tế theo khung phân tích
  tại \texttt{docs/deliverables/analysis/framework.md} (\texttt{Wallet},
  \texttt{Transaction}, \texttt{Transfer}, \texttt{Category},
  \texttt{Debt}, \texttt{SavingAccount}, \texttt{FundCertificate},
  \texttt{FinancialGoal}, \texttt{MonthlyStat},
  \texttt{FinancialLedger}).
\end{enumerate}

{\def\LTcaptype{none}
\begin{longtable}[]{|>{\raggedright\arraybackslash}p{0.07\linewidth}|>{\raggedright\arraybackslash}p{0.13\linewidth}|>{\raggedright\arraybackslash}p{0.10\linewidth}|>{\raggedright\arraybackslash}p{0.18\linewidth}|>{\raggedright\arraybackslash}p{0.13\linewidth}|>{\raggedright\arraybackslash}p{0.19\linewidth}|}
\hline
\everycr{\noalign{\ifnum0=`}\fi\hrule height 0.4pt\futurelet\reserved@a\LT@@tabularcr}
\textbf{Mã mục tiêu} & \textbf{Phân hệ nghiệp vụ} & \textbf{Mã yêu cầu} & \textbf{Tên chức năng} & \textbf{Ca sử dụng tham chiếu} & \textbf{Lớp thực thể miền OOA (Domain Model)} \\
\hline
\endfirsthead
\hline
\textbf{Mã mục tiêu} & \textbf{Phân hệ nghiệp vụ} & \textbf{Mã yêu cầu} & \textbf{Tên chức năng} & \textbf{Ca sử dụng tham chiếu} & \textbf{Lớp thực thể miền OOA (Domain Model)} \\
\hline
\endhead
\hline
\endfoot
\hline
\endlastfoot
G2, G1 & Quản lý ví tiền & FR-WAL-01 & Khởi tạo ví tiền mới & UC01 &
\texttt{Wallet} \\
G2 & Quản lý ví tiền & FR-WAL-02 & Xem danh sách và chi tiết số dư động
& UC01 & \texttt{Wallet} \\
G2 & Quản lý ví tiền & FR-WAL-03 & Cập nhật thông tin ví tiền & UC01 &
\texttt{Wallet} \\
G2 & Quản lý ví tiền & FR-WAL-04 & Thống kê cơ cấu tài sản theo ví &
UC01 & \texttt{Wallet} \\
G1, G5 & Ghi nhận thu chi & FR-TXN-01 & Tạo mới giao dịch thu nhập &
UC02 & \texttt{Transaction}, \texttt{Wallet}, \texttt{Category} \\
G1, G5 & Ghi nhận thu chi & FR-TXN-02 & Tạo mới giao dịch chi tiêu &
UC02 & \texttt{Transaction}, \texttt{Wallet}, \texttt{Category} \\
G1 & Ghi nhận thu chi & FR-TXN-03 & Tra cứu và lọc lịch sử giao dịch &
UC02 & \texttt{Transaction}, \texttt{Wallet}, \texttt{Category} \\
G1, G5 & Ghi nhận thu chi & FR-TXN-04 & Chỉnh sửa giao dịch thu chi &
UC02 & \texttt{Transaction}, \texttt{Wallet}, \texttt{Category} \\
G1, G5 & Ghi nhận thu chi & FR-TXN-05 & Xóa giao dịch thu chi & UC02 &
\texttt{Transaction}, \texttt{Wallet} \\
G1, G5 & Chuyển tiền nội bộ & FR-TRF-01 & Chuyển tiền giữa hai ví nội bộ
& UC03 & \texttt{Transfer}, \texttt{Wallet} \\
G1, G5 & Chuyển tiền nội bộ & FR-TRF-02 & Ghi nhận phí chuyển tiền phát
sinh & UC03 & \texttt{Transfer}, \texttt{Transaction},
\texttt{Wallet} \\
G1, G5 & Chuyển tiền nội bộ & FR-TRF-03 & Lịch sử và hủy chuyển tiền nội
bộ & UC03 & \texttt{Transfer}, \texttt{Wallet} \\
G1 & Danh mục phân loại & FR-CAT-01 & Khởi tạo và quản lý cấu trúc phân
loại 3 cấp & UC04 & \texttt{Category} \\
G1 & Danh mục phân loại & FR-CAT-02 & Tùy biến danh mục người dùng &
UC04 & \texttt{Category} \\
G1 & Danh mục phân loại & FR-CAT-03 & Ràng buộc toàn vẹn khi sửa và xóa
danh mục & UC04 & \texttt{Category}, \texttt{Transaction} \\
G3, G2 & Quản lý công nợ & FR-DEBT-01 & Ghi nhận khoản cho vay (Lend) &
UC05 & \texttt{Debt}, \texttt{Wallet}, \texttt{Transaction} \\
G3, G2 & Quản lý công nợ & FR-DEBT-02 & Ghi nhận khoản đi vay (Borrow) &
UC05 & \texttt{Debt}, \texttt{Wallet}, \texttt{Transaction} \\
G3, G2 & Quản lý công nợ & FR-DEBT-03 & Thu hồi nợ cho vay (Collect) &
UC05 & \texttt{Debt}, \texttt{Wallet}, \texttt{Transaction} \\
G3, G2 & Quản lý công nợ & FR-DEBT-04 & Trả nợ đi vay (Repay) & UC05 &
\texttt{Debt}, \texttt{Wallet}, \texttt{Transaction} \\
G3, G2 & Quản lý công nợ & FR-DEBT-05 & Xóa nợ khó đòi và miễn trừ nợ
(Write-off) & UC05 & \texttt{Debt} \\
G3, G2 & Quản lý sổ tiết kiệm & FR-SAV-01 & Mở sổ tiết kiệm mới & UC06 &
\texttt{SavingAccount}, \texttt{Wallet}, \texttt{Transaction} \\
G3, G5 & Quản lý sổ tiết kiệm & FR-SAV-02 & Tính toán tiền lãi dự kiến
và theo dõi kỳ hạn & UC06 & \texttt{SavingAccount} \\
G3, G2, G5 & Quản lý sổ tiết kiệm & FR-SAV-03 & Tất toán sổ tiết kiệm
đúng hạn & UC06 & \texttt{SavingAccount}, \texttt{Wallet},
\texttt{Transaction} \\
G3, G2, G5 & Quản lý sổ tiết kiệm & FR-SAV-04 & Tất toán sổ tiết kiệm
trước hạn & UC06 & \texttt{SavingAccount}, \texttt{Wallet},
\texttt{Transaction} \\
G3, G2 & Chứng chỉ quỹ đầu tư & FR-INV-01 & Ghi nhận lệnh mua chứng chỉ
quỹ & UC07 & \texttt{FundCertificate}, \texttt{Wallet},
\texttt{Transaction} \\
G3, G2 & Chứng chỉ quỹ đầu tư & FR-INV-02 & Cập nhật giá trị tài sản
ròng (NAV Update) & UC07 & \texttt{FundCertificate} \\
G3, G2, G5 & Chứng chỉ quỹ đầu tư & FR-INV-03 & Định giá thị trường và
tính lãi lỗ chưa thực hiện & UC07 & \texttt{FundCertificate} \\
G3, G2, G5 & Chứng chỉ quỹ đầu tư & FR-INV-04 & Ghi nhận bán chứng chỉ
quỹ & UC07 & \texttt{FundCertificate}, \texttt{Wallet},
\texttt{Transaction} \\
G4, G2 & Quản lý quỹ mục tiêu & FR-FND-01 & Thiết lập quỹ tài chính mục
tiêu & UC08 & \texttt{FinancialGoal} \\
G4, G2 & Quản lý quỹ mục tiêu & FR-FND-02 & Phân bổ tài sản vào quỹ mục
tiêu & UC08 & \texttt{FinancialGoal}, \texttt{Wallet},
\texttt{SavingAccount}, \texttt{FundCertificate} \\
G4, G2 & Quản lý quỹ mục tiêu & FR-FND-03 & Theo dõi tiến độ tích lũy
quỹ & UC08 & \texttt{FinancialGoal} \\
G4, G2 & Quản lý quỹ mục tiêu & FR-FND-04 & Điều chỉnh và tái phân bổ
quỹ mục tiêu & UC08 & \texttt{FinancialGoal} \\
G1, G5 & Thống kê và báo cáo & FR-REP-01 & Bảng thống kê dòng tiền trượt
12 tháng & UC09 & \texttt{MonthlyStat}, \texttt{Transaction} \\
G1 & Thống kê và báo cáo & FR-REP-02 & Báo cáo cơ cấu thu chi theo nhóm
và danh mục & UC09 & \texttt{MonthlyStat}, \texttt{Category},
\texttt{Transaction} \\
G2, G5 & Thống kê và báo cáo & FR-REP-03 & Báo cáo tổng tài sản và giá
trị tài sản ròng & UC09 & \texttt{FinancialLedger}, \texttt{Wallet},
\texttt{SavingAccount}, \texttt{FundCertificate}, \texttt{Debt} \\
G3 & Thống kê và báo cáo & FR-REP-04 & Sao kê công nợ chi tiết theo đối
tượng & UC09 & \texttt{Debt} \\
G2, G5 & Quản lý dữ liệu & FR-DAT-01 & Xuất dữ liệu tài chính (Export) &
UC10 & \texttt{FinancialLedger} \\
G2, G5 & Quản lý dữ liệu & FR-DAT-02 & Nhập dữ liệu và phục hồi (Import
\& Restore) & UC10 & \texttt{FinancialLedger} \\
G5 & Quản lý dữ liệu & FR-DAT-03 & Kiểm tra tính toàn vẹn dữ liệu (Data
Integrity Audit) & UC10 & \texttt{FinancialLedger} \\
\end{longtable}
}

\subsubsection{8.2 Ma trận phương pháp kiểm chứng yêu cầu (Verification
Methods
Matrix)}\label{82-ma-trux1eadn-phux1b0ux1a1ng-phuxe1p-kiux1ec3m-chux1ee9ng-yuxeau-cux1ea7u-verification-methods-matrix}

Theo tiêu chuẩn ISO/IEC/IEEE 29148:2018 §7, các phương pháp kiểm chứng
được phân định thành 4 loại hình chuẩn mực:

\begin{itemize}
\tightlist
\item
  \textbf{I (Inspection: Soát xét kỹ thuật):} Kiểm tra trực quan mã
  nguồn, tài liệu thiết kế, rà soát tĩnh (Code Review, Static Code
  Analysis) mà không thực thi chương trình.
\item
  \textbf{T (Test: Kiểm thử tự động):} Thực thi chương trình hoặc các
  module với tập dữ liệu đầu vào xác định, đối chiếu tự động kết quả
  thực tế với kết quả kỳ vọng (Unit test JUnit 5, Integration test).
\item
  \textbf{D (Demonstration: Trình diễn hoạt động):} Vận hành chức năng
  trực tiếp trên giao diện tương tác người dùng để quan sát hành vi,
  luồng xử lý và kết quả đầu ra theo kịch bản sử dụng thực tế.
\item
  \textbf{A (Analysis: Phân tích định lượng):} Đánh giá các chỉ số kỹ
  thuật thông qua công cụ phân tích dữ liệu, tính toán thống kê,
  profiling hiệu năng và đo lường độ bao phủ mã nguồn (JMH, VisualVM,
  JaCoCo).
\end{itemize}

\paragraph{8.2.1 Ma trận phương pháp kiểm chứng yêu cầu chức năng (39
FRs)}\label{821-ma-trux1eadn-phux1b0ux1a1ng-phuxe1p-kiux1ec3m-chux1ee9ng-yuxeau-cux1ea7u-chux1ee9c-nux103ng-39-frs}

{\def\LTcaptype{none}
\begin{longtable}[]{|>{\raggedright\arraybackslash}p{0.07\linewidth}|>{\raggedright\arraybackslash}p{0.16\linewidth}|>{\raggedright\arraybackslash}p{0.10\linewidth}|>{\raggedright\arraybackslash}p{0.18\linewidth}|>{\raggedright\arraybackslash}p{0.27\linewidth}|}
\hline
\everycr{\noalign{\ifnum0=`}\fi\hrule height 0.4pt\futurelet\reserved@a\LT@@tabularcr}
\textbf{Mã yêu cầu} & \textbf{Tên chức năng} & \textbf{Phương pháp chính} & \textbf{Môi trường và công cụ kiểm chứng} & \textbf{Tiêu chí nghiệm thu cốt lõi} \\
\hline
\endfirsthead
\hline
\textbf{Mã yêu cầu} & \textbf{Tên chức năng} & \textbf{Phương pháp chính} & \textbf{Môi trường và công cụ kiểm chứng} & \textbf{Tiêu chí nghiệm thu cốt lõi} \\
\hline
\endhead
\hline
\endfoot
\hline
\endlastfoot
FR-WAL-01 & Khởi tạo ví tiền mới & T, D & JUnit 5 / JavaFX UI & Tạo ví
thành công với tên duy nhất, loại ví hợp lệ và số dư ban đầu không âm \\
FR-WAL-02 & Xem danh sách và chi tiết số dư động & T, D & JUnit 5 /
JavaFX UI & Hiển thị chính xác danh sách ví và số dư động theo đúng công
thức BR-FIN-02 \\
FR-WAL-03 & Cập nhật thông tin ví tiền & T, D & JUnit 5 / JavaFX UI &
Cập nhật tên ví và trạng thái; cấm đổi trực tiếp số dư khả dụng \\
FR-WAL-04 & Thống kê cơ cấu tài sản theo ví & T, D & JUnit 5 / JavaFX UI
& Tính toán chính xác tỷ trọng phần trăm số dư từng ví trên tổng tài sản
thanh khoản \\
FR-TXN-01 & Tạo mới giao dịch thu nhập & T, D & JUnit 5 / JavaFX UI & Số
tiền \textgreater{} 0, danh mục Type = Income; số dư ví nhận tăng đúng
bằng số tiền thu \\
FR-TXN-02 & Tạo mới giao dịch chi tiêu & T, D & JUnit 5 / JavaFX UI & Số
tiền \textgreater{} 0, danh mục Type = Expense, số dư ví chi đủ thanh
toán; trừ số dư ví chính xác \\
FR-TXN-03 & Tra cứu và lọc lịch sử giao dịch & T, D & JUnit 5 / JavaFX
UI & Lọc chính xác theo khoảng ngày, ví, danh mục, từ khóa; tính đúng
tổng thu chi \\
FR-TXN-04 & Chỉnh sửa giao dịch thu chi & T, D & JUnit 5 / JavaFX UI &
Hoàn tác số tiền cũ trên ví cũ, áp dụng số tiền mới trên ví mới; số dư
ví nhất quán \\
FR-TXN-05 & Xóa giao dịch thu chi & T, D & JUnit 5 / JavaFX UI & Xóa
giao dịch thành công và phục hồi số dư ví chính xác theo quy tắc đảo
ngược \\
FR-TRF-01 & Chuyển tiền giữa hai ví nội bộ & T, D & JUnit 5 / JavaFX UI
& Ghi sổ kép đối ứng: ví nguồn trừ đúng bằng ví đích cộng; bảo toàn tổng
tiền \\
FR-TRF-02 & Ghi nhận phí chuyển tiền phát sinh & T, D & JUnit 5 / JavaFX
UI & Ví nguồn bị trừ thêm khoản phí; tạo giao dịch chi phí chuyển khoản
tự động \\
FR-TRF-03 & Lịch sử và hủy chuyển tiền nội bộ & T, D & JUnit 5 / JavaFX
UI & Hủy chuyển tiền hoàn lại tiền cho cả hai ví và xóa phí phát sinh
liên quan \\
FR-CAT-01 & Khởi tạo và quản lý cấu trúc phân loại 3 cấp & T, I & JUnit
5 / Inspection & Cây danh mục 3 cấp (Type -\textgreater{} Group
-\textgreater{} Category) đúng chuẩn phân cấp \\
FR-CAT-02 & Tùy biến danh mục người dùng & T, D & JUnit 5 / JavaFX UI &
Thêm mới danh mục con thành công; tên danh mục không trùng trong cùng
một nhóm \\
FR-CAT-03 & Ràng buộc toàn vẹn khi sửa và xóa danh mục & T, D & JUnit 5
/ JavaFX UI & Chặn xóa danh mục đang có giao dịch tham chiếu; hỗ trợ gom
danh mục an toàn \\
FR-DEBT-01 & Ghi nhận khoản cho vay (Lend) & T, D & JUnit 5 / JavaFX UI
& Trừ tiền ví cho vay; ghi nhận dư nợ phải thu của đối tác; cập nhật sổ
nợ \\
FR-DEBT-02 & Ghi nhận khoản đi vay (Borrow) & T, D & JUnit 5 / JavaFX UI
& Cộng tiền vào ví nhận; ghi nhận nghĩa vụ nợ phải trả của đối tác \\
FR-DEBT-03 & Thu hồi nợ cho vay (Collect) & T, D & JUnit 5 / JavaFX UI &
Cộng tiền vào ví; giảm trừ dư nợ phải thu; chuyển trạng thái PAID khi
hết nợ \\
FR-DEBT-04 & Trả nợ đi vay (Repay) & T, D & JUnit 5 / JavaFX UI & Trừ
tiền ví chi trả; giảm trừ nghĩa vụ nợ phải trả; kiểm tra số tiền trả
không vượt nợ \\
FR-DEBT-05 & Xóa nợ khó đòi và miễn trừ nợ & T, D & JUnit 5 / JavaFX UI
& Đưa dư nợ về 0 với trạng thái WRITTEN\_OFF; số dư các ví tiền giữ
nguyên không đổi \\
FR-SAV-01 & Mở sổ tiết kiệm mới & T, D & JUnit 5 / JavaFX UI & Trừ tiền
ví nguồn gửi; ghi nhận sổ tiết kiệm với gốc, kỳ hạn, lãi suất và ngày
đáo hạn \\
FR-SAV-02 & Tính toán tiền lãi dự kiến và theo dõi kỳ hạn & T, A & JUnit
5 / Analysis & Tính đúng tiền lãi dự kiến theo công thức BR-FIN-05; phân
loại trạng thái đáo hạn chuẩn \\
FR-SAV-03 & Tất toán sổ tiết kiệm đúng hạn & T, D & JUnit 5 / JavaFX UI
& Cộng cả gốc và đủ lãi có kỳ hạn vào ví nhận; chuyển trạng thái sổ sang
CLOSED \\
FR-SAV-04 & Tất toán sổ tiết kiệm trước hạn & T, D & JUnit 5 / JavaFX UI
& Cộng gốc và tính lãi không kỳ hạn theo số ngày thực gửi; tất toán an
toàn \\
FR-INV-01 & Ghi nhận lệnh mua chứng chỉ quỹ & T, D & JUnit 5 / JavaFX UI
& Trừ tiền ví mua; tính chính xác số lượng chứng chỉ quỹ khớp lệnh với
scale 4 \\
FR-INV-02 & Cập nhật giá trị tài sản ròng (NAV Update) & T, D & JUnit 5
/ JavaFX UI & Lưu lịch sử biến động NAV theo ngày; cập nhật giá trị thị
trường của danh mục quỹ \\
FR-INV-03 & Định giá thị trường và tính lãi lỗ chưa thực hiện & T, A &
JUnit 5 / Analysis & Tính đúng giá trị hiện tại theo NAV mới nhất, tỷ
suất sinh lời và lãi lỗ chưa thực hiện \\
FR-INV-04 & Ghi nhận bán chứng chỉ quỹ & T, D & JUnit 5 / JavaFX UI &
Giảm số lượng chứng chỉ quỹ sở hữu; tính tiền thu ròng sau thuế phí và
cộng vào ví \\
FR-FND-01 & Thiết lập quỹ tài chính mục tiêu & T, D & JUnit 5 / JavaFX
UI & Khởi tạo mục tiêu với số tiền đích và hạn chót; trạng thái
IN\_PROGRESS \\
FR-FND-02 & Phân bổ tài sản vào quỹ mục tiêu & T, D & JUnit 5 / JavaFX
UI & Gán tài sản (tiền mặt, sổ tiết kiệm, quỹ) vào mục tiêu; kiểm soát
không phân bổ vượt 100\% \\
FR-FND-03 & Theo dõi tiến độ tích lũy quỹ & T, D & JUnit 5 / JavaFX UI &
Tính đúng tỷ lệ phần trăm hoàn thành; kích hoạt trạng thái REACHED khi
đạt 100\% \\
FR-FND-04 & Điều chỉnh và tái phân bổ quỹ mục tiêu & T, D & JUnit 5 /
JavaFX UI & Điều chỉnh số tiền mục tiêu hoặc tái cấu trúc danh mục tài
sản phân bổ an toàn \\
FR-REP-01 & Bảng thống kê dòng tiền trượt 12 tháng & T, A & JUnit 5 /
Analysis & Trượt đúng 12 tháng liên tục; tổng hợp chính xác tổng thu,
tổng chi, dòng tiền ròng và tỷ lệ tiết kiệm \\
FR-REP-02 & Báo cáo cơ cấu thu chi theo nhóm và danh mục & T, D & JUnit
5 / JavaFX UI & Tính đúng tổng số tiền và tỷ trọng phần trăm chi tiêu
của từng nhóm nhu cầu và danh mục \\
FR-REP-03 & Báo cáo tổng tài sản và giá trị tài sản ròng & T, A & JUnit
5 / Analysis & Tổng hợp chuẩn xác: Tổng tài sản = Tiền mặt + Tiết kiệm +
Đầu tư + Nợ phải thu - Nợ phải trả \\
FR-REP-04 & Sao kê công nợ chi tiết theo đối tượng & T, D & JUnit 5 /
JavaFX UI & Liệt kê chi tiết toàn bộ lịch sử vay, trả, dư nợ còn lại và
trạng thái theo từng đối tác \\
FR-DAT-01 & Xuất dữ liệu tài chính (Export) & T, I & JUnit 5 /
Inspection & Xuất toàn bộ dữ liệu ra tệp JSON chuẩn UTF-8; bảo toàn
nguyên vẹn kiểu số BigDecimal \\
FR-DAT-02 & Nhập dữ liệu và phục hồi (Import \& Restore) & T, D & JUnit
5 / JavaFX UI & Kiểm tra tính hợp lệ dữ liệu 4 bước trước khi nạp; khôi
phục hoàn toàn trạng thái hệ thống \\
FR-DAT-03 & Kiểm tra tính toàn vẹn dữ liệu (Data Integrity Audit) & T, A
& JUnit 5 / Analysis & Thuật toán kiểm toán quét toàn bộ sổ cái; phát
hiện 100\% các sai lệch số dư nếu có \\
\end{longtable}
}

\paragraph{8.2.2 Ma trận phương pháp kiểm chứng yêu cầu phi chức năng (6
NFRs)}\label{822-ma-trux1eadn-phux1b0ux1a1ng-phuxe1p-kiux1ec3m-chux1ee9ng-yuxeau-cux1ea7u-phi-chux1ee9c-nux103ng-6-nfrs}

{\def\LTcaptype{none}
\begin{longtable}[]{|>{\raggedright\arraybackslash}p{0.07\linewidth}|>{\raggedright\arraybackslash}p{0.13\linewidth}|>{\raggedright\arraybackslash}p{0.10\linewidth}|>{\raggedright\arraybackslash}p{0.18\linewidth}|>{\raggedright\arraybackslash}p{0.13\linewidth}|>{\raggedright\arraybackslash}p{0.19\linewidth}|}
\hline
\everycr{\noalign{\ifnum0=`}\fi\hrule height 0.4pt\futurelet\reserved@a\LT@@tabularcr}
\textbf{Mã yêu cầu} & \textbf{Tên yêu cầu phi chức năng} & \textbf{Phân loại chất lượng (ISO 25010)} & \textbf{Phương pháp chính} & \textbf{Công cụ và kỹ thuật đo lường} & \textbf{Ngưỡng định lượng nghiệm thu} \\
\hline
\endfirsthead
\hline
\textbf{Mã yêu cầu} & \textbf{Tên yêu cầu phi chức năng} & \textbf{Phân loại chất lượng (ISO 25010)} & \textbf{Phương pháp chính} & \textbf{Công cụ và kỹ thuật đo lường} & \textbf{Ngưỡng định lượng nghiệm thu} \\
\hline
\endhead
\hline
\endfoot
\hline
\endlastfoot
NFR-01 & Độ chính xác và toàn vẹn dữ liệu tài chính & Functional
Accuracy / Data Integrity & T, I & JUnit 5, Checkstyle, ArchUnit & Sai
số tích lũy = 0.00 VNĐ; 0\% kiểu float/double; vi phạm BR-FIN-02 =
0.00\% \\
NFR-02 & Hiệu năng và độ trễ phản hồi & Perfor\allowbreak mance Efficiency / Time
Behaviour & T, A & JMH, VisualVM, System.nanoTime & Ghi sổ \textless{}
100ms; lọc dữ liệu \textless{} 300ms; báo cáo \textless{} 500ms; RAM
\textless{} 256MB; khởi động \textless{} 2s \\
NFR-03 & Độ tin cậy và lưu trữ an toàn & Reliability / Fault Tolerance
\& Recoverability & T, I & Fault Injection Test, Process Kill & Tỷ lệ
hỏng tệp = 0.00\%; ghi tệp nguyên tử ATOMIC\_MOVE; lưu trữ \(\ge 5\) bản
sao lưu \\
NFR-04 & Khả năng bảo trì và kiến trúc hướng đối tượng & Maintain\allowbreak ability
/ Modularity \& Testability & A, I & JaCoCo, SonarQube, Metrics Reloaded
& Test coverage dòng \(\ge 80\%\), nhánh \(\ge 75\%\); độ phức tạp
\(V(G) \le 10\); 0 cyclic dependency \\
NFR-05 & Tính khả dụng và công thái học & Usability / Learnability \&
Error Protection & D, I & Usability Walkthrough, Error Catalog & Ghi
nhận chi tiêu \(\le 3\) tương tác; 100\% thông báo lỗi tiếng Việt rõ
ràng có gợi ý xử lý \\
NFR-06 & Tính tương thích và môi trường thực thi & Portability /
Adaptability \& Installability & T, I & Cross-Platform CI
(Win/\allowbreak macOS/\allowbreak Linux) & 100\% pass kiểm thử trên cả 3 hệ điều hành; 100\%
tệp dùng UTF-8; 0 native library \\
\end{longtable}
}

\end{document}
